\documentclass[11pt,a4paper]{scrartcl}
\usepackage[margin=22mm,headheight=18pt]{geometry}
\usepackage{fontspec}
\usepackage{xeCJK}
\IfFontExistsTF{Songti SC}{\setCJKmainfont{Songti SC}}{\setCJKmainfont{FandolSong-Regular.otf}[BoldFont=FandolSong-Bold.otf]}
\IfFontExistsTF{Heiti SC}{\setCJKsansfont{Heiti SC}}{\setCJKsansfont{FandolHei-Regular.otf}[BoldFont=FandolHei-Bold.otf]}
\usepackage{amsmath,amssymb,mathtools,mathrsfs,bm,cancel}
\usepackage{mathpazo}
\setmainfont{texgyrepagella-regular.otf}[BoldFont=texgyrepagella-bold.otf,ItalicFont=texgyrepagella-italic.otf,BoldItalicFont=texgyrepagella-bolditalic.otf]
\usepackage[normalem]{ulem}
\usepackage[svgnames,table]{xcolor}
\usepackage[most]{tcolorbox}
\usepackage{tikz}
\usetikzlibrary{arrows.meta,positioning,calc,matrix,patterns,decorations.markings,shapes.geometric}
\usepackage{enumitem,booktabs,longtable,array,tabularx,multicol}
\usepackage{microtype,fancyhdr}
\definecolor{ChapterBlue}{RGB}{20,60,120}
\definecolor{NoteRed}{RGB}{179,0,0}
\definecolor{AnnotationBlue}{RGB}{30,65,135}
\setkomafont{section}{\Large\bfseries\color{ChapterBlue}}
\setkomafont{subsection}{\large\bfseries\color{ChapterBlue}}
\setlist{nosep,leftmargin=*,topsep=3pt}
\setlength{\parindent}{0pt}
\setlength{\parskip}{5pt plus 1pt}
\linespread{1.12}
\allowdisplaybreaks
\emergencystretch=3em
\widowpenalty=10000
\clubpenalty=10000
\pagestyle{fancy}
\fancyhf{}
\fancyhead[C]{\small\nouppercase{\leftmark}}
\fancyfoot[C]{\thepage}
\renewcommand{\headrulewidth}{0.3pt}
% Original scan boundaries are internal anchors only; no visible page labels.
\newcommand{\sourcepage}[1]{\par\phantomsection\label{source:#1}}
\newif\ifchapterstart
\newcommand{\noteschapter}[1]{\clearpage\section{#1}\global\chapterstarttrue}
\newcommand{\notessection}[1]{\ifchapterstart\global\chapterstartfalse\else\clearpage\fi\subsection{#1}}
\renewcommand{\sectionmark}[1]{\markboth{}{}}
\renewcommand{\subsectionmark}[1]{\markboth{\thesubsection\quad #1}{}}
\renewcommand{\contentsname}{目录}
\setcounter{tocdepth}{2}
\setcounter{secnumdepth}{2}
\newcommand{\unclear}[2]{\textnormal{\textcolor{NoteRed}{［无法辨认 #1：#2］}}}
\newcommand{\ann}[1]{{\color{AnnotationBlue}#1}}
\newcommand{\redann}[1]{{\color{NoteRed}#1}}
\newcommand{\circled}[1]{\mbox{\textcircled{\scriptsize #1}}}
\newcommand{\R}{\mathbb{R}}
\newcommand{\N}{\mathbb{N}}
\newcommand{\Z}{\mathbb{Z}}
\newcommand{\Q}{\mathbb{Q}}
\newcommand{\C}{\mathbb{C}}
\newcommand{\F}{\mathbb{F}}
\newcommand{\V}{\mathcal{V}}
\newcommand{\E}{\mathcal{E}}
\newcommand{\abs}[1]{\left|#1\right|}
\newcommand{\norm}[1]{\left\|#1\right\|}
\newcommand{\inner}[2]{\langle #1,#2\rangle}
\newcommand{\floor}[1]{\lfloor #1\rfloor}
\newcommand{\ceil}[1]{\lceil #1\rceil}
\DeclareMathOperator{\rank}{rank}
\DeclareMathOperator{\tr}{tr}
\DeclareMathOperator{\Span}{span}
\DeclareMathOperator{\diag}{diag}
\DeclareMathOperator{\im}{Im}
\DeclareMathOperator{\id}{id}
\DeclareMathOperator{\Hom}{Hom}
\DeclareMathOperator{\End}{End}
\DeclareMathOperator{\Ker}{Ker}
\DeclareMathOperator{\sgn}{sgn}
\newtcolorbox{noteblock}[1]{enhanced,breakable,colback=ChapterBlue!2,colframe=ChapterBlue!60,title={#1},fonttitle=\bfseries,boxrule=0.4pt,arc=1mm}
\newenvironment{proof}{\par\textbf{Proof.}\ }{\par}
\tikzset{vertex/.style={circle,draw,inner sep=2pt},edge/.style={thick},dedge/.style={thick,-Stealth}}

% Clickable contents and PDF outline.
\usepackage[unicode,hidelinks,bookmarksnumbered,bookmarksopen,linktoc=all]{hyperref}
\hypersetup{pdftitle={抽象代数},pdfsubject={完整转录，含可点击目录},pdfauthor={芬达不带气儿}}

% Coordinate system: x and y are per-mille of original 1415x2000 preview image.
% Approximate each original baseline as (pixel_x/1.415,pixel_y/2).
% Width is also in per-mille of the original image width.
\newlength{\handunit}\setlength{\handunit}{0.194mm}
\newlength{\handyunit}\setlength{\handyunit}{0.265mm}
\newsavebox{\handbox}
\newcommand{\handtext}[6]{%
  \node[anchor=base west,inner sep=0pt,outer sep=0pt] at (#1,#2)
    {{\normalfont\fontsize{#4}{#4}\selectfont\color{#5}%
      \sbox{\handbox}{#6}%
      \ifdim\wd\handbox>#3\handunit\resizebox{#3\handunit}{!}{\usebox{\handbox}}\else\usebox{\handbox}\fi}};}
\newcommand{\handline}[5][black]{\handtext{#2}{#3}{#4}{10.5}{#1}{#5}}
\newcommand{\handtitle}[4]{\handtext{#1}{#2}{#3}{13}{black}{\bfseries #4}}
\newcommand{\handnote}[5][AnnotationBlue]{\handtext{#2}{#3}{#4}{7.2}{#1}{#5}}
\newcommand{\handarrow}[5][AnnotationBlue]{\draw[->,color=#1,line width=.35pt] (#2,#3)--(#4,#5);}
\newcommand{\handdiagram}[5]{\node[anchor=north west,inner sep=0pt] at (#1,#2)
 {\resizebox{#3\handunit}{#4\handyunit}{\begin{tikzpicture}[x=1cm,y=1cm]#5\end{tikzpicture}}};}
\newcommand{\handchapterentry}[1]{\refstepcounter{section}\setcounter{subsection}{0}\addcontentsline{toc}{section}{\protect\numberline{\thesection}#1}}
\newcommand{\handsectionentry}[1]{\refstepcounter{subsection}\addcontentsline{toc}{subsection}{\protect\numberline{\thesubsection}#1}}
\newcommand{\handregister}[1]{\ifcsname handchapter@#1\endcsname\handchapterentry{\csname handchapter@#1\endcsname}\fi
 \ifcsname handsection@#1\endcsname\handsectionentry{\csname handsection@#1\endcsname}\fi}
\newenvironment{handpage}[1]{\clearpage\thispagestyle{fancy}\markboth{\csname handheader@#1\endcsname}{}\null\handregister{#1}%
 \begin{tikzpicture}[remember picture,overlay]
 \begin{scope}[shift={([xshift=8mm,yshift=-13mm]current page.north west)},x=0.194mm,y=-0.265mm]
}{\end{scope}\end{tikzpicture}\par}

\expandafter\def\csname handchapter@1\endcsname{域、线性空间}
\expandafter\def\csname handsection@1\endcsname{定义和例子}
\expandafter\def\csname handheader@1\endcsname{1.1 定义和例子}
\expandafter\def\csname handheader@2\endcsname{1.1 定义和例子}
\expandafter\def\csname handheader@3\endcsname{1.1 定义和例子}
\expandafter\def\csname handsection@4\endcsname{域同态}
\expandafter\def\csname handheader@4\endcsname{1.2 域同态}
\expandafter\def\csname handsection@5\endcsname{域的特征}
\expandafter\def\csname handheader@5\endcsname{1.3 域的特征}
\expandafter\def\csname handheader@6\endcsname{1.3 域的特征}
\expandafter\def\csname handsection@7\endcsname{域的扩张}
\expandafter\def\csname handheader@7\endcsname{1.4 域的扩张}
\expandafter\def\csname handheader@8\endcsname{1.4 域的扩张}
\expandafter\def\csname handheader@9\endcsname{1.4 域的扩张}
\expandafter\def\csname handheader@10\endcsname{1.4 域的扩张}
\expandafter\def\csname handheader@11\endcsname{1.4 域的扩张}
\expandafter\def\csname handheader@12\endcsname{1.4 域的扩张}
\expandafter\def\csname handsection@13\endcsname{代数闭包}
\expandafter\def\csname handheader@13\endcsname{1.5 代数闭包}
\expandafter\def\csname handsection@14\endcsname{Galois 群}
\expandafter\def\csname handheader@14\endcsname{1.6 Galois 群}
\expandafter\def\csname handheader@15\endcsname{1.6 Galois 群}
\expandafter\def\csname handheader@16\endcsname{1.6 Galois 群}
\expandafter\def\csname handchapter@17\endcsname{环}
\expandafter\def\csname handsection@17\endcsname{定义，例子}
\expandafter\def\csname handheader@17\endcsname{2.1 定义，例子}
\expandafter\def\csname handheader@18\endcsname{2.1 定义，例子}
\expandafter\def\csname handheader@19\endcsname{2.1 定义，例子}
\expandafter\def\csname handheader@20\endcsname{2.1 定义，例子}
\expandafter\def\csname handheader@21\endcsname{2.1 定义，例子}
\expandafter\def\csname handheader@22\endcsname{2.1 定义，例子}
\expandafter\def\csname handheader@23\endcsname{2.1 定义，例子}
\expandafter\def\csname handsection@24\endcsname{同态的基本概念}
\expandafter\def\csname handheader@24\endcsname{2.2 同态的基本概念}
\expandafter\def\csname handheader@25\endcsname{2.2 同态的基本概念}
\expandafter\def\csname handsection@26\endcsname{子环、理想、商环}
\expandafter\def\csname handheader@26\endcsname{2.3 子环、理想、商环}
\expandafter\def\csname handheader@27\endcsname{2.3 子环、理想、商环}
\expandafter\def\csname handheader@28\endcsname{2.3 子环、理想、商环}
\expandafter\def\csname handsection@29\endcsname{特征}
\expandafter\def\csname handheader@29\endcsname{2.4 特征}
\expandafter\def\csname handheader@30\endcsname{2.4 特征}
\expandafter\def\csname handsection@31\endcsname{同态基本定理}
\expandafter\def\csname handheader@31\endcsname{2.5 同态基本定理}
\expandafter\def\csname handheader@32\endcsname{2.5 同态基本定理}
\expandafter\def\csname handheader@33\endcsname{2.5 同态基本定理}
\expandafter\def\csname handheader@34\endcsname{2.5 同态基本定理}
\expandafter\def\csname handheader@35\endcsname{2.5 同态基本定理}
\expandafter\def\csname handheader@36\endcsname{2.5 同态基本定理}
\expandafter\def\csname handheader@37\endcsname{2.5 同态基本定理}
\expandafter\def\csname handheader@38\endcsname{2.5 同态基本定理}
\expandafter\def\csname handheader@39\endcsname{2.5 同态基本定理}
\expandafter\def\csname handsection@40\endcsname{整环与整除性}
\expandafter\def\csname handheader@40\endcsname{2.6 整环与整除性}
\expandafter\def\csname handheader@41\endcsname{2.6 整环与整除性}
\expandafter\def\csname handheader@42\endcsname{2.6 整环与整除性}
\expandafter\def\csname handheader@43\endcsname{2.6 整环与整除性}
\expandafter\def\csname handheader@44\endcsname{2.6 整环与整除性}
\expandafter\def\csname handheader@45\endcsname{2.6 整环与整除性}
\expandafter\def\csname handheader@46\endcsname{2.6 整环与整除性}
\expandafter\def\csname handheader@47\endcsname{2.6 整环与整除性}
\expandafter\def\csname handheader@48\endcsname{2.6 整环与整除性}
\expandafter\def\csname handheader@49\endcsname{2.6 整环与整除性}
\expandafter\def\csname handheader@50\endcsname{2.6 整环与整除性}
\expandafter\def\csname handheader@51\endcsname{2.6 整环与整除性}
\expandafter\def\csname handheader@52\endcsname{2.6 整环与整除性}
\expandafter\def\csname handheader@53\endcsname{2.6 整环与整除性}
\expandafter\def\csname handsection@54\endcsname{多项式环}
\expandafter\def\csname handheader@54\endcsname{2.7 多项式环}
\expandafter\def\csname handheader@55\endcsname{2.7 多项式环}
\expandafter\def\csname handheader@56\endcsname{2.7 多项式环}
\expandafter\def\csname handheader@57\endcsname{2.7 多项式环}
\expandafter\def\csname handheader@58\endcsname{2.7 多项式环}
\expandafter\def\csname handchapter@59\endcsname{群}
\expandafter\def\csname handsection@59\endcsname{群与群作用}
\expandafter\def\csname handheader@59\endcsname{3.1 群与群作用}
\expandafter\def\csname handheader@60\endcsname{3.1 群与群作用}
\expandafter\def\csname handheader@61\endcsname{3.1 群与群作用}
\expandafter\def\csname handheader@62\endcsname{3.1 群与群作用}
\expandafter\def\csname handsection@63\endcsname{群同态}
\expandafter\def\csname handheader@63\endcsname{3.2 群同态}
\expandafter\def\csname handheader@64\endcsname{3.2 群同态}
\expandafter\def\csname handsection@65\endcsname{子群、正规子群、商群}
\expandafter\def\csname handheader@65\endcsname{3.3 子群、正规子群、商群}
\expandafter\def\csname handheader@66\endcsname{3.3 子群、正规子群、商群}
\expandafter\def\csname handheader@67\endcsname{3.3 子群、正规子群、商群}
\expandafter\def\csname handheader@68\endcsname{3.3 子群、正规子群、商群}
\expandafter\def\csname handheader@69\endcsname{3.3 子群、正规子群、商群}
\expandafter\def\csname handheader@70\endcsname{3.3 子群、正规子群、商群}
\expandafter\def\csname handheader@71\endcsname{3.3 子群、正规子群、商群}
\expandafter\def\csname handsection@72\endcsname{同态基本定理}
\expandafter\def\csname handheader@72\endcsname{3.4 同态基本定理}
\expandafter\def\csname handheader@73\endcsname{3.4 同态基本定理}
\expandafter\def\csname handheader@74\endcsname{3.4 同态基本定理}
\expandafter\def\csname handheader@75\endcsname{3.4 同态基本定理}
\expandafter\def\csname handheader@76\endcsname{3.4 同态基本定理}
\expandafter\def\csname handsection@77\endcsname{群作用}
\expandafter\def\csname handheader@77\endcsname{3.5 群作用}
\expandafter\def\csname handheader@78\endcsname{3.5 群作用}
\expandafter\def\csname handheader@79\endcsname{3.5 群作用}
\expandafter\def\csname handsection@80\endcsname{Sylow 定理}
\expandafter\def\csname handheader@80\endcsname{3.6 Sylow 定理}
\expandafter\def\csname handheader@81\endcsname{3.6 Sylow 定理}
\expandafter\def\csname handheader@82\endcsname{3.6 Sylow 定理}
\expandafter\def\csname handheader@83\endcsname{3.6 Sylow 定理}
\expandafter\def\csname handheader@84\endcsname{3.6 Sylow 定理}
\expandafter\def\csname handheader@85\endcsname{3.6 Sylow 定理}
\expandafter\def\csname handheader@86\endcsname{3.6 Sylow 定理}
\expandafter\def\csname handsection@87\endcsname{对称群}
\expandafter\def\csname handheader@87\endcsname{3.7 对称群}
\expandafter\def\csname handheader@88\endcsname{3.7 对称群}
\expandafter\def\csname handheader@89\endcsname{3.7 对称群}
\expandafter\def\csname handheader@90\endcsname{3.7 对称群}
\expandafter\def\csname handheader@91\endcsname{3.7 对称群}
\expandafter\def\csname handheader@92\endcsname{3.7 对称群}
\expandafter\def\csname handheader@93\endcsname{3.7 对称群}
\expandafter\def\csname handheader@94\endcsname{3.7 对称群}
\expandafter\def\csname handsection@95\endcsname{交换群}
\expandafter\def\csname handheader@95\endcsname{3.8 交换群}
\expandafter\def\csname handheader@96\endcsname{3.8 交换群}
\expandafter\def\csname handheader@97\endcsname{3.8 交换群}
\expandafter\def\csname handheader@98\endcsname{3.8 交换群}
\expandafter\def\csname handheader@99\endcsname{3.8 交换群}
\expandafter\def\csname handheader@100\endcsname{3.8 交换群}
\expandafter\def\csname handheader@101\endcsname{3.8 交换群}
\expandafter\def\csname handheader@102\endcsname{3.8 交换群}
\expandafter\def\csname handheader@103\endcsname{3.8 交换群}
\expandafter\def\csname handheader@104\endcsname{3.8 交换群}
\expandafter\def\csname handheader@105\endcsname{3.8 交换群}
\expandafter\def\csname handheader@106\endcsname{3.8 交换群}
\expandafter\def\csname handsection@107\endcsname{可解群}
\expandafter\def\csname handheader@107\endcsname{3.9 可解群}
\expandafter\def\csname handheader@108\endcsname{3.9 可解群}
\expandafter\def\csname handheader@109\endcsname{3.9 可解群}
\expandafter\def\csname handheader@110\endcsname{3.9 可解群}
\begin{document}
\hypersetup{pageanchor=false}
\begin{titlepage}
\thispagestyle{empty}
\null
\begin{tikzpicture}[remember picture,overlay]
\node[anchor=center,inner sep=0pt] at ([yshift=-0.382\paperheight]current page.north)
  {{\sffamily\bfseries\fontsize{40}{50}\selectfont 抽象代数}};
\node[anchor=center,inner sep=0pt] at ([yshift=-0.618\paperheight]current page.north)
  {{\rmfamily\fontsize{24}{32}\selectfont By 芬达不带气儿}};
\end{tikzpicture}
\end{titlepage}
\pagenumbering{roman}
\hypersetup{pageanchor=true}
\pagestyle{plain}
\renewcommand{\contentsname}{目录}
\tableofcontents
\clearpage
\pagenumbering{arabic}
\pagestyle{fancy}
% ===== Line-layout page-001.tex =====
\begin{handpage}{1}
\handtitle{72}{124}{245}{\S 1. 域、线性空间}
\handline[NoteRed]{312}{120}{480}{（用线性空间的观点看待域扩张很重要）}
\handtitle{72}{153}{310}{\S 1.1 定义和例子}
\handline{72}{183}{790}{\textbf{Def.}\quad $F:$ 集合\quad $0_F,1_F\in F\ (0_F\ne1_F)$\quad $F$ 上有两个二元运算：$+,\cdot$.}
\handline{72}{210}{650}{$+ : F\times F\to F\qquad(x,y)\mapsto x+y$}
\handline{72}{238}{650}{$\cdot : F\times F\to F\qquad(x,y)\mapsto x\cdot y$}
\handline{72}{268}{700}{满足\quad\circled{1} 加法结合律\qquad $(x+y)+z=x+(y+z)$}
\handline{72}{295}{580}{\circled{2} 加法交换律\qquad $x+y=y+x$}
\handline{72}{325}{600}{\circled{3} 加法单位元\qquad $x+0_F=0_F+x=x$}
\handline{72}{354}{650}{\circled{4} 加法逆存在\qquad $\forall x\in F,\ \exists y\in F\ \text{s.t.}\ x+y=0_F$}
\handline[AnnotationBlue]{72}{388}{805}{这样的 $y$ 是唯一的：若 $y_1,y_2$ 为 $x$ 的加法逆，则 $y_1=y_1+0_F=y_1+(x+y_2)=(y_1+x)+y_2=0_F+y_2=y_2$}
\handline[AnnotationBlue]{72}{416}{650}{称这个 $y$ 为 $x$ 的加法逆元，记为 $-x$}
\handline{72}{446}{600}{\circled{5} 乘法结合律\qquad $(x\cdot y)\cdot z=x\cdot(y\cdot z)$}
\handline{72}{477}{580}{\circled{6} 乘法交换律\qquad $x\cdot y=y\cdot x$}
\handline{72}{510}{620}{\circled{7} 乘法单位元\qquad $x\cdot1_F=1_F\cdot x=x$}
\handline{72}{543}{620}{\circled{8} 乘法逆存在\quad 记 $F^*=F\setminus\{0_F\}$}
\handline{72}{574}{590}{$\forall x\in F^*\quad\exists y\in F\qquad x\cdot y=y\cdot x=1_F$}
\handline[AnnotationBlue]{72}{605}{655}{类似可证 $y$ 唯一，称为 $x$ 的乘法逆元，记为 $x^{-1}$}
\handline{72}{636}{650}{\circled{9} 分配律\qquad $x(y+z)=xy+xz\qquad(x+y)z=xz+yz$}
\handline{72}{675}{200}{\textbf{Prop.}}
\handline{72}{706}{260}{\circled{1}\quad $0_F\cdot0_F=0_F$}
\handline{72}{737}{310}{$0\cdot(0+1)=0\cdot1=0$}
\handline{72}{768}{425}{$\therefore\ 0\cdot0+0\cdot1=0\cdot0+0=0\quad\therefore\ 0\cdot0=0$}
\handline{72}{802}{440}{\circled{2}\quad $\forall x\in F,\ x\cdot0_F=0_F$.}
\handline{72}{835}{790}{$x\cdot0=x(0+0)=x\cdot0+x\cdot0\qquad\text{令 }y=-(x\cdot0)\ \Rightarrow\ x\cdot0+y=x\cdot0+y+x\cdot0$}
\handline{72}{865}{240}{$\therefore\ x\cdot0=0$.}
\end{handpage}

% ===== Line-layout page-002.tex =====
\begin{handpage}{2}
\handline{72}{94}{835}{\textbf{Ex.}\quad\circled{1}\quad $\Q(\sqrt2)=\{x=a+b\sqrt2\mid a,b\in\Q\}\quad\Q(\sqrt[3]2)=\{x=a+b\sqrt[3]2+c(\sqrt[3]2)^2\mid a,b,c\in\Q\}$.}
\handline{72}{123}{530}{更一般地，$\alpha\in\C$，$\Q(\alpha)$.}
\handline{72}{153}{690}{设 $\alpha$ 在 $\Q$ 上的极小多项式为 $f(x)$，令 $n=\deg(f(x))$.}
\handline{72}{184}{690}{则 $\Q(\alpha)=\{x_1+x_2\alpha+\cdots+x_n\alpha^{n-1}\mid x_i\in\Q,1\le i\le n\}$.}
\handline{72}{215}{720}{除了逆元存在之外的其他条件很容易验证.}
\handline[AnnotationBlue]{72}{245}{730}{下面考虑 $g(\alpha)=x_1+x_2\alpha+\cdots+x_n\alpha^{n-1}$ 的逆元}
\handline{72}{275}{690}{$g(x)\in\Q[x]\quad\deg(g(x))=n-1$，而 $\deg(f(x))=n$}
\handline{72}{305}{740}{$\therefore g(x),f(x)$ 互素\quad$\therefore(g(x),f(x))=1\quad\exists u(x),v(x)\in\Q[x]\ \text{s.t.}\ ug+vf=1$.}
\handline{72}{345}{765}{$\displaystyle\therefore\ \frac1{g(\alpha)}=\frac{u(\alpha)}{u(\alpha)g(\alpha)}=\frac{u(\alpha)}{u(\alpha)g(\alpha)+v(\alpha)f(\alpha)}=\frac{u(\alpha)}1=u(\alpha)$ 为 $g(\alpha)$ 逆元.}
\handline{72}{405}{750}{\circled{2} 考虑 $\Q(\sqrt2,\sqrt3)$ 与 $\Q(\sqrt2+\sqrt3)$}
\handline{72}{435}{810}{$\sqrt2\cdot\sqrt3=\sqrt6,\ \sqrt2\cdot\sqrt6=2\sqrt3,\ \sqrt3\cdot\sqrt6=3\sqrt2\Rightarrow\Q(\sqrt2,\sqrt3)=\{x_1+x_2\sqrt2+x_3\sqrt3+x_4\sqrt6\mid x_1,x_2,x_3,x_4\in\Q\}$.}
\handline[NoteRed]{72}{468}{780}{下面从线性空间的角度证明 $\Q(\sqrt2,\sqrt3)=\Q(\sqrt2+\sqrt3)$}
\handline[AnnotationBlue]{72}{499}{750}{令 $\alpha=\sqrt2+\sqrt3$，$\Rightarrow\alpha^2=5+2\sqrt6$，$\alpha^3=11\sqrt2+9\sqrt3$}
\handline[AnnotationBlue]{72}{533}{790}{$\alpha$ 在 $\Q$ 中的极小多项式为 $f(x)=x^4-10x^2+1=0$，$\deg(f(x))=4$.}
\handline[AnnotationBlue]{72}{570}{800}{$\therefore\Q(\sqrt2+\sqrt3)=\{x_1+x_2\alpha+x_3\alpha^2+x_4\alpha^3\mid x_1,x_2,x_3,x_4\in\Q\}$.}
\handline[AnnotationBlue]{72}{606}{85}{而}
\handline[AnnotationBlue]{130}{606}{280}{$1=1$}
\handline[AnnotationBlue]{72}{635}{300}{$\displaystyle\sqrt2=-\frac92\alpha+\frac12\alpha^3$}
\handline[AnnotationBlue]{72}{666}{300}{$\displaystyle\sqrt3=\frac{11}2\alpha-\frac12\alpha^3$}
\handline[AnnotationBlue]{72}{697}{300}{$\displaystyle\sqrt6=-\frac52+\frac{\alpha^2}2$}
\handline[AnnotationBlue]{392}{606}{300}{$1=1$}
\handline[AnnotationBlue]{329}{635}{300}{$\alpha=\sqrt2+\sqrt3$}
\handline[AnnotationBlue]{327}{666}{300}{$\alpha^2=5+2\sqrt6$}
\handline[AnnotationBlue]{327}{697}{330}{$\alpha^3=11\sqrt2+9\sqrt3$.}
\handline[AnnotationBlue]{72}{752}{825}{$\Q(\sqrt2,\sqrt3)=\Span_{\Q}\{1,\sqrt2,\sqrt3,\sqrt6\}.\quad\Q(\sqrt2+\sqrt3)=\Span_{\Q}\{1,\alpha,\alpha^2,\alpha^3\}$.}
\handline[AnnotationBlue]{72}{788}{820}{且两个空间的基可以互相线性表出；故为同一线性空间}
\handline[AnnotationBlue]{72}{821}{735}{$\therefore\Q(\sqrt2,\sqrt3)=\Q(\sqrt2+\sqrt3)$.}
\handline[NoteRed]{72}{859}{800}{后面可以证明，在一定条件下，有限的代数扩张都是单扩张.}
\handline[NoteRed]{72}{889}{430}{（可分扩张）}
\end{handpage}

% ===== Line-layout page-003.tex =====
\begin{handpage}{3}
\handline{72}{93}{760}{\circled{3}\quad $\F_p=\{\bar0,\bar1,\bar2,\ldots,\overline{p-1}\}\quad\bmod p$ 剩余类\quad（$p$ 素）}
\handline{72}{126}{720}{\textbf{Prop.}\quad $\forall\bar k\in\F_p^*\ \exists\bar l\in\F_p^*\ \text{s.t.}\ \bar k\cdot\bar l=\bar1$.}
\handline[AnnotationBlue]{72}{152}{640}{该命题可以从两个角度来看待.}
\handline[NoteRed]{72}{183}{600}{\circled{1} 经典数论，Bézout 定理；}
\handline[NoteRed]{72}{214}{740}{$\forall\bar k\in\F_p^*,\quad(k,p)=1\quad\therefore\exists l,m\in\N\ \text{s.t.}\ lk+mp=1\Rightarrow\bar k\bar l=\bar1$}
\handline[NoteRed]{72}{241}{720}{\circled{2} 群作用：定义 $\varphi_k:\F_p\to\F_p$}
\handline[NoteRed]{72}{271}{470}{$\bar y\mapsto\overline{ky}\qquad\forall k\in\{1,2,\ldots,p-1\}$.}
\handline[NoteRed]{72}{304}{660}{显然 $\varphi_k(\bar0)=\bar0$.}
\handline[NoteRed]{72}{336}{650}{若 $\varphi_k(\bar y)=\bar0$ 则 $\bar y=\bar0$：}
\handline[NoteRed]{72}{366}{745}{$\varphi_k(\bar y)=\overline{ky}=\bar0\quad\therefore ky\equiv0\pmod p\quad\therefore p\mid ky\quad\therefore(k,p)=1\quad\therefore p\mid y\quad\therefore y\equiv0\pmod p$}
\handline[NoteRed]{72}{399}{720}{$\therefore\varphi_k$ 是单射，而 $\varphi_k$ 是 $\F_p$ 到自身的映射，故也为满射.}
\handline[NoteRed]{72}{430}{720}{$\therefore\bar1$ 有原像，即 $\exists\bar l\in\F_p\ \text{s.t.}\ \varphi_k(\bar l)=\bar1$，$\bar k\cdot\bar l=\bar1$}
\handline{72}{463}{780}{\circled{4} 考虑 $K=\C(x,\sqrt{x^3+2})$ 来自代数几何的例子.令 $y=\sqrt{x^3+2}\quad y^2=x^3+2$.}
\handline{72}{493}{660}{$=\C(x)(y)$，可以类似于 $\Q(\sqrt2)$ 处理.}
\handline{72}{525}{780}{$K=\{R_1(x)+R_2(x)y\mid R_1,R_2\in\C[x]\}$，是一条代数曲线上的亚纯函数域.}
\handline[AnnotationBlue]{72}{557}{730}{是 $\C[x]$ 上的 $2$ 维线性空间.}
\handline[NoteRed]{72}{591}{780}{$\C,\R$ 是 $\Q$ 上的无穷维空间（找不到 $\Q$ 上的有限基），$\C$ 是 $\R$ 上的 $2$ 维线性空间}
\handline{72}{625}{795}{\circled{5}\quad $\F_2=\{\bar0,\bar1\}$，$2$ 元有限域.下面构造 $4$ 元有限域 $\F_4=\F_2(\alpha)=\{x+y\alpha\mid x,y\in\F_2\}$.}
\handline[NoteRed]{72}{657}{760}{下面来看 $\alpha$ 是如何构造的.}
\handline[AnnotationBlue]{72}{690}{790}{考虑 $\F_2[x]$ 上的 $2$ 次不可约多项式：$\F_2[x]$ 上的全体 $2$ 次多项式只有 $4$ 个：}
\handline[AnnotationBlue]{72}{723}{790}{\circled{1} $x^2=x\cdot x$ 可约\quad\circled{2} $x^2+1=x^2-1=(x+1)(x-1)$ 可约\quad\circled{3} $x^2+x=x(x+1)$ 可约}
\handline[AnnotationBlue]{72}{755}{720}{\circled{4}\quad $x^2+x+1$，只有它不可约.}
\handline[NoteRed]{72}{787}{770}{令 $\alpha^2+\alpha+1=0\quad\therefore\alpha^2=-\alpha-1=\alpha+1$.}
\handline[NoteRed]{72}{821}{800}{$\Rightarrow\F_4=\{x+y\alpha\mid x,y\in\F_2,\alpha^2=\alpha+1\}=\{0,1,\alpha,\alpha+1=\alpha^2\}=\{0,1,\alpha,\alpha^2\}$.}
\handline[NoteRed]{72}{850}{800}{由于 $\alpha\cdot\alpha^2=\alpha(\alpha+1)=\alpha^2+\alpha=2\alpha+1=1\quad\therefore\alpha,\alpha^2$ 存在逆，从而 $\F_4$ 是域.}
\handline[NoteRed]{72}{885}{856}{学习群论之后，可以证明：对素数 $p$ 及正整数 $d$，$q=p^d$，一定存在 $q$ 元有限域 $\F_q$（Galois 域）}
\handline[NoteRed]{72}{917}{760}{且 $\F_q$ 在同构意义下唯一}
\handline[AnnotationBlue]{72}{950}{820}{\textbf{Ex.}\quad $\F_9=\F_3(\alpha)$（$\alpha^2=2$，模 $3$ 意义下）\quad$\F_8=\F_2(\alpha)$（$\alpha^3=1+\alpha$，模 $2$ 意义下）.}
\end{handpage}

% ===== Line-layout page-004.tex =====
\begin{handpage}{4}
\handtitle{72}{104}{550}{\S 1.2 域同态}
\handline{72}{132}{190}{\textbf{Def.}}
\handline{72}{163}{795}{(1) $V_1,V_2$ 是 $F$ 上线性空间，$\varphi:V_1\to V_2\quad\alpha\mapsto\varphi(\alpha)$.}
\handline{72}{193}{470}{若 $\varphi(\alpha+\beta)=\varphi(\alpha)+\varphi(\beta)$，$k\in F\quad\varphi(k\alpha)=k\varphi(\alpha)$}
\handline{72}{220}{690}{则称 $\varphi$ 是从 $V_1$ 到 $V_2$ 的同态（线性映射）}
\handline{72}{251}{780}{(2) $F_1,F_2$ 是两个域.\quad$\varphi:F_1\to F_2$，满足}
\handline{72}{282}{790}{\circled{1} $\varphi(0_{F_1})=0_{F_2}$\quad\circled{2} $\varphi(1_{F_1})=1_{F_2}$\quad\circled{3} $\varphi(x+y)=\varphi(x)+\varphi(y)$\quad\circled{4} $\varphi(xy)=\varphi(x)\varphi(y)$}
\handline{72}{309}{720}{则称 $\varphi$ 是 $F_1$ 到 $F_2$ 的同态.}
\handline[AnnotationBlue]{72}{341}{775}{\textit{Remark:}\quad\circled{1}\circled{2}\circled{3}\circled{4}\quad$\Rightarrow\varphi(-x)=-\varphi(x)$，$\varphi(x^{-1})=(\varphi(x))^{-1}$.}
\handline{72}{371}{790}{\textbf{Thm.}\quad 若 $\varphi:F_1\to F_2$ 是同态，则 $\varphi$ 是单射.\quad$F_2\ne\{0\}$.}
\handline{72}{402}{770}{\textbf{Pf.}\quad 若 $\varphi(x_1)=\varphi(x_2)\Rightarrow\varphi(x_1-x_2)=\varphi(x_1)-\varphi(x_2)=0_{F_2}$.}
\handline{72}{430}{680}{若 $x_1-x_2\ne0_{F_1}$ 则 $\exists(x_1-x_2)^{-1}\in F_1$}
\handline{72}{462}{725}{$0_{F_2}=\varphi(x_1-x_2)\cdot\varphi((x_1-x_2)^{-1})=\varphi(1_{F_1})=1_{F_2}$，矛盾！}
\handline{72}{491}{650}{$\therefore x_1=x_2$，即 $\varphi$ 为单射.}
\handline[AnnotationBlue]{72}{521}{790}{域之间的同态一定是单同态，$F_1$ 称为 $F_2$ 的子域.}
\handline{72}{554}{805}{\textbf{Def.}\quad\circled{1} 若 $\varphi:F_1\to F_2$ 同态，是满射，则称 $\varphi$ 是同构.}
\handline{72}{582}{750}{\circled{2} 若 $F_1=F_2=F$，则称 $\varphi$ 是 $F$ 的自同构.}
\handline{72}{615}{780}{\circled{3} 设 $\sigma:F\to F$ 是自同构，则 $E=\{x\in F\mid\sigma(x)=x\}$ 称为 $\sigma$ 的不动域.}
\handline{72}{645}{760}{\textbf{Ex.}\quad\circled{1}\quad $\varphi:\C\to\C$}
\handline{72}{678}{715}{$x+yi\mapsto x-iy$\quad 则 $\varphi$ 是自同构，且不动域为 $\R$.}
\handline{72}{710}{480}{\circled{2}\quad $\Q(\sqrt2)$ 与 $\Q(\sqrt3)$ 之间不存在同态.}
\handline[NoteRed]{72}{745}{755}{若 $\exists\varphi:\Q(\sqrt2)\to\Q(\sqrt3)$.}
\handline[NoteRed]{72}{781}{765}{设 $\alpha=\varphi(\sqrt2)=x+\sqrt3y$，$x,y\in\Q$}
\handline[NoteRed]{72}{812}{800}{则 $\alpha^2=\varphi^2(\sqrt2)=\varphi(2)=\varphi(2\times1)=2\varphi(1)=2=(x+\sqrt3y)^2=x^2+3y^2+2\sqrt3xy$}
\handline[NoteRed]{72}{847}{780}{$\displaystyle\Rightarrow\begin{cases}x^2+3y^2=2,\\2xy=0,\end{cases}$\quad 但在 $\Q$ 上无解，$\therefore$ 不存在这样的 $\varphi$.}
\handline[AnnotationBlue]{72}{929}{355}{$\Q(\sqrt2),\Q(\sqrt3)$ 作为线性空间}
\handline[AnnotationBlue]{72}{952}{340}{是线性同构的，}
\handline[AnnotationBlue]{72}{974}{340}{但不是 $\Q$-代数同构.}
\end{handpage}

% ===== Line-layout page-005.tex =====
\begin{handpage}{5}
\handtitle{72}{78}{720}{\S1.3 域的特征}
\handline{72}{114}{820}{\textbf{Def.} 设 $F$ 是一个域且 $0\ne1$. 定义映射 $N\colon\N\to F$，$n\mapsto N(n)$.}
\handline{72}{141}{740}{满足 $N(0)=0_F$ \redann{\circled{1}}，且 $N(n+1)=N(n)+1_F$ \redann{\circled{2}}.}
\handline{72}{168}{740}{若 $N$ 是单射，则称 $F$ 的特征为 $0$，$\operatorname{char}F=0$.}
\handline{72}{196}{740}{若 $N$ 不是单射，则存在最小的正整数 $p$ s.t. $N(p)=0$，则称 $F$ 特征为 $p$，$\operatorname{char}F=p$.}
\handline[AnnotationBlue]{72}{224}{810}{这里采用了定义自然数的方法（Peano 公理）来得到 $N$ 这个映射的定义.}
\handline[AnnotationBlue]{72}{252}{820}{我们要定义 $N$ 的加法、乘法，即 $N(n+m)=N(n)+N(m)$，$N(nm)=N(n)N(m)$，$N(n-m)=N(n)-N(m)$.}
\handline[NoteRed]{72}{281}{805}{以 $N(n+m)=N(n)+N(m)$ 为例，将其视为关于 $m$ 的一个命题 $P(m)$.}
\handline[NoteRed]{72}{309}{810}{$P(0)\colon\quad N(n+0)=N(n)+N(0)=N(n)+0_F=N(n)$，成立.}
\handline[NoteRed]{72}{367}{810}{若 $P(m)$ 成立. $P(m+1)\colon\ N(n+m+1)\overset{\text{利用 }P(m)}=N(n+1)+N(m)\overset{\text{利用\circled{2}}}=N(n)+N(m)+1_F$}
\handline[NoteRed]{72}{399}{580}{$N(n)+N(m+1)\overset{\text{利用\circled{2}}}=N(n)+N(m)+1_F$}
\handline[NoteRed]{72}{429}{640}{$\therefore N(n+m+1)=N(n)+N(m+1)$，$P(m+1)$ 成立.}
\handline[NoteRed]{72}{460}{760}{类似也定义乘法与减法.}
\handline[AnnotationBlue]{72}{493}{810}{若 $N$ 不单，$\exists n,m\in\N$ s.t. $N(n)=N(m)$，不妨设 $n>m$.}
\handline[AnnotationBlue]{72}{524}{710}{$\therefore N(n-m)=N(n)-N(m)=0$. 于是 $M=\{q\in\N\mid N(q)=0\}$ 非空.}
\handline[AnnotationBlue]{72}{556}{685}{由 $\N$ 的良序性 $\Rightarrow M$ 有最小元 $p$，$N(p)=0$. 则称 $p=\operatorname{char}F$.}
\handline{72}{590}{800}{\textbf{Prop.} $F$ 是域 $\ne\{0\}$，则无零因子.}
\handline{72}{623}{420}{\textbf{Prop.} 若 $\operatorname{char}F=p>0$，则 $p$ 为素数.}
\handline{72}{657}{750}{\textbf{Pf.} 若 $p=qr$，$q,r\in\N$，$0<q,r<p$.}
\handline{72}{689}{680}{$\therefore N(p)=N(qr)=N(q)N(r)=0$.}
\handline{72}{721}{710}{由 $F$ 无零因子 $\Rightarrow N(q)=0$ 或 $N(r)=0$，与 $p$ 最小性矛盾.}
\handline{72}{760}{825}{\textbf{Prop.} 若 $\operatorname{char}F=0$，则 $\Q$ 是 $F$ 的子域. \ann{$\longrightarrow$ 理解为存在单同态.}}
\handline{72}{792}{748}{若 $\operatorname{char}F=p>0$，则 $\F_p$ 是 $F$ 的子域. \ann{（每个域都是 $\Q$ 或 $\F_p$ 的扩张.）}}
\handline{72}{830}{835}{\textbf{Prop.} 若 $\varphi\colon E\to F$ 域同态，则 $\operatorname{char}E=\operatorname{char}F$.}
\handline{72}{864}{785}{\textbf{Pf.} \circled{1} $\operatorname{char}E=0\Rightarrow\Q\subseteq E\Rightarrow\Q\subseteq F\Rightarrow\operatorname{char}F=0$.}
\handline{72}{897}{765}{\circled{2} $\operatorname{char}E=p>0\Rightarrow\underbrace{1_E+\cdots+1_E}_{p\text{ 个}}=0$，$\therefore\varphi(1_E+\cdots+1_E)=\underbrace{1_F+\cdots+1_F}_{p\text{ 个}}=0_F$.}
\handline{72}{953}{743}{而由于 $p$ 是素的，故没有更小的 $p'$ s.t. $p'$ 个 $1_F$ 相加为 $0_F$，$\therefore\operatorname{char}F=p$.}
\end{handpage}

% ===== Line-layout page-006.tex =====
\begin{handpage}{6}
\handline{72}{99}{790}{\textbf{Ex.} 若 $\operatorname{char}F=p>0$，$\forall x,y\in F$.}
\handline{72}{127}{720}{$(x+y)^p=x^p+\binom p1x^{p-1}y+\cdots+\binom p1xy^{p-1}+y^p$.}
\handline{72}{157}{700}{$p$ 素 $\Rightarrow p\mid\binom pk$（$1\le k\le p-1$）.}
\handline{72}{185}{680}{$\therefore(x+y)^p=x^p+y^p$.}
\handline[AnnotationBlue]{72}{213}{726}{（中间项都被加了 $p$ 的倍数次，为 $0$.）}
\handline{72}{247}{700}{考虑 $\sigma\colon F\to F$}
\handline{72}{274}{600}{$x\mapsto x^p$.}
\handline{72}{304}{754}{\circled{1} $\sigma(0)=0$ \circled{2} $\sigma(1)=1$ \circled{3} $\sigma(xy)=(xy)^p=x^py^p=\sigma(x)\sigma(y)$ \circled{4} $\sigma(x+y)=(x+y)^p=x^p+y^p=\sigma(x)+\sigma(y)$.}
\handline{72}{337}{750}{$\therefore\sigma$ 是一个自同态.}
\handline{72}{369}{757}{若 $F$ 为一个有限域，则 $\sigma$ 是同态 $\Rightarrow\sigma$ 是单射 $\overset{|F|<\infty}\Longrightarrow$ 满射 $\Rightarrow\sigma$ 是自同构.}
\handline{72}{399}{730}{此时 $\sigma$ 称为 Frobenius 自同构.}
\handline[AnnotationBlue]{72}{431}{745}{（在一个特征为 $p$ 的有限域上，“求 $p$ 次幂”是一个自同构.）}
\handline{72}{464}{740}{eg. $\sigma\colon\F_4\to\F_4$.\quad $\F_4=\F_2(\alpha)$，$\operatorname{char}\F_4=\operatorname{char}\F_2=2$.}
\handline{72}{491}{260}{$x\mapsto x^2$}
\handline{301}{491}{380}{$\alpha^2=\alpha+1$.}
\handline{72}{548}{254}{$\begin{array}{|c|c|c|c|c|}\hline x&0&1&\alpha&\alpha^2\\\hline x^2&0&1&\alpha^2&\alpha\\\hline\end{array}$}
\handline{358}{526}{410}{$(\alpha^2)^2=(\alpha+1)^2=\alpha^2+1=\alpha+2=\alpha$.}
\handline{72}{626}{700}{注意 $\sigma(0)=0$，$\sigma(1)=1$.}
\handline{72}{660}{710}{$\therefore\sigma$ 的不动域为 $\F_2$.}
\handline[AnnotationBlue]{72}{699}{730}{与 Fermat 小定理的联系：}
\handline[AnnotationBlue]{72}{732}{730}{考虑 $\sigma\colon\F_p\to\F_p$}
\handline[AnnotationBlue]{72}{762}{600}{$x\mapsto x^p$}
\handline[AnnotationBlue]{72}{797}{700}{则 $\sigma$ 为恒等映射.}
\handline[AnnotationBlue]{72}{839}{700}{$\Rightarrow a^p\equiv a\pmod p$}
\handline[AnnotationBlue]{72}{874}{650}{$a^{p-1}\equiv1\pmod p,\quad(a,p)=1$.}
\end{handpage}

% ===== Line-layout page-007.tex =====
\begin{handpage}{7}
\handtitle{72}{80}{700}{\S1.4 域的扩张}
\handline{72}{109}{807}{\textbf{Def.} \circled{1} 若 $F$ 是 $E$ 的子域，则称 $E$ 是 $F$ 的扩域（扩张），记作 $E/F$.}
\handline{72}{138}{743}{\circled{2} 若 $E$ 是 $F$ 的扩域，则 \redann{\fbox{\color{black}$E$ 可视为 $F$ 上的线性空间.}}}
\handline{72}{167}{739}{记 $[E:F]=\dim_F(E)$，称为 $E/F$ 的次数. 若 $[E:F]<+\infty$ 则称 $E/F$ 为有限扩张.}
\handline{72}{197}{820}{\textbf{Def.} 设 $E/F$，$S\subseteq E$，记 $E$ 中包含 $S$ 和 $F$ 最小的子域为 $F(S)$.}
\handline{72}{225}{735}{$F(S)$ 称为在 $F$ 上由 $S$ 生成的 $E$ 中的子域.}
\handline{72}{254}{590}{若 $S$ 为有限集，且 $E=F(S)$，则称 $E/F$ 是有限生成的.}
\handline{72}{286}{590}{若 $E$ 的任一有限子集 $S$，都有 $F(S)\ne E$，则称 $E$ 是无限生成的.}
\handline{72}{320}{818}{\textbf{Ex.} \circled{1} $\Q(\sqrt2)$，$\Q(\sqrt2,\sqrt3)$ 都是有限生成扩张. $S_1=\{\sqrt2\}$，$S_2=\{\sqrt2,\sqrt3\}$.}
\handline{72}{350}{739}{$[\Q(\sqrt2):\Q]=2$，$[\Q(\sqrt2,\sqrt3):\Q]=4$. $\therefore$ 它们也是有限扩张.}
\handline{72}{378}{765}{\circled{2} $\R(x)$：$\R$ 上有理函数域，由未知元 $x$，在 $\R$ 上有限生成.}
\handline{72}{407}{737}{但 $[\R(x):\R]=+\infty$，故不是有限生成扩张.}
\handline[AnnotationBlue]{72}{440}{760}{\circled{3} \textbf{Thm.} $\Q$ 与 $\F_p$ 均无真子域，即 $\Q$ 与 $\F_p$ 是素域.}
\handline{72}{469}{750}{$1^\circ$：若 $E\subseteq\Q\Rightarrow0,1\in E$.}
\handline{72}{500}{710}{加法封闭 $\Rightarrow\N\subseteq E$.}
\handline{72}{531}{715}{加法逆存在 $\Rightarrow\Z\subseteq E$.}
\handline{72}{562}{721}{乘法封闭，逆存在 $\Rightarrow\Q\subseteq E$. $\therefore E=\Q$.}
\handline{72}{601}{768}{$2^\circ$：记 $E$ 为 $\F_p$ 子域，则 $E,\F_p$ 均为有限集，$\therefore[\F_p:E]<+\infty$.}
\handline{72}{636}{765}{取 $\F_p$ 的 $[\F_p:E]$ 个基，\ann{（$\F_p$ 视为 $E$-线性空间）}每个基前系数有 $|E|$ 种选择.}
\handline{72}{670}{769}{$\therefore|\F_p|=|E|^{[\F_p:E]}$. 而 $|\F_p|=p$ 素 $\Rightarrow|E|=p$，$[\F_p:E]=1$. $\therefore E=\F_p$.}
\handline{72}{701}{750}{$\therefore\Q$ 与 $\F_p$ 是素域.}
\handline[AnnotationBlue]{72}{733}{775}{下面给出两个不是有限生成的例子.}
\handline{72}{766}{775}{\circled{4} $E=\Q\left(2^{\frac1{2^k}}\mid k=1,2,3,\ldots\right)$.}
\handline{72}{802}{784}{$E_0=\Q$，$E_1=\Q(\sqrt2)$，$E_2=\Q(\sqrt2,\sqrt[4]2)=E_1(\sqrt[4]2)$，$E_3=\Q(\sqrt2,\sqrt[4]2,\sqrt[8]2)=E_2(\sqrt[8]2)$.}
\handline{72}{838}{781}{$E_n=E_{n-1}\left(2^{\frac1{2^n}}\right)$，$\Q=E_0\subset E_1\subset E_2\cdots\subset E_n\subset\cdots\subset E$，$E=\bigcup_{k=0}^{\infty}E_k$.}
\handline{72}{870}{770}{$1^\circ$ 证明 $E$ 是一个域.}
\handline{72}{911}{785}{$\forall a,b\in E$，$k,l\in\N$ s.t. $a\in E_k$，$b\in E_l$. $\Rightarrow a,b\in E_{\max(k,l)}$.}
\handline{72}{943}{785}{$E_{\max(k,l)}$ 是域，$\therefore a,b$ 可作运算，$\therefore E$ 是一个域.}
\end{handpage}

% ===== Line-layout page-008.tex =====
\begin{handpage}{8}
\handline{72}{110}{400}{$2^\circ$ 证明 $E/\Q$ 不是有限生成扩张.}
\handline{72}{142}{370}{设 $S$ 为 $E$ 的任一有限子集.}
\handline{72}{168}{270}{$S=\{a_1,a_2,\ldots,a_n\}$.}
\handline{72}{196}{500}{$\forall a_i\in S\quad1\le i\le n\quad\exists k_i\in\N\ \text{s.t.}\ a_i\in E_{k_i}$.}
\handline{72}{225}{380}{令 $N=\max_{1\le i\le n}\{k_i\}$ 则 $S\subseteq E_N$.}
\handline{72}{280}{610}{$\Q\cup S\subseteq E_N\quad\therefore\Q(S)\subseteq E_N\varsubsetneqq E\quad\therefore\Q(S)\ne E$.}
\handline{72}{308}{340}{$\therefore E/\Q$ 不是有限生成扩张.}
\handline{72}{339}{460}{\circled{5}\quad$E=\Q(\sqrt2,\sqrt3,\sqrt5,\ldots,\sqrt p,\ldots)$，$p$ 素.}
\handline[AnnotationBlue]{72}{368}{620}{但是要证 $\sqrt2,\sqrt3,\sqrt5,\ldots,\sqrt p,\ldots$ 在 $\Q$ 上线性无关.}
\handline[AnnotationBlue]{72}{395}{350}{（将 $\R$ 视为 $\Q$ 上线性空间）.}
\handline{72}{432}{610}{Thm.\quad 有限扩张一定是有限生成扩张，但反之未必.}
\handline{72}{463}{610}{Pf.\quad 设 $E/F$.\quad$[E:F]=n\quad\therefore E=\Span_F(e_1,e_2,\ldots,e_n)$.}
\handline{72}{492}{490}{$\therefore E=F(e_1,e_2,\ldots,e_n)$ 为有限生成扩张.}
\handline[AnnotationBlue]{72}{523}{430}{有限生成扩张不一定是有限扩张.}
\handline[AnnotationBlue]{72}{552}{830}{反例：\circled{1} $\Q[x]$.$x$ 是超元.$[\Q[x]:\Q]=+\infty$.但 $\Q[x]/\Q$ 是有限扩张.}
\handline[AnnotationBlue]{72}{584}{490}{由于 $\pi$ 是超越数.$\forall p\in\Q[x]$.\quad$p(\pi)\ne0$.}
\handline[AnnotationBlue]{72}{615}{320}{$\therefore\Q(\pi)\cong\Q[x]$.}
\handline[AnnotationBlue]{72}{647}{640}{甚至 $\Q(\pi)$ 不仅是有限生成，还是单扩张，但不是有限扩张.}
\handline{72}{684}{230}{Def.\quad 设 $E/F$.}
\handline{72}{716}{815}{\circled{1} 若 $u\in E$.\ $\exists f\ne0\in F[x]\ \text{s.t.}\ f(u)=0$.则称 $u$ 在 $F$ 上代数（$u$ 是 $F$ 上代数元）.}
\handline{72}{748}{620}{\circled{2} 若 $\forall u\in E$.$u$ 在 $F$ 上代数，则称 $E$ 是 $F$ 的代数扩张.}
\handline{72}{783}{820}{\circled{3} 若 $\exists u\in E$，s.t. $u$ 不是任何 $F[x]$ 中多项式的根，则称 $E$ 是 $F$ 的超越扩张.}
\handline[AnnotationBlue]{72}{824}{570}{Ex.\quad$\Q(\sqrt2)/\Q$，$\Q(\sqrt3)/\Q$ 是代数扩张.}
\handline[AnnotationBlue]{72}{858}{480}{$\Q(\pi)/\Q$，$\Q[x]/\Q$ 是超越扩张.}
\end{handpage}

% ===== Line-layout page-009.tex =====
\begin{handpage}{9}
\handline[NoteRed]{72}{96}{660}{下面研究有限扩张、有限生成扩张、代数扩张之间的关系.}
\handline{72}{130}{700}{Thm.\quad 设 $E/F$.$\alpha,\beta$ 是 $F$ 上代数元.则 $\alpha+\beta,\alpha\beta$ 也是 $F$ 上代数元.}
\handline{72}{158}{340}{Pf.（高代证法，利用结式）}
\handline{72}{187}{620}{$\alpha,\beta$ 为 $F$ 上代数元，则 $\exists f,g\in F[x]\ \text{s.t.}\ f(\alpha)=g(\beta)=0$.}
\handline{72}{215}{420}{$h(y):=R_x(f(x),g(y-x))\in F[y]$.}
\handline{72}{247}{420}{$h(\alpha+\beta)=R_x(f(x),g(\alpha+\beta-x))$.}
\handline{72}{276}{640}{而 $f(x)$ 与 $g(\alpha+\beta-x)$ 有相同根 $\alpha$，$\therefore R_x(f(x),g(\alpha+\beta-x))=0$.}
\handline{72}{304}{470}{$\therefore h(\alpha+\beta)=0\quad\therefore\alpha+\beta$ 为 $F$ 上代数元.}
\handline{72}{337}{610}{$k(y):=R_x\bigl(f(x),x^dg(\frac yx)\bigr)\in F[y]$.\quad$d=\deg(g(x))$.}
\handline{72}{371}{430}{$k(\alpha\beta)=R_x\bigl(f(x),x^dg(\frac{\alpha\beta}x)\bigr)$.}
\handline{72}{400}{650}{而 $f(x)$ 与 $x^dg(\frac{\alpha\beta}x)$ 有相同根 $\alpha$，$\therefore R_x\bigl(f(x),x^dg(\frac{\alpha\beta}x)\bigr)=0$.}
\handline{72}{431}{460}{$\therefore k(\alpha\beta)=0\quad\therefore\alpha\beta$ 为 $F$ 上代数元.}
\handline[AnnotationBlue]{72}{463}{280}{证法二：对称多项式.}
\handline[AnnotationBlue]{72}{494}{260}{$f(\alpha)=0$.\quad$g(\beta)=0$.}
\handline[AnnotationBlue]{72}{527}{390}{$f(x)=(x-\alpha_1)(x-\alpha_2)\cdots(x-\alpha_n)$.}
\handline[AnnotationBlue]{72}{558}{390}{$g(x)=(x-\beta_1)(x-\beta_2)\cdots(x-\beta_m)$.}
\handline[AnnotationBlue]{72}{589}{800}{$x_{ij}=\alpha_i+\beta_j\quad h(x)=\prod_{i,j}(x-x_{ij})$.证明 $h(x)$ 系数在 $\alpha,\beta$ 的任意置换下不变.}
\handline{72}{621}{640}{\ann{证法三：抽代证法.}\redann{（提前用一下域扩张的维数公式）}}
\handline{72}{660}{856}{\ann{Lemma.\quad$\alpha\in E$.$\alpha$ 是 $F$ 上代数元 $\Longleftrightarrow F(\alpha)/F$ 是有限扩张，$[F(\alpha):F]<\infty$.}\redann{（trivial）}}
\handline[AnnotationBlue]{72}{695}{650}{考虑 $F(\alpha,\beta)=F(\alpha)(\beta)$.则 $\beta\in F(\alpha,\beta)$ 是 $F(\alpha)$ 上代数元.}
\handline[AnnotationBlue]{72}{726}{380}{$\therefore[F(\alpha,\beta):F(\alpha)]<\infty$.}
\handline[AnnotationBlue]{72}{761}{730}{$\therefore$ 由维数公式 $[F(\alpha,\beta):F]=[F(\alpha,\beta):F(\alpha)][F(\alpha):F]<\infty$.}
\handline{72}{798}{812}{\ann{$\therefore F(\alpha,\beta)/F$ 是有限扩张 $\Rightarrow F(\alpha,\beta)/F$ 是代数扩张}\redann{（即下一个定理）}.}
\handline[AnnotationBlue]{72}{832}{590}{$\therefore\alpha+\beta,\alpha\beta\in F(\alpha,\beta)$ 是 $F$ 上代数元.}
\end{handpage}

% ===== Line-layout page-010.tex =====
\begin{handpage}{10}
\handline{72}{94}{500}{Thm.\quad 有限扩张一定是代数扩张，反之未必.}
\handline{72}{127}{330}{Pf.\quad 设 $E/F$，$[E:F]=n$.}
\handline{72}{157}{630}{$\forall u\in E$.\quad$1,u,u^2,\ldots,u^n$ 视为 $F$-线性空间 $(E)$ 上 $n+1$ 个向量.}
\handline{72}{185}{360}{而 $\dim_F(E)=n\quad\therefore$ 必线性相关.}
\handline{72}{218}{620}{$\exists b_0,b_1,\ldots,b_n\in F$ 不全为 $0$.\ s.t. $b_0+b_1u+\cdots+b_nu^n=0$.}
\handline{72}{248}{490}{$\therefore u$ 是 $F$ 的代数元 $\therefore E/F$ 是代数扩张.}
\handline[AnnotationBlue]{72}{276}{330}{代数扩张不一定是有限扩张.}
\handline[AnnotationBlue]{72}{307}{400}{反例：$E=\Q(2^{\frac1{2^k}}\mid k=1,2,3,\ldots)$.}
\handline[AnnotationBlue]{72}{339}{740}{$\forall\alpha\in E$.\ $\exists N\in\N$.\ s.t. $\alpha\in E_N$.\quad$E_N/\Q$ 是有限扩张 $\Rightarrow$ 是代数扩张.}
\handline[AnnotationBlue]{72}{370}{500}{$\therefore\alpha$ 是 $\Q$ 上代数元 $\therefore E/\Q$ 是代数扩张.}
\handline[AnnotationBlue]{72}{400}{300}{但 $E/\Q$ 不为有限扩张.}
\handline[NoteRed]{72}{436}{810}{（若 $E/\Q$ 为有限扩张，则 $E/\Q$ 为有限生成扩张，但之前分析出 $E/\Q$ 不能有限生成）.}
% 按用户要求删除关系图；保留原图区域 x=197,y=473,w=535,h=152，供手绘.
\handline[NoteRed]{72}{676}{450}{反例：\circled{1}：$E=\Q(2^{\frac1{2^k}}\mid k=1,2,3,\ldots)$.}
\handline[NoteRed]{72}{713}{270}{\circled{2}：$\Q(\pi)$，$\Q[x]$.}
\handline[NoteRed]{72}{749}{430}{\circled{3}：$\Q(\pi)$，$\Q[x]$\quad$\pi$ 是超越数.}
\handline[NoteRed]{72}{787}{390}{\circled{4}：$E=\Q(2^{\frac1{2^k}}\mid k=1,2,3,\ldots)$.}
\handline[NoteRed]{72}{861}{730}{下面我们来看 \circled{1} \circled{2} 这两个箭头，在什么条件下可以成立.}
\end{handpage}

% ===== Line-layout page-011.tex =====
\begin{handpage}{11}
\handline{72}{80}{450}{Thm.\quad 有限生成的代数扩张是有限扩张.}
\handline{72}{110}{710}{加强形式：对于 $E/F$，T.F.A.E.\ \ann{（the following are equivalent）}.}
\handline{72}{139}{220}{\circled{1} $E/F$ 是有限扩张.}
\handline{72}{165}{440}{\circled{2} $E=F(u_1,u_2,\ldots,u_n)$，$u_i$ 为 $F$ 上代数元.}
\handline{72}{194}{240}{先做一些准备工作.}
\handline{72}{222}{710}{Def.\quad 设 $E/F$.若有 $E$ 的子域 $K$.有 $F\subseteq K$，则称 $K$ 为 $E/F$ 的一个中间域.}
\handline{72}{251}{470}{Ex.\quad$\Q(\sqrt2,\sqrt3)/\Q$ 的中间域为 $\Q(\sqrt2)$，$\Q(\sqrt3)$.}
\handdiagram{748}{257}{180}{56}{
\node(q) at (0,0) {$\Q$};\node(a) at (-1.5,1) {$\Q(\sqrt2)$};\node(b) at (1.5,1) {$\Q(\sqrt3)$};\node(c) at (0,2) {$\Q(\sqrt2,\sqrt3)$};
\draw[->,>=Stealth](q)--(a);\draw[->,>=Stealth](q)--(b);\draw[->,>=Stealth](a)--(c);\draw[->,>=Stealth](b)--(c);
}
\handline{72}{277}{300}{$\Q(\sqrt2),\Q(\sqrt3)$ 之间没有同态.}
\handline{72}{337}{300}{Lemma. 1.（维数公式）.}
\handline{72}{366}{810}{设 $E/F$ 是有限扩张.$K$ 为 $E/F$ 中间域，则 $E/K,K/F$ 也是有限扩张，且 $[E:F]=[E:K][K:F]$.}
\handline{72}{397}{780}{Pf.\quad 首先 $E$ 是 $F$-线性空间.且为有限维.$E/K,K/F$ 是子空间，故也是有限维.}
\handline{72}{426}{350}{$\therefore E/K,K/F$ 是有限扩张.}
\handline{72}{457}{285}{设 $[E:K]=m$，$[K:F]=n$.}
\handline{72}{491}{370}{$u_1,u_2,\ldots,u_n\in K$ 是 $K/F$ 的一组基.}
\handline{72}{522}{400}{$v_1,v_2,\ldots,v_m\in E$ 是 $E/K$ 的一组基.}
\handline{72}{558}{490}{$\forall\beta\in E$.\ $\exists\alpha_1,\alpha_2,\ldots,\alpha_m\in K$.\ s.t. $\beta=\sum_{i=1}^m\alpha_iv_i$.}
\handline{72}{589}{630}{$\forall\alpha_i\in K\ 1\le i\le m$，$\exists a_{i1},a_{i2},\ldots,a_{in}\in F\ \text{s.t.}\ \alpha_i=\sum_{j=1}^na_{ij}u_j$.}
\handline{72}{625}{610}{$\Rightarrow\beta=\sum_{i=1}^m(\sum_{j=1}^na_{ij}u_j)v_i=\sum_{i=1}^m\sum_{j=1}^na_{ij}u_jv_i\quad(u_jv_i\in E)$.}
\handline{72}{660}{610}{$\Rightarrow E\subseteq\Span_F(u_jv_i\mid\substack{1\le i\le m\\1\le j\le n})$.\quad$\therefore\dim_F(E)\le mn$.}
\handline{72}{691}{340}{下证 $u_jv_i$ $F$-线性无关.}
\handline{72}{726}{365}{设 $\sum C_{ij}u_jv_i=0$\quad$C_{ij}\in F$.}
\handline{72}{762}{570}{$\therefore\sum_{i=1}^m(\sum_{j=1}^nC_{ij}u_j)v_i=0$\quad 而 $v_i$ 为 $E/K$ 一组基.}
\handline{72}{822}{600}{$\therefore\sum_{j=1}^nC_{ij}u_j=0\quad(\forall1\le i\le m)$.\quad 而 $u_j$ 为 $K/F$ 一组基.}
\handline{72}{882}{640}{$\therefore C_{ij}=0\quad\forall1\le i\le m,1\le j\le n\quad\therefore u_jv_i$ $F$-线性无关.}
\handline{72}{930}{710}{$\therefore mn\le\dim_F(E)\quad\therefore\dim_F(E)=mn\quad\text{i.e. }[E:F]=[E:K][K:F]$.}
\end{handpage}

% ===== Line-layout page-012.tex =====
\begin{handpage}{12}
\handline{72}{87}{450}{Cor.\quad 若 $[E:F]$ 为素数，则 $E/F$ 无非平凡中间域.}
\handline{72}{115}{390}{Lemma. 2.\quad 单代数扩张是有限扩张.}
\handline{72}{142}{420}{Pf.\quad 设 $E=F(u)$.$u$ 是 $F$ 上代数元.}
\handline{72}{169}{500}{设 $u$ 在 $F$ 上的极小多项式为 $f$.则 $f$ 在 $F$ 上不可约.}
\handline{72}{198}{705}{设 $\deg f=n$.\ $\Rightarrow E=F(u)=\Span_F(1,u,\ldots,u^{n-1})\Rightarrow\dim_F(E)=n<\infty$.}
\handline{72}{223}{240}{i.e. $[E:F]=n$.}
\handline[AnnotationBlue]{72}{250}{430}{（极小多项式的次数即为 $E/F$ 的维数）.}
\handline[AnnotationBlue]{72}{279}{510}{Remark：\circled{1} 极小多项式：$1^\circ$ 唯一，$2^\circ$ 不可约.}
\handline[AnnotationBlue]{72}{307}{723}{\circled{2} $1,u,\ldots,u^{n-1}$ 线性无关，否则 $\exists g$ s.t. $g(u)=0$ 且 $0<\deg(g)\le n-1<n$.}
\handline[AnnotationBlue]{72}{335}{270}{与 $f$ 的极小性矛盾.}
\handline[AnnotationBlue]{72}{365}{627}{$E$ 中每个元素可以由 $1,u,\ldots,u^{n-1}$ 线性表出.回忆 $\Q(\sqrt2+\sqrt3)$.}
\handline{72}{395}{250}{回到定理的证明：}
\handline{72}{426}{700}{Pf.\quad\circled{1}$\Rightarrow$\circled{2}.设 $[E:F]=n$.取一组基 $u_1,u_2,\ldots,u_n$ 则 $E=F(u_1,\ldots,u_n)$.}
\handline{72}{457}{635}{因为有限扩张 $\Rightarrow$ 代数扩张.$\therefore u_i$ 为 $F$ 上代数元\quad$1\le i\le n$.}
\handline{72}{489}{570}{\circled{2}$\Rightarrow$\circled{1}\quad 设 $E=F(u_1,\ldots,u_n)$，$u_i$ 为 $F$ 上代数元\ $1\le i\le n$.}
\handline{72}{518}{370}{考虑 $F=F_0\subset F_1\subset F_2\cdots F_n=E$.}
\handline{72}{548}{350}{$F_k=F(u_1,u_2,\ldots,u_k)\quad1\le k\le n$.}
\handline{72}{579}{490}{则 $F_k/F_{k-1}$ 是单代数扩张，故为有限扩张.}
\handline{72}{610}{685}{由维数公式 $\Rightarrow[E:F]=\prod_{i=0}^{n-1}\dim_{F_i}(E_i)<\infty\quad\therefore E/F$ 是有限扩张.}
\handline{72}{641}{856}{Thm.（代数扩张的传递性）.若 $F\subseteq K\subseteq E$，$E/K,K/F$ 是代数扩张，则 $E/F$ 是代数扩张.}
\handline{72}{674}{475}{Pf.\quad$\forall\alpha\in E$.\ $\exists f\ne0\in K[x]\ \text{s.t.}\ f(\alpha)=0$.}
\handline{72}{704}{610}{设 $f(x)=x^n+a_1x^{n-1}+a_2x^{n-2}+\cdots+a_{n-1}x+a_n$.\ $a_i\in K$.\ $1\le i\le n$.}
\handline{72}{736}{525}{由 $K/F$ 是代数扩张.$\therefore a_i$ 是 $F$ 上代数元 $1\le i\le n$.}
\handline{72}{767}{625}{设 $K'=F(a_1,a_2,\ldots,a_n)$ 为有限生成的代数扩张 $\Rightarrow$ 有限扩张.}
\handline{72}{804}{774}{而 $f\in K'[x]\quad\therefore\alpha$ 在 $K'$ 上也代数.$\therefore K'(\alpha)/K'$ 是单代数扩张 $\Rightarrow$ 有限扩张.}
\handline{72}{838}{856}{维数公式 $\Rightarrow K'(\alpha)/F$ 也是有限扩张.而 $F(\alpha)$ 是 $K'(\alpha)/F$ 中间域 $\therefore F(\alpha)/F$ 也是有限扩张.}
\handline{72}{874}{789}{从而 $F(\alpha)/F$ 是代数扩张.从而 $\alpha$ 是 $F$ 上代数元.$\therefore E/F$ 是代数扩张.}
\handline[NoteRed]{72}{904}{505}{Remark：这个定理是前面一个定理的充分应用.}
\handline[NoteRed]{72}{934}{795}{且证 $\alpha$ 是 $F$ 代数元，只要证 $F(\alpha)/F$ 是代数扩张，将 $E/F$ 简化为每个 $F(\alpha)/F$ 代数.}
\end{handpage}

% ===== Line-layout page-013.tex =====
\begin{handpage}{13}
\handtitle{72}{113}{230}{§1.5 代数闭包.}
\handline{72}{141}{460}{Def.\quad 设 $E/F$，$K=\{\alpha\in E\mid\alpha$ 在 $F$ 上代数$\}$.}
\handline{72}{164}{615}{则 $K$ 为 $E/F$ 的中间域，称为 $F$ 在 $E$ 中的（相对）代数闭包.}
\handline[AnnotationBlue]{72}{196}{300}{“闭包”的“闭”体现在哪？}
\handline[NoteRed]{72}{223}{783}{若 $K'$ 是 $E/K$ 的中间域，且 $K'/K$ 代数，由于 $K/F$ 代数 $\Longrightarrow K'/F$ 代数.}
\handline[NoteRed]{72}{250}{785}{再结合代数闭包定义.$\forall\alpha\in K'$，$\alpha$ 在 $F$ 上代数.$\therefore\alpha\in K\quad\therefore K'\subset K\quad\therefore K'=K$.}
\handline{72}{281}{560}{Def.\quad(1) 若 $K$ 没有真代数扩张，则称 $K$ 为代数闭域；}
\handline{72}{307}{744}{(2) 若 $K/F$ 是代数扩张且 $K$ 是代数闭域，则称 $K$ 是 $F$ 的一个绝对代数闭包.}
\handline{72}{337}{726}{Ex.\quad\circled{1} $E_1=\R$，$F_1=\Q$.令 $K_1=\{$实代数数$\}$，则 $K_1$ 为 $F_1$ 在 $E_1$ 中的代数闭包.}
\handline{72}{365}{673}{\circled{2} $E_2=\C$，$F_2=\Q$.令 $K_2=\{$复代数数$\}$，则 $K_2$ 为 $F_2$ 在 $E_2$ 中的代数闭包.}
\handline{72}{395}{414}{\circled{3} $\C$ 是一个代数闭域.\ann{（代数基本定理）}}
\handline{72}{425}{654}{Thm.\quad$\Q$ 在 $\C$ 中的相对代数闭包，就是 $\Q$ 的一个绝对代数闭包.}
\handline[AnnotationBlue]{72}{454}{370}{下面直接证一个更一般的形式：}
\handline{72}{484}{766}{Thm.\quad$E$ 是代数闭域.$F$ 是 $E$ 的子域.$K$ 是 $F$ 在 $E$ 中的相对代数闭包，}
\handline{72}{513}{350}{则 $K$ 是 $F$ 的绝对代数闭包.}
\handline{72}{543}{615}{Pf.\quad 显然 $K/F$ 是代数扩张，下证 $K$ 是代数闭域.}
\handline{72}{574}{627}{若 $\exists K'$.\ s.t. $K'/K$ 为代数扩张.$\forall\alpha\in K'$，$\alpha$ 在 $K$ 上代数.}
\handline{72}{605}{647}{设 $\alpha$ 在 $K[x]$ 上的极小多项式为 $f(x)$.由 $K\subseteq E\Rightarrow f\in E[x]$.}
\handline{72}{636}{631}{由 $E$ 是代数闭域 $\Rightarrow f$ 在 $E$ 上可彻底分解为一次因式乘积，}
\handline{72}{671}{527}{$f(x)=(x-\alpha_1)(x-\alpha_2)\cdots(x-\alpha_n)$.\quad$\alpha_i\in E\quad1\le i\le n$.}
\handline{72}{704}{675}{$\therefore f(\alpha_i)=0$.$\alpha_i$ 在 $K$ 上代数.$\Rightarrow K(\alpha_i)/K$ 为代数扩张.}
\handline{72}{738}{575}{由代数扩张传递性 $\Rightarrow K(\alpha_i)/F$ 是代数扩张.}
\handline{72}{771}{385}{$\therefore\alpha_i$ 在 $F$ 上代数，$\therefore\alpha_i\in K$.}
\handline{72}{807}{630}{某个 $\alpha_i$ 即为 $\alpha$，$\therefore\alpha\in K\quad\therefore K'\subset K\quad\therefore K'=K$.}
\handline{72}{842}{350}{$\therefore K$ 没有真代数扩张.}
\handline{72}{875}{569}{$\therefore K$ 为代数闭域.$K$ 是 $F$ 的绝对代数闭包.}
\handline{72}{910}{856}{Thm.\quad 对任一域 $F$，存在 $F$ 的绝对代数闭包，并且在同构意义下唯一.\ann{（Artin）}}
\handline[AnnotationBlue]{72}{938}{255}{见姚慕生附录.}
\end{handpage}

% ===== Line-layout page-014.tex =====
\begin{handpage}{14}
\handtitle{72}{126}{250}{§1.6 Galois 群}
\handline{72}{156}{803}{Def.\quad$E$ 是一个域，$\operatorname{Aut}(E):=\{\sigma\in S(E)\mid\begin{aligned}\sigma(0)&=0,\ \sigma(1)=1,\ \sigma(\alpha+\beta)=\sigma(\alpha)+\sigma(\beta)\\\sigma(\alpha\beta)&=\sigma(\alpha)\sigma(\beta)\end{aligned}\}$.}
\handline{72}{199}{610}{其中 $S(E)$ 为 $E$ 上对称群\quad$S(E)=\{\sigma\colon E\to E\mid\sigma$ 是双射$\}$.}
\handline{72}{231}{420}{$\operatorname{Aut}(E)$ 称为域 $E$ 上的自同构群.}
\handline{72}{260}{662}{$\operatorname{Gal}(E/F):=\{\sigma\in\operatorname{Aut}(E)\mid\sigma|_F=\id_F\}$ 称为 $E/F$ 的 Galois 群.}
\handline[NoteRed]{72}{289}{476}{Remark：$\sigma|_F=\id_F$ 与 $\sigma(F)=F$ 的区别.}
\handline[NoteRed]{72}{321}{741}{$(\Longleftrightarrow)\ \forall x\in F,\ \sigma(x)=x.\qquad(\Longleftrightarrow)\ \forall x\in F,\ \sigma(x)\in F\ \wedge\ \forall y\in F,\exists x\in F\ \text{s.t.}\ \sigma(x)=y$.}
\handline{72}{348}{773}{Thm.\quad 若 $E/F$ 是有限扩张，则 $\operatorname{Gal}(E/F)$ 是有限群.\ann{（域 $\longrightarrow$ 群）}}
\handline{72}{379}{380}{Pf.\quad 设 $E=\Span_F(u_1,u_2,\ldots,u_n)$.}
\handline{72}{408}{590}{$u_i$ 是 $F$ 上代数元（$1\le i\le n$）（有限扩张一定是代数扩张）.}
\handline{72}{437}{570}{对于 $u_i$，$\exists$ 极小多项式 $f_i(x)\in F[x]$.\ s.t. $f_i(u_i)=0$.}
\handline{72}{468}{774}{设 $\sigma\in\operatorname{Aut}(E)$.且 $\sigma|_F=\id_F$.设 $f_i(x)=a_nx^n+a_{n-1}x^{n-1}+\cdots+a_1x+a_0\quad a_j\in F$.}
\handline{72}{500}{415}{由 $\sigma|_F=\id_F\quad\therefore\sigma(a_j)=a_j\quad(1\le j\le n)$.}
\handline{72}{530}{727}{$\therefore\sigma(f_i(x))=a_n\sigma^n(x)+a_{n-1}\sigma^{n-1}(x)+\cdots+a_1\sigma(x)+a_0=f_i(\sigma(x))$.}
\handline{72}{562}{340}{令 $X_i=\{\alpha\in E\mid f_i(\alpha)=0\}$.}
\handline{72}{597}{757}{对 $\alpha\in X_i$\quad$f_i(\sigma(\alpha))=\sigma(f_i(\alpha))=\sigma(0)=0\Rightarrow\sigma(\alpha)\in X_i$.}
\handline{72}{629}{355}{于是 $\sigma|_{X_i}$ 是 $X_i$ 的一个置换.}
\handline{72}{660}{450}{令 $X=\bigcup_{i=1}^nX_i$，则 $\sigma|_X$ 是 $X$ 的一个置换.}
\handline{72}{692}{814}{而 $E=\Span_F(u_1,\ldots,u_n)$.\ $\forall\alpha\in E$.$\alpha$ 可唯一表示为 $u_1,\ldots,u_n$ 的 $F$-线性组合.}
\handline{72}{726}{808}{$\therefore$ 一旦确定 $\sigma(u_1),\ldots,\sigma(u_n)$ 的值.$\forall\alpha\in E$，$\sigma(\alpha)$ 的值唯一确定.}
\handline{72}{759}{491}{而当 $\sigma|_X$ 确定，$\sigma(u_i)$ 也随之确定.}
\handline{72}{793}{704}{而 $X$ 是有限集\quad$\therefore\sigma|_X$ 在 $X$ 上的置换方式为有限个.}
\handline{72}{827}{365}{$\therefore\operatorname{Gal}(E/F)$ 为有限群.}
\handline[NoteRed]{72}{865}{479}{关注 $\sigma$ 对极小多项式的根的置换.}
\end{handpage}

% ===== Line-layout page-015.tex =====
\begin{handpage}{15}
\handline{72}{91}{500}{\textbf{Def.} 设 $E$ 是一个域，$G\le\operatorname{Aut}(E)$.}
\handline{72}{117}{405}{定义 $\operatorname{Inv}(G)=\{\alpha\in E\mid\forall\sigma\in G,\ \sigma(\alpha)=\alpha\}$}
\handline{72}{148}{590}{$\operatorname{Inv}(G)$ 是 $E$ 的子域，称为 $E$ 的不动域.}
\handline{72}{176}{742}{\textbf{Thm.} \ann{(Artin's Lemma)（某种程度上可视为上一个定理的逆）（群$\to$域）}}
\handline{72}{203}{798}{$E$：域，$G\le\operatorname{Aut}(E)$ 且 $|G|<\infty$. 记 $F=\operatorname{Inv}(G)$，则 $[E:F]\le|G|$，$E/F$ 为有限扩张.}
\handline[NoteRed]{72}{230}{383}{这里只证“$\le$”，事实上可改进为“$=$”.}
\handline{72}{262}{560}{\textbf{Pf.} 设 $G=\{\eta_1,\eta_2,\ldots,\eta_n\}$，$|G|=n$.}
\handline{72}{289}{599}{下证当 $m>n$ 时，任意 $m$ 个 $E$ 中的数 $F$-线性相关.}
\handline{72}{320}{533}{$\forall u_1,\ldots,u_m\in E$，不妨设全不为 $0$（否则已经线性相关）.}
\handline{72}{350}{465}{则 $\eta_j(u_i)$，$1\le j\le n$，$1\le i\le m$ 是 $mn$ 个定义好的数.}
\handline{72}{381}{415}{构造齐次线性方程组.\quad $\eta_j(u_i)\in E$.}
\handline{72}{439}{562}{$\begin{cases}
\eta_1(u_1)x_1+\eta_1(u_2)x_2+\cdots+\eta_1(u_m)x_m=0,\\
\eta_2(u_1)x_1+\eta_2(u_2)x_2+\cdots+\eta_2(u_m)x_m=0,\\
\cdots\\
\eta_n(u_1)x_1+\eta_n(u_2)x_2+\cdots+\eta_n(u_m)x_m=0.
\end{cases}$}
\handline{72}{500}{704}{共 $n$ 个方程，$m$ 个未知数，$m>n$，故必有非零解 $(x_1,\ldots,x_m)\in E^m$.}
\handline{72}{529}{542}{设 $\bar X=(x_1,\ldots,x_m)$ 是所有非零解中含 $0$ 最多的.}
\handline{72}{561}{445}{不妨设 $x_1\ne0$，则由齐次性不妨设 $x_1=1$.}
\handline{72}{590}{460}{下证 $(1,x_2,\ldots,x_m)$ 中 $x_2,\ldots,x_m\in F$.}
\handline{72}{621}{785}{若 $\exists x_j$，$2\le j\le m$，s.t. $x_j\notin F$，不妨设 $x_2\notin F$，所以 $\exists\eta\in G$ s.t. $\eta(x_2)\ne x_2$.}
\handline{72}{649}{800}{显然 $\eta(\bar X)=(\eta(1),\eta(x_2),\ldots,\eta(x_m))=(1,\eta(x_2),\ldots,\eta(x_m))$ 仍为一组非零解.}
\handline{72}{679}{400}{我们有以下事实：}
\handline{72}{711}{810}{由于 $\eta\in G\le\operatorname{Aut}(E)$，$\eta$ 为双射，因而 $\eta(\bar X)$ 与 $\bar X$ 中零的个数相同，且位置对应.}
\handline{72}{741}{814}{考虑 $\eta(\bar X)-\bar X$，则也为方程组的解，且第 $1$ 分量为 $1-1=0$，第 $2$ 分量 $\eta(x_2)-x_2\ne0$.}
\handline{72}{780}{570}{其余分量：若在 $\bar X$ 中为 $0$，则在 $\eta(\bar X)-\bar X$ 中也为 $0$.}
\handline{72}{812}{480}{若在 $\bar X$ 中不为 $0$，则在 $\eta(\bar X)-\bar X$ 中可能变为 $0$.}
\handline{72}{845}{735}{从而 $\eta(\bar X)-\bar X$ 中零的个数比 $\bar X$ 至少多 $1$，与 $\bar X$ 取法矛盾.}
\handline{72}{879}{817}{$\Rightarrow(1,x_2,\ldots,x_m)\in F^m$；$0=\eta_1(u_1)x_1+\cdots+\eta_1(u_m)x_m=\eta_1(x_1u_1+\cdots+x_mu_m)=0$.}
\handline{72}{909}{794}{由 $\eta_1$ 为双射 $\Rightarrow u_1+x_2u_2+\cdots+x_mu_m=0$. 所以 $u_1,\ldots,u_m$ $F$-线性相关.}
\handline{72}{943}{545}{所以 $[E:F]=\dim_F(E)\le n=|G|$.}
\end{handpage}

% ===== Line-layout page-016.tex =====
\begin{handpage}{16}
\handline{72}{83}{510}{有了上面两个定理，我们可以问两个问题.}
\handline{72}{114}{481}{由第一个定理知，有限扩张的 Galois 群是有限群，}
\handline{72}{140}{839}{Artin 引理则说明一个域的自同构群的有限子群的不动域到原来域的扩张是有限扩张.}
\handline{72}{170}{440}{于是可以问两个“拉回”问题.}
\handline{72}{197}{420}{Q1：域 $\to$ 群 $\to$ 域.}
\handline{72}{222}{710}{$[E:F]<\infty\Rightarrow\operatorname{Gal}(E/F)$ 有限 $\le\operatorname{Aut}(E)\Rightarrow F'=\operatorname{Inv}(\operatorname{Gal}(E/F))$.}
\handline{72}{251}{747}{$[E:F']<\infty$. 问 $F$ 与 $F'$ 是什么关系？$F\subseteq F'$ 是显然的. \redann{$F=F'$？}}
\handline{72}{280}{390}{Q2：群 $\to$ 域 $\to$ 群.}
\handline{72}{310}{786}{$E$ 是域，$G$ 有限 $\le\operatorname{Aut}(E)$，$F=\operatorname{Inv}(G)$. $\Rightarrow[E:F]<\infty\Rightarrow\operatorname{Gal}(E/F)$ 有限.}
\handline{72}{339}{763}{问 $\operatorname{Gal}(E/F)$ 与 $G$ 什么关系？$G\subseteq\operatorname{Gal}(E/F)$ 显然. \redann{$G=\operatorname{Gal}(E/F)$？}}
\handline[NoteRed]{72}{366}{450}{结论：Q2 正确，称为 Artin 定理.}
\handline[NoteRed]{72}{395}{340}{Q1 错误，下面有反例.}
\handline[NoteRed]{72}{424}{747}{Galois 理论就是研究什么样的有限扩张 $E/F$，满足 $\operatorname{Inv}(\operatorname{Gal}(E/F))=F$.}
\handline{72}{456}{416}{\textbf{Ex.} \circled{1} $E=\Q(\sqrt[3]{2})$，$F=\Q$.}
\handline{72}{485}{717}{计算 $\operatorname{Gal}(E/F)$. 首先 $\sqrt[3]{2}$ 在 $\Q$ 中极小多项式为 $x^3-2$，次数为 $3$.}
\handline{72}{514}{590}{所以 $E=\{x_1+x_2\alpha+x_3\alpha^2\mid\alpha=\sqrt[3]{2},\ x_1,x_2,x_3\in\Q\}$.}
\handline{72}{543}{770}{设 $\sigma\in\operatorname{Aut}(E)$，$\sigma(\alpha)=\beta\in E$，$\sigma(\alpha^3)=\sigma(2)=2\sigma(1)=2\Rightarrow\beta^3=2$.}
\handline{72}{562}{180}{且 $\sigma|_\Q=\id_\Q$.}
\handline{72}{585}{544}{所以 $\beta=\sqrt[3]{2}=\alpha$，所以 $\sigma(\alpha)=\alpha$，$\sigma(\alpha^2)=\alpha^2$.}
\handline{72}{615}{738}{$\Rightarrow\sigma(x_1+x_2\alpha+x_3\alpha^2)=x_1\sigma(1)+x_2\sigma(\alpha)+x_3\sigma(\alpha^2)=x_1+x_2\alpha+x_3\alpha^2$.}
\handline{72}{645}{750}{所以 $\sigma=\id\Rightarrow\operatorname{Gal}(E/F)=\{\id_E\}$. 所以 $\operatorname{Inv}(\{\id_E\})=E\ne F$.}
\handline[AnnotationBlue]{72}{673}{697}{原因是 $x^3-2$ 还有两个根 $\alpha\omega$ 与 $\alpha\omega^2$ 不在 $E$ 中，$E$ 不够大.}
\handline{72}{703}{807}{\ann{若把 $\alpha\omega,\alpha\omega^2$ 加进去，则 $\operatorname{Inv}(\operatorname{Gal}(E/F))=F$.}\redann{在 Galois 理论中会学到.}}
\handline{72}{735}{653}{\circled{2} $F=\F_2(T)$，$T$ 是未定元. $F$ 是 $\F_2$ 上有理函数域.}
\handline{72}{767}{430}{$E=F(\sqrt T)=\operatorname{Span}_F(1,S)$，$S=\sqrt T$.}
\handline{72}{839}{195}{乘法表 $\begin{array}{c|cc}\cdot&1&S\\\hline1&1&S\\S&S&T\end{array}$}
\handline{309}{803}{479}{设 $\sigma\in\operatorname{Aut}(E)$ 且 $\sigma|_F=\id_F$. 设 $\sigma(S)=u$.}
\handline{322}{836}{547}{$u^2=\sigma(S^2)=\sigma(T)=T$. 设 $u=a+bS$，$(a,b\in F)$.}
\handline{302}{869}{626}{$\Rightarrow u^2=a^2+2abS+b^2S^2=a^2+b^2S^2=a^2+b^2T=T\Rightarrow(1+b^2)T=a^2$.}
\handline{72}{897}{695}{结合两边次数 $\Rightarrow1-b^2=a^2=0\Rightarrow b=1,a=0\Rightarrow u=S$，所以 $\sigma(S)=S$.}
\handline{72}{930}{668}{所以 $\sigma=\id_E$，所以 $\operatorname{Gal}(E/F)=\{\id_E\}$，$\operatorname{Inv}(\operatorname{Gal}(E/F))=E\ne F$.}
\end{handpage}

% ===== Line-layout page-017.tex =====
\begin{handpage}{17}
\handtitle{72}{109}{330}{$\S2$ 环与模}
\handtitle{72}{143}{380}{$\S2.1$ 定义，例子.}
\handline{72}{172}{553}{\textbf{Def.} (1) $R$ 是一个集合，$+:R\times R\to R$，$\cdot:R\times R\to R$}
\handline{72}{198}{470}{满足：\circled{1} $\forall a,b,c\in R$，$(a+b)+c=a+(b+c)$.}
\handline{72}{226}{320}{\circled{2} $\forall a,b\in R$，$a+b=b+a$}
\handline{72}{254}{360}{\circled{3} $\exists0_R\in R$，$\forall a\in R$，$a+0_R=0_R+a=a$}
\handline{72}{282}{372}{\circled{4} $\forall a\in R$，$\exists b\in R$ s.t. $a+b=b+a=0_R$}
\handline{72}{310}{348}{\circled{5} $\forall a,b,c\in R$，$(ab)c=a(bc)$}
\handline{72}{339}{536}{\circled{6} $\forall a,b,c\in R$，$a(b+c)=ab+ac$，$(b+c)a=ba+ca$.}
\handline{72}{369}{355}{则 $(R,+,\cdot)$ 称为一个环.}
\handline{72}{399}{768}{(2) 若 $\exists1_R\in R$，s.t. $\forall a\in R$，$a\cdot1_R=1_R\cdot a=a$，则 $1_R$ 唯一，称为 $R$ 中单位元.}
\handline{72}{428}{360}{$(R,+,\cdot)$ 此时称为幺环.}
\handline{72}{459}{625}{(3) 若 $\forall a,b\in R$，$ab=ba$，则称 $(R,+,\cdot)$ 为交换环.}
\handline[AnnotationBlue]{72}{489}{712}{环与域最大的差别是，不一定所有非零元都存在乘法逆.}
\handline[AnnotationBlue]{72}{520}{422}{含不含幺、交不交换是次要的.}
\handline{72}{557}{130}{\textbf{Ex.}}
\handline{72}{590}{479}{\circled{1} 所有的域都是环，如 $\Q,\R,\C$ 等.}
\handline{72}{622}{584}{\circled{2} $\Z$，$\Z_m=\{\bar0,\bar1,\ldots,\overline{m-1}\}$（$m\in\omega$；剩余类环）.}
\handline{72}{653}{440}{\circled{3} $F[x]$，（$F$ 是一个域）多项式环.}
\handline{72}{686}{550}{\circled{4} $V$ 是域 $F$ 上 $n$ 维线性空间，则 $V\cong F^n$.}
\handline{72}{719}{595}{$\mathcal L(V)$ 为 $V$ 上线性变换全体构成的线性空间.}
\handline{72}{752}{529}{则 $(\mathcal L(V),+,\circ)\cong M_n(F)$ \ann{（非交换环）}}
\handline{72}{827}{400}{\circled{5} $R=m\Z$（$m\in\omega$，$m\ge2$）}
\handline{72}{861}{798}{$ma+mb:=m(a+b)$，$(ma)\cdot(mb):=m(mab)$. 则 $R$ 是一个不含幺的环.}
\end{handpage}

% ===== Line-layout page-018.tex =====
\begin{handpage}{18}
\handline{72}{104}{503}{\circled{6} 有关代数数论与环论起源的例子：}
\handline{72}{135}{644}{$\Z[i]$ Gauss 整环，$\Z(\zeta_m)$，$\zeta_m=e^{\frac{2\pi i}{m}}$，$m$ 次单位根.}
\handline{72}{164}{684}{$\Z[\sqrt D]$，$D$ 中无平方因子，可以为负. 与 Pell 方程 $x^2-Dy^2=N$ 有关.}
\handline[AnnotationBlue]{72}{192}{799}{Kummer 证明了 $\Z[\sqrt D]$ 上没有唯一分解定理，反驳了 Cauchy 等人对费马大定理的证明，}
\handline{72}{221}{610}{$i,\sqrt D,\zeta_m$ 都是首一整系数多项式的根，称为代数整数.}
\handline{72}{250}{439}{由此引出一些代数数论的问题.}
\handline{72}{280}{335}{\circled{7} $\bar X$：集合，$R$：环.}
\handline{72}{312}{400}{$R^{\bar X}:=\{f:\bar X\to R\}$. \ann{函数环.}}
\handline{72}{342}{565}{$(f+g)(x):=f(x)+g(x)$，$(f\cdot g)(x):=f(x)\cdot g(x)$.}
\handline{72}{369}{300}{\circled{8} 来自群表示论.}
\handline{72}{400}{543}{设 $G$ 是有限群，$F$ 是一个域. 令 $R=\operatorname{Span}_F(G)$.}
\handline{72}{431}{670}{$a,b\in R$，$a=\sum_{g\in G}a_gg$（$a_g\in F$），$b=\sum_{h\in G}b_hh$（$b_h\in F$）.}
\handline{72}{459}{290}{$+$：按一般加法定义.}
\handline{72}{488}{380}{$\cdot$：$a\cdot b:=\sum(a_g\cdot b_h)(g\cdot h)$}
\handline[NoteRed]{487}{488}{296}{（乘法还可以用卷积定义，}
\handline[NoteRed]{72}{520}{343}{比如多项式环就用卷积来定义乘法）.}
\handline{72}{555}{658}{则 $(R,+,\cdot)$ 是一个环，且 $1_R=e_G\in G$. 即 $R$ 是幺环.}
\handline{72}{585}{440}{记 $R=F(G)$，称为 $G$ 的群代数.}
\handline{72}{616}{594}{\circled{9} Quiver 表示（有向、有限的图，不含定向圈）}
\handline{72}{648}{602}{道路 (path)：$1^\circ$ 每个顶点定义了长度为 $0$ 的道路.}
\handline{72}{679}{423}{$2^\circ$ 每个箭头定义了长度为 $1$ 的道路.}
\handline{72}{711}{710}{$3^\circ$ 若一个道路 $\ell_1$ 的头与另一个道路 $\ell_2$ 的尾相同，可定义新道路 $\ell=\ell_2\cdot\ell_1$.}
\handline{72}{747}{470}{设 $F$：域，$Q$ 是一个 Quiver 表示.}
\handline{72}{780}{488}{$R:=F[Q]=\operatorname{Span}_F\{Q\text{ 的所有道路}\}$.}
\handline{72}{815}{431}{若 $\ell_1$ 头与 $\ell_2$ 尾不同，则 $\ell_2\cdot\ell_1:=0$.}
\handline{72}{846}{422}{则 $R=F[Q]$ 是 $Q$ 的群代数.}
\handline{72}{883}{737}{$1_R=\sum_{i=1}^{|Q|}V_i$，$V_i$ 为 $Q$ 中的顶点，\ann{（即长度为 $0$ 的道路本体）}.}
\end{handpage}

% ===== Line-layout page-019.tex =====
\begin{handpage}{19}
\handline{72}{127}{591}{\circled{10} 四元数 (Hamilton) $\mathbb H=\operatorname{Span}_{\R}(1,i,j,k)$.}
\handline{72}{158}{596}{$i^2=j^2=k^2=-1$，$ij=-ji=k$，$jk=-kj=i$，$ki=-ik=j$}
\handline{72}{187}{654}{$\mathbb H$ 不交换，含幺，非零元可逆\ann{（有些书称这种环为体）}}
\handline[AnnotationBlue]{72}{215}{353}{体 $+$ 乘法交换 $\Rightarrow$ 域}
\handline{72}{245}{639}{可逆：$(a+bi+cj+dk)(a-bi-cj-dk)=a^2+b^2+c^2+d^2\in\R$}
\handline{72}{279}{435}{所以 $(a+bi+cj+dk)^{-1}=\dfrac{a-bi-cj-dk}{a^2+b^2+c^2+d^2}$.}
\handline{72}{307}{250}{四元数起源：}
\handline{72}{338}{547}{$\varphi:\C\to M_2(\R)$，则 $\varphi$ 是环的单同态.}
\handline{72}{367}{534}{$a+bi\mapsto\left(\begin{smallmatrix}a&-b\\b&a\end{smallmatrix}\right)$，$(\R^2,\varphi)$ 是一个 $\R$-代数.}
\handline{72}{396}{433}{$\left(\begin{smallmatrix}a&-b\\b&a\end{smallmatrix}\right)$ 即为旋转与伸缩的复合.}
\handline{72}{426}{782}{Hamilton 想在 $\R^3$ 上定义一个类似的运算来描述刚体运动\redann{（但失败）}}
\handline{72}{454}{808}{\textbf{Def.} $V$ 是 $\R$-线性空间，若有双线性映射 $\mu:V\times V\to V$，则称 $(V,\varphi)$ 为一个 $\R$-代数.}
\handline[AnnotationBlue]{72}{484}{554}{$\mu$ 可以加各种条件：如自反、对称、传递……}
\handline[AnnotationBlue]{72}{516}{573}{如 $\R^3$ 上的叉积，具有反称性 $a\times b=-b\times a$.}
\handline[AnnotationBlue]{72}{547}{686}{同时满足 Jacobi 恒等式：$(u\times v)\times w+(v\times w)\times u+(w\times u)\times v=0$}
\handline[AnnotationBlue]{72}{578}{730}{\textbf{Def.} 在域 $F$ 的向量空间 $V$ 中，若有双线性映射 $\mu:V\times V\to V$，}
\handline[AnnotationBlue]{72}{609}{688}{且 $\mu$ 满足反称性、Jacobi 恒等式，则 $V$ 称为李代数 (Lie)}
\handline[NoteRed]{72}{643}{731}{\textbf{Thm.} 设 $V$ 是一个 $\R$-代数，含有单位元 $e$，且 $\dim V=2k+1$，$k\ge1$，$\in\omega$，}
\handline[NoteRed]{72}{674}{314}{则 $(V,\mu)$ 不是域.}
\handline[NoteRed]{72}{708}{770}{\textbf{Pf.} $\forall\alpha\in V$，与 $e$ $\R$-线性无关，考虑 $T_\alpha:V\to V$，$T_\alpha\in\mathcal L(V)$}
\handline[NoteRed]{72}{735}{219}{$\beta\mapsto\mu(\alpha,\beta)$}
\handline[NoteRed]{72}{772}{624}{由于 $\dim V$ 为奇，故 $T_\alpha$ 必有实特征值 $\lambda\in\R$.}
\handline[NoteRed]{72}{805}{641}{$\exists\beta\ne0$ s.t. $\mu(\alpha,\beta)=\lambda\beta=\lambda\mu(e,\beta)=\mu(\lambda e,\beta)$}
\handline[NoteRed]{72}{841}{300}{所以 $\mu(\alpha-\lambda e,\beta)=0$.}
\handline[NoteRed]{72}{874}{376}{而 $\alpha-\lambda e$、$\beta$ 均为非零元.}
\handline[NoteRed]{72}{905}{452}{所以 $(V,\mu)$ 存在零因子，故不是域.}
\handline[NoteRed]{72}{944}{713}{也就是 $\R^3$ 上没有类似于 $(\R^2,\varphi)$ 这样的 $\R$-代数.}
\end{handpage}

% ===== Line-layout page-020.tex =====
\begin{handpage}{20}
\handline{72}{100}{591}{但 Hamilton 发现 $\R^4$ 上可以有这样子的 $\R$-代数.}
\handline{72}{131}{578}{在 $\R^4$ 中将 $\R^4$ 视为 $\C^2\Rightarrow\mathbb H=\{(z,w)\mid z,w\in\C\}$.}
\handline{72}{158}{573}{$\varphi:\C^2\to M_2(\C)$，则 $(\R^4,\varphi)$ 是一个 $\R$-代数.}
\handline{72}{188}{541}{$(z,w)\mapsto\begin{pmatrix}z&-w\\\bar w&\bar z\end{pmatrix}$，且 $\mathbb H\cong\{\begin{pmatrix}z&-w\\\bar w&\bar z\end{pmatrix}\mid z,w\in\C\}$.}
\handline{72}{237}{595}{令 $1=\begin{pmatrix}1&0\\0&1\end{pmatrix}$，$I=\begin{pmatrix}i&0\\0&-i\end{pmatrix}$，$J=\begin{pmatrix}0&-1\\1&0\end{pmatrix}$，$K=\begin{pmatrix}0&-i\\-i&0\end{pmatrix}$.}
\handline{72}{273.408}{416}{则 $\forall q\in\{\begin{pmatrix}z&-w\\\bar w&\bar z\end{pmatrix}\mid z,w\in\C\}$}
\handline{72}{319}{666}{$1\cdot q=q$，$I\cdot J=-J\cdot I=K$，$J\cdot K=-K\cdot J=I$，$K\cdot I=-I\cdot K=J$.}
\handline[AnnotationBlue]{72}{351}{586}{可以用类似的方法定义八元数、十六元数，……}
\handline[AnnotationBlue]{72}{382}{665}{统称为 Clifford 代数. 二维 Clifford 代数等价于复数.}
\handline[AnnotationBlue]{72}{411}{469}{三维 Clifford 代数等价于四元数.}
\handline[AnnotationBlue]{72}{442}{735}{在相对论（时空代数）中，Clifford 代数可以更自然地描述}
\handline[AnnotationBlue]{72}{472}{682}{物理定律，比如将麦克斯韦方程组简化为一个方程.}
\handline[AnnotationBlue]{72}{502}{540}{尤其擅长简化高维空间的几何运算.}
\handline[AnnotationBlue]{72}{534}{805}{三十二元数应用于量子力学（$\gamma$ 矩阵，Dirac 方程）、Spin Geometry（自旋几何）.}
\handline[AnnotationBlue]{72}{567}{390}{Index theorem（指标定理）}
\handline[NoteRed]{72}{597}{758}{\textbf{Thm.} 设 $\dim V=2$，$(V,\mu)$ 是一个交换、结合且含幺的 $\R$-代数且为域.}
\handline{72}{627}{639}{\redann{则 $(V,\mu)\cong\C$}\ann{（揭示了 $\C$ 某种程度上的唯一性）}}
\handline[NoteRed]{72}{657}{624}{\textbf{Pf.} 设 $e$ 是单位元，$e\ne0$，任取 $f\in V$，与 $e$ 线性无关.}
\handline[NoteRed]{72}{688}{474}{考虑 $(e,f)$ 的乘法表. 取 $f_1=f-\frac y2e$}
\handline[NoteRed]{72}{765}{158}{$\begin{array}{c|cc}\mu&e&f\\\hline e&e&f\\f&f&xe+yf\end{array}$}
\handline[NoteRed]{72}{813}{150}{（$x,y\in\R$）}
\handline[NoteRed]{72}{720}{277}{则 $ef_1=f_1e=f_1$.}
\handline[NoteRed]{309}{751}{322}{$f_1\cdot f_1=(f-\frac y2e)(f-\frac y2e)$}
\handline[NoteRed]{315}{782}{479}{$=f\cdot f-yf+\frac{y^2}4e=xe+\frac{y^2}4e=\widetilde xe$.}
\handline[NoteRed]{258}{819}{670}{设 $\widetilde x=x+\frac{y^2}4$. 若 $\widetilde x\ge0$，则取 $f_2=f_1+\sqrt{\widetilde x}\,e$，$f_3=f_1-\sqrt{\widetilde x}\,e$.}
\handline[NoteRed]{72}{849}{732}{则 $f_2\cdot f_3=f_1\cdot f_1-\widetilde xe=0$，所以 $f_1,f_3$ 是零因子，与域矛盾.}
\handline[NoteRed]{72}{887}{520}{所以 $\widetilde x<0$，取 $f_4=f_1/\sqrt{-\widetilde x}$，则 $f_4\cdot f_4=-e$.}
\handline[NoteRed]{72}{917}{678}{定义 $\varphi:(V,\mu)\to\C$，则 $\varphi$ 是同构，所以 $(V,\mu)\cong\C$.}
\handline[NoteRed]{72}{932}{116}{$e\mapsto1$}
\handline[NoteRed]{72}{947}{123}{$f_4\mapsto i$}
\end{handpage}

% ===== Line-layout page-021.tex =====
\begin{handpage}{21}
\handline{72}{157}{700}{\textbf{Def.} (1) $R$：环，$M$：集合，$+:M\times M\to M$，$\cdot:R\times M\to M$.}
\handline{72}{185}{560}{若 $+$ 与 $\cdot$ 满足 \circled{1} $(x+y)+z=x+(y+z)$ \circled{2} $x+y=y+x$}
\handline{72}{212}{450}{\circled{3} $\exists0_M\in M$ s.t. $\forall x\in M$，$0_M+x=x+0_M=x$}
\handline{72}{242}{523}{\circled{4} $\forall x\in M$，$\exists y\in M$ s.t. $x+y=0_M$. 记 $y=-x$.}
\handline{72}{270}{368}{\circled{5} $\forall a,b\in R$，$x\in M$，$a(bx)=(ab)x$}
\handline{72}{299}{460}{\circled{6} $a(x+y)=ax+ay$，$a\in R$，$x,y\in M$.}
\handline{72}{329}{440}{\circled{7} $(a+b)x=ax+bx$，$a,b\in R$，$x\in M$}
\handline{72}{361}{708}{则称 $M$ 是一个左 $R$ 模. 右 $R$ 模同理\ann{（后面默认讨论左 $R$ 模）}}
\handline{72}{391}{390}{若 $R$ 可交换，则左右 $R$ 模等价.}
\handline{72}{421}{455}{(2) 若 $R$ 是幺环，且 $1_Rx=x$，$\forall x\in M$ \ann{(*)}}
\handline{72}{452}{312}{则称 $M$ 是一个幺模.}
\handline{72}{481}{325}{(3) 若 $R,S$ 是两个环.}
\handline{72}{512}{515}{定义两种数乘 $\cdot_R:R\times M\to M$，$\cdot_S:M\times S\to M$}
\handline{72}{544}{551}{$M$ 既是左 $R$ 模又是右 $S$ 模，称为 $(R,S)$ 双模.}
\handline{72}{576}{802}{\textbf{Ex.} \circled{1} 模的定义有 $7$ 条公理，而域 $F$ 上线性空间的定义有 $8$ 条公理.}
\handline{72}{605}{370}{多了一条 $1_F\cdot\alpha=\alpha$，即 \ann{(*)}.}
\handline{72}{636}{618}{因此 $F$-线性空间是一个左 $F$ 模，且是一个幺模.}
\handline[AnnotationBlue]{72}{666}{643}{线性空间是用来理解模这一概念的重要例子.}
\handline{72}{697}{667}{\circled{2} $\Z_m$，$+:\Z_m\times\Z_m\to\Z_m$，$\cdot:\Z\times\Z_m\to\Z_m$}
\handline{72}{726}{445}{$(\bar a,\bar b)\mapsto\overline{a+b}$，\hspace{35mm}$(a,\bar x)\mapsto\overline{ax}$}
\handline{72}{767}{330}{则 $\Z_m$ 是一个左 $\Z$ 模.}
\handline{72}{803}{478}{\redann{$\star$} 事实上，所有 Abel 群与 $\Z$ 模等价.}
\handline{72}{837}{600}{$G$：交换群，$+$：按 $G$ 中加法，$\cdot:\Z\times G\to G$}
\handline{72}{884.311}{332}{$(n,g)\mapsto ng=\begin{cases}\underbrace{g+\cdots+g}_{n\text{ 个}},&n>0,\\0,&n=0,\\-|n|g,&n<0.\end{cases}$}
\handline[NoteRed]{72}{933}{740}{这一观点在研究有限 Abel 群结构与有限生成模结构时会用到.}
\end{handpage}

% ===== Line-layout page-022.tex =====
\begin{handpage}{22}
\handline{72}{96}{450}{\circled{3} $R=\R[T]$，$M$ 是一个 $\R$-线性空间.}
\handline{72}{119}{470}{首先加法是 $M$ 作为线性空间自带的.}
\handline{72}{145}{480}{我们要问如何定义乘法使 $M$ 成为 $R$ 模}
\handline{72}{174}{230}{$\cdot:R\times M\longrightarrow M$}
\handline{72}{202}{300}{$(f(T),x)\longmapsto f(T)(x)$.}
\handline{72}{228}{620}{而如果考虑 $f(T)=T$，就要问 $T$ 作用在 $x$ 上是什么意思.}
\handline{72}{254}{500}{而 $T:M\to M$ 就可以取成一个线性变换.}
\handline{72}{285}{640}{于是就可以取 $M$ 上线性变换全体作为 $\R$ 上多项式的未定元.}
\handline{72}{311}{590}{而 $f(T)$ 就是线性变换的多项式.\quad$T^n:=T\circ T\circ\cdots\circ T$}
\handline[AnnotationBlue]{72}{368}{665}{于是研究一元多项式环上的模等价于研究 $M$ 上的线性变换.}
\handline[AnnotationBlue]{72}{399}{635}{后面关于模的分类定理就可以应用在这个例子上.}
\handline[AnnotationBlue]{72}{425}{735}{就发现线性代数中的若当标准型理论是分类定理的特殊情况.}
\handline{72}{457}{525}{\circled{4} $\widetilde R=M_n(R)$，$R$ 是交换幺环，$\widetilde M=R^n$.}
\handline{72}{487}{225}{$\cdot:\widetilde R\times\widetilde M\longrightarrow\widetilde M$.}
\handline{72}{518}{620}{就定义为 $\widetilde R$ 中的矩阵在 $\widetilde M$ 中向量上的作用，得到 $\widetilde M$ 中的向量.}
\handline{72}{550}{325}{则 $\widetilde M$ 是一个 $\widetilde R$ 模.}
\handline{72}{580}{750}{\circled{5} $U:\R^n$ 中的开集，$A=C^\infty(U)$，$D:U$ 上的 $C^\infty$ 微分算子环.}
\handline{72}{624}{705}{$\displaystyle D=\left\{\sum_{k(i_1\cdots i_k)}\frac{a_{i_1i_2\cdots i_k}(x)\partial^k}{\partial x_{i_1}\partial x_{i_2}\cdots\partial x_{i_k}}\ \middle|\ i_1,\ldots,i_k\text{ 是 }1\cdots k\text{ 的排列},\ k\in\N\right\}$.}
\handline{72}{676}{505}{$D$ 是非交换的.\quad\ann{例如，只考虑 $1$ 维情形}}
\handline[AnnotationBlue]{72}{713}{440}{定义 $x:C^\infty(U)\longrightarrow C^\infty(U)$}
\handline[AnnotationBlue]{72}{745}{210}{$f\longmapsto xf$}
\handline[AnnotationBlue]{72}{783}{385}{$\partial=\dfrac d{dx}:C^\infty(U)\longrightarrow C^\infty(U)$}
\handline[AnnotationBlue]{72}{824}{225}{$f\longmapsto\dfrac d{dx}f$.}
\handline[AnnotationBlue]{72}{855}{455}{$x\partial\ne\partial x$，即 $x\,df\ne d(xf)$.}
\handline{72}{886}{700}{定义 $D$ 上乘法：$D\times D\longrightarrow D$\qquad 定义数乘 $D\times A\longrightarrow A$}
\handline{72}{918}{550}{$(D_1,D_2)\longmapsto D_1\circ D_2$\qquad\qquad $(D,f)\longmapsto D(f)$.}
\handline{72}{956}{275}{则 $A$ 是左 $D$ 模.}
\end{handpage}

% ===== Line-layout page-023.tex =====
\begin{handpage}{23}
\handline{72}{131}{655}{另一种看法：$A\subseteq D$，视 $A$ 为零阶算子构成的子环.}
\handline{72}{161}{330}{定义数乘 \circled{1} $D\times A\longrightarrow D$}
\handline{72}{185}{410}{$(D,f)\longmapsto D\circ f$\qquad $D$ 是右 $A$ 模}
\handline{72}{214}{215}{\circled{2} $A\times D\longrightarrow D$}
\handline{72}{242}{400}{$(f,D)\longmapsto f\circ D$\qquad $D$ 是左 $A$ 模.}
\handline{72}{274}{370}{$\therefore D$ 是 $(A,A)$ 双模.}
\end{handpage}

% ===== Line-layout page-024.tex =====
\begin{handpage}{24}
\handtitle{72}{94}{375}{§2.2 同态的基本概念.}
\handline{72}{121}{310}{\textbf{Def.}\quad (a) $R,S$ 为两个环.}
\handline{72}{147}{605}{$\varphi:R\longrightarrow S$ 满足\quad$\varphi(a+b)=\varphi(a)+\varphi(b)$\quad\redann{$\Rightarrow\varphi(0_R)=0_S$}}
\handline{72}{171}{205}{$\varphi(ab)=\varphi(a)\varphi(b)$.}
\handline{72}{199}{230}{则称 $\varphi$ 是环同态.}
\handline{72}{229}{445}{若 $R,S$ 均为幺环，还要求 $\varphi(1_R)=1_S$.}
\handline{72}{259}{405}{(b) $R$ 是一个环，$M,N$ 是 $R$ 的两个模}
\handline{72}{285}{155}{$\varphi:M\longrightarrow N$}
\handline{72}{314}{260}{$\varphi(m+n)=\varphi(m)+\varphi(n)$.}
\handline{72}{343}{475}{$\varphi(mn)=\varphi(m)\varphi(n)$\quad 则称 $\varphi$ 是一个模同态.}
\handline[AnnotationBlue]{72}{376}{545}{特别地：由于域都是环，故域同态都是环同态.}
\handline[AnnotationBlue]{72}{405}{735}{域上线性空间都是域上的模，故线性映射都是模同态.}
\handline{72}{434}{530}{\textbf{Ex.}\quad\circled{1} $\varphi:\Z\longrightarrow\Z_m$\qquad $\varphi$ 是环同态.}
\handline{72}{463}{185}{$n\longmapsto n\pmod m$}
\handline[NoteRed]{72}{497}{345}{但 $\Z_m\to\Z$ 不存在环同态.}
\handline[NoteRed]{72}{528}{510}{因 $\Z_m,\Z$ 均为幺环.若存在同态 $\varphi:\Z_m\to\Z$}
\handline[NoteRed]{72}{563}{390}{则 $\varphi(\overline1)=1\Rightarrow\varphi(\overline m)=m$}
\handline[NoteRed]{72}{593}{415}{而 $\varphi(\overline m)=\varphi(\overline0)=0\ne m$，矛盾!}
\handline[NoteRed]{72}{625}{330}{$\therefore\Z_m$ 与 $\Z$ 不存在环同态.}
\handline{72}{657}{460}{\circled{2} $\varphi:\Z_{mn}\to\Z_m$ 存在自然同态.}
\handline{72}{693}{420}{但同样 $\Z_m\to\Z_{mn}$ 不存在同态.}
\handline{72}{723}{730}{\circled{3} $R=F[x]$，$F:$ 域，$\alpha\in F$，定义 $\operatorname{ev}_\alpha:R\longrightarrow F$\quad\ann{(evaluation)}}
\handline{72}{753}{170}{$f\longmapsto f(\alpha)$.}
\handline{72}{793}{485}{则 $\operatorname{ev}_\alpha$ 是环同态.\ann{（还是 $R$-代数同态.）}}
\handline{72}{831}{550}{进一步，设 $E/F$ 扩域，$\alpha\in E$，$\operatorname{ev}_\alpha:R\longrightarrow F$}
\handline{72}{861}{170}{$f\longmapsto f(\alpha)$}
\handline[AnnotationBlue]{72}{906}{570}{考虑 $\operatorname{ev}_\alpha$，实际上是在考虑 $F(\alpha)$ 这个单扩张.}
\handline[NoteRed]{72}{811}{160}{从而可以用环}
\handline[NoteRed]{710}{837}{155}{同态的观}
\handline[NoteRed]{262}{867}{150}{点来看待}
\handline[NoteRed]{723}{899}{145}{域扩张.}
\handarrow[NoteRed]{643.0}{890}{690.0}{847}
\end{handpage}

% ===== Line-layout page-025.tex =====
\begin{handpage}{25}
\handline{72}{155}{585}{\textbf{Def.}\quad$\varphi:R\to S$，$R,S$ 为两个环，$\varphi$ 是环同态.}
\handline{72}{184}{760}{\circled{1} $\varphi$ 是单射，$\varphi$ 称为单同态 \ann{(monomorphism / injective homomorphism)}}
\handline{72}{212}{760}{\circled{2} $\varphi$ 是满射，$\varphi$ 称为满同态 \ann{(epimorphism / surjective homomorphism)}}
\handline{72}{239}{495}{\circled{3} $\varphi$ 是双射，$\varphi$ 称为同构 \ann{(isomorphism)}}
\handline[AnnotationBlue]{72}{271}{610}{\textbf{Fact:}\quad$\varphi:R\to S$ 是同构，则 $\varphi^{-1}:S\to R$ 也是同构}
\handline{72}{300}{290}{\textbf{Def.}\quad$\varphi:R\to S$ 同态.}
\handline{72}{329}{345}{$\ker\varphi:=\{a\in R\mid\varphi(a)=0_S\}$}
\handline{72}{358}{370}{$\im\varphi:=\{b\in S\mid\exists a\in R,b=\varphi(a)\}$}
\handline[AnnotationBlue]{72}{385}{420}{\textbf{Fact:}\quad\circled{1} $\varphi$ 是单同态 $\iff\ker\varphi=\{0_R\}$}
\handline[AnnotationBlue]{72}{413}{310}{\circled{2} $\varphi$ 是满同态 $\iff\im\varphi=S$}
\handline{72}{447}{245}{\textbf{Ex.}\quad$\varphi:\Z\to\Z_m$}
\handline{72}{478}{305}{$\ker\varphi=m\Z$，$\im\varphi=\Z_m$.}
\handline{72}{514}{760}{\textbf{Def.}\quad 若 $S\subseteq R$，$\forall a\in S,b\in S$，$\Rightarrow a+b\in S,ab\in S,-a\in S,-b\in S$.}
\handline{72}{544}{245}{则 $S$ 是 $R$ 的子环.}
\handline[AnnotationBlue]{72}{577}{750}{\textbf{Fact:}\quad$\varphi:R\to S$ 同态，则 $\ker\varphi$ 是 $R$ 的子环，$\im\varphi$ 是 $S$ 的子环.}
\handline{72}{616}{545}{\textbf{Lemma.}\quad 若 $a\in\ker\varphi$，$b\in R$，则 $ab\in\ker\varphi$.}
\handline{72}{645}{695}{\textbf{Pf.}\quad$\varphi(ab)=\varphi(a)\cdot\varphi(b)=0_S\cdot\varphi(b)=0_S$，$\therefore ab\in\ker\varphi$.}
\handline[AnnotationBlue]{72}{678}{405}{即 $\ker\varphi$ 是 $R$ 的一个理想 (ideal)}
\end{handpage}

% ===== Line-layout page-026.tex =====
\begin{handpage}{26}
\handtitle{72}{120}{340}{§2.3 子环、理想、商环}
\handline{72}{153}{750}{\textbf{Def.}\quad$R$ 是一个环，$S\subseteq R$，若 $a\in S,b\in S\Rightarrow\left\{\begin{smallmatrix}a+b\in S\\ab\in S\\-a\in S\\-b\in S\end{smallmatrix}\right.$ 则称 $S$ 为 $R$ 的子环.}
\handline[AnnotationBlue]{72}{200}{610}{显然，子环包含加法单位 $0$.\quad$a\in S\Rightarrow-a\in S\Rightarrow0=a+(-a)\in S$.}
\handline{72}{230}{675}{\textbf{Ex.}\quad\circled{1} $\Z$ 是 $\Q$ 的子环.\circled{2} $F$ 是域，$F$ 是 $F[x]$ 的子环.\circled{3} $\C$ 是 $\mathbb H$ 的子环.}
\handline{72}{255}{285}{\circled{4} $\varphi:R\to S$ 环同态.}
\handline{72}{281}{440}{$\ker\varphi$ 是 $R$ 的子环，$\im\varphi$ 是 $S$ 的子环.}
\handline{72}{309}{550}{\textbf{Def.}\quad$I\subseteq R$ 是子环，$\forall a\in I,b\in R\Rightarrow ab,ba\in I$}
\handline{72}{338}{695}{则 $I$ 称为 $R$ 的（双边）理想.\ann{（更一般地可以定义左理想和右理想）}}
\handline{72}{368}{525}{\textbf{Ex.}\quad\circled{1} $\varphi:R\to S$ 环同态，$\ker\varphi$ 是理想.}
\handline[AnnotationBlue]{72}{393}{620}{（事实上任何理想 $I$ 都是某个 $\varphi$ 的核，$\varphi$ 取为典则投影即可）}
\handline{72}{425}{350}{\circled{2} $R,\{0_R\}$ 称为 $R$ 的平凡理想}
\handline{72}{454}{485}{\circled{3} 若 $R$ 是幺环，$I$ 是理想，且 $1_R\in I$，则 $I=R$.}
\handline{72}{483}{445}{$\forall a\in R$，$1_R\in I\Rightarrow a=a\cdot1_R\in I\Rightarrow I=R$}
\handline{72}{575}{375}{\circled{4} 若 $R$ 是域，则 $R$ 只有平凡理想.}
\handline{72}{607}{385}{设 $I$ 是 $R$ 的理想，且 $I\ne\{0_R\}$.}
\handline{72}{637}{380}{$\therefore\exists a\ne0_R\in I$，由 $R$ 是域，$\therefore a^{-1}$ 存在.}
\handline{72}{667}{650}{$a\in I$，$a^{-1}\in R\Rightarrow1_R=aa^{-1}\in I\Rightarrow R=(1_R)\subseteq I$，$\therefore I=R$.}
\handline[AnnotationBlue]{72}{701}{435}{\textbf{Fact:}\quad 任意个理想之交仍为理想.}
\handline{72}{735}{730}{\textbf{Def.}\quad 设 $S$ 是环 $R$ 的子集.定义由 $S$ 生成的理想为：$R$ 中所有}
\handline{72}{769}{715}{包含 $S$ 的理想之交，记为 $(S)$\quad\ann{$(S)$ 是 $R$ 中包含 $S$ 的最小理想}}
\handline{72}{803}{560}{若 $S$ 为有限集，则称 $(S)$ 是有限生成的理想.}
\handline{72}{838}{805}{一个生成元\ann{（也称主理想）}：$\displaystyle I_a=\left\{\sum_{i=1}^kf_iag_i\ \middle|\ k\in\N,f_i,g_i\in R\right\}$，\ann{也记作 $(a)$}}
\handline{72}{872}{485}{若 $R$ 是交换环，则 $(a)=\{ag\mid g\in R\}$.}
\handline{72}{916}{856}{$m$ 个生成元：$\displaystyle I_{(a_1,a_2,\ldots,a_m)}=(a_1,a_2,\ldots,a_m)=\left\{\sum_{j=1}^m\sum_{i=1}^{k_j}f_{ij}a_jg_{ij}\ \middle|\ k_j\in\N,f_{ij},g_{ij}\in R\right\}$}
\handline[AnnotationBlue]{72}{960}{810}{无限个生成元生成的理想，要结合泛函与拓扑来定义，这里略.}
\end{handpage}

% ===== Line-layout page-027.tex =====
\begin{handpage}{27}
\handline{72}{111}{365}{\textbf{Def.}\quad$R:$ 环，$I$ 是 $R$ 中的理想.}
\handline{72}{138}{395}{定义二元关系“$\sim$”：$a\sim b\iff b-a\in I$}
\handline{72}{166}{305}{\circled{1} $a\sim a$：$0_R=a-a\in I$}
\handline{72}{190}{610}{\circled{2} $a\sim b\Rightarrow b\sim a$：$a\sim b\iff b-a\in I\Rightarrow a-b\in I\iff b\sim a$.}
\handline{72}{219}{710}{\circled{3} $a\sim b,b\sim c\Rightarrow a\sim c$：$a\sim b,b\sim c\Rightarrow b-a,c-b\in I\Rightarrow c-a\in I\Rightarrow a\sim c$.}
\handline{72}{247}{610}{$\therefore$“$\sim$”是 $R$ 上的等价关系.“$\sim$”将 $R$ 划分为若干等价类.}
\handline{72}{275}{425}{记这些等价类构成的集合为 $R/I$.}
\handline{72}{307}{410}{$a\in R$ 所在的等价类记为 $a+I$.}
\handline{72}{336}{350}{定义 $R/I$ 上的加法与乘法：}
\handline{72}{364}{345}{$(a+I)+(b+I):=(a+b)+I$}
\handline{72}{392}{330}{$(a+I)\cdot(b+I):=ab+I$}
\handline{72}{423}{590}{$0_{R/I}=I$\quad 易验证 $(R/I,+,\cdot,0_{R/I})$ 构成一个环，}
\handline{72}{453}{295}{称为 $R$ 模 $I$ 的商环.}
\handline[AnnotationBlue]{72}{486}{655}{注意，我们需要检查加法和乘法是否是良定义的.}
\handline[AnnotationBlue]{72}{517}{610}{即运算的结果必须不依赖于代表元的选取.}
\handline[NoteRed]{72}{551}{630}{对于 $a'\in a+I$，$b'\in b+I$，则 $a'+I=a+I$，$b'+I=b+I$}
\handline[NoteRed]{72}{583}{470}{$\exists c,d\in I$ s.t. $a'=a+c$，$b'=b+d$}
\handline[NoteRed]{72}{616}{480}{\circled{1} 加法良定：$(a+I)+(b+I)=a+b+I$}
\handline[NoteRed]{72}{646}{300}{$(a'+I)+(b'+I)=a'+b'+I$.}
\handline[NoteRed]{72}{684}{440}{而 $a'+b'=a+b+c+d$，$c+d\in I$.}
\handline[NoteRed]{72}{716}{280}{$\therefore a'+b'+I=a+b+I$.}
\handline[NoteRed]{72}{754}{455}{\circled{2} 乘法良定：$(a+I)(b+I)=ab+I$}
\handline[NoteRed]{72}{784}{280}{$(a'+I)(b'+I)=a'b'+I$.}
\handline[NoteRed]{72}{820}{500}{而 $a'b'=(a+c)(b+d)=ab+ad+cb+cd$}
\handline[NoteRed]{72}{863}{720}{$I$ 是理想，$c,d\in I\Rightarrow ad,cb,cd\in I\Rightarrow ad+cb+cd\in I$}
\handline[NoteRed]{72}{901}{285}{$\therefore a'b'+I=ab+I$}
\end{handpage}

% ===== Line-layout page-028.tex =====
\begin{handpage}{28}
\handline{72}{120}{720}{\textbf{Def.}\quad 设 $R$ 是一个环，$I$ 是 $R$ 的理想.定义 $\pi:R\longrightarrow R/I$，称为典则投影.}
\handline{72}{137}{120}{$a\longmapsto a+I$}
\handline[AnnotationBlue]{72}{172}{690}{$\pi$ 是满同态.若 $\pi(a)=0_{R/I}$，则 $a+I=I\Rightarrow a\in I$，$\therefore\ker\pi\subseteq I$.}
\handline[AnnotationBlue]{72}{199}{555}{$\forall a\in I$，$\pi(a)=a+I=I=0_{R/I}\Rightarrow a\in\ker\pi$，$\therefore I\subseteq\ker\pi$}
\handline[AnnotationBlue]{72}{226}{590}{$\therefore I=\ker\pi$.\quad\redann{从而说明，每一个理想 $I$ 都是其典则投影的核.}}
\handline{72}{261}{585}{\textbf{Ex.}\quad\circled{1} $\Z_m=\Z/m\Z$\quad$\pi:\Z\to\Z_m$，$\ker\pi=m\Z$}
\handline{72}{290}{320}{\circled{2} $\Q(\sqrt2)\cong\Q[x]/(x^2-2)$}
\handline{72}{318}{540}{$p(x)+(x^2-2)=q(x)+(x^2-2)\iff x^2-2\mid p(x)-q(x)$.}
\handline{72}{349}{530}{对于 $\forall p(x)\in\Q[x]$，考虑与 $(x^2-2)$ 的带余除法.}
\handline{72}{378}{650}{$p(x)$ 所在的等价类由余式 $r(x)$ 唯一决定，且 $\deg r(x)\leq1$.}
\handline{72}{408}{675}{于是 $\Q[x]/(x^2-2)$ 中元素可表示为 $a+bx+(x^2-2)$，$a,b\in\Q$.}
\handline{72}{439}{365}{$\varphi:\Q(\sqrt2)\longrightarrow\Q[x]/(x^2-2)$}
\handline{72}{468}{320}{$a+b\sqrt2\longmapsto a+bx+(x^2-2)$.}
\handline{72}{497}{240}{则 $\varphi$ 为同构映射.}
\handline{72}{530}{425}{类似地 $\C=\R(i)\cong\R[x]/(x^2+1)$.}
\handline{72}{562}{420}{\circled{3} $R/R\cong\{0\}$，$R/\{0_R\}\cong R$.}
\handline{72}{590}{595}{\circled{4} 若 $R$ 交换，则 $R/I$ 交换；若 $R$ 含幺，则 $R/I$ 含幺.}
\handline{72}{625}{470}{\textbf{Def.}\quad$R:$ 环，$M:R$ 模，$N:M$ 的子模.}
\handline{72}{655}{360}{同样有 $M/N$ 为 $M$ 模 $N$ 的商模.}
\handline{72}{690}{620}{$\pi:M\longrightarrow N$：$x\longmapsto x+N$\quad 典则投影，且为模的满同态.}
\handline{72}{722}{610}{\textbf{Ex.}\quad$R$ 是一个交换环，则 $R$ 是 $R$ 模.\quad$(R\times R\to R)$.}
\handline{72}{753}{440}{若 $R$ 有一个子模 $I$，则 $I$ 是 $R$ 的理想.}
\handline{72}{788}{410}{反之 $R$ 的理想 $I$ 是 $R$ 的子模.}
\handline{72}{825}{575}{商环 $R/I$ 也是商模：\quad$R\times R/I\longrightarrow R/I$.}
\handline{72}{866}{290}{$(a,b+I)\longmapsto ab+I$.}
\end{handpage}

% ===== Line-layout page-029.tex =====
\begin{handpage}{29}
\handtitle{72}{74}{720}{\S2.4 特征}
\handline{72}{109}{779}{\textbf{Def.} \circled{1} $R$：环. 若 $\exists a,b\ne0\in R$ s.t. $ab=0$，则 $a,b$ 称为 $R$ 的零因子.}
\handline{72}{137}{721}{没有零因子的交换幺环称为整环. 整环满足消去律（$a\ne0$，$ab=ac\Rightarrow b=c$）.}
\handline{72}{164}{739}{\circled{2} 理想都是主理想的整环称为主理想整环. \ann{（PID）}}
\handline{72}{194}{780}{\textbf{Ex.} 一些不是整环的例子：}
\handline{72}{222}{724}{\circled{1} $U\subseteq\R^n$ 为开集.}
\handline{72}{251}{727}{令 $R=C^\infty(U)$，$U$ 上的光滑函数环.}
\handline{72}{279}{730}{取 $f,g$ s.t. $\operatorname{supp}f\cap\operatorname{supp}g=\varnothing\Rightarrow f\cdot g\equiv0$.}
\handline{72}{332}{739}{\circled{2} $A,B\in M_n(\R)$.}
\handline{72}{360}{714}{$AB=0\ \nRightarrow\ A=0\text{ 或 }B=0$. 如 $A=\begin{pmatrix}0&1\\0&0\end{pmatrix}$，$A^2=0$.}
\handline{72}{391}{802}{\textbf{Ex.} 关于主理想整环的例子.}
\handline{72}{421}{713}{$\Z$ 是主理想整环：}
\handline{72}{451}{716}{首先 $\Z$ 是交换幺环. 设 $I$ 是 $\Z$ 的一个理想，且 $I\ne\{0\}$.}
\handline{72}{481}{706}{$a\in I\Rightarrow-a\in I$.}
\handline{72}{512}{719}{$\therefore$ 考虑 $\min_{m\in I}\{|m|>0\}$，记为 $m_0$.}
\handline{72}{541}{702}{断言 $I=m_0\Z$.}
\handline{72}{572}{704}{若 $\exists n\in I$ s.t. $n\ne km_0$，$k\in\Z$，即 $m_0\nmid n$.}
\handline{72}{604}{698}{设 $d=\gcd(m_0,n)$，则 $1\le d<m_0$.}
\handline{72}{636}{705}{由 Bézout 定理 $\Rightarrow\exists u,v\in\Z$ s.t. $um_0+vn=d$.}
\handline{72}{672}{687}{$\therefore d\in I$，且 $|d|<|m_0|$，与 $m_0$ 最小性矛盾. $\therefore I=m_0\Z$.}
\handline{72}{704}{728}{$\therefore\Z$ 是主理想整环.}
\handline{72}{745}{836}{\textbf{Def.} $R$：环，$\varphi\colon\Z\to R$. \hfill $\varphi$ 是自然同态.}
\handline{72}{824}{400}{$n\mapsto\begin{cases}\underbrace{1_R+1_R+\cdots+1_R}_{n\text{ 个}},&n>0,\\0,&n=0,\\-\bigl(\underbrace{1_R+\cdots+1_R}_{-n\text{ 个}}\bigr),&n<0.\end{cases}$}
\handline{72}{884}{618}{则满足 $\varphi(m+n)=\varphi(m)+\varphi(n)$，$\varphi(mn)=\varphi(m)\varphi(n)$.}
\handline{72}{918}{815}{则 $\ker\varphi$ 是 $\Z$ 的一个理想，$\Z$ 是主理想整环，$\therefore\exists m\in\Z$ s.t. $\ker\varphi=m\Z$.}
\handline{72}{958}{803}{$m$ 称为环 $R$ 的特征.}
\end{handpage}

% ===== Line-layout page-030.tex =====
\begin{handpage}{30}
\handline[AnnotationBlue]{72}{141}{772}{$m$ 是 $\ker\varphi$ 中的最小正数.}
\handline[AnnotationBlue]{72}{168}{772}{但特征这一概念对一般的环而言没有什么用.}
\handline[AnnotationBlue]{72}{194}{495}{对域来说，若 $F$ 的特征为 $p$（$p$ 素），则可以保证}
\handline[AnnotationBlue]{575}{190}{243}{$\forall a\in F,\ pa=0$.}
\handline[AnnotationBlue]{575}{213}{244}{且当 $k\in\Z$，$p\nmid k$ 时 $ka\ne0$.}
\handline[AnnotationBlue]{72}{224}{487}{但一般的环没有这样的性质.}
\handline[AnnotationBlue]{72}{255}{751}{除非这个环具有与域相似的特性，比如这个环是幺环.}
\handline{72}{283}{778}{\textbf{Def.} $R$：幺环，$\varphi\colon\Z\to R$ 自然同态，定义同前.}
\handline{72}{311}{680}{$p$ 是使得 $\varphi(p)=0$ 的最小正整数.}
\handline[AnnotationBlue]{72}{343}{780}{但这种定义还是不够好. 比如 $\Z_6=\{\bar0,\bar1,\bar2,\bar3,\bar4,\bar5\}$.}
\handline[AnnotationBlue]{72}{375}{750}{$p=6$，但 $6$ 并不是使得 $k\cdot\bar2=0$ 最小的 $k$.}
\handline[AnnotationBlue]{72}{405}{726}{事实上 $k$ 的值依赖于元素的选取.}
\handline[AnnotationBlue]{72}{433}{726}{比如对 $\bar2$ 来说 $k=3$，$3\cdot\bar2=0$.}
\handline[AnnotationBlue]{72}{462}{686}{对 $\bar3$ 来说 $k=2$，$2\cdot\bar3=0$.}
\handline[AnnotationBlue]{72}{490}{713}{但 $2,3$ 均不等于特征 $6$.}
\handline[AnnotationBlue]{72}{519}{728}{对于域来讲若特征为 $p$，则 $k$ 是不依赖于元素的选取的，}
\handline[AnnotationBlue]{72}{550}{712}{每个元素 $x$ 对应的 $k$ 都是特征 $p$.}
\handline[AnnotationBlue]{72}{581}{750}{出现问题的原因在于 $\Z_6$ 这个环含有零因子.}
\handline[AnnotationBlue]{72}{613}{738}{因此对它定义特征没有意义. 因此，一般情况下}
\handline[AnnotationBlue]{72}{644}{724}{当我们谈论一个环的特征时，只谈论整环的特征.}
\handline{72}{678}{845}{\textbf{Thm.} $R$：整环. 若 $\exists a\in R$，$a\ne0$，$m\in\N$，$m\ne0$，s.t. $ma=0$，则 $\exists p$ 素数 s.t. $\forall b\in R$，$pb=0$.}
\handline{72}{713}{849}{\textbf{Pf.} \circled{1} $\forall b\in R$：$0=0\cdot b=(ma)\cdot b=(a+a+\cdots+a)b=ab+\cdots+ab=a\cdot(b+\cdots+b)=a(mb)$.}
\handline{72}{745}{710}{$a\in R$，$a\ne0$，$R$ 是整环，$a\cdot(mb)=0$，$\therefore mb=0$.}
\handline{72}{777}{794}{\circled{2} 对于 $1_R\in R$，由 \circled{1}，$m\cdot1_R=0$. 令 $K=\{k\mid k\cdot1_R=0,\ k\in\N_+\}$.}
\handline{72}{808}{773}{$K$ 是 $\N_+$ 的子集，故存在最小元，记为 $p$，$\therefore p\cdot1_R=0$.}
\handline{72}{840}{781}{若 $p$ 为合数，$p=qr$，$1<q,r<p$，$\therefore qr\cdot1_R=q(r\cdot1_R)=0$.}
\handline[AnnotationBlue]{72}{872}{777}{此时再将 $q$ 视为原来的 $m$，由 $0<r<p$，$p$ 为 $K$ 中最小元 $\Rightarrow r\notin K$.}
\handline{72}{910}{783}{$\therefore r\cdot1_R\ne0$. \ann{视 $r\cdot1_R$ 为原来的 $a$，则由 \circled{1} $\Rightarrow\forall b\in R$，$q\cdot b=0$.}}
\handline[AnnotationBlue]{72}{944}{764}{再取 $b=1_R\Rightarrow q\cdot1_R=0$，与 $p$ 最小性矛盾，$\therefore p$ 为素数.}
\handline{72}{971}{774}{$\therefore\forall b\in R$，$pb=b+\cdots+b=1_Rb+\cdots+1_Rb=(1_R+\cdots+1_R)b=(p\cdot1_R)b$}
\handline{72}{991}{151}{$=0\cdot b=0$.}
\end{handpage}

% ===== Line-layout page-031.tex =====
\begin{handpage}{31}
\handtitle{72}{122}{204}{\S2.5 同态基本定理}
\handline[AnnotationBlue]{285}{118}{400}{（顺序为 第一 $\to$ 对应 $\to$ 第三 $\to$ 第二）}
\handline{72}{150}{792}{\textbf{Thm.}（第一同构定理）$\varphi\colon R\to S$ 同态，则 $R/\ker\varphi\cong\im\varphi$.}
\handline{72}{177}{729}{特别地若 $\varphi$ 是满同态，则 $R/\ker\varphi\cong S$.}
\handline{72}{207}{805}{\textbf{Pf.} 令 $I=\ker\varphi$，则 $I$ 是 $R$ 的理想.}
\handline{72}{234}{754}{定义映射 $\psi\colon R/I\to\im\varphi$}
\handline{72}{261}{592}{$a+I\mapsto\varphi(a)$}
\handline{72}{294}{786}{\circled{1} 检查 $\psi$ 是良定义的，即不依赖于代表元的选取.}
\handline{72}{323}{762}{若 $a'\in a+I$，$a'-a\in I=\ker\varphi$，$\varphi(a'-a)=0$.}
\handline{72}{351}{738}{$\therefore\psi(a'+I)=\varphi(a')=\varphi(a'-a+a)=\varphi(a'-a)+\varphi(a)=\varphi(a)=\psi(a+I)$.}
\handline{72}{381}{793}{\circled{2} $\psi$ 单：求 $\ker\psi$.}
\handline{72}{408}{656}{$\psi(a+I)=0\Longleftrightarrow\varphi(a)=0\Longleftrightarrow a\in\ker\varphi=I\Longleftrightarrow a+I=0_{R/I}$.}
\handline{72}{440}{682}{$\therefore\ker\psi=\{0_{R/I}\}$，$\therefore\psi$ 是单射.}
\handline{72}{472}{803}{\circled{3} $\psi$ 满：$\forall b\in\im\varphi$，$\exists a\in R$ s.t. $b=\varphi(a)$.}
\handline{72}{503}{680}{$\therefore\psi(a+I)=\varphi(a)=b$，$\therefore\psi$ 是满射.}
\handline{72}{537}{806}{\circled{4} 保持加法：}
\handline{72}{570}{734}{$\psi(a+I)+\psi(b+I)=\varphi(a)+\varphi(b)=\varphi(a+b)=\psi(a+b+I)=\psi((a+I)+(b+I))$.}
\handline{72}{600}{806}{\circled{5} 保持乘法：}
\handline{72}{633}{744}{$\psi(a+I)\psi(b+I)=\varphi(a)\varphi(b)=\varphi(ab)=\psi(ab+I)=\psi((a+I)(b+I))$.}
\handline{72}{669}{794}{$\therefore\psi$ 是一个同构映射，$R/\ker\varphi\cong\im\varphi$.}
\handline[NoteRed]{72}{702}{784}{（由交换图知，$\psi$ 实际上是唯一的.）}
\handline[AnnotationBlue]{72}{735}{370}{环同态的典则分解：}
\handline[AnnotationBlue]{529}{735}{350}{交换图：}
\handdiagram{139}{772}{283}{167}{
\node (r) at (0,2) {$R$};\node (s) at (4,2) {$S$};\node (q) at (0,0) {$R/\ker\varphi$};\node (im) at (4,0) {$\im\varphi$};
\draw[->,color=AnnotationBlue] (r)--node[above]{$\varphi$}(s);\draw[->>,color=AnnotationBlue] (r)--(q);\draw[->,color=AnnotationBlue] (q)--node[above]{$\eta$}(im);\draw[->,color=AnnotationBlue] (im)--(s);}
\handdiagram{588}{761}{292}{155}{
\node (r) at (0,2) {$R$};\node (s) at (3,2) {$S\ (\im\varphi)$};\node (q) at (0,0) {$R/\ker\varphi$};
\draw[->,color=AnnotationBlue] (r)--node[above]{$\varphi$}(s);\draw[->,color=AnnotationBlue] (r)--node[left]{$\pi$}(q);\draw[->,color=AnnotationBlue] (q)--node[right]{$\psi$}(s);\node at (1,1.3) {$\circlearrowright$};}
\handline[AnnotationBlue]{72}{954}{336}{即 $\varphi=\psi\circ\pi$.}
\end{handpage}

% ===== Line-layout page-032.tex =====
\begin{handpage}{32}
\handline{72}{116}{813}{设 $I,J$ 是环 $R$ 的理想，我们可以提出以下问题.}
\handline{72}{146}{150}{$\bullet$ 若 $I\subseteq J$：}
\handline{235}{143}{337}{\circled{1} $I$ 是 $J$ 的理想吗？}
\handline{72}{170}{351}{\circled{2} $R/I$，$R/J$，$J/I$ 之间有什么关系？}
\handline[AnnotationBlue]{585}{153}{255}{$\longrightarrow$ 第三同构定理.}
\handline{72}{202}{177}{$\bullet$ 若 $I,J$ 互不包含：}
\handline{260}{199}{204}{\circled{1} $I\cap J$ 是 $R$ 的理想}
\handline{72}{226}{170}{$I\cap J\subseteq I,J$}
\handline{72}{254}{204}{\circled{2} $I+J$ 是 $R$ 的理想}
\handline{72}{282}{170}{$I,J\subseteq I+J$}
\handline{276}{219}{145}{$I/(I\cap J)$，$J/(I\cap J)$}
\handline{306}{245}{150}{$(I+J)/I$，$(I+J)/J$}
\handline{72}{267}{152}{之间有什么关系？}
\handline[AnnotationBlue]{432}{228}{170}{$\longrightarrow$ 第二同构定理.}
\handline{72}{318}{827}{\textbf{Thm.}（对应定理）$\varphi\colon R\to S$ 满同态，$J=\ker\varphi$ 是 $R$ 的理想.}
\handline{72}{347}{555}{则对 $R$ 中任一包含 $J$ 的子环 $R'$，$S'=\varphi(R')$ 也是 $S$ 的子环.}
\handline{72}{377}{555}{反之，对 $S$ 的任一子环 $S'$，$R'=\varphi^{-1}(S')$ 是 $R$ 中包含 $J$ 的子环.}
\handline{72}{407}{554}{即 $R$ 中包含 $J$ 的子环与 $S$ 中的子环存在着一一对应.}
\handline{72}{438}{818}{且上述表述中将“子环”改为“理想”，结论也成立，且 $R/R'\cong S/S'$（$R',S'$ 是理想）.}
\handline{72}{472}{831}{\textbf{Pf.} \circled{1} $R'$ 是子环 $\Rightarrow S'=\varphi(R')$ 是子环.}
\handline{72}{505}{748}{$S_1,S_2\in S'$，$\exists r_1,r_2\in R'$ s.t. $S_1=\varphi(r_1)$，$S_2=\varphi(r_2)$.}
\handline{72}{538}{752}{$S_1+S_2=\varphi(r_1)+\varphi(r_2)=\varphi(r_1+r_2)\in\varphi(R')=S'$.}
\handline{72}{570}{758}{$S_1S_2=\varphi(r_1)\varphi(r_2)=\varphi(r_1r_2)\in\varphi(R')=S'$.}
\handline{72}{601}{759}{$-S_1=-\varphi(r_1)=\varphi(0)-\varphi(r_1)=\varphi(0-r_1)=\varphi(-r_1)\in\varphi(R')=S'$.}
\handline{72}{637}{783}{\circled{2} $S'$ 是子环 $\Rightarrow R'=\varphi^{-1}(S')$ 是包含 $J$ 的子环.}
\handline{72}{672}{758}{$r_1,r_2\in R'$，$\exists S_1,S_2\in S'$ s.t. $S_1=\varphi(r_1)$，$S_2=\varphi(r_2)$.}
\handline{72}{706}{758}{$S_1+S_2=\varphi(r_1+r_2)\Rightarrow r_1+r_2\in R'$，$S_1S_2=\varphi(r_1r_2)\Rightarrow r_1r_2\in R'$.}
\handline{72}{740}{763}{$-S_1=\varphi(0-r_1)=\varphi(-r_1)\Rightarrow-r_1\in R'$. $\therefore R'$ 是子环.}
\handline{72}{775}{772}{由 $S'$ 是子环 $\Rightarrow0_S\in S'\Rightarrow\varphi^{-1}(0_S)=\ker\varphi=J\subseteq\varphi^{-1}(S')=R'$.}
\handline{72}{809}{764}{$\therefore R'$ 是包含 $J$ 的子环.}
\handline[AnnotationBlue]{72}{841}{774}{将子环改为理想，结论仍然成立.}
\handline[AnnotationBlue]{72}{880}{781}{但是不可以写同理可证或类似可证.}
\handline[AnnotationBlue]{72}{921}{781}{因为由“$R'$ 是理想”推“$\varphi(R')$ 是理想”依赖于 $\varphi$ 是满同态.}
\handline[AnnotationBlue]{72}{953}{774}{下面看子环改为理想的证明.}
\end{handpage}

% ===== Line-layout page-033.tex =====
\begin{handpage}{33}
\handline{72}{94}{778}{\circled{3} $R'$ 是理想 $\Rightarrow S'=\varphi(R')$ 是理想.}
\handline{72}{123}{747}{$\forall S_1\in S'$，$\forall S_2\in S$，要证 $S_1S_2,S_2S_1\in S'$.}
\handline{72}{151}{745}{$\exists r_1\in R'$，s.t. $\varphi(r_1)=S_1$. \redann{由 $\varphi$ 是满同态，$\exists r_2\in R$ s.t. $\varphi(r_2)=S_2$.}}
\handline{72}{179}{532}{$R'$ 是理想 $\Rightarrow r_1r_2,r_2r_1\in R'$.}
\handline{72}{208}{546}{$\therefore S_1S_2=\varphi(r_1r_2)$，$S_2S_1=\varphi(r_2r_1)\in\varphi(R')=S'$.}
\handline{72}{236}{586}{\circled{4} $S'$ 是 $S$ 的理想 $\Rightarrow R'=\varphi^{-1}(S')$ 是包含 $J$ 的理想.}
\handline{72}{266}{559}{$S'$ 是理想 $\Rightarrow0_S\in S'\Rightarrow\varphi^{-1}(0_S)=\ker\varphi=J\subseteq R'$.}
\handline{72}{296}{766}{只需证 $R'$ 是 $R$ 的理想.}
\handline{72}{326}{748}{$\forall r_1\in R'$，$\forall r_2\in R$，要证 $r_1r_2\in R'$.}
\handline{72}{359}{752}{若 $r_1r_2\notin R'=\varphi^{-1}(S')$，i.e. $\varphi(r_1r_2)\notin S'$.}
\handline{72}{391}{742}{而 $\varphi(r_1r_2)=\varphi(r_1)\varphi(r_2)$，$\varphi(r_1)\in S'$，$\varphi(r_2)\in S$，而 $S'$ 是理想.}
\handline{72}{420}{724}{$\therefore\varphi(r_1r_2)\in S'$，矛盾！$\therefore r_1r_2\in R'$.}
\handline{72}{450}{551}{$\therefore R'$ 是 $R$ 中包含 $J$ 的理想.}
\handline[AnnotationBlue]{72}{482}{557}{\circled{1}\circled{2}\circled{3}\circled{4} 中只有 \circled{3} 需要满同态的条件.}
\handline{72}{515}{595}{\circled{5} $R'$ 是 $R$ 的理想，令 $S'=\varphi(R')$，$S'$ 是 $S$ 的理想.}
\handline{72}{545}{568}{则 $R/R'\cong S/S'$.}
\handline{72}{579}{574}{定义 $\psi\colon R/R'\to S/S'$}
\handline{72}{610}{466}{$a+R'\mapsto\varphi(a)+S'$.}
\handline{72}{642}{583}{$\bullet$ 良定义：若 $a_1+R'=a_2+R'$，$a_2-a_1\in R'$.}
\handline{72}{678}{480}{$\varphi(a_2)-\varphi(a_1)=\varphi(a_2-a_1)\in\varphi(R')=S'$.}
\handline{72}{713}{494}{$\therefore\varphi(a_2)+S'=\varphi(a_1)+S'$}
\handline{72}{751}{498}{i.e. $\psi(a_1+R')=\psi(a_2+R')$.}
\handline{72}{790}{616}{下面要证 $\psi$ 是同构，在此之前，先证 \redann{$\varphi^{-1}(S')=R'$}.}
\handline{72}{828}{766}{$\forall a\in R'$，$\varphi(a)\in\varphi(R')=S'$，$\therefore a\in\varphi^{-1}(S')$，$\therefore R'\subseteq\varphi^{-1}(S')$.}
\handline{72}{864}{778}{$\forall a\in\varphi^{-1}(S')$，$\varphi(a)\in S'=\varphi(R')$，$\therefore\exists b\in R'$ s.t. $\varphi(a)=\varphi(b)$.}
\handline{72}{902}{785}{$\therefore\varphi(a-b)=0\Rightarrow a-b\in\ker\varphi=\redann{J\subseteq R'}$，$\therefore a-b\in R'$.}
\handline{72}{938}{789}{$\therefore a=(a-b)+b$，$a-b,b\in R'$，$\therefore a\in R'$，$\therefore\varphi^{-1}(S')\subseteq R'$.}
\handline{72}{970}{300}{$\therefore\varphi^{-1}(S')=R'$.}
\end{handpage}

% ===== Line-layout page-034.tex =====
\begin{handpage}{34}
\handline{72}{119}{805}{$\bullet$ 单射：若 $\psi(a+R')=0_{S/S'}=0+S'$}
\handline{72}{149}{706}{则 $\varphi(a)+S'=0+S'$，$\therefore\varphi(a)\in S'$.}
\handline{72}{180}{696}{$\therefore a\in\redann{\underline{\varphi^{-1}(S')=R'}\quad\star}$.}
\handline{72}{211}{671}{$\therefore a+R'=0+R'=0_{R/R'}$，$\therefore\ker\psi=\{0_{R/R'}\}$.}
\handline{72}{242}{713}{$\therefore\psi$ 是单射.}
\handline{72}{274}{814}{$\bullet$ 满射：$\forall b+S'\in S/S'$，$b\in S$. 由 $\varphi$ 满射知 $\exists a\in R'$ s.t. $b=\varphi(a)$.}
\handline{72}{306}{716}{$\therefore\psi(a+R')=\varphi(a)+S'=b+S'$.}
\handline{72}{338}{716}{$\therefore\psi$ 是满射.}
\handline{72}{366}{829}{$\bullet$ 同态：}
\handline{72}{397}{761}{加法：$\psi(a_1+R'+a_2+R')=\psi((a_1+a_2)+R')$}
\handline{72}{428}{682}{$=\varphi(a_1+a_2)+S'=\varphi(a_1)+\varphi(a_2)+S'$}
\handline{72}{459}{671}{$=(\varphi(a_1)+S')+(\varphi(a_2)+S')$}
\handline{72}{490}{667}{$=\psi(a_1+R')+\psi(a_2+R')$.}
\handline{72}{524}{764}{乘法：$\psi((a_1+R')(a_2+R'))=\psi(a_1a_2+R')$}
\handline{72}{555}{681}{$=\varphi(a_1a_2)+S'=\varphi(a_1)\varphi(a_2)+S'$}
\handline{72}{587}{677}{$=(\varphi(a_1)+S')(\varphi(a_2)+S')$}
\handline{72}{619}{668}{$=\psi(a_1+R')\psi(a_2+R')$.}
\handline{72}{651}{777}{$\therefore\psi$ 是同态映射.}
\handline{72}{687}{789}{综上所述，$\psi$ 是 $R/R'\to S/S'$ 的同构映射.}
\handline{72}{722}{670}{$\therefore R/R'\cong S/S'$.}
\handline[NoteRed]{72}{758}{824}{在验证加法时，用到了 $(a_1+R')+(a_2+R')=(a_1+a_2)+R'$}
\handline[NoteRed]{72}{797}{544}{$(\varphi(a_1)+S')+(\varphi(a_2)+S')=(\varphi(a_1)+\varphi(a_2))+S'$.}
\handline[NoteRed]{72}{829}{815}{事实上这两步没有用到 $R',S'$ 是理想这个条件}
\handline[NoteRed]{72}{866}{827}{而是用了 $R',S'$ 作为加法子群是正规子群.}
\handline[NoteRed]{72}{901}{818}{验证乘法时，用到了 $(a_1+R')(a_2+R')=a_1a_2+R'$}
\handline[NoteRed]{72}{943}{561}{$(\varphi(a_1)+S')(\varphi(a_2)+S')=(\varphi(a_1)\varphi(a_2))+S'$.}
\handline[NoteRed]{72}{977}{823}{这两步用了 $R',S'$ 是理想这一条件.}
\end{handpage}

% ===== Line-layout page-035.tex =====
\begin{handpage}{35}
\handline{72}{97}{812}{\textbf{Thm.}（第三同构定理）$R,S$：环，$\varphi\colon R\to S$ 满同态，$J=\ker\varphi$.}
\handline{72}{126}{800}{对于每个包含于 $J$ 的 $R$ 的理想 $I$，存在唯一满同态}
\handline{72}{155}{783}{$\widehat\varphi\colon R/I\to S$ s.t. $\varphi=\widehat\varphi\circ\pi_I$，且 $\ker\widehat\varphi=J/I$.}
\handline{72}{185}{427}{于是 $(R/I)/(J/I)\cong R/J$.}
\handline[AnnotationBlue]{513}{178}{303}{即存在唯一的 $\widehat\varphi$，使下图交换.}
\handdiagram{717}{193}{211}{120}{\node(r)at(0,2){$R$};\node(s)at(3,2){$S$};\node(q)at(0,0){$R/I$};\draw[->,color=AnnotationBlue](r)--node[above]{$\varphi$}(s);\draw[->,color=AnnotationBlue](r)--node[left]{$\pi_I$}(q);\draw[->,color=AnnotationBlue](q)--node[right]{$\widehat\varphi$}(s);\node at(1,1.3){$\circlearrowright$};}
\handline{72}{238}{447}{\textbf{Pf.} 首先 $I$ 是 $J$ 的理想.}
\handline{72}{268}{373}{由于 $I$ 是 $R$ 的理想，}
\handline{72}{300}{409}{故 $\forall a\in I$，$b\in J$，$ab,ba\in I$，$\therefore I$ 是 $J$ 的理想.}
\handline{72}{330}{770}{$\therefore J/I$ 是良定义的.}
\handline{72}{362}{770}{定义 $\widehat\varphi\colon R/I\to S$}
\handline{72}{388}{670}{$a+I\mapsto\varphi(a)$.}
\handline{72}{425}{768}{下证 $\widehat\varphi$ 是满同态，且 $\varphi=\widehat\varphi\circ\pi_I$，$\ker\widehat\varphi=J/I$.}
\handline{72}{458}{787}{$\bullet$ 良定义：$a+I=a'+I$，$a-a'\in I$.}
\handline{72}{491}{691}{$\widehat\varphi(a+I)-\widehat\varphi(a'+I)=\varphi(a)-\varphi(a')=0$.}
\handline{72}{523}{791}{$\bullet$ 满：$\forall b\in S$，由 $\varphi$ 满 $\Rightarrow\exists a\in R$ s.t. $b=\varphi(a)$.}
\handline{72}{558}{719}{$\therefore b=\varphi(a)=\widehat\varphi(a+I)$.}
\handline{72}{585}{710}{$\therefore\widehat\varphi$ 是满射.}
\handline{72}{623}{794}{$\bullet$ 同态：$\widehat\varphi((a+I)(b+I))=\widehat\varphi(ab+I)=\varphi(ab)=\varphi(a)\varphi(b)=\widehat\varphi(a+I)\widehat\varphi(b+I)$}
\handline{72}{659}{678}{$\widehat\varphi((a+I)+(b+I))=\widehat\varphi(a+b+I)=\varphi(a+b)=\varphi(a)+\varphi(b)=\widehat\varphi(a+I)+\widehat\varphi(b+I)$.}
\handline{72}{694}{799}{$\bullet$ $\ker$：$\forall a+I\in\ker\widehat\varphi$.}
\handline{72}{731}{674}{$\widehat\varphi(a+I)=\varphi(a)=0\Rightarrow a\in J$，$\therefore a+I\in J/I$，$\therefore\ker\widehat\varphi\subset J/I$.}
\handline{72}{767}{672}{$\forall a+I\in J/I\Rightarrow a\in J\Rightarrow\widehat\varphi(a+I)=\varphi(a)=0$，$\therefore J/I\subset\ker\widehat\varphi$.}
\handline{72}{803}{670}{$\therefore\ker\widehat\varphi=J/I$.}
\handline{72}{842}{806}{而 $\widehat\varphi\circ\pi_I(a)=\widehat\varphi(a+I)=\varphi(a)$，$\therefore\varphi=\widehat\varphi\circ\pi_I$.}
\handline{72}{874}{799}{对 $\widehat\varphi$ 用第一同构定理 $\Rightarrow (R/I)/(J/I)\cong S\cong R/J$.}
\handline{72}{914}{820}{唯一性：设 $\widehat\varphi_1,\widehat\varphi_2$ 满足 $\varphi=\widehat\varphi_1\circ\pi_I$，$\varphi=\widehat\varphi_2\circ\pi_I$.}
\handline{72}{948}{782}{$\forall a\in R$，$\varphi(a)=\widehat\varphi_1(a+I)=\widehat\varphi_2(a+I)$，$\therefore\widehat\varphi_1=\widehat\varphi_2$.}
\handline{72}{990}{780}{$\therefore\widehat\varphi$ 是唯一的.}
\end{handpage}

% ===== Line-layout page-036.tex =====
\begin{handpage}{36}
\handline{72}{94}{300}{Def.\quad$I,J$ 是 $R$ 的理想.}
\handline{72}{123}{550}{$I+J:=\{x+y\mid x\in I,y\in J\}$\quad$I\cap J:=I\cap J$（作为集合）.}
\handline[AnnotationBlue]{72}{149}{290}{$I+J,I\cap J$ 是 $R$ 的理想.}
\handline[AnnotationBlue]{72}{179}{525}{$\forall x\in I,y\in J,z\in R$.\quad$(x+y)z=xz+yz\in I+J$.}
\handline[AnnotationBlue]{72}{210}{360}{$\forall x\in I\cap J,z\in R\quad xz\in I\cap J$.}
\handline{72}{239}{290}{Thm.（第二同构定理）.}
\handline{72}{268}{662}{$\varphi\colon I/(I\cap J)\longrightarrow(I+J)/J\quad$则 $\varphi$ 是同构，$I/(I\cap J)\cong(I+J)/J$.}
\handline{72}{295}{180}{$x+I\cap J\mapsto x+J$}
\handline{72}{325}{670}{Pf.\ $\bullet$ 良定义：$x_1+I\cap J=x_2+I\cap J\quad x_1-x_2\in I\cap J\subseteq J$.}
\handline{72}{354}{520}{$\therefore x_1+J=x_2+J\quad\therefore\varphi(x_1+I\cap J)=\varphi(x_2+I\cap J)$.}
\handline{72}{386}{340}{$\bullet$ 单射：若 $\varphi(x+I\cap J)=0$}
\handline{72}{418}{268}{则 $x+J=0\quad\therefore x\in J$}
\handline[NoteRed]{72}{446}{361}{在定义 $x+I\cap J$ 时，自动要求 $x\in I$.}
\handline{72}{480}{380}{$\therefore x\in I\cap J\quad\therefore x+I\cap J=0_{I/(I\cap J)}$.}
\handline{72}{513}{425}{$\therefore\ker\varphi=\{0_{I/(I\cap J)}\}\quad\therefore\varphi$ 是单射.}
\handline{72}{544}{420}{$\bullet$ 满射：$\forall z+J\in(I+J)/J$，$z\in I+J$}
\handline{72}{578}{300}{$\exists x\in I,y\in J\ \text{s.t.}\ z=x+y$.}
\handline{72}{612}{549}{$\therefore z+J=x+J\quad\therefore\varphi(x+I\cap J)=x+J=z+J$.}
\handline{72}{644}{198}{$\therefore\varphi$ 是满射.}
\handline{72}{674}{230}{$\bullet$ 同态：$\forall a,b\in I$.}
\handline{72}{710}{650}{$\varphi((a+I\cap J)(b+I\cap J))=\varphi(ab+I\cap J)=ab+J=(a+J)(b+J)$.}
\handline{72}{741}{300}{$=\varphi(a+I\cap J)\varphi(b+I\cap J)$}
\handline{72}{781}{670}{$\varphi((a+I\cap J)+(b+I\cap J))=\varphi(a+b+I\cap J)=a+b+J=(a+J)+(b+J)$}
\handline{72}{815}{355}{$=\varphi(a+I\cap J)+\varphi(b+I\cap J)$.}
\handline{72}{848}{210}{$\therefore\varphi$ 是同态.}
\handline{72}{882}{243}{综上，$\varphi$ 是同构.}
\end{handpage}

% ===== Line-layout page-037.tex =====
\begin{handpage}{37}
\handline{72}{87}{400}{Ex.\quad\circled{1} $\varphi\colon\Z\to\Q\quad\varphi$ 不是满同态.}
\handline{72}{114}{360}{$\Q$ 是域，故理想只有 $\{0\},\Q$.}
\handline{72}{141}{495}{$\Z$ 是主理想整环：理想有形 $m\Z$，$m\in\Z$.}
\handline{72}{167}{420}{$\Z$ 有无穷多个理想，$\Q$ 只有 $2$ 个理想.}
\handline{72}{195}{710}{$\Z$ 与 $\Q$ 理想之间不存在一一对应，结合对应定理知 $\varphi$ 不是满同态.}
\handline{72}{225}{747}{进一步在对应定理的证明中若 $\varphi$ 不是满的，则不一定把理想映为理想.}
\handline{72}{254}{454}{在这里，$\varphi(m\Z)\ne\{0\}$ 且 $\varphi(m\Z)\varsubsetneqq\Q$}
\handline{72}{283}{392}{故 $\varphi(m\Z)$ 一定不是 $\Q$ 的理想.}
\handline{72}{313}{311}{\circled{2}\quad$R=\Z$，$I=2\Z$，$J=3\Z$.}
\handline{72}{343}{390}{$I+J=\{2n+3m\mid m,n\in\Z\}=\Z=R$.}
\handline{72}{371}{192}{$I\cap J=6\Z$.}
\handline{72}{397}{217}{由第二同构定理}
\handline{72}{427}{250}{$I/(I\cap J)\cong(I+J)/J$}
\handline{72}{465}{453}{$I/(I\cap J)=2\Z/6\Z\cong\Z/3\Z=(I+J)/J$.}
\handline[AnnotationBlue]{72}{494}{550}{由 Bézout 定理可知 $p\Z+q\Z=\Z\Longleftrightarrow(p,q)=1$.}
\handline[AnnotationBlue]{72}{523}{405}{由此引出一般环上互素的概念：}
\handline{72}{558}{630}{Def.\quad$I,J$ 是 $R$ 的理想.若 $I+J=R$，则称 $I,J$ 互素.}
\handline{72}{591}{468}{若 $R$ 是幺环，则上述定义 $\Longleftrightarrow1_R\in I+J$}
\handline{72}{625}{328}{$\Longleftrightarrow1_R=x+y\ (x\in I,y\in J)$}
\handline{72}{658}{810}{Thm.（中国剩余定理.\ann{Chinese Remainder Theorem CRT}）\redann{（二元版）}.}
\handline{72}{687}{660}{设 $I,J$ 是 $R$ 的互素理想，则 $R/(I\cap J)\cong(R/I)\times(R/J)$.}
\handline{72}{721}{455}{Pf.\quad 定义 $\varphi\colon R\longrightarrow(R/I)\times(R/J)$.}
\handline{72}{760}{440}{$x=y+z\mapsto(y+I,z+J)\qquad y\in I,z\in J$.}
\handline{72}{791}{580}{$\bullet$ 良定义：$x=y_1+z_1=y_2+z_2$，$y_1,y_2\in I$，$z_1,z_2\in J$.}
\handline{72}{827}{563}{则 $y_1-y_2=z_2-z_1\in I\cap J\quad\therefore y_1-y_2\in I$，$z_1-z_2\in J$.}
\handline{72}{857}{380}{$\therefore y_1+I=y_2+I$，$z_1+J=z_2+J$}
\handline{72}{894}{270}{$\therefore\varphi$ 是良定义的.}
\end{handpage}

% ===== Line-layout page-038.tex =====
\begin{handpage}{38}
\handline{72}{106}{657}{$\bullet$ 满射：$\forall z+I\in R/I$，$z\in R$，$z=z_1+z_2\quad z_1\in I,z_2\in J$.}
\handline{72}{135}{526}{$\forall y+J\in R/J$，$y\in R$，$y=y_1+y_2\quad y_1\in I,y_2\in J$.}
\handline{72}{163}{333}{$\therefore z+I=z_2+I\quad y+J=y_1+J$.}
\handline{72}{191}{241}{取 $x=y_1+z_2\in J$.}
\handline{72}{218}{466}{则 $\varphi(x)=(z_2+I,y_1+J)=(z+I,y+J)$.}
\handline{72}{249}{478}{易验证，$\varphi$ 是一个同态.$\therefore\varphi$ 是满同态.}
\handline{72}{278}{200}{若 $x\in\ker\varphi$.}
\handline{72}{309}{625}{$\varphi(x)=(x+I,x+J)=(0,0)\quad\therefore x\in I,x\in J\quad\therefore x\in I\cap J$.}
\handline{72}{335}{210}{$\therefore\ker\varphi\subseteq I\cap J$}
\handline{72}{363}{398}{若 $x\in I\cap J$，则 $x+I=0$，$x+J=0$}
\handline{72}{392}{572}{$\therefore\varphi(x)=(0,0)\quad\therefore x\in\ker\varphi\quad\therefore I\cap J\subseteq\ker\varphi$}
\handline{72}{422}{223}{$\therefore\ker\varphi=I\cap J$.}
\handline{72}{454}{367}{$\therefore R/\ker\varphi\cong(R/I)\times(R/J)$}
\handline{72}{487}{365}{即 $R/(I\cap J)\cong(R/I)\times(R/J)$.}
\handline{72}{539}{735}{Def.\quad 设 $I,J$ 是 $R$ 的理想，$IJ:=\{\sum_{i=1}^mx_iy_i\mid x_i\in I,y_i\in J,1\le i\le m,m\in\omega\}$.}
\handline[AnnotationBlue]{72}{568}{375}{为什么用 $IJ=\{xy\mid x\in I,y\in J\}$ 来定义？}
\handline[AnnotationBlue]{72}{598}{705}{若采用这种定义方式，则 $x_1y_1+x_2y_2$ 不一定可以写成 $xy$（$x\in I,y\in J$）的形式.}
\handline[NoteRed]{72}{628}{707}{显然 $IJ\subset I\cap J$，例如 $R=\Z$，$I=J=2\Z$，则 $IJ=4\Z\subseteq I\cap J=2\Z$}
\handline[NoteRed]{72}{660}{347}{下面看何时会有 $IJ=I\cap J$}
\handline{72}{694}{538}{Lemma.\quad 设 $R$ 是交换幺环，$I_1,I_2,\ldots,I_n$ 两两互素}
\handline{72}{725}{614}{则 (1) $I_1\cdots I_n=I_1\cap\cdots\cap I_n$\quad(2) $I_1\cap I_2\cdots\cap I_{n-1}$ 与 $I_n$ 互素.}
\handline{72}{754}{591}{Pf.\quad$n=2$ 时.$I_1I_2\subseteq I_1\cap I_2$ 显然，下证 $I_1\cap I_2\subseteq I_1I_2$.}
\handline{72}{785}{527}{由 $I_1,I_2$ 互素 $\Rightarrow1_R\in I_1+I_2\quad\therefore1_R=y+z$，$y\in I_1,z\in I_2$.}
\handline{72}{826}{692}{$\forall x\in I_1\cap I_2\quad x=x\cdot1_R=x(y+z)=xy+xz=yx+xz\in I_1I_2$}
\handline{72}{858}{235}{$\therefore I_1\cap I_2\subseteq I_1I_2$}
\handline{72}{896}{444}{$\therefore(1)\ I_1I_2=I_1\cap I_2\quad(2)\ I_1$ 与 $I_2$ 互素.}
\handline[NoteRed]{72}{928}{338}{注：由 $I,J$ 互素可推 $IJ=I\cap J$}
\handline[NoteRed]{72}{956}{303}{但 $IJ=I\cap J\not\Rightarrow I,J$ 互素.}
\end{handpage}

% ===== Line-layout page-039.tex =====
\begin{handpage}{39}
\handline{72}{86}{390}{$n>2$ 时，假设 $n-1$ 时成立，下证 $n$ 时成立.}
\handline{72}{113}{351}{$1\le k\le n-1$：$I_k$ 与 $I_n$ 互素\quad$\therefore I_n+I_k=R$}
\handline{72}{144}{566}{$\therefore R=R\cdots R=\prod_{k=1}^{n-1}(I_n+I_k)=I_n+I_nf(I_1,\ldots,I_{n-1})+\prod_{k=1}^{n-1}I_k$.}
\handline{72}{199}{300}{其中 $f$ 是一个 $n-1$ 元多项式.}
\handline{72}{229}{358}{$\therefore\prod_{k=1}^{n-1}(I_n+I_k)\subseteq I_n+I_1I_2\cdots I_{n-1}$}
\handline{72}{258}{386}{显然有 $I_n+I_1\cdots I_{n-1}\subseteq R=\prod_{k=1}^{n-1}(I_n+I_k)$}
\handline{72}{288}{224}{$\therefore R=I_n+I_1\cdots I_{n-1}$}
\handline{72}{317}{252}{$\therefore I_n$ 与 $I_1\cdots I_{n-1}$ 互素.}
\handline{72}{347}{690}{由归纳假设 $\Rightarrow I_1\cdots I_{n-1}=I_1\cap\cdots\cap I_{n-1}\quad\therefore I_n$ 与 $I_1\cap\cdots\cap I_{n-1}$ 互素.}
\handline{72}{379}{651}{由 $n=2$ 的情形 $\Rightarrow I_n(I_1\cdots I_{n-1})=I_n\cap(I_1\cdots I_{n-1})=I_1\cap\cdots\cap I_n$.}
\handline{72}{410}{198}{Thm.（CRT 多元版）}
\handline{72}{438}{520}{$R$：交换幺环，$I_1,\ldots,I_n$：$R$ 的两两互素的理想.}
\handline{72}{469}{416}{则 $R/(I_1\cap\cdots\cap I_n)\cong(R/I_1)\times\cdots\times(R/I_n)$.}
\handline{72}{501}{229}{Pf.\quad$n=2$ 时已知.}
\handline{72}{532}{355}{假设 $n-1$ 时成立，下证 $n$ 时成立.}
\handline{72}{565}{336}{由引理 $I_n$ 与 $I_1\cap\cdots\cap I_{n-1}$ 互素}
\handline{72}{598}{725}{$\therefore R/(I_n\cap(I_1\cap\cdots\cap I_{n-1}))\cong(R/(I_1\cap\cdots\cap I_{n-1}))\times(R/I_n)\cong(R/I_1)\times\cdots\times(R/I_n)$.}
\handline[AnnotationBlue]{72}{632}{712}{Ex.\quad$R=\Z$，$I_1=3\Z$，$I_2=5\Z$，$I_3=7\Z$.\quad$I_1,I_2,I_3$ 两两互素.}
\handline[AnnotationBlue]{72}{663}{453}{则 $R/(I_1\cap I_2\cap I_3)\cong(R/I_1)\times(R/I_2)\times(R/I_3)$.}
\handline[AnnotationBlue]{72}{695}{463}{即 $\Z/105\Z\cong(\Z/3\Z)\times(\Z/5\Z)\times(\Z/7\Z)$.}
\handline[AnnotationBlue]{72}{744}{560}{对于同余方程组 $\begin{cases}n\equiv r_1\pmod3\\n\equiv r_2\pmod5\\n\equiv r_3\pmod7\end{cases}\Rightarrow n\equiv{}?\pmod{105}$.}
\handline[AnnotationBlue]{72}{816}{476}{求基解 $\begin{cases}e_1\equiv1\ (3)\\e_1\equiv0\ (5)\\e_1\equiv0\ (7)\end{cases}\quad\begin{cases}e_2\equiv0\ (3)\\e_2\equiv1\ (5)\\e_2\equiv0\ (7)\end{cases}\quad\begin{cases}e_3\equiv0\ (3)\\e_3\equiv0\ (5)\\e_3\equiv1\ (7)\end{cases}$.}
\handline[AnnotationBlue]{72}{868}{558}{$\Rightarrow e_1\equiv70\ (105)\quad e_2\equiv21\ (105)\quad e_3\equiv15\ (105)$.}
\handline[AnnotationBlue]{72}{904}{448}{$\therefore n\equiv70r_1+21r_2+15r_3\pmod{105}$.}
\end{handpage}

% ===== Line-layout page-040.tex =====
\begin{handpage}{40}
\handtitle{72}{93}{235}{§2.6 整环与整除性.}
\handline{72}{119}{715}{环 $\supseteq$ 幺环 $\supseteq$ 交换幺环 $\supseteq$ 整环 $\supseteq$ 整闭整环 $\supseteq\mathrm{GCD}\supseteq\mathrm{UFD}\supseteq\mathrm{PID}\supseteq\mathrm{ED}\supseteq$ 域}
\handline{72}{147}{741}{此外还有一些没有明显包含关系的环：Bézout 整环，Dedekind 整环，Noether 环.}
\handline{72}{176}{100}{一、整环}
\handline{72}{207}{774}{Def.\quad$R$：交换幺环，$1_R\ne0_R$，若 $R$ 没有零因子则称 $R$ 是一个整环.\ann{（Integral Domain）}}
\handline{72}{235}{795}{Def.\quad 对于整环 $R$.若 $m\in\omega$，$a\in R$，$a\ne0$，$ma=0$，则 $\exists$ 素数 $p$ s.t. $\forall x\in R$，$px=0$.}
\handline{72}{261}{610}{记 $\operatorname{char}R=p$ 称为 $R$ 的特征.$a$ 称为周期元，$m$ 称为周期}
\handline{72}{289}{390}{若 $R$ 没有周期元则记 $\operatorname{char}R=0$.}
\handline{72}{319}{510}{Ex.\quad\circled{1} $\Z$，$\C[x]$，域，$\F_p=\Z/p\Z$（$p$ 素）都是整环}
\handline{72}{347}{370}{\circled{2} $\Z_m=\Z/m\Z$ 不是整环（$m$ 是合数）}
\handline{72}{377}{460}{\circled{3} $\Z[i]\cong\Z[x]/(x^2+1)$ 是整环，称为高斯整环}
\handline{72}{404}{399}{\circled{4} $\C(x,y)/(y^2-(4x^3-g_2x-g_3))$ 是整环}
\handline{72}{461}{741}{\circled{5} 设 $R=\{\sum_{n\ge0}a_nx^n\mid a_n\in\R$，有正收敛半径$\}$\ann{（即 $0$ 附近的实解析函数）}是整环.}
\handline{72}{491}{533}{\circled{6} $R=C^\infty(U)$，其中 $U$ 是 $\R^n$ 中的开集，$R$ 不是整环.}
\handline{72}{523}{536}{取 $f,g\in R\ \text{s.t.}\ \operatorname{supp}f\cap\operatorname{supp}g=\varnothing$ 则 $f\cdot g=0$.}
\handline{72}{555}{247}{\circled{7} $\R[x]/(x^2-1)$ 不是整环}
\handline{72}{587}{651}{$\R[x]/(x^2-1)=\R[x]/((x+1)(x-1))\cong(\R[x]/(x+1))\times(\R[x]/(x-1))\cong\R\times\R$.}
\handline[AnnotationBlue]{72}{620}{615}{形如 $R=R_1\times R_2$ 的环一定不是整环（$0_{R_1}\ne1_{R_1}$，$0_{R_2}\ne1_{R_2}$）.}
\handline[AnnotationBlue]{72}{655}{475}{$\because(0_{R_1},1_{R_2})\cdot(1_{R_1},0_{R_2})=(0_{R_1},0_{R_2})=0_R$.}
\handline{72}{684}{300}{而 $\R[x]/(x^2+1)\cong\C$ 是整环.}
\handline{72}{718}{785}{Def.\quad 设 $R$ 是非零交换幺环，$I$ 是 $R$ 的理想.若 $R/I$ 是整环，则称 $I$ 是素理想.}
\handline{72}{750}{650}{Thm.\quad$I$ 是素理想 $\Longleftrightarrow$ “$I\ne R$ 且 $ab\in I$ 则 $a\in I$ 或 $b\in I$”.}
\handline{72}{783}{640}{Pf.\quad“$\Rightarrow$” $I$ 是素理想.$\therefore R/I$ 是整环.$ab\in I\Rightarrow ab+I=0_{R/I}$.}
\handline{72}{821}{765}{而 $ab+I=(a+I)(b+I)=0_{R/I}\quad\therefore a+I=0_{R/I}$ 或 $b+I=0_{R/I}$，即 $a\in I$ 或 $b\in I$.}
\handline{72}{850}{310}{“$\Leftarrow$” $I\ne R\Rightarrow R/I$ 非零.}
\handline{72}{895}{780}{$ab\in I\Longleftrightarrow ab+I=0\Longleftrightarrow(a+I)(b+I)=0\quad\Rightarrow a\in I$ 或 $b\in I\Longleftrightarrow a+I=0$ 或 $b+I=0$}
\handline{72}{925}{367}{$\therefore R/I$ 是整环.$\therefore I$ 是素理想.}
\handline[AnnotationBlue]{72}{957}{388}{特别地，若 $R$ 是整环，则 $\{0_R\}$ 是素理想.}
\end{handpage}

% ===== Line-layout page-041.tex =====
\begin{handpage}{41}
\handline{72}{124}{742}{Def.\quad$R$：交换幺环，$m$ 是理想.若 $R/m$ 是一个域，则称 $m$ 是 $R$ 的一个极大理想.}
\handline[AnnotationBlue]{72}{151}{400}{域是整环，故极大理想一定是素理想.}
\handline{72}{180}{689}{Thm.\quad$m$ 是极大理想 $\Longleftrightarrow m\ne R$ 且任何包含 $m$ 的理想只能是 $m$ 或 $R$.}
\handline{72}{209}{427}{Pf.\quad“$\Rightarrow$” $m$ 是极大理想.$\therefore R/m$ 是域}
\handline{72}{235}{328}{$\therefore R/m$ 的理想只有 $0$ 和 $R/m$}
\handline{72}{261}{392}{令 $\varphi\colon R\longrightarrow R/m\qquad\varphi$ 是满同态.}
\handline{72}{289}{136}{$a\mapsto a+m$}
\handline{72}{317}{188}{则 $\ker\varphi=m$}
\handline{72}{348}{419}{$\therefore$ 由对应定理.$R$ 中包含 $\ker\varphi=m$ 的理想，}
\handline{72}{374}{282}{与 $R/m$ 中的理想一一对应.}
\handline{72}{405}{599}{$\therefore R$ 中包含 $m$ 的理想只有 $\varphi^{-1}(R/m)=R$ 和 $\varphi^{-1}(0)=m$}
\handline{72}{435}{621}{“$\Leftarrow$” 若 $R/m$ 不是域.$\therefore\exists a\in R/m\ \text{s.t.}\ a\ne0$ 且 $a$ 不可逆.}
\handline{242}{467}{402}{考虑 $J=(a)\lhd R/m$ 是由 $a$ 生成的主理想.}
\handline{245}{497}{625}{首先，$J\ne0$.其次由于 $a$ 不可逆.$\therefore1_R+R/m\notin J\quad\therefore J\ne R/m$.}
\handline[AnnotationBlue]{72}{464}{130}{非域的环}
\handline[AnnotationBlue]{72}{497}{150}{一定有非平凡理想}
\handdiagram{40}{447}{20}{50}{\draw (1,3)--(.7,3)..controls(.2,3)and(.6,1.6)..(0,1.5)..controls(.6,1.4)and(.2,0)..(.7,0)--(1,0);}
\handline{72}{527}{571}{$\therefore$ 由对应定理，存在 $R$ 中包含 $m$ 的理想 $\widetilde J$ 与 $(a)$ 对应}
\handline{72}{558}{293}{且 $\widetilde J\ne m$，$\widetilde J\ne R$，矛盾！}
\handline{72}{590}{210}{$\therefore R/m$ 是域.}
\handline{72}{623}{833}{Ex.\quad$R=\Z[x]$\quad 令 $I=(x)=xR=\{f(x)=a_1x+a_2x^2+\cdots+a_nx^n\mid n\in\omega,a_i\in\Z\}$}
\handline{72}{654}{377}{则 $R/I=\Z$ 是整环但不是域}
\handline{72}{687}{511}{$\therefore I=(x)$ 是 $R$ 的素理想但不是极大理想.}
\handline{72}{723}{580}{令 $m_p=(x,p)$ 是由 $x$ 和 $p$ 生成的理想，其中 $p$ 为素数.}
\handline{72}{759}{542}{则 $R/m_p\cong(R/I)/(m_p/I)\cong\Z/p\Z=\F_p$ 是一个域}
\handline{72}{791}{377}{$\therefore m_p=(x,p)$ 是 $R$ 的极大理想.}
\end{handpage}

% ===== Line-layout page-042.tex =====
\begin{handpage}{42}
\handline{72}{67}{242}{已知 $\Z\subseteq\Q$，而 $\Q$ 是域.}
\handline{72}{96}{520}{于是可以问，对于一般的环 $R$，能否找到包含它的域？}
\handline{72}{128}{660}{Def.\quad$R$：整环，$S=R^*=R\setminus\{0\}$，则 \circled{1} $a,b\in S\Rightarrow a\cdot b\in S$\quad\circled{2} $1_R\in S$}
\handline[AnnotationBlue]{72}{157}{759}{（对于一般的幺环 $R$，若 $\exists S\subseteq R$（$S$ 不一定是 $R^*$）满足 \circled{1}\circled{2}，则称 $S$ 是 $R$ 的一个乘法系统）}
\handline{72}{184}{626}{在 $R\times S$ 上定义关系“$\sim$”\qquad$(r_1,s_1)\sim(r_2,s_2)\Longleftrightarrow r_1s_2=r_2s_1$}
\handline{72}{213}{280}{易验证它也是一个等价关系}
\handline{72}{242}{316}{令 $F=R\times S/\sim$ 将 $(r,s)$ 记为 $\frac rs$}
\handline{72}{270}{746}{$+:\ \frac{r_1}{s_1}+\frac{r_2}{s_2}:=\frac{r_1s_2+r_2s_1}{s_1s_2}\qquad\cdot:\ \frac{r_1}{s_1}\cdot\frac{r_2}{s_2}:=\frac{r_1r_2}{s_1s_2}\qquad1_F:=\frac{1_R}{1_R}\quad0_F:=\frac{0_R}{1_R}$}
\handline{72}{302}{522}{易验证 $(F,+,\cdot,1_F,0_F)$ 是一个域，称为 $R$ 的分式域}
\handline[AnnotationBlue]{72}{335}{790}{Remark：\circled{1} 对于不是整环的 $R$.则 $S$ 上的“$\sim$”不能这样定义.\sout{因为有零因子}，传递性证不出来}
\handline[AnnotationBlue]{72}{363}{819}{\circled{2} 事实上对于“$+$”和“$\cdot$”还要检验良定义：$1^\circ$ 与代表元无关；$2^\circ$ $s_1\in S,s_2\in S\Rightarrow s_1s_2\in S$}
\handline[AnnotationBlue]{72}{405}{580}{\circled{3} $S$ 除了取成 $R^*$ 还可以取成 $\begin{cases}S=R/p\quad(p\text{ 是素理想})\\S=\{1,a,a^2,\ldots\}\quad(a\ne0)\end{cases}$}
\handline[AnnotationBlue]{72}{447}{477}{同样按上述步骤作得到一个域，称为商域.}
\handline{72}{481}{532}{Thm.\ann{（分式域的泛性质）}\redann{分式域是包含 $R$ 的最小域}}
\handline{72}{511}{687}{$R$ 是整环，$F$ 是 $R$ 的分式域，若有另一个域 $F'$，及单的环同态 $\varphi\colon R\to F'$}
\handline{72}{541}{820}{则存在唯一域同态 $\psi\colon F\to F'\ \text{s.t.}\ \varphi=\psi\circ i_R$，其中 $i_R\colon R\to F'$ 为嵌入映射}
\handline{753}{576}{94}{$r\mapsto\frac r{1_R}$}
\handline{72}{586}{210}{Pf.\quad$\psi\colon F\to F'$}
\handline{72}{613}{116}{$\frac rs\mapsto\frac{\varphi(r)}{\varphi(s)}$}
\handline{72}{641}{727}{\circled{1} 良定义：$1^\circ$ 由 $\varphi$ 是单同态.$\therefore\forall s\ne0_F$ 有 $\varphi(s)\ne0_{F'}\quad\therefore\frac{\varphi(r)}{\varphi(s)}$ 有意义}
\handline{72}{679}{764}{$2^\circ$ 若 $\frac{r_1}{s_1}=\frac{r_2}{s_2}$ 则 $\psi(\frac{r_1}{s_1})-\psi(\frac{r_2}{s_2})=\frac{\varphi(r_1)}{\varphi(s_1)}-\frac{\varphi(r_2)}{\varphi(s_2)}=\frac{\varphi(r_1s_2-r_2s_1)}{\varphi(s_1s_2)}=0_{F'}$}
\handline{72}{707}{278}{$\therefore\psi$ 与代表元选取无关}
\handline{72}{747}{856}{\circled{2} 同态：$1^\circ\psi(\frac{r_1}{s_1}+\frac{r_2}{s_2})=\psi(\frac{r_1s_2+r_2s_1}{s_1s_2})=\frac{\varphi(r_1s_2+r_2s_1)}{\varphi(s_1s_2)}=\frac{\varphi(r_1)\varphi(s_2)+\varphi(r_2)\varphi(s_2)}{\varphi(s_1)\varphi(s_2)}=\frac{\varphi(r_1)}{\varphi(s_1)}+\frac{\varphi(r_2)}{\varphi(s_2)}$}
\handline{72}{783}{856}{$2^\circ\psi(\frac{r_1}{s_1}\cdot\frac{r_2}{s_2})=\psi(\frac{r_1r_2}{s_1s_2})=\frac{\varphi(r_1r_2)}{\varphi(s_1s_2)}=\frac{\varphi(r_1)\varphi(r_2)}{\varphi(s_1)\varphi(s_2)}=\psi(\frac{r_1}{s_1})\cdot\psi(\frac{r_2}{s_2})\qquad\qquad=\psi(\frac{r_1}{s_1})+\psi(\frac{r_2}{s_2})$}
\handline{72}{817}{766}{$3^\circ\quad\psi(0_F)=\psi(\frac{0_R}{1_R})=\frac{\varphi(0_R)}{\varphi(1_R)}=0_{F'}\qquad4^\circ\quad\psi(1_F)=\psi(\frac{1_R}{1_R})=\frac{\varphi(1_R)}{\varphi(1_R)}=1_{F'}$}
\handline{72}{850}{546}{\circled{3} 唯一性：设 $\eta\colon F\to F'$ 也是域同态，且 $\varphi=\eta\circ i_R$}
\handline{72}{888}{714}{$\forall\frac rs\in F$，$\psi(\frac rs)=\frac{\varphi(r)}{\varphi(s)}=\frac{\eta\circ i_R(r)}{\eta\circ i_R(s)}=\frac{\eta(r/1_R)}{\eta(s/1_R)}=\eta(\frac r{1_R})\cdot\eta(\frac{1_R}s)=\eta(\frac rs)$.}
\handline{72}{935}{100}{$\therefore\eta=\psi$.}
\end{handpage}

% ===== Line-layout page-043.tex =====
\begin{handpage}{43}
\handline{72}{105}{400}{\textbf{Def.} $R$：整环，$F$ 是 $R$ 的分式域.}
\handline{72}{132}{521}{\circled{1} $\bar R=\{\alpha\in F\mid\alpha\text{ 满足 }f(\alpha)=0,\ f\in R[x]\text{ 且首一}\}$}
\handline{72}{158}{344}{称为 $R$ 在 $F$ 中的整闭包}
\handline{72}{188}{400}{\circled{2} 若 $\bar R=R$，则称 $R$ 是整闭整环.}
\handline{72}{221}{323}{\textbf{Ex.} \circled{1} $R=\Z$，$F=\Q$}
\handline{72}{248}{321}{显然 $R\subseteq\bar R$. 下证 $\bar R\subseteq R$.}
\handline{72}{282}{423}{$\forall\alpha\in\bar R$，$\alpha=\frac pq$，$p,q\in\Z$，$(p,q)=1$.}
\handline{72}{313}{536}{且 $\alpha^n+a_{n-1}\alpha^{n-1}+\cdots+a_1\alpha+a_0=0$，$a_i\in\Z$（$0\le i\le n-1$）}
\handline{72}{342}{583}{所以 $(\frac pq)^n+\cdots+a_1\frac pq+a_0=0$，所以 $p^n+\cdots+a_1pq^{n-1}+a_0q^n=0$}
\handline{72}{377}{768}{所以 $q\mid p^n$ 而 $(p,q)=1$，所以 $q=1$，所以 $\alpha=p\in\Z$，所以 $\bar R\subseteq R$，所以 $R=\bar R$，即 $\Z$ 是整闭整环.}
\handline{72}{410}{469}{\circled{2} $R=\Z[\sqrt5]$，$F=\Q[\sqrt5]$. \ann{（来自代数数论）}}
\handline{72}{439}{350}{$\alpha=\frac{\sqrt5-1}2\in F$，但 $\alpha\notin R$.}
\handline{72}{470}{655}{而 $\alpha^2+\alpha-1=0$，所以 $\alpha\in\bar R$，所以 $\bar R\ne R$，所以 $\Z[\sqrt5]$ 不是整闭整环.}
\handline{72}{506}{527}{\circled{3} $R=\C[t^2,t^3]=\C[x,y]/(x^3-y^2)$. \ann{（来自代数几何）}}
\handline{72}{540}{475}{$\alpha=\frac{t^3}{t^2}=\frac yx\in F$，但 $\frac yx$ 不是多项式，所以 $\alpha\notin R$}
\handline{72}{573}{641}{取 $f(T)=T^2-x$，$f(\alpha)=\frac{y^2}{x^2}-x=\frac{y^2-x^3}{x^2}=0$，$f\in R[T]$ 且首一}
\handline{72}{603}{465}{所以 $\alpha\in\bar R$，所以 $\C[t^2,t^3]$ 不是整闭整环.}
\handtitle{72}{640}{230}{二．整除性}
\handline{72}{675}{817}{\textbf{Def.} $R$：非零交换幺环，$a,b\in R$，若 $\exists c\in R$ s.t. $a=bc$，则称 $b$ 整除 $a$，记作 $b\mid a$}
\handline[AnnotationBlue]{72}{703}{585}{从理想的观点看 $b\mid a\Longleftrightarrow a\in(b)\Longleftrightarrow(a)\subseteq(b)$}
\handline{72}{744}{650}{\textbf{Def.} $R$：整环，$a\in R$，若 $a$ 在 $R$ 中有乘法逆，则称 $a$ 是 $R$ 中的单位}
\handline{72}{774}{766}{记 $R$ 中所有单位的全体为 $U(R)$ 或 $R^\times$，则 $R^\times$ 是一个乘法 Abel 群，称为单位群}
\handline[AnnotationBlue]{72}{809}{617}{单位：unit，而单位元 $1_R$：unity 或 identity，注意区分}
\handline{72}{847}{378}{\textbf{Ex.} \circled{1} $R=\Z$，$R^\times=\{1,-1\}$}
\handline{72}{880}{385}{\circled{2} $R=\Z[i]$，$R^\times=\{\pm1,\pm i\}$}
\handline{72}{917}{405}{令 $\alpha=x+yi$，$\beta=z+wi$，$x,y,z,w\in\Z$，}
\handline{72}{950}{383}{$\alpha\beta=1\Rightarrow$ 解得 $R^\times=\{\pm1,\pm i\}$}
\end{handpage}

% ===== Line-layout page-044.tex =====
\begin{handpage}{44}
\handline{72}{98}{525}{\circled{3} $R=\Z[\sqrt3]$. 设 $\alpha=x+y\sqrt3$，$\beta=z+w\sqrt3$}
\handline{72}{138}{350}{$\alpha\beta=1\Longleftrightarrow\begin{cases}xz+3yw=1,\\xw+yz=0\end{cases}$}
\handline{72}{184}{500}{观察到 $(2+\sqrt3)(2-\sqrt3)=1\Rightarrow(2+\sqrt3)^n(2-\sqrt3)^n=1$.}
\handline{72}{215}{500}{于是 $(2+\sqrt3)^n=x_n+y_n\sqrt3$，$(2-\sqrt3)^n=x_n-y_n\sqrt3$ 都是单位.}
\handline{72}{246}{311}{且 $x_n^2-3y_n^2=1$ \ann{（Pell 方程）}}
\handline[NoteRed]{72}{276}{512}{这里可以拓展到代数数论的 Dirichlet 单位定理}
\handline{72}{312}{500}{\textbf{Def.} $R$：整环，在 $R$ 上定义关系，称为相伴.}
\handline{72}{342}{460}{若 $a,b\ne0$，$a\sim b\overset{\mathrm{def}}\Longleftrightarrow\exists u\in R^\times$ s.t. $a=ub$}
\handline[AnnotationBlue]{72}{371}{340}{易验证，相伴是一个等价关系}
\handline{72}{410}{400}{\textbf{Lemma.} $a\sim b\Longleftrightarrow a\mid b$ 且 $b\mid a$}
\handline{72}{441}{342}{\textbf{Pf.} “$\Rightarrow$” $a=ub\Rightarrow b\mid a$}
\handline{72}{473}{255}{$b=u^{-1}a\Rightarrow a\mid b$}
\handline{72}{516}{714}{“$\Leftarrow$” $\begin{aligned}a\mid b&\Rightarrow b=ac\\b\mid a&\Rightarrow a=bd\end{aligned}$ $\Rightarrow a=bd=acd$. 由 $a\ne0$ 及整环 $\Rightarrow cd=1$，所以 $a\sim b$.}
\handline[AnnotationBlue]{72}{573}{415}{理想的观点：$a\sim b\Longleftrightarrow(a)=(b)$}
\handline[AnnotationBlue]{72}{611}{482}{在相伴意义下，整除关系是一个偏序关系}
\handline[AnnotationBlue]{72}{641}{607}{\circled{1} $a\mid a$；\circled{2} $a\mid b$，$b\mid c\Rightarrow a\mid c$；\circled{3} $a\mid b$，$b\mid a\Rightarrow a\sim b$.}
\handline{72}{680}{617}{\textbf{Def.} \circled{1} $a,b\in R$，不为零，$a\mid b$，$b\nmid a$，则称 $a$ 是 $b$ 的真因子.}
\handline{72}{714}{666}{\circled{2} $a\in R$，$a\ne0$，若 $a$ 不能分解为真因子的乘积，则称 $a$ 是不可约元.}
\handline{72}{749}{573}{\circled{3} $a\in R$，$a\ne0$，若 $(a)$ 是素理想，则称 $a$ 是一个素元.}
\handline{72}{786}{440}{在初等数论中，不可约元与素元是一回事.}
\handline[NoteRed]{72}{825}{506}{问题：不可约元与素元的关系是什么？}
\end{handpage}

% ===== Line-layout page-045.tex =====
\begin{handpage}{45}
\handline{72}{88}{270}{先看两个例子.}
\handline{72}{119}{558}{\textbf{Ex.} \circled{1} $R=\C[x,y]/(y^2-x^3)=\C[t^2,t^3]$. \ann{（来自代数几何）}}
\handline{72}{149}{480}{设 $\bar x=x+(y^2-x^3)$，则 $\bar x$ 不可约，但不是素的.}
\handline{72}{178}{250}{设 $\bar x=\bar P\cdot\bar Q$}
\handline{72}{207}{762}{\ann{一方面：}$x+A(x,y)(y^2-x^3)=(P(x,y)+B(x,y)(y^2-x^3))(Q(x,y)+C(x,y)(y^2-x^3))$}
\handline{72}{237}{447}{令 $x=t^2$，$y=t^3\Rightarrow t^2=P(t^2,t^3)Q(t^2,t^3)$}
\handline{72}{267}{661}{若 $P,Q$ 均不为常数，则 $\deg(P),\deg(Q)\ge2$，$\deg(PQ)\ge4$，矛盾！}
\handline{72}{296}{604}{所以 $P$ 或 $Q$ 为非零常数，$\Rightarrow P$ 或 $Q$ 不为真因子 $\Rightarrow\bar x$ 不可约.}
\handline{72}{327}{326}{\ann{另一方面：}若 $(\bar x)$ 是素理想}
\handline{72}{357}{595}{考虑 $\bar y\cdot\bar y=\bar y^2=\bar x^3\in\bar x(x^2)\subseteq(\bar x)\Rightarrow\bar y\in(\bar x)$，所以 $\bar x\mid\bar y$}
\handline{72}{391}{687}{设 $\bar y=\bar q\cdot\bar x$，$y+A(x,y)(y^2-x^3)=(q(x,y)+B(x,y)(y^2-x^3))(x+C(x,y)(y^2-x^3))$}
\handline{72}{423}{562}{令 $x=t^2$，$y=t^3\Rightarrow t^3=t^2q(t^2,t^3)$，即 $t=q(t^2,t^3)$}
\handline{72}{455}{626}{而 $\deg q=0$ 或 $\ge2$，矛盾！所以 $(\bar x)$ 不是素理想，$\bar x$ 不是素元.}
\handline{72}{491}{408}{\circled{2} $R=\Z[\sqrt{-3}]$. \ann{（来自代数数论）}}
\handline{72}{522}{364}{$a=2$，则 $2$ 不可约，但不是素的}
\handline{72}{565}{570}{\ann{一方面，}设 $2=(a+b\sqrt{-3})(c+d\sqrt{-3})\Longleftrightarrow\begin{cases}ac-3bd=2,\\ad+bc=0\end{cases}$}
\handline{72}{625}{564}{而 $(a-b\sqrt{-3})(c-d\sqrt{-3})=ac-3bd-\sqrt{-3}(ad+bc)=2$}
\handline{72}{665}{370}{所以 $(a^2+3b^2)(c^2+3d^2)=4$}
\handline{72}{714}{420}{$\Rightarrow\begin{cases}a^2+3b^2=1,\\c^2+3d^2=4\end{cases}$ 或 $\begin{cases}a^2+3b^2=2,\\c^2+3d^2=2\end{cases}$ 或 $\begin{cases}a^2+3b^2=4,\\c^2+3d^2=1\end{cases}$}
\handline{72}{767}{432}{$a^2+3b^2=1\Rightarrow a=\pm1,b=0\Rightarrow a+b\sqrt{-3}\in R^\times$}
\handline{72}{805}{345}{$a^2+3b^2=2\Rightarrow a=\pm\sqrt2\notin R$}
\handline{72}{840}{435}{$a^2+3b^2=4\Rightarrow c=\pm1,d=0\Rightarrow c+d\sqrt{-3}\in R^\times$}
\handline{72}{876}{219}{所以 $2$ 不可约}
\handline{72}{911}{520}{\ann{另一方面，}$(1+\sqrt{-3})(1-\sqrt{-3})=4=2\times2\in(2)$}
\handline{72}{950}{674}{但 $1+\sqrt{-3},1-\sqrt{-3}\notin(2)$，故 $(2)$ 不是素理想，所以 $2$ 不是素元.}
\end{handpage}

% ===== Line-layout page-046.tex =====
\begin{handpage}{46}
\handline{72}{115}{682}{\textbf{Def.} $R$ 是整环，$q\in R$，$q\ne0$，若 $(q)\subseteq(a)$ 则必有 $(a)=(q)$ 或 $(a)=(1_R)=R$}
\handline{72}{146}{360}{则称 $(q)$ 是极大主理想.}
\handline{72}{177}{405}{\textbf{Lemma.} $R$ 是整环，$q\in R$，$q\ne0$，$q\notin R^\times$}
\handline{72}{204}{361}{则 $q$ 不可约 $\Longleftrightarrow(q)$ 是极大主理想.}
\handline{72}{236}{355}{\textbf{Pf.} “$\Rightarrow$” 若 $q$ 不可约.}
\handline{72}{265}{438}{设 $(q)\subseteq(a)\Rightarrow a\mid q\Rightarrow q=ac$. 由于 $q$ 不可约}
\handline{72}{294}{245}{所以 $a\in R^\times$ 或 $c\in R^\times$，}
\handline{72}{326}{675}{若 $a\in R^\times$，则 $(a)=(1_R)=R$；若 $c\in R^\times$，则 $(a)=(q)$. 所以 $(q)$ 是极大主理想}
\handline{72}{357}{330}{“$\Leftarrow$” 若 $(q)$ 是极大主理想.}
\handline{72}{387}{471}{设 $q=ab$，则 $(q)\subseteq(a)\Rightarrow(a)=R$ 或 $(a)=(q)$.}
\handline{72}{418}{590}{若 $(a)=R$，则 $1_R\in R=(a)$，$\exists a'\in R$，s.t. $aa'=1_R$，所以 $a\in R^\times$}
\handline{72}{447}{624}{若 $(a)=(q)$，则 $a\in(q)$，$\exists a'\in R$ s.t. $a=qa'\Rightarrow a'b=1$，所以 $b\in R^\times$}
\handline{72}{480}{242}{所以 $q$ 是不可约元.}
\handline{72}{513}{392}{\textbf{Thm.} $R$：整环，则素元一定是不可约元.}
\handline{72}{545}{385}{\textbf{Pf.} 设 $p$ 是素元，$(p)$ 是素理想.}
\handline{72}{576}{410}{若 $(p)\subseteq(a)$，所以 $\exists c\in R$ s.t. $p=ac$}
\handline{72}{609}{285}{由素性 $\Rightarrow p\mid a$ 或 $p\mid c$}
\handline{72}{641}{381}{若 $p\mid a$，则 $(a)\subseteq(p)$，所以 $(a)=(p)$}
\handline{72}{676}{644}{若 $p\mid c$，则 $\exists d\in R$ s.t. $c=pd$，所以 $p=pad$，所以 $ad=1$，$a\in R^\times$}
\handline{72}{710}{216}{所以 $(a)=(1_R)=R$}
\handline{72}{746}{289}{所以 $(p)$ 是极大主理想}
\handline{72}{781}{253}{所以 $p$ 是不可约元.}
\end{handpage}

% ===== Line-layout page-047.tex =====
\begin{handpage}{47}
\handline[NoteRed]{72}{94}{438}{下面探究什么情况下不可约元也是素元}
\handline{72}{124}{706}{\textbf{Def.} \circled{1} $R$ 是整环，$a,b\in R$ 非零，若 $d\in R$，$d\mid a$ 且 $d\mid b$，则称 $d$ 是 $a,b$ 的一个公因子}
\handline{72}{151}{660}{\circled{2} 若 $d$ 是 $a,b$ 的公因子，若 $a,b$ 的另一公因子 $\forall c$，有 $c\mid d$，则称 $d$ 是 $a,b$ 的最大公因子}
\handline{72}{180}{263}{记为 $d=\gcd(a,b)$}
\handline{72}{207}{526}{\circled{3} 若 $\forall a,b\in R$，$\gcd(a,b)$ 存在，则称 $R$ 是 GCD 整环.}
\handline{72}{237}{492}{\textbf{Lemma.} 若 $d_1,d_2$ 都是 $a,b$ 的 g.c.d.，则 $d_1\sim d_2$.}
\handline{72}{265}{310}{即 g.c.d. 在相伴意义下唯一.}
\handline[AnnotationBlue]{72}{297}{754}{初等数论中，有 B\'ezout 定理，同样可以问在 GCD 环中，是否也有 B\'ezout 定理？}
\handline{72}{324}{650}{\ann{即 $d=\gcd(a,b)$. 是否 $\exists u,v\in R$，s.t. $ua+vb=d$？}\redann{$\Rightarrow(a,b)=(d)$}}
\handline{72}{356}{753}{\textbf{Def.} $R$ 是整环，若 $\forall a,b\in R$，$(a,b)$ 是 $R$ 的主理想，则称 $R$ 为 B\'ezout 整环.}
\handline{72}{386}{406}{\textbf{Prop.} B\'ezout 整环 $\Rightarrow$ GCD 环.}
\handline{72}{416}{641}{\textbf{Pf.} 设 $R$ 是 B\'ezout 整环，设 $\forall a,b\in R$，$\exists d\in R$ s.t. $(d)=(a,b)$.}
\handline{72}{446}{593}{$a\in(a,b)\Rightarrow d\mid a$，$b\in(a,b)\Rightarrow d\mid b$，所以 $d$ 是 $a,b$ 的公因子.}
\handline{72}{477}{592}{设 $c$ 为 $a,b$ 的公因子，$\Rightarrow(a)\subseteq(c)$，$(b)\subseteq(c)\Rightarrow(a,b)\subseteq(c)$}
\handline{72}{509}{483}{所以 $(d)\subseteq(c)$，所以 $c\mid d$，所以 $d=\gcd(a,b)$.}
\handline{72}{547}{410}{\textbf{Ex.} \circled{1} GCD 环 $\nRightarrow$ B\'ezout 整环.}
\handline{72}{577}{340}{$R=\C[x,y]$，则 $R$ 是 GCD 环}
\handline{72}{607}{632}{$\gcd(x,y)=1$，但 $\forall u(x,y),v(x,y)\in R$，$u(x,y)x+v(x,y)y\ne1$}
\handline{72}{637}{332}{所以 $R$ 不是 B\'ezout 整环.}
\handline{72}{669}{565}{\circled{2} $R=\Z[\sqrt{-3}]$ 既不是 GCD 环，又不是 B\'ezout 整环.}
\handline{72}{701}{520}{$1^\circ$ 取 $a=4=(1+\sqrt{-3})(1-\sqrt{-3})$，$b=2(1+\sqrt{-3})$}
\handline{72}{734}{469}{$2\nmid1+\sqrt{-3}$，$2$ 不可约 $\Rightarrow\gcd(2,1+\sqrt{-3})\sim1$}
\handline{72}{767}{240}{若 $a,b$ 有 g.c.d.}
\handline{72}{801}{446}{则 $\gcd(a,b)\sim2\gcd(2,1+\sqrt{-3})\sim2$}
\handline{72}{835}{691}{但 $1+\sqrt{-3}$ 为 $a,b$ 的一个公因子，$1+\sqrt{-3}\nmid2$，所以矛盾，故 $a,b$ 无 g.c.d..}
\handline{72}{868}{699}{$2^\circ$ 由于 $\gcd(2,1+\sqrt{-3})=1$，设 $2(x+y\sqrt{-3})+(1+\sqrt{-3})(z+w\sqrt{-3})=1$}
\handline{72}{901}{394}{$\Rightarrow2x-2y-4w=1$，奇偶矛盾！}
\handline{72}{927}{386}{所以 $\forall u,v\in R$，$2u+(1+\sqrt{-3})v\ne1$}
\handline{72}{950}{280}{故 $R$ 不是 B\'ezout 整环.}
\end{handpage}

% ===== Line-layout page-048.tex =====
\begin{handpage}{48}
\handline{72}{117}{350}{\textbf{Prop.} 设 $R$ 是 GCD 环，则}
\handline{72}{148}{370}{\circled{1} $R$ 中任意 $n$ 个元素都有 g.c.d..}
\handline{72}{177}{440}{\circled{2} $\gcd(\gcd(a,b),c)\sim\gcd(a,\gcd(b,c))$}
\handline{72}{205}{354}{\circled{3} $\gcd(ca,cb)\sim c\gcd(a,b)$}
\handline{72}{234}{514}{\circled{4} $\gcd(a,b)\sim1$，$\gcd(a,c)\sim1\Rightarrow\gcd(a,bc)\sim1$.}
\handline{72}{266}{357}{\textbf{Pf.} \circled{1} 设 $a_1,\ldots,a_n\in R$}
\handline{72}{295}{568}{令 $d_1=\gcd(a_1,a_2)$，$d_2=\gcd(d_1,a_3)$，……，$d_{n-1}=\gcd(d_{n-2},a_n)$}
\handline{72}{325}{328}{下证 $d_{n-1}=\gcd(a_1,\ldots,a_n)$}
\handline{72}{354}{418}{首先 $d=\gcd(d_{n-2},a_n)\Rightarrow d\mid d_{n-2}$，$d\mid a_n$}
\handline{72}{384}{565}{而 $d_{n-2}=\gcd(d_{n-3},a_{n-1})\Rightarrow d_{n-2}\mid d_{n-3}$，$d_{n-2}\mid a_{n-1}\Rightarrow d\mid a_{n-1}$}
\handline{72}{416}{474}{$\cdots\cdots\Rightarrow d\mid a_1$，所以 $d$ 是 $a_1,\ldots,a_n$ 的公因子.}
\handline{72}{447}{340}{设 $c$ 是 $a_1,\ldots,a_n$ 的另一公因子.}
\handline{72}{480}{560}{$c\mid a_1$，$c\mid a_2\Rightarrow c\mid d_1$，又 $c\mid a_3\Rightarrow c\mid d_2\Rightarrow\cdots\Rightarrow c\mid d_{n-1}=d$}
\handline{72}{511}{290}{所以 $d=\gcd(a_1,\ldots,a_n)$}
\handline{72}{544}{600}{\circled{2} $\gcd(\gcd(a,b),c)\sim\gcd(a,b,c)\sim\gcd(a,\gcd(b,c))$}
\handline{72}{576}{454}{\circled{3} 设 $d=\gcd(a,b)$，$e=\gcd(ca,cb)$}
\handline{72}{608}{536}{$d\mid a$，$d\mid b\Rightarrow cd\mid ca$，$cd\mid cb\Rightarrow cd\mid e$，所以 $e=cdq$}
\handline{72}{642}{564}{所以 $ca=ex=cdqx\Rightarrow dqx=a$，所以 $dq\mid a$，同理 $dq\mid b$}
\handline{72}{675}{658}{所以 $dq\mid d$，所以 $q\in R^\times$，所以 $e\sim cd$，即 $\gcd(ca,cb)\sim c\gcd(a,b)$}
\handline{72}{709}{704}{\circled{4} $\gcd(a,bc)\sim\gcd(\gcd(a,ac),bc)\sim\gcd(a,\gcd(ac,bc))$}
\handline{72}{741}{484}{$\sim\gcd(a,c\gcd(a,b))\sim\gcd(a,c)\sim1$}
\handline{72}{777}{447}{\textbf{Thm.} 若 $R$ 是 GCD 环，则不可约元也是素元.}
\handline{72}{810}{327}{\textbf{Pf.} 设 $p$ 是不可约元.}
\handline{72}{843}{350}{下证“若 $p\nmid a$，$p\nmid b$，则 $p\nmid ab$”.}
\handline{72}{876}{587}{首先 $p\nmid a\Longleftrightarrow\gcd(p,a)\sim1$，$p\nmid b\Longleftrightarrow\gcd(p,b)\sim1$}
\handline{72}{910}{363}{所以 $\gcd(p,ab)\sim1\Longleftrightarrow p\nmid ab$}
\handline{72}{947}{455}{所以若 $p\mid ab$ 则 $p\mid a$ 或 $p\mid b$，所以 $p$ 是素元.}
\end{handpage}

% ===== Line-layout page-049.tex =====
\begin{handpage}{49}
\handline{72}{126}{424}{\textbf{Def.} $R$：整环，$a\in R$，非零、非单位}
\handline{72}{154}{633}{若 $a$ 可分解为 $R$ 中有限个不可约元的乘积，且在相伴意义下分解方式唯一.}
\handline{72}{184}{365}{即 $a=p_1p_2\cdots p_s=q_1q_2\cdots q_t$}
\handline{72}{211}{420}{则 $s=t$ 且 $\exists\sigma\in S_t$ s.t. $q_i\sim p_{\sigma(i)}$}
\handline{72}{240}{407}{则称 $R$ 是一个唯一分解整环 (UFD).}
\handline{72}{275}{383}{\textbf{Ex.} \circled{1} $\Z$，$\C[x]$ 都是 UFD 的}
\handline{72}{303}{320}{\circled{2} $\Z[\sqrt{-5}]$ 不是 UFD 的.}
\handline{72}{334}{443}{显然 $2,3,1+\sqrt{-5},1-\sqrt{-5}$ 都是不可约元.}
\handline{72}{363}{419}{以 $3$ 为例，设 $3=(x+y\sqrt{-5})(z+w\sqrt{-5})$}
\handline{72}{394}{348}{取模 $9=(x^2+5y^2)(z^2+5w^2)$}
\handline{72}{423}{383}{若 $x^2+5y^2=3\Rightarrow y=0\Rightarrow x^2=3$，矛盾}
\handline{72}{452}{596}{所以 $x^2+5y^2=9$ 或 $1$，$x^2+5y^2=1\Rightarrow y=0\Rightarrow x^2=1\Rightarrow x+y\sqrt{-5}\in R^\times$}
\handline{72}{482}{389}{$x^2+5y^2=9\Rightarrow z^2+5w^2=1\Rightarrow z+w\sqrt{-5}\in R^\times$}
\handline{72}{516}{249}{所以 $3$ 是不可约元.}
\handline{72}{549}{724}{但是 $6=2\times3=(1+\sqrt{-5})(1-\sqrt{-5})$ 都是不可约元之积. 故 $\Z[\sqrt{-5}]$ 不是 UFD}
\handline[NoteRed]{72}{585}{799}{延伸：但 $\Z[\sqrt{-5}]$ 中的理想可以唯一分解为素理想之积，为 Dedekind 整环.}
\handline[NoteRed]{72}{619}{746}{令 $P_2=(2,1+\sqrt{-5})=(2,1-\sqrt{-5})$，$P_3=(3,1+\sqrt{-5})$，$P'_3=(3,1-\sqrt{-5})$ 都是素理想.}
\handline[NoteRed]{72}{653}{273}{则 $(6)=P_2^2P_3P'_3$}
\handline{72}{691}{307}{\textbf{Thm.} UFD $\Rightarrow$ GCD.}
\handline{72}{723}{404}{\textbf{Pf.} 设 $a,b\in R$，$a,b\ne0$ 且 $a,b\notin R^\times$}
\handline{72}{757}{339}{$a=p_1^{\alpha_1}\cdots p_k^{\alpha_k}$，$b=q_1^{\beta_1}\cdots q_\ell^{\beta_\ell}$}
\handline{72}{789}{636}{将 $\{p_1,\ldots,p_k,q_1,\ldots,q_\ell\}$ 在模去相伴关系后得到 $\{r_1,r_2,\ldots,r_m\}$.}
\handline{72}{823}{390}{则 $a\sim r_1^{c_1}\cdots r_m^{c_m}$，$b\sim r_1^{d_1}\cdots r_m^{d_m}$}
\handline{72}{861}{363}{则令 $e_i=\min(c_i,d_i)$，$i=1,2,\ldots,m$}
\handline{72}{894}{469}{$d=r_1^{e_1}\cdots r_m^{e_m}$. 可验证 $d=\gcd(a,b)$.}
\end{handpage}

% ===== Line-layout page-050.tex =====
\begin{handpage}{50}
\handline{72}{116}{325}{反之，\redann{GCD 环未必是 UFD 的.}}
\handline{72}{145}{280}{$R=\C[x^r\mid r\in\Q,r>0]$.}
\handline{72}{175}{590}{$R$ 中元素形如 $f=a_0+a_1x^{r_1}+\cdots+a_nx^{r_n}$，$a_i\in\C,r_j\in\Q_+$.}
\handline{72}{203}{255}{\circled{1} $R$ 不是 UFD 的.}
\handline{72}{233}{710}{考虑 $x\in R$.$x=x^{1/2}\cdot x^{1/2}=x^{1/2}\cdot x^{1/4}x^{1/4}=\cdots=x^{1/2}x^{1/4}\cdots x^{1/2^n}\cdot x^{1/2^n}=\cdots$}
\handline{72}{266}{510}{这个分解可以无限进行下去，故 $R$ 不是 UFD 的.}
\handline{72}{296}{240}{\circled{2} $R$ 是 GCD 的.}
\handline{72}{328}{155}{设 $f,g\in R$.}
\handline{72}{359}{720}{$f=a_0+a_1x^{\alpha_1}+\cdots+a_nx^{\alpha_n}$，$\alpha_i=\dfrac{p_i}{q_i}$，$p_i,q_i\in\Z_+$，$i=0,1,2,\ldots,n$}
\handline{72}{392}{730}{$g=b_0+b_1x^{\beta_1}+\cdots+b_mx^{\beta_m}$，$\beta_j=\dfrac{s_j}{r_j}$，$s_j,r_j\in\Z_+$，$j=0,1,2,\ldots,m$.}
\handline{72}{422}{605}{通分，取 $N=\operatorname{Lcm}(q_1,\ldots,q_n,r_1,\ldots,r_m)$\quad\ann{（最小公倍数）}}
\handline{72}{452}{430}{则 $\alpha_i=\dfrac{A_i}{N}$，$\beta_j=\dfrac{B_j}{N}$，$A_i,B_j\in\Z_+$.}
\handline{72}{486}{650}{$\therefore f=a_0+a_1(x^{1/N})^{A_1}+\cdots+a_n(x^{1/N})^{A_n}$，$\therefore f\in\C[x^{1/N}]$}
\handline{72}{518}{515}{同理 $g\in\C[x^{1/N}]$，而 $\C[x^{1/N}]$ 是 GCD 的}
\handline{72}{550}{510}{$\therefore d=\gcd(f,g)$ 存在，$\therefore R$ 是 GCD 环.}
\handline{72}{580}{440}{对于这个环而言，整体上看上去是很大的.}
\handline{72}{615}{770}{但是单拎出两个元素又能包含在一个看上去不那么大的，“有限”的小环里.}
\handline[NoteRed]{72}{651}{705}{下面就要问，这种“局部小、整体大”的环，能否让它整体“有限”.}
\handline[NoteRed]{72}{686}{570}{对 GCD 环，加一些什么限制，可以得到 UFD 环?}
\end{handpage}

% ===== Line-layout page-051.tex =====
\begin{handpage}{51}
\handline{72}{89}{470}{\textbf{Def.}\quad 设 $R$ 是一个整环，若 $R$ 中不存在如下形式的序列}
\handline{72}{114}{375}{$a_1,a_2,\ldots$，其中 $a_{k+1}$ 是 $a_k$ 的真因子.}
\handline{72}{140}{255}{则称 $R$ 满足因子链条件.}
\handline[AnnotationBlue]{72}{168}{620}{理想的观点：设 $I_k=(a_k)$，则 $a_{k+1}\mid a_k$，$a_k\nmid a_{k+1}\iff I_k\subsetneq I_{k+1}$.}
\handline[AnnotationBlue]{72}{195}{510}{定义 $\iff R$ 中不存在主理想升链 $I_1\subsetneq I_2\subsetneq\cdots\subsetneq I_k\cdots$}
\handline[AnnotationBlue]{72}{219}{610}{称为 ACCP 条件 (ascending chain condition on principal ideas)}
\handline[AnnotationBlue]{72}{246}{270}{中文：主理想升链条件.}
\handline{72}{274}{480}{上述定义还等价于：若 $\exists$ 主理想升链 $I_1\subseteq I_2\cdots I_n\cdots$}
\handline{72}{301}{330}{则 $\exists m\in\N$，s.t. $n>m$ 时 $I_n=I_m$.}
\handline[AnnotationBlue]{72}{332}{455}{显然，$R=\C[x^r\mid r\in\Q_+]$ 不满足 ACCP 条件.}
\handline[AnnotationBlue]{72}{365}{420}{因为 $(x)\subsetneq(x^{1/2})\subsetneq\cdots\subsetneq(x^{1/2^n})\subsetneq\cdots$}
\handline[NoteRed]{72}{389}{685}{下面证明 $\mathrm{UFD}\iff\mathrm{GCD}+\mathrm{ACCP}$.但我们会证明一个更强的结果}
\handline{72}{420}{625}{\textbf{Def.}\quad 若在整环 $R$ 中，不可约元都是素元，则称 $R$ 满足素性条件.}
\handline{72}{454}{575}{\textbf{Thm.}\quad 设 $R$ 是整环，则 $R$ 是 UFD 的 $\iff\begin{cases}R\text{ 满足 ACCP}\\R\text{ 满足素性条件}\end{cases}$}
\handline{72}{485.527}{595}{\textbf{Pf.}\quad“$\Rightarrow$”设 $a_1,a_2,\ldots$ 是一条因子链，$a_{k+1}$ 是 $a_k$ 的真因子.}
\handline{72}{514}{455}{$\therefore a_1=a_2q_1=a_3q_2q_1=\cdots=a_{N+1}q_N\cdots q_1$.\quad\ann{(*)}}
\handline{72}{549}{305}{其中 $q_1,\ldots,q_N$ 不是 $R$ 中单位.}
\handline{72}{579}{600}{另一方面 $R$ 是 UFD，$\Rightarrow a_1=p_1\cdots p_s$，$s$ 依赖于 $a_1$，与分解方式无关.}
\handline{72}{612}{575}{而由 \ann{(*)} $a_1$ 的不可约因子至少有 $N$ 个，而 $N$ 是任意的，故矛盾!}
\handline{72}{644}{210}{$\therefore R$ 满足 ACCP}
\handline{72}{679}{545}{而 $\mathrm{UFD}\Rightarrow\mathrm{GCD}\Rightarrow$ 素性条件，故 $R$ 满足素性条件.}
\handline{72}{712}{335}{“$\Leftarrow$”设 $a\in R$ 且 $a\ne0,a\notin R^\times$.}
\handline{72}{747}{655}{不妨设 $a$ 不可约，则 $a=p_1p_2$，若 $p_1,p_2$ 有一个可约，则继续分解}
\handline{72}{782}{600}{一直分解下去，一定会在有限步停止，否则与 ACCP 矛盾!}
\handline{72}{817}{760}{故 $a$ 存在不可约分解，下证唯一性.若 $a=p_1\cdots p_s=q_1\cdots q_t$，$p_i,q_j$ 不可约.}
\handline{72}{852}{770}{由素性条件 $\Rightarrow p_i,q_j$ 是素元.$p_1\mid q_1\cdots q_t$，$\therefore\exists j_1$ s.t. $p_1\mid q_{j_1}$，由 $q_{j_1}$ 不可约 $\Rightarrow p_1\sim q_{j_1}$.}
\handline{72}{885}{490}{消去 $p_1,q_{j_1}$，$\Rightarrow\exists j_2$ s.t. $p_2\mid q_{j_2}\Rightarrow p_2\sim q_{j_2}$.}
\handline{72}{916}{755}{重复上述步骤，若最后 $p_k\cdots p_s\sim1$，则 $p_k,\ldots,p_s$ 是单位，矛盾!$\therefore s\leq t$.}
\handline{72}{943}{740}{同理 $s\geq t$，故 $s=t$，且调序后 $p_i\sim q_i$，从而唯一性得证.}
\end{handpage}

% ===== Line-layout page-052.tex =====
\begin{handpage}{52}
\handline{72}{115}{730}{\textbf{Def.}\quad 设 $R$ 是整环，若 $R$ 的理想都是主理想，则称 $R$ 是主理想整环 \ann{(PID)}}
\handline{72}{146}{265}{\textbf{Ex.}\quad\circled{1} $\Z$ 是主理想整环.}
\handline{72}{171}{255}{\circled{2} $\C[x]$ 是主理想整环.}
\handline{72}{198}{840}{\redann{Non-Ex.}\ \circled{3} $\Z[x]$ 不是主理想整环.$(2,x)$ 不是主理想.设 $I=(2,x)=\{2f(x)+xg(x)\mid f,g\in\Z[x]\}$.}
\handline{72}{227}{365}{若存在 $p(x)\in\Z[x]$ s.t. $I=(p(x))$}
\handline{72}{255}{330}{显然 $I$ 中多项式的常数项为偶数}
\handline{72}{284}{430}{由 $2\in I\Rightarrow p(x)\mid2\Rightarrow p(x)=\pm1$ 或 $\pm2$.}
\handline{72}{315}{615}{若 $p(x)=\pm1$，则 $I=\Z[x]$，但 $\Z[x]$ 中包含常数项为奇数的多项式.}
\handline{72}{344}{510}{若 $p(x)=\pm2$，则 $I=(2)$，而 $x\in I\Rightarrow2\mid x$，矛盾!}
\handline{72}{378}{475}{\circled{4} $\C[x,y]$ 不是主理想整环，$(x,y)$ 不是主理想.}
\handline{72}{411}{790}{\circled{5} $\Z[\sqrt{-5}]$ 不是主理想整环，$(2,1+\sqrt{-5})$ 不是主理想.\redann{（$\Z[\sqrt{-5}]$ 也不是 UFD 的）}}
\handline{72}{444}{770}{\textbf{Thm.}\quad$\mathrm{PID}\Rightarrow\mathrm{UFD}$\quad\ann{（反之不对.$\C[x,y]$ 是 UFD 的，但 $(x,y)$ 不是主理想）}}
\handline{72}{473}{285}{\textbf{Pf.}\quad\circled{1} $\mathrm{PID}\Rightarrow\mathrm{ACCP}$.}
\handline{72}{501}{315}{设 $(a_1)\subseteq(a_2)\subseteq\cdots\subseteq(a_n)\cdots$}
\handline{72}{533}{490}{令 $\displaystyle I=\bigcup_{n\geq1}^{\infty}(a_n)$.\quad$\forall x\in I$，$\exists k\in\N$ s.t. $x\in(a_k)$.}
\handline{72}{563}{450}{$\forall y\in R$，$xy\in(a_k)\subseteq I$，$\therefore I$ 是理想.}
\handline{72}{597}{435}{由 $R$ 是 PID 的，$\therefore\exists a\in R$ s.t. $I=(a)$.}
\handline{72}{638}{615}{$a\in I$，$\therefore\exists N$ s.t. $a\in(a_N)\Rightarrow(a)\subseteq(a_N)$，而 $(a_N)\subseteq I$.}
\handline{72}{669}{600}{$\therefore I=(a_N)$，且 $n\geq N$ 时 $(a_n)=(a_N)$.于是 $R$ 满足 ACCP.}
\handline{72}{701}{255}{\circled{2} $\mathrm{PID}\Rightarrow\mathrm{GCD}$.}
\handline{72}{737}{435}{证一个更强结论，$\mathrm{PID}\Rightarrow$ Bézout 整环.}
\handline{72}{775}{590}{$\forall x,y\in R$，由 $\mathrm{PID}\Rightarrow\exists z\in R$ s.t. $(x)+(y)=(z)$.}
\handline{72}{808}{470}{$\therefore R$ 是 Bézout 整环，$\therefore R$ 是 GCD 的.}
\handline{72}{854}{390}{于是 $\mathrm{PID}\Rightarrow\begin{cases}\mathrm{ACCP}\\\mathrm{GCD}\end{cases}\iff\mathrm{UFD}$.}
\end{handpage}

% ===== Line-layout page-053.tex =====
\begin{handpage}{53}
\handline{72}{96}{460}{\textbf{Def.}\quad 设 $R$ 是一个整环，若 $\exists\nu:R\setminus\{0\}\to\N_+$ 满足}
\handline{72}{125}{555}{$\forall a\in R,b\in R,b\ne0$，$\exists q\in R,r\in R$ s.t. $a=bq+r$\quad\redann{（带余除法）}}
\handline{72}{149}{250}{其中 $r=0$ 或 $\nu(r)<\nu(b)$}
\handline{72}{175}{530}{则称 $R$ 是一个欧几里得整环 \ann{(Euclidean Domain; ED)}}
\handline{72}{202}{485}{$\nu$ 称为 $R$ 的欧氏赋值 \ann{(Euclidean valuation)}}
\handline{72}{232}{265}{\textbf{Ex.}\quad\circled{1} $R=\Z$，$\nu(n)=|n|$.}
\handline{72}{258}{280}{\circled{2} $R=\C[x]$，$\nu(f)=\deg f$.}
\handline{72}{286}{380}{\circled{3} $R=\Z[i]$，\ann{Gauss 整环}，$\nu(z)=|z|^2$.}
\handline{72}{315}{385}{设 $a=x+yi,b=z+wi\in\Z[i]$，$b\ne0$.}
\handline{72}{344}{545}{考虑 $\dfrac ab=\dfrac{x+yi}{z+wi}=\dfrac{xz+yw+(yz-xw)i}{z^2+w^2}=u+vi$，其中 $u,v\in\Q$.}
\handline{72}{374}{365}{设 $u,v$ 的四舍五入为 $\widetilde u,\widetilde v\in\Z$.}
\handline{72}{405}{385}{即 $u=\widetilde u+\varepsilon$，$v=\widetilde v+\eta$，$|\varepsilon|\leq\frac12$，$|\eta|\leq\frac12$.}
\handline{72}{436}{285}{$\therefore\dfrac ab=(\widetilde u+\widetilde vi)+(\varepsilon+\eta i)$.}
\handline{72}{466}{305}{$\therefore a=(\widetilde u+\widetilde vi)b+(\varepsilon+\eta i)b$.}
\handline{72}{496}{570}{令 $q=\widetilde u+\widetilde vi$，$r=(\varepsilon+\eta i)b$，则 $q\in\Z[i]$，$r=a-bq\in\Z[i]$.}
\handline{72}{527}{425}{且 $\nu(r)=|r|^2=(\varepsilon^2+\eta^2)|b|^2\leq\frac12|b|^2<\nu(b)$}
\handline{72}{560}{245}{$\therefore\Z[i]$ 是 ED 的.}
\handline{72}{594}{825}{\textbf{Thm.}\quad$\mathrm{ED}\Rightarrow\mathrm{PID}$.\ann{（$\Rightarrow\mathrm{UFD}\Rightarrow\mathrm{GCD}$）}\redann{（反之不对，$\Z[\frac{1+\sqrt{-19}}2]$ 是 PID，但不是 ED，1949）}}
\handline{72}{627}{640}{\textbf{Pf.}\quad 设 $I$ 是 $R$ 的理想，不妨设 $I\ne0$，故 $I$ 至少有一个非零元.}
\handline{72}{661}{580}{考虑 $\nu(I\setminus\{0\})\subseteq\N_+$，则 $\nu(I\setminus\{0\})$ 有最小元，设为 $d>0$.}
\handline{72}{695}{525}{设 $b\in I\setminus\{0\}$ s.t. $\nu(b)=d$.\quad$b\in I\Rightarrow(b)\subseteq I$.}
\handline{72}{732}{370}{$\forall a\in I$，$a=bq+r$，$\nu(r)<\nu(b)$.}
\handline{72}{766}{530}{而 $r=a-bq\in I$.由于 $\nu(b)$ 最小，故只能 $r=0$.}
\handline{72}{801}{520}{$\therefore a=bq$，$\therefore I=(b)$，$\therefore R$ 是 PID 的.}
\handline[NoteRed]{72}{835}{750}{任何一个 ED 环，都可以通过辗转相除法来求得最大公因数.}
\handline[AnnotationBlue]{72}{873}{560}{事实上，带余除法中的商 $q$ 和余数 $r$，不一定唯一.}
\handline[AnnotationBlue]{72}{910}{800}{比如在 $\Z$ 中，对 $5$ 和 $3$ 进行带余除法，$5=3+2=2\cdot3-1$.$(q,r)=(1,2)$ 或 $(2,-1)$.}
\handline{72}{946}{835}{Rhai、Jodeit 等人证明了：在 ED 中，“$q,r$ 唯一”$\iff R$ 是域或“域上一元多项式环”.}
\end{handpage}

% ===== Line-layout page-054.tex =====
\begin{handpage}{54}
\handtitle{72}{102}{250}{§2.7 多项式环}
\handline{72}{130}{740}{\textbf{Def.}\quad\circled{1} 设 $R$ 是交换幺环且 $1_R\ne0_R$.$f:\N\to R:n\mapsto a_n$ 且 $\exists d\in\N$ s.t. $n>d$ 时 $a_n=0$.}
\handline{72}{160}{440}{所有这样子的映射构成 $R[x]$，称为多项式环.}
\handline{72}{187}{415}{$+:(n\mapsto a_n)+(n\mapsto b_n):=(n\mapsto a_n+b_n)$.}
\handline{72}{214}{440}{$\textstyle\cdot:(n\mapsto a_n)\cdot(n\mapsto b_n):=n\mapsto\sum_{k=0}^na_kb_{n-k}$}
\handline{72}{242}{475}{$1_{R[x]}:=\left(\begin{smallmatrix}0\mapsto1_R\\n\mapsto0_R\quad n\geq1\end{smallmatrix}\right)$\quad$0_{R[x]}:=(n\mapsto0_R)$.}
\handline{72}{273}{570}{\circled{2} 多元多项式环：$R[x_1,x_2,\ldots,x_n]:=R[x_1,\ldots,x_{n-1}][x_n]$.}
\handline{72}{304}{560}{\circled{3} $R[[x]]:=\{f:\N\to R:n\mapsto a_n\}$ 称为形式幂级数环.}
\handline{72}{335}{350}{\textbf{Def.}\quad 设 $R$ 是交换幺环，$f(x)\in R[x]$.}
\handline{72}{362}{665}{$f(x)=a_0+a_1x+\cdots+a_nx^n$.若 $a_n\ne0$，则 $\deg f:=n$.若 $f\equiv0$，则 $\deg f:=-\infty$.}
\handline{72}{398}{560}{\textbf{Lemma.}\quad 设 $R$ 是整环，$f,g\in R[x]$，则 $\deg(fg)=\deg f+\deg g$.}
\handline[AnnotationBlue]{72}{425}{425}{（若 $R$ 不是整环，则 $\deg(fg)\leq\deg f+\deg g$）}
\handline{72}{461}{370}{\textbf{Cor.}\quad 若 $R$ 是整环，则 $R[x]$ 是整环}
\handline{72}{493}{765}{\textbf{Pf.}\quad$fg=0\Rightarrow\deg f+\deg g=-\infty$，$\therefore\deg f=-\infty$ 或 $\deg g=-\infty$，即 $f=0$ 或 $g=0$}
\handline{72}{526}{735}{\textbf{Lemma.}\quad$F$ 是域，$\deg$ 是 $R=F[x]$ 上的欧几里得赋值，从而 $F[x]$ 是 ED 的.}
\handline{72}{557}{515}{\textbf{Pf.}\quad$\deg:R\setminus\{0\}\to\N$，对于 $f\in R,g\in R\setminus\{0\}$，}
\handline{72}{590}{530}{$f=a_nx^n+\cdots+a_1x+a_0$，$g=b_mx^m+\cdots+b_1x+b_0$，$b_m\ne0$.}
\handline{72}{623}{580}{\circled{1} 若 $n<m$，则取 $q=0,r=f$.$f=qg+r$，$\deg(r)<\deg(g)$.}
\handline{72}{656}{600}{\circled{2} 若 $n\geq m$，令 $\widetilde f=f-(\frac{a_n}{b_m}x^{n-m})g=\widetilde a_{m-1}x_{m-1}+\cdots+\widetilde a_1x+\widetilde a_0$}
\handline{72}{684}{220}{化归到情形\circled{1}.}
\handline[NoteRed]{72}{728}{845}{$F$ 必须是域.若取 $\Z[x]$，$(2,x)$ 不是主理想，$\therefore\Z[x]$ 不是 PID 的，$\therefore\Z[x]$ 不是 ED 的.}
\handline[NoteRed]{72}{762}{555}{从证明中也可看出，需要 $b_m^{-1}$ 存在，于是需要 $F$ 是域.}
\handline{72}{797}{510}{\textbf{Cor.}\quad\circled{1} $F[x]$ 是 ED 的 $\Rightarrow\mathrm{PID}\Rightarrow\mathrm{UFD}\Rightarrow\mathrm{GCD}$.}
\handline{72}{824}{180}{$\mathrm{PID}\Rightarrow\text{Bézout}$.}
\handline{72}{868}{715}{\circled{2} $F[x]$ 中的不可约元是素元（即满足素性条件）.\quad$\bigl(\mathrm{UFD}\iff\begin{cases}\mathrm{ACCP}\\\text{素性条件}\end{cases}\bigr)$.}
\handline{72}{909}{650}{\circled{3} $f\in F[x]$ 不可约 $\iff(f)$ 是极大理想 $\iff F[x]/f$ 是域.}
\end{handpage}

% ===== Line-layout page-055.tex =====
\begin{handpage}{55}
\handline{72}{90}{250}{\textbf{Thm.}\quad\ann{（多项式的泛性质）}}
\handline{72}{116}{490}{设 $R,S$ 是交换幺环.$\varphi:R\to S$ 是同态.设 $u_1,\ldots,u_n\in S$.}
\handline{72}{145}{730}{则存在唯一环同态 $\widetilde\varphi:R[x_1,\ldots,x_n]\to S$ s.t. $\widetilde\varphi|_R=\varphi$，$\widetilde\varphi(x_i)=u_i$，$1\leq i\leq n$.}
\handline{72}{175}{665}{\textbf{Pf.}\quad$n=1$ 时，定义 $\widetilde\varphi(a_0+a_1x+\cdots+a_nx^n):=\varphi(a_0)+\varphi(a_1)u+\cdots+\varphi(a_n)u^n$.}
\handline{72}{200}{180}{则 $\widetilde\varphi$ 满足条件.}
\handline{72}{229}{490}{设 $n-1$ 时 $\widetilde\varphi$ 存在，记为 $\widetilde\varphi_{n-1}:R[x_1,\ldots,x_{n-1}]\to S$}
\handline{72}{259}{625}{对于 $n$ 的情况.$R[x_1,\ldots,x_n]=R[x_1,\ldots,x_{n-1}][x_n]$，记 $\widetilde R=R[x_1,\ldots,x_{n-1}]$.}
\handline{72}{290}{420}{应用 $n=1$，对于 $\widetilde\varphi_{n-1}:\widetilde R\to S$，以及 $u_n\in S$.}
\handline{72}{321}{560}{存在 $\widetilde\varphi_n:\widetilde R[x_n]\to S$ s.t. $\widetilde\varphi_n|_{\widetilde R}=\widetilde\varphi_{n-1}$，$\widetilde\varphi_n(x_n)=u_n$.}
\handline{72}{354}{185}{$\widetilde\varphi_n$ 即为所求.}
\handline{72}{382}{340}{唯一性：若 $\widetilde\varphi,\widehat\varphi$ 满足条件.}
\handline{72}{414}{445}{$\forall f\in R[x_1,\ldots,x_n]$，$f=\sum a_{i_1,\ldots,i_n}x_1^{i_1}\cdots x_n^{i_n}$.}
\handline{72}{447}{660}{$\widetilde\varphi(f)=\sum\widetilde\varphi(a_{i_1,\ldots,i_n})\widetilde\varphi(x_1^{i_1}\cdots x_n^{i_n})=\sum\varphi(a_{i_1,\ldots,i_n})u_1^{i_1}\cdots u_n^{i_n}$}
\handline{72}{479}{465}{$=\sum\widehat\varphi(a_{i_1,\ldots,i_n})\widehat\varphi(x_1^{i_1}\cdots x_n^{i_n})=\widehat\varphi(f)$.}
\handline{72}{509}{155}{$\therefore\widetilde\varphi=\widehat\varphi$.}
\handline{72}{544}{420}{\textbf{Cor.}\quad$R[x_1,\ldots,x_n]$ 与 $R[y_1,\ldots,y_n]$ 同构.}
\handline{72}{580}{630}{\textbf{Pf.}\quad 取 $S=R[y_1,\ldots,y_n]$，考虑嵌入映射 $\varphi:R\to S:a\mapsto a$.}
\handline{72}{613}{620}{于是 $\exists$ 同态 $\widetilde\varphi:R[x_1,\ldots,x_n]\to S$ s.t. $\widetilde\varphi|_R=\varphi$，$\widetilde\varphi(x_i)=y_i$.}
\handline{72}{643}{600}{再取 $S'=R[x_1,\ldots,x_n]$.嵌入映射 $\psi:R\to S':a\mapsto a$.}
\handline{72}{675}{650}{同理 $\exists$ 同态 $\widetilde\psi:R[y_1,\ldots,y_n]\to S'$ s.t. $\widetilde\psi|_R=\psi$，$\widetilde\psi(y_i)=x_i$.}
\handline{72}{711}{495}{下面考虑 $\widetilde\varphi\circ\widetilde\psi:R[y_1,\ldots,y_n]\to R[y_1,\ldots,y_n]=S$}
\handline{72}{750}{425}{对于嵌入映射 $h:R\to S:a\mapsto a$.}
\handline{72}{784}{665}{存在唯一同态 $\widetilde h:R[y_1,\ldots,y_n]\to S$ s.t. $\widetilde h|_R=h$，$\widetilde h(y_i)=y_i$.}
\handline{72}{817}{430}{显然 $\id$ 满足，由唯一性 $\widetilde h=\id$.}
\handline{72}{856}{635}{易验证 $\widetilde\varphi\circ\widetilde\psi=\widetilde h$，$\therefore\widetilde\varphi\circ\widetilde\psi=\id$\quad 同理 $\widetilde\psi\circ\widetilde\varphi=\id$}
\handline{72}{893}{470}{于是 $R[x_1,\ldots,x_n]$ 与 $R[y_1,\ldots,y_n]$ 同构.}
\end{handpage}

% ===== Line-layout page-056.tex =====
\begin{handpage}{56}
\handline{72}{118}{545}{\textbf{Def.}\quad\circled{1} $R$ 是 UFD 的，对于 $f\in R[x]$，$f(x)=a_0+a_1x+\cdots+a_nx^n$}
\handline{72}{146}{425}{$d:=\gcd_R(a_0,a_1,\ldots,a_n)$ 称为 $f$ 的最大公因子.}
\handline{72}{173}{625}{\circled{2} 若 $\gcd(f)=1$，则称 $f$ 是本原多项式.\redann{（不可约多项式一定是本原多项式）}}
\handline[AnnotationBlue]{72}{204}{525}{对于 $f\in R[x]$，设 $d=\gcd(f)$，$f=df_1$，则 $f_1$ 是本原多项式.}
\handline{72}{235}{350}{\textbf{Lemma 1}\quad$R:$ UFD 环，$F$ 是 $R$ 的分式域}
\handline{72}{264}{635}{$f\in F[x]\setminus\{0\}$，则 $\exists r\in F\setminus\{0\}$ s.t. $f=rf_1$，$f_1$ 是 $R[x]$ 中的本原多项式.}
\handline{72}{295}{530}{且上述分解在相差一个\redann{$R$ 中的单位}的意义下是唯一的.}
\handline{72}{330}{720}{\textbf{Pf.}\quad 设 $f=a_0+a_1x+\cdots+a_nx^n$，$a_i\in F$，$a_i=\dfrac{c_i}{b_i}$，$c_i\in R,b_i\in R\setminus\{0\}$.}
\handline{72}{360}{510}{$R$ 是 UFD $\Rightarrow\mathrm{GCD}\Rightarrow\mathrm{LCM}$ 存在.令 $b=\operatorname{LCM}(b_0,\ldots,b_n)$.}
\handline{72}{393}{635}{则 $bf\in R[x]$.令 $c=\gcd(bf)$，则 $f=cf_1$，$f_1$ 为本原多项式.}
\handline{72}{425}{230}{令 $r=\dfrac cb$，$\therefore f=rf_1$.}
\handline{72}{456}{640}{下证唯一性：若 $f=r_1f_1=r_2f_2$，$r_1,r_2\in F\setminus\{0\}$，$f_1,f_2$ 为本原多项式.}
\handline{72}{493}{565}{设 $r_1=\dfrac{c_1}{d_1}$，$r_2=\dfrac{c_2}{d_2}\Rightarrow d_2c_1f_1=d_1c_2f_2=g\in R[x]$.}
\handline{72}{523}{465}{设 $d=\gcd(g)$，则 $d_2c_1\sim d_1c_2\sim\gcd(g)=d$.}
\handline{72}{557}{495}{于是 $d_2c_1=d_1c_2u$，$u\in R^\times\Rightarrow r_1=r_2u$，$u\in R^\times$}
\handline{72}{588}{355}{同时 $f_2=f_1u$，$\therefore r_1\sim r_2,f_1\sim f_2$.}
\handline{72}{628}{856}{\textbf{Cor.}\quad 设 $f,g\in R[x]$，是本原的.且在 $F[x]$ 中相伴\ann{（$F$ 是 $R$ 的分式域）}，则在 $R[x]$ 中也相伴.}
\handline{72}{660}{700}{\textbf{Pf.}\quad$f,g$ 在 $F[x]$ 中相伴 $\Rightarrow f=gu$.令 $h=f=gu\in F[x]$.}
\handline{72}{696}{555}{则 $h$ 可分解为 $h=1\cdot f=g\cdot u$，$f,g$ 是本原的.}
\handline{72}{735}{790}{由唯一性知 $1$ 和 $u$ 相差一个 $R$ 中单位，$\therefore u\in R^\times$，$\therefore f,g$ 在 $R[x]$ 中也相伴.}
\end{handpage}

% ===== Line-layout page-057.tex =====
\begin{handpage}{57}
\handline{72}{99}{780}{\textbf{Lemma 2.（Gauss）}}
\handline{72}{128}{768}{本原多项式的乘积仍是本原多项式.}
\handline{72}{157}{795}{\textbf{Pf.} 设 $f,g$ 本原，$h=f\cdot g$.}
\handline{72}{185}{741}{若 $h$ 不是本原的，$\gcd(h)\ne1$，$\therefore\exists R$ 中素元 $p$，s.t. $p\mid\gcd(h)$.}
\handline{72}{216}{744}{考虑 $\bar R=R/(p)$，$(p)$ 是素理想，$\therefore\bar R$ 是整环，$\therefore\bar R[x]$ 是整环.}
\handline{72}{247}{742}{定义环同态 $\varphi\colon R[x]\to\bar R[x]$}
\handline{72}{274}{611}{$\sum a_ix^i\mapsto\sum(a_i+(p))x^i$.}
\handline{72}{308}{641}{则 $\varphi(h)=\varphi(f)\cdot\varphi(g)$.}
\handline{72}{339}{672}{而由于 $p\mid\gcd(h)$，$\therefore\varphi(h)=0$，$\therefore\varphi(f)\cdot\varphi(g)=0$.}
\handline{72}{373}{700}{由于 $\bar R[x]$ 是整环，$\therefore\varphi(f)=0$ 或 $\varphi(g)=0$.}
\handline{72}{406}{705}{$\therefore p\mid\gcd f$ 或 $p\mid\gcd g$，与 $f,g$ 都本原矛盾！}
\handline{72}{471}{818}{\textbf{Lemma 3.（Gauss）}}
\handline{72}{508}{791}{$f\in R[x]$，$\deg f>0$，$F$ 为 $R$ 的分式域.}
\handline{72}{540}{795}{若 $f$ 在 $R[x]$ 中不可约，则 $f$ 在 $F[x]$ 中不可约.}
\handline{72}{575}{812}{\textbf{Pf.} 设 $f=g_1\cdot g_2$，$g_1,g_2\in F[x]$，且 $g_1,g_2$ 不是 $F[x]$ 中的单位.}
\handline{72}{610}{741}{设 $g_1=r_1f_1$，$g_2=r_2f_2$，$r_1,r_2\in F\setminus\{0\}$，$f_1,f_2\in R[x]$ 是本原的.}
\handline{72}{645}{730}{$\therefore f=(r_1r_2)(f_1f_2)$，$f_1\cdot f_2$ 本原.}
\handline{72}{683}{760}{由于 $f$ 在 $R[x]$ 中不可约，故 $f$ 也是本原多项式.}
\handline{72}{718}{693}{$1\cdot f=(r_1r_2)\cdot(f_1f_2)$.}
\handline{72}{756}{730}{由分解的唯一性知 $f\sim f_1\cdot f_2$.}
\handline{72}{794}{716}{与 $f$ 在 $R[x]$ 中不可约矛盾.}
\handline{72}{831}{727}{$\therefore f$ 在 $F[x]$ 中不可约.}
\end{handpage}

% ===== Line-layout page-058.tex =====
\begin{handpage}{58}
\handline{72}{117}{810}{\textbf{Thm.} 若 $R$ 是 UFD 的，则 $R[x]$ 也是 UFD 的.}
\handline{72}{146}{810}{\textbf{Pf.} 设 $f\in R[x]$，则 $f=df_1$，$d=\gcd(f)$，$f_1$ 是本原的.}
\handline{72}{179}{760}{存在性：$d\in R$ 有唯一分解 $d=p_1\cdots p_s$.}
\handline{72}{204}{391}{若 $f_1$ 有分解，$f_1=f_2\cdot f_3$，则 $f_2,f_3$ 是本原的.}
\handline{72}{238}{719}{继续分解下去，由于 $\deg f_1<+\infty$，故一定会终止.}
\handline{72}{267}{716}{$f_1=q_1\cdots q_t$，其中 $q_1,\ldots,q_t$ 在 $R[x]$ 中不可约.}
\handline{72}{301}{723}{于是 $f$ 有不可约分解 $f=df_1=p_1\cdots p_t\cdot q_1\cdots q_s$.}
\handline{72}{332}{770}{唯一性：}
\handline{72}{363}{713}{\circled{1} 若 $f$ 是本原的，$\Rightarrow f=q_1\cdots q_n=q'_1\cdots q'_k$，$q_i,q'_j$ 是 $R[x]$ 中不可约多项式.}
\handline{72}{396}{709}{由引理3 $\Rightarrow q_i,q'_j$ 是 $F[x]$ 中不可约多项式.}
\handline{72}{431}{708}{由于 $F[x]$ 是 UFD 的，$\therefore f$ 在 $F[x]$ 中的分解是唯一的.}
\handline{72}{464}{707}{$\therefore n=k$，且通过调序后，可使 $q_i\sim q'_i$.}
\handline{72}{497}{706}{$\therefore f$ 在 $R[x]$ 中的分解是唯一的.}
\handline{72}{533}{716}{\circled{2} 若 $f$ 不是本原的，$f=(p_1\cdots p_s)(q_1\cdots q_n)=(p'_1\cdots p'_t)(q'_1\cdots q'_k)$.}
\handline{72}{566}{712}{由引理1 $\Rightarrow p_1\cdots p_s\sim p'_1\cdots p'_t$}
\handline{72}{598}{565}{$q_1\cdots q_n\sim q'_1\cdots q'_k$.}
\handline{72}{639}{688}{由 $R$ 是 UFD 的，及情形 \circled{1} 知分解是唯一的.}
\handline{72}{676}{706}{$\therefore R[x]$ 是 UFD 的.}
\handline{72}{718}{841}{\textbf{Cor.} 若 $R$ 是 UFD 的，则 $R[x_1,\ldots,x_n]=R[x_1,\ldots,x_{n-1}][x_n]$ 也是 UFD 的.}
\end{handpage}

% ===== Line-layout page-059.tex =====
\begin{handpage}{59}
\handtitle{72}{93}{740}{\S3 群}
\handtitle{72}{125}{741}{\S3.1 群与群作用}
\handline{72}{153}{760}{\textbf{Def.} $X$：集合. $S(X)$：$X$ 上的双射全体.}
\handline{72}{181}{685}{$\circ$：复合；$^{-1}$：逆.}
\handline{72}{211}{724}{\circled{1} $(S(X),\circ,{}^{-1},\id_X)$ 称为 $X$ 的对称群.}
\handline{72}{239}{721}{\circled{2} 若 $|X|=n$，则 $S(X)=S_n$ 称为 $n$ 元对称群.}
\handline{72}{271}{721}{\circled{3} 若 $G\subseteq S(X)$ 满足：\circled{1} $\id_X\in G$}
\handline{72}{300}{515}{\circled{2} $g_1,g_2\in G\Rightarrow g_1\circ g_2\in G$}
\handline{72}{330}{518}{\circled{3} $g\in G\Rightarrow g^{-1}\in G$.}
\handline{72}{362}{586}{则称 $G$ 是一个置换群. \ann{（置换群是对称群的子群.）}}
\handline{72}{395}{739}{\circled{4} 若 $X$ 上有某种代数或几何的结构，}
\handline{72}{423}{720}{且 $G$ 中的双射保持这些结构}
\handline{72}{453}{719}{则称 $G$ 为 $X$ 上的一个变换群.}
\handline{72}{493}{800}{\textbf{Def.} \circled{1} 设 $G$ 是一个集合，存在 $\circ\colon G\times G\to G$，$(\ )^{-1}\colon G\to G$ 及 $e\in G$}
\handline{72}{527}{692}{s.t. $1^\circ$ $\sigma\circ(\tau\circ\mu)=(\sigma\circ\tau)\circ\mu$，$2^\circ$ $\sigma\circ e=e\circ\sigma=\sigma$，$3^\circ$ $\sigma\circ\sigma^{-1}=\sigma^{-1}\circ\sigma=e$.}
\handline{72}{559}{691}{则称 $(G,\circ,e)$ 是一个群. \ann{（易知 $e,\sigma^{-1}$ 是唯一的.）}}
\handline{72}{596}{718}{\circled{2} $G$：群，$X$：集合.}
\handline{72}{629}{472}{若有 $\mu\colon G\times X\to X$，满足 $(1)\ ex=x$}
\handline{72}{662}{414}{$(g,x)\mapsto gx$，\quad $(2)\ g(hx)=(gh)x$}
\handline{72}{702}{514}{则称 $\mu$ 定义了一个 $G$ 在 $X$ 上的（左）作用.}
\handline{72}{733}{530}{类似可定义右作用}
\handline{72}{769}{472}{$\mu\colon X\times G\to X$，满足 $(1)\ xe=x$}
\handline{72}{810}{413}{$(x,g)\mapsto xg$，\quad $(2)\ (xh)g=x(hg)$.}
\handline{72}{851}{856}{\textbf{Def.} 设 $G,H$ 是两个群. 若 $\varphi\colon G\to H$ 满足 $\varphi(g_1g_2)=\varphi(g_1)\varphi(g_2)$，则称 $\varphi$ 是一个同态.}
\handline[AnnotationBlue]{72}{889}{827}{取 $g_1=g_2=e_G\Rightarrow\varphi(e_G)=\varphi(e_G)\varphi(e_G)\Rightarrow\varphi(e_G)=e_H$.}
\handline[AnnotationBlue]{72}{933}{834}{取 $g_1=g$，$g_2=g^{-1}\Rightarrow e_H=\varphi(g)\varphi(g^{-1})\Rightarrow\varphi(g^{-1})=\varphi(g)^{-1}$.}
\end{handpage}

% ===== Line-layout page-060.tex =====
\begin{handpage}{60}
\handline{72}{98}{807}{\textbf{Lemma.} 若 $\mu\colon G\times X\to X$ 定义了 $G$ 在 $X$ 上的一个左作用}
\handline{72}{128}{681}{则 $\varphi\colon G\to S(X)$}
\handline{72}{156}{597}{$g\mapsto\mu_g$，\qquad 是一个群同态.}
\handline{72}{189}{727}{反之，若给出一个群同态，则可定义 $G$ 在 $X$ 上的作用.}
\handline{72}{223}{782}{\textbf{Pf.} $\varphi(g_1g_2)=\mu_{g_1g_2}$}
\handline{72}{251}{701}{$\varphi(g_1)\varphi(g_2)=\mu_{g_1}\cdot\mu_{g_2}$}
\handline{72}{279}{682}{$\mu_{g_1g_2}(x)=(g_1g_2)(x)$，\quad $\mu_{g_1}\circ\mu_{g_2}=g_1(g_2x)$.}
\handline{72}{311}{701}{由 $\mu$ 是左作用，$\therefore\varphi(g_1g_2)=\varphi(g_1)\varphi(g_2)$.}
\handline{72}{345}{708}{反之也是然.}
\handline[AnnotationBlue]{72}{376}{787}{该引理仅适用于左作用，对于右作用有三种解决方案：}
\handline{72}{410}{145}{\circled{1} 定义反同态}
\handline{72}{439}{153}{$\varphi(g_1g_2)=\varphi(g_2)\varphi(g_1)$.}
\handline{218}{409}{216}{\circled{2} 定义反群，反乘法}
\handline{260}{439}{182}{$G^{\mathrm{op}}$，$g_1\circ_{\mathrm{op}}g_2:=g_2\cdot g_1$.}
\handline{438}{408}{401}{\circled{3} $\mu\colon X\times G\to X$}
\handline{470}{435}{377}{定义 $\mu^{-1}\colon G\times X\to X$，把右转化为左.}
\handline{72}{466}{376}{$(g,x)\mapsto\mu(x,g^{-1})=xg^{-1}$.}
\handline{72}{511}{820}{\textbf{Def.} \circled{1} 若 $(G,\circ)$ 满足交换律，则称为 Abel 群.}
\handline{72}{540}{747}{\circled{2} 若 $G$ 为有限集，则称为有限群.}
\handline{72}{581}{822}{\textbf{Prop.} \circled{1} $e$ 唯一；\circled{2} $\forall g\in G$，$g^{-1}$ 唯一；\circled{3} $(g^{-1})^{-1}=g$；\circled{4} $(gh)^{-1}=h^{-1}g^{-1}$.}
\handline[NoteRed]{72}{615}{815}{注：$(\ )^{-1}$ 是 $G\to G$ 的反同态.}
\handline[NoteRed]{72}{650}{752}{矩阵的 $(\ )^{-1}$，$(\ )^*$，$(\ )^{\mathrm T}$ 三个运算可视为矩阵乘法环到自身的反同态.}
\handline{72}{691}{817}{\textbf{Ex.} \circled{1} 域、线性空间、环、模在加法运算下构成交换群.}
\handline{72}{754}{728}{整环的单位 $R^\times$，域的非零元在乘法运算下构成群.}
\handline{72}{792}{758}{\circled{2} 有限群.}
\handline{72}{829}{737}{$|G|=1\Rightarrow G=\{e\}$.}
\handline{72}{867}{741}{$|G|=2\Rightarrow G=\{e,a\}$，$a^2=a$ 或 $e$.}
\handline{72}{903}{733}{若 $a^2=a$，则 $aa^{-1}=a^2a^{-1}\Rightarrow a=e$，$\therefore a^2=e$.}
\handline{72}{943}{135}{乘法表：}
\handline{72}{975.718}{155}{$\begin{array}{c|cc}\cdot&e&a\\\hline e&e&a\\a&a&e\end{array}$}
\end{handpage}

% ===== Line-layout page-061.tex =====
\begin{handpage}{61}
\handline{72}{112}{748}{$|G|=3$，$G=\{e,a,a^2\}$，$a^3=e$.}
\handline{72}{141}{649}{$G\cong\Z_3\cong\Z/3\Z$.}
\handline{72}{170}{751}{$|G|=4$，有两种群.}
\handline{72}{199}{677}{$G_1=\{e,a,a^2,a^3\}$，$a^4=e$，$G_1\cong\Z_4$.}
\handline{72}{227}{675}{$G_2=\{e,a,b,ab\}$，$a^2=b^2=e$，$ab=ba$，$G_2\cong\Z_2\times\Z_2$.}
\handline{72}{258}{758}{$|G|=5$，$G\cong\Z_5\cong\Z/5\Z$.}
\handline{72}{286}{762}{$|G|=6$，有两种群.}
\handline{72}{316}{682}{$G_1\cong\Z/6\Z\cong(\Z/2\Z)\times(\Z/3\Z)$ \redann{$(\mathrm{CRT})$}.}
\handline{72}{346}{690}{$G_1$ 有两种写法：$G_1=\{e,a,a^2,a^3,a^4,a^5\}$，$a^6=e$；}
\handline{72}{375}{546}{$G_1=\{e,a,b,b^2,ab,ab^2\}$，$a^2=b^3=e$，$ab=ba$.}
\handline{72}{409}{689}{$G_2\cong S_3=\{(123),(231),(312),(132),(321),(213)\}$.}
\handline[NoteRed]{72}{444}{698}{$G_1$ 交换，$G_2$ 非交换，是阶数最小的非交换群.}
\handline{72}{475}{773}{$\cdots$}
\handline{72}{512}{792}{\circled{3} 李群.}
\handline{72}{546}{768}{$V$：$\R$ 或 $\C$-线性空间.}
\handline{72}{584}{768}{$\mathrm{GL}(V)=\{\varphi\colon V\to V\mid\varphi\text{ 线性、可逆}\}$，一般线性群，General Linear Group.}
\handline{72}{619}{769}{$\mathrm{SL}(V)=\{\varphi\in\mathrm{GL}(V)\mid\det(\varphi)=1\}$，特殊线性群，Special Linear Group.}
\handline{72}{654}{775}{若 $V$ 是带内积的 $\R$-线性空间，}
\handline{72}{691}{775}{$\mathrm O(V)=\{\varphi\in\mathrm{GL}(V)\mid\langle\varphi(\alpha),\varphi(\alpha)\rangle=\langle\alpha,\alpha\rangle\}$，正交变换群.}
\handline{72}{728}{777}{$\mathrm{SO}(V)=\mathrm O(V)\cap\mathrm{SL}(V)$，特殊正交变换群.}
\handline{72}{765}{785}{内积是一种正定、对称、双线性型.}
\handline{72}{800}{789}{定义一个反称、非退化的双线性型 $\omega\colon V\times V\to\R$.}
\handline{72}{839}{789}{$(V,\omega)$ 称为辛空间.}
\handline{72}{881}{790}{$\mathrm{Sp}(V,\omega)=\{\varphi\in\mathrm{GL}(V)\mid\omega(\varphi(\alpha),\varphi(\beta))=\omega(\alpha,\beta)\}$，称为辛群.}
\handline[AnnotationBlue]{72}{921}{798}{辛，Symplectic.}
\end{handpage}

% ===== Line-layout page-062.tex =====
\begin{handpage}{62}
\handline{72}{91}{821}{\circled{4} 二面体群：正 $n$ 边形的对称群 $D_{2n}$.}
\handline{72}{122}{794}{$n=3$：反射 $3$ 个：$s,sr,sr^2$. $s^2=e$，$r^3=e$.}
\handline{72}{151}{695}{旋转 $3$ 个：$e,r,r^2$.}
\handline{72}{180}{701}{由于反射两次回到原位置.}
\handline{72}{209}{699}{$\therefore(sr)(sr)=e\Rightarrow rsr=s$.}
\handline{72}{240}{802}{$n=4$：反射 $4$ 个：$s,sr,sr^2,sr^3$. $s^2=e$，$r^4=e$，$rsr=s$.}
\handline{72}{269}{707}{旋转 $4$ 个：$e,r,r^2,r^3$.}
\handline{72}{301}{802}{一般地，设正 $n$ 边形顶点为 $e^{\frac{2\pi i}{n}k}$，$k=0,1,2,\ldots,n-1$.}
\handline{72}{332}{706}{定义反射 $s\colon z\mapsto\bar z$，$\therefore s^2=\id=e$.}
\handline{72}{361}{661}{旋转 $r\colon z\mapsto e^{\frac{2\pi i}{n}}z$，$r^n=\id=e$.}
\handline{72}{396}{690}{$rsr(z)=rs(e^{\frac{2\pi i}{n}}z)=r(e^{-\frac{2\pi i}{n}}\bar z)=\bar z=s(z)$，$\therefore rsr=s$.}
\handline{72}{433}{708}{设 $r,s\in G$.}
\handline{72}{464}{731}{考虑 $s^{k_1}r^{l_1}s^{k_2}r^{l_2}\cdots s^{k_m}r^{l_m}$，$k_i,l_i\in\Z$.}
\handline{72}{500}{728}{由 $s^2=e$，$r^n=e\Rightarrow s^{-1}=s$，$r^{-1}=r^{n-1}$，$\therefore$ 不妨设 $k_i,l_i\in\N_+$.}
\handline{72}{533}{722}{由 $s^2=e$ 不妨设 $k_i=1$.}
\handline{72}{569}{720}{又 $rsr=s$，$\therefore rs=sr^{-1}$，$\therefore r^ks=sr^{-k}$. 故可借此将 $s,r$ 调序.}
\handline{72}{606}{734}{$\therefore\text{原式}=s^m\cdot r^{l_m-l_{m-1}-\cdots-l_1}$，再将 $l_m-l_{m-1}\cdots-l_1$ 化正.}
\handline{72}{638}{660}{$=s^m\cdot r^k$.}
\handline{72}{673}{719}{最终元素可分为两类：}
\handline{72}{708}{673}{带 $s$，反射：$s,sr,sr^2,\ldots,sr^{n-1}$.}
\handline{72}{744}{690}{不带 $s$，旋转：$e,r,r^2,\ldots,r^{n-1}$.}
\handline{72}{790}{701}{$\therefore D_{2n}=G=\langle r,s\mid s^2=r^n=e,\ rsr=s\rangle=\left\{\begin{array}{l}s,sr,\ldots,sr^{n-1},\\e,r,\ldots,r^{n-1}\end{array}\right\}$.}
\handline[AnnotationBlue]{72}{826}{814}{二面体群还可以用两个反射来生成.}
\handline[AnnotationBlue]{72}{864}{796}{若令 $t=sr$，由 $rsr=s\Rightarrow srsr=s^2\Rightarrow t^2=s^2=e$.}
\handline[AnnotationBlue]{72}{895}{671}{而 $r=s^{-1}t=st$，$\therefore(st)^n=e$.}
\handline[AnnotationBlue]{72}{943}{742}{$\therefore G=\langle s,t\mid s^2=t^2=(st)^n=e\rangle$. 也称为反射群.}
\end{handpage}

% ===== Line-layout page-063.tex =====
\begin{handpage}{63}
\handtitle{72}{98}{757}{\S3.2 群同态}
\handline{72}{125}{762}{定义略.}
\handline{72}{156}{765}{\textbf{Ex.}}
\handline{72}{187}{762}{\circled{1} $H$ 是 $G$ 的子群，则 $i\colon H\to G$ 是群同态，$i$ 是嵌入映射.}
\handline{72}{215}{762}{\circled{2} $V$ 是域 $F$ 上的线性空间，$G=\mathrm{GL}(V)$，则 $\det\colon G\to F\setminus\{0\}$，$\varphi\mapsto\det(\varphi)$}
\handline{72}{243}{736}{是从 $G$ 到 $F\setminus\{0\}$ 的群同态，$F\setminus\{0\}$ 视为乘法群.}
\handline{72}{274}{739}{$\ker(\det)=\{\varphi\in G\mid\det(\varphi)=1\}=\mathrm{SL}(V)$，即特殊线性群.}
\handline{72}{304}{766}{\circled{3} $G$ 是一个群，$\varphi_g\colon G\to G$，令 $\operatorname{Aut}(G)$ 为 $G$ 的自同构群.}
\handline{72}{332}{566}{$h\mapsto ghg^{-1}$}
\handline{72}{371}{730}{$\varphi_g$ 是 $G$ 的一个自同构，即 $\varphi_g\in\operatorname{Aut}(G)$. 且 $(\varphi_g)^{-1}=\varphi_{g^{-1}}$.}
\handline{72}{402}{732}{令 $\Phi\colon G\to\operatorname{Aut}(G)$，$g\mapsto\varphi_g$，则 $\Phi$ 是一个群同态.}
\handline[AnnotationBlue]{72}{435}{723}{$\operatorname{Inn}(G):=\im\Phi$，称为 $G$ 的内自同构群，它是 $\operatorname{Aut}(G)$ 的子群.}
\handline[NoteRed]{72}{468}{727}{重点：看 $\ker\Phi$ 是什么：}
\handline{72}{504}{722}{$\ker\Phi=\{g\in G\mid\varphi_g=\id\}$.}
\handline{72}{537}{716}{$g\in\ker\Phi\Longleftrightarrow\forall h\in G$，$\varphi_g(h)=h$，即 $gh=hg$. \ann{（$g$ 与 $G$ 中所有元素交换.）}}
\handline[AnnotationBlue]{72}{571}{727}{$Z(G):=\{g\in G\mid\forall h\in G,\ gh=hg\}=\ker\Phi$，称为 $G$ 的中心.}
\handline[AnnotationBlue]{72}{606}{726}{$g\in Z(G)$ 称为中心元.}
\handline{72}{638}{725}{$G=\mathrm{GL}(V)$ 的中心元是纯量变换 $\lambda I_V$（$\lambda\in F$）.}
\handline{72}{674}{775}{\circled{4} $G$：有限群，$S(G)$ 是 $G$ 的对称群.}
\handline{72}{710}{728}{$\forall a\in G$，定义 $t_a\colon G\to G$，$x\mapsto ax$ \redann{（左乘）}.}
\handline{72}{746}{731}{则 $t_a$ 是 $G$ 上双射，且 $(t_a)^{-1}=t_{a^{-1}}$，$\therefore t_a\in S(G)$.}
\handline{72}{783}{217}{考虑 $T\colon G\to S(G)$}
\handline{72}{818}{119}{$a\mapsto t_a$.}
\handline{337}{783}{463}{$\forall a,b\in G$，$\forall x\in G$}
\handline{243}{815}{459}{$t_a(t_b(x))=t_a(bx)=abx=(ab)x=t_{ab}(x)$.}
\handline{72}{855}{712}{$\therefore T(a)T(b)=T(ab)$，$\therefore T$ 是一个同态.}
\handline{72}{894}{730}{且 $a\in\ker T\Longleftrightarrow\forall x\in G$，$t_a(x)=ax=x\Longleftrightarrow a=e$，$\therefore\ker T=\{e\}$.}
\handline{72}{931}{724}{$\therefore T$ 是一个单同态，从而 \ann{$G$ 可借助 $T$ 来嵌入 $S(G)$.}}
\handline[AnnotationBlue]{72}{977}{713}{故 $G$ 是 $S(G)$ 的子群，即 $G$ 可看作 $G$ 自身的一个置换群.}
\end{handpage}

% ===== Line-layout page-064.tex =====
\begin{handpage}{64}
\handline{72}{93}{810}{\circled{5} 指数映射 $\exp:(\R,+)\longrightarrow(\R^\times,\cdot)$，是一个群同态.}
\handline{72}{119}{534}{$x\longmapsto e^x$}
\handline{72}{151}{621}{$e:\R\longrightarrow S^1=\{(x,y)\in\R^2\mid x^2+y^2=1\}$ 也是一个群同态.}
\handline{72}{177}{545}{$\theta\longmapsto(\cos\theta,\sin\theta)$}
\handline{72}{215}{818}{\circled{6} $m\in\Z$，$\Z_m\cong\Z/m\Z$，有限交换群，且为循环群.}
\handline{72}{244}{783}{若 $\varphi:\Z_m\to\Z_n$ 是群同态，\ann{只有一个生成元，故只需确定 $\varphi(\overline1)$ 的值.}}
\handline{72}{276}{776}{首先肯定有 $\varphi(\overline0)=\overline0$. 下设 $\varphi(\overline1)=\overline a\Rightarrow\varphi(\overline m)=m\overline a=\varphi(\overline0)=\overline0$.}
\handline[NoteRed]{72}{304}{774}{问：$\Z_n$ 中哪些 $a$ 满足 $m\overline a=\overline0$？}
\handline{72}{337}{779}{\circled{1} 若 $(m,n)=1$，$\exists u,v\in\Z$ s.t. $um+vn=1$，$\therefore a=u(ma)+v(na)$.}
\handline{72}{371}{756}{$\therefore\varphi(\overline1)=\overline a=u(\overline{ma})+v(\overline{na})=\overline0$，$\therefore\varphi$ 是零同态.}
\handline{72}{406}{786}{\circled{2} 若 $(m,n)=d>1$. 设 $n=dn'$. 则 $a$ 可取 $n',2n',\ldots,kn',\ldots$. $m\overline a=m\overline{kn'}=dm'\overline{kn'}=m'k\overline n=\overline0$.}
\handline{72}{436}{783}{且 $\forall k\in\omega$，$d\overline a=d\overline{kn'}=k\overline{dn'}=k\overline n=\overline0$.}
\handline[AnnotationBlue]{72}{469}{782}{类似的想法可用于证明两个 Abel 群不同构.}
\handline{72}{501}{784}{$1^\circ$ $\Z_2\times\Z_2$ 与 $\Z_4$ 不同构；$2^\circ$ $\varphi:\Z_m\to\Z$ 只能是零同态.}
\handline{72}{534}{785}{$1^\circ$ $\Z_2\times\Z_2=\{1,a,b,ab\}$（$a^2=b^2=1$，$ab=ba$），$\Z_4=\{1,c,c^2,c^3\}$（$c^4=1$）.}
\handline{72}{568}{771}{设 $\varphi:\Z_2\times\Z_2\to\Z_4$ 是群同态，则 $\varphi(1)=1$.}
\handline{72}{601}{772}{而 $1=\varphi(1)=\varphi(a^2)=(\varphi(a))^2$，而只有 $1^2=1$、$(c^2)^2=1$，$\therefore\varphi(a)=1$ 或 $c^2$.}
\handline{72}{634}{771}{同理 $\varphi(b)=1$ 或 $c^2$，$\therefore$ 共有 $4$ 种同态.}
\handline{72}{761}{220}{$\begin{array}{c|cccc}&1&a&b&a\cdot b\\\hline\varphi_1&1&1&1&1\\\varphi_2&1&1&c^2&c^2\\\varphi_3&1&c^2&1&c^2\\\varphi_4&1&c^2&c^2&1\end{array}$}
\handline[AnnotationBlue]{72}{673}{489}{显然，$\{\varphi_1,\varphi_2,\varphi_3,\varphi_4\}$ 也构成群}
\handline[AnnotationBlue]{386}{708}{460}{且同构于 $\Z_2\times\Z_2$.}
\handline[NoteRed]{352}{746}{495}{一般地，设 $G,H$ 是交换群，则 $\Hom(G,H)$ 也交换.}
\handline[NoteRed]{354}{783}{492}{$\Hom(G,H)$ 为 $G\to H$ 的群同态的全体.}
\handline[NoteRed]{345}{811}{500}{若 $G,H$ 不交换，则同态的和不一定是同态}
\handline[NoteRed]{72}{834}{235}{从而 $\Hom(G,H)$ 构不成群.}
\handline{72}{850}{701}{$2^\circ$ 设 $\varphi:\Z_m\to\Z$ 是同态，则 $\varphi(\overline0)=0$.}
\handline{72}{886}{668}{设 $\varphi(\overline1)=a$}
\handline{72}{930}{676}{则 $\varphi(\overline m)=ma=0$，由 $\Z$ 是整环，$m\ne0\Rightarrow a=0$.}
\handline{72}{961}{681}{$\therefore\varphi$ 是零同态.}
\end{handpage}

% ===== Line-layout page-065.tex =====
\begin{handpage}{65}
\handtitle{72}{101}{771}{\S3.3 子群、正规子群、商群}
\handline{72}{132}{776}{\textbf{Def.} $G$：群，$H\subseteq G$，若 \circled{1} $g_1,g_2\in H\Rightarrow g_1g_2\in H$；\circled{2} $g\in H\Rightarrow g^{-1}\in H$.}
\handline{72}{160}{711}{则称 $H$ 是 $G$ 的子群，记为 $H\le G$.}
\handline{72}{190}{746}{(i) 等价定义：$g_1,g_2\in H\Rightarrow g_1\cdot g_2^{-1}\in H$.}
\handline{72}{220}{749}{(ii) 若 $G$ 是有限群则 \circled{2} 可省去：$g\in H$，$g^2,g^3,\ldots,g^n,\ldots\in H$.}
\handline{72}{249}{713}{由于有限，$\therefore\exists r,s\in\omega$，s.t. $g^r=g^s$.}
\handline{72}{279}{712}{不妨设 $r>s$，$\Rightarrow g^{r-s-1}=g^{-1}\in H$.}
\handline{72}{311}{790}{\textbf{Ex.} \circled{1} 子域，子空间，子环，子模.}
\handline{72}{341}{737}{\circled{2} $\varphi:G\to H$ 群同态，则 $\ker\varphi\le G$，$\im\varphi\le H$.}
\handline{72}{373}{706}{若 $G'\le G$ 则 $\varphi(G')\le H$.}
\handline{72}{403}{709}{若 $H'\le H$ 则 $\varphi^{-1}(H')\le G$.}
\handline{72}{436}{740}{\circled{3} $\operatorname{GL}(V)$ 的各种子群.}
\handline{72}{468}{745}{\circled{4} 设 $H_\alpha$（$\alpha\in\Lambda$）是 $G$ 的子群，则 $\bigcap_{\alpha\in\Lambda}H_\alpha$ 也是 $G$ 的子群.}
\handline{72}{504}{749}{\circled{5} $G,H$ 是群，$K:=G\times H$；}
\handline{72}{532}{696}{$(g_1,h_1)\cdot(g_2,h_2):=(g_1g_2,h_1h_2)$，$e_K=(e_G,e_H)$.}
\handline{72}{570}{340}{考虑嵌入映射：$i_G:G\to K$}
\handline{72}{599}{122}{$g\mapsto(g,e_H)$}
\handline{432}{570}{353}{$i_H:H\to K$}
\handline{263}{600}{306}{$h\mapsto(e_G,h)$.}
\handline{72}{641}{709}{则 $G,H$ 可视为 $K$ 的子群.}
\handline{72}{683}{752}{\circled{6} $\operatorname{SO}(2,\R)\le\operatorname{SL}(2,\R)\le\operatorname{SL}(2,\C)$.}
\handline{72}{721}{714}{再看 $\operatorname{SL}(2,\Z)=\{\left(\begin{smallmatrix}a&b\\c&d\end{smallmatrix}\right)\mid ad-bc=1,\ a,b,c,d\in\Z\}$.}
\handline{72}{775}{708}{$\begin{pmatrix}a&b\\c&d\end{pmatrix}^{-1}=\frac1{ad-bc}\begin{pmatrix}d&-b\\-c&a\end{pmatrix}=\begin{pmatrix}d&-b\\-c&a\end{pmatrix}\in\operatorname{SL}(2,\Z)$.}
\handline{72}{827}{723}{$\Rightarrow\operatorname{SL}(2,\Z)$ 也是子群.}
\handline{72}{865}{720}{考虑 $\varphi:\operatorname{SL}(2,\Z)\to\operatorname{SL}(2,\Z_m)$.}
\handline{72}{914}{706}{$\ker(\varphi)=\left\{\begin{pmatrix}a&b\\c&d\end{pmatrix}\in M_2(\Z)\ \middle|\ \begin{gathered}a\equiv d\equiv1\pmod m,\\c\equiv b\equiv0\pmod m\end{gathered}\right\}$}
\handline{72}{947}{709}{是 $\operatorname{SL}(2,\Z)$ 的子群，称为模 $m$ 的主同余子群.}
\handline{72}{987}{713}{学习模形式时会与之打交道.}
\end{handpage}

% ===== Line-layout page-066.tex =====
\begin{handpage}{66}
\handline{72}{105}{824}{\textbf{Def.} $S\subseteq G$，$\mathcal H=\{H\le G\mid S\subseteq H\}$.}
\handline{72}{137}{749}{$\langle S\rangle:=\bigcap_{H\in\mathcal H}H$ 称为 $G$ 中由 $S$ 生成的子群.}
\handline{72}{169}{744}{若 $G=\langle S\rangle$，且 $S$ 为有限集，则称 $G$ 是有限生成的.}
\handline{72}{201}{827}{\textbf{Def.} $H\le G$，$\varphi:H\times G\to G$ 为左作用.}
\handline{72}{231}{594}{$(h,x)\mapsto hx$}
\handline{72}{269}{756}{在 $G$ 中定义关系“$x\sim y$”$\Longleftrightarrow\exists h\in H$，s.t. $y=hx$.}
\handline{72}{300}{756}{易验证“$\sim$”是等价关系. 考虑 $G/\sim$.}
\handline{72}{332}{781}{\circled{1} $[x]=\{y\in G\mid y\sim x\}=\{hx\mid h\in H\}=Hx$，称为 $H$ 的右陪集.}
\handline{72}{365}{752}{类似地可从右作用出发定义左陪集.}
\handline{72}{398}{785}{\circled{2} $H\backslash G=\{Hx\mid x\in G\}$，$G/H=\{xH\mid x\in G\}$，分别称为右、左商集.}
\handline{72}{434}{849}{\textbf{Prop.} 两个陪集要么相等，要么不交.}
\handline{72}{467}{841}{\textbf{Pf.} 若 $xH\cap yH\ne\varnothing$.}
\handline{72}{499}{764}{设 $a\in xH\cap yH$，$\therefore\exists h_1,h_2\in H$ s.t. $a=xh_1=yh_2$.}
\handline{72}{533}{742}{$\Rightarrow x=y(h_2\cdot h_1^{-1})\Rightarrow x\sim y\Rightarrow xH=yH$.}
\handline[AnnotationBlue]{72}{572}{780}{令 $X=G/H$，$\mu:G\times X\to X$}
\handline[AnnotationBlue]{72}{605}{574}{$(g,xH)\mapsto(gx)H$.}
\handline[AnnotationBlue]{72}{642}{776}{\circled{1} $\mu$ 不依赖于代表元的选取.}
\handline[AnnotationBlue]{72}{682}{744}{若 $xH=yH$，则 $x\sim y$，$\Rightarrow gx\sim gy\Rightarrow(gx)H=(gy)H$.}
\handline[AnnotationBlue]{72}{719}{782}{\circled{2} $(e,xH)\mapsto(ex)H=xH$.}
\handline[AnnotationBlue]{72}{758}{783}{\circled{3} $g(hx)H=(gh)xH$.}
\handline[AnnotationBlue]{72}{796}{778}{$\therefore\mu$ 是 $G$ 在 $X$ 上的左作用.}
\handline[AnnotationBlue]{72}{832}{771}{而 $X=G/H$ 称为 $G$ 模 $H$ 的左齐性空间.}
\handline[AnnotationBlue]{72}{870}{775}{右齐性空间同理可定义.}
\end{handpage}

% ===== Line-layout page-067.tex =====
\begin{handpage}{67}
\handline[AnnotationBlue]{72}{104}{768}{若要在商集上定义一个运算，$H$ 要满足什么条件？}
\handline{72}{137}{755}{设 $X=G/H$，定义乘法 $\cdot:X\times X\to X$}
\handline{72}{166}{496}{$(xH,yH)\mapsto(xy)H$.}
\handline{72}{198}{761}{不依赖于代表元选取：若 $xH=x'H$，$yH=y'H$，则 $(xy)H=(x'y')H$.}
\handline{72}{229}{700}{$\exists h_1,h_2\in H$ s.t. $x'=xh_1$，$y'=yh_2$.}
\handline{72}{259}{655}{$\therefore x'y'=xh_1yh_2$，又需要 $\exists h_3\in H$ s.t. $x'y'=xyh_3$.}
\handline{72}{290}{642}{$\therefore xh_1yh_2=xyh_3$，$\therefore h_1y=yh_3h_2^{-1}=yh_4$（$h_4=h_3h_2^{-1}$）.}
\handline{72}{320}{641}{或者记为 $h_4=y^{-1}h_1y$.}
\handline{72}{355}{661}{用陪集的观点看 $\forall y\in G,h_1\in H$，必须要有 $h_1y\in yH$\ann{（或 $y^{-1}h_1y\in H$）}.}
\handline{72}{388}{660}{由 $h_1$ 的任意性 $\Rightarrow Hy\subseteq yH$\ann{（或 $y^{-1}Hy\subseteq H$）}.}
\handline{72}{422}{798}{\textbf{Prop.} 若 $H\le G$，则 TFAE：}
\handline{72}{455}{300}{\circled{1} $\forall g\in G$，$gH\subseteq Hg$ \ann{$(Hg\subseteq gH)$}}
\handline{383}{455}{407}{\circled{1}' $\forall g\in G$，$gH=Hg$ \ann{（左右陪集相等）}}
\handline{72}{487}{300}{\circled{2} $\forall g\in G$，$gHg^{-1}\subseteq H$ \ann{$(H\subseteq gHg^{-1})$}}
\handline{383}{487}{407}{\circled{2}' $\forall g\in G$，$gHg^{-1}=H$ \ann{$(g^{-1}Hg=H)$}.}
\handline{72}{526}{802}{\textbf{Pf.} 只证 \circled{1}$\Rightarrow$\circled{1}'. 若 $\forall g\in G$，$gH\subseteq Hg$，$\Rightarrow gHg^{-1}\subseteq H\Rightarrow(g^{-1})^{-1}Hg^{-1}\subseteq H$.}
\handline{72}{561}{549}{令 $\widetilde g=g^{-1}$，$\Rightarrow\forall\widetilde g\in G$，$(\widetilde g)^{-1}H\widetilde g\subseteq H\Rightarrow H\widetilde g\subseteq H\widetilde g$.}
\handline{72}{594}{546}{$\therefore\forall g\in G$，$Hg\subseteq gH$，$\therefore gH=Hg$.}
\handline{72}{635}{830}{\textbf{Def.} 若 $H\le G$，满足 \circled{1}\circled{1}'\circled{2}\circled{2}' 之一，则称 $H$ 为 $G$ 的正规子群\ann{（normal subgroup）}}
\handline{72}{669}{738}{记为 $H\lhd G$.}
\handline[AnnotationBlue]{72}{706}{812}{显然，只有在正规子群的商集上才能定义乘法运算.}
\handline{72}{747}{821}{\textbf{Def.} 设 $H\lhd G$，在 $G/H$ 上定义 $\cdot:G/H\times G/H\to G/H$，}
\handline{72}{781}{721}{$(xH)(yH)=(xy)H$，$e_{G/H}:=H$，$(xH)^{-1}:=x^{-1}H$.}
\handline{72}{819}{741}{则 $(G/H,\cdot)$ 是一个群，称为 $G$ 模 $H$ 的商群.}
\handline{72}{863}{826}{\textbf{Prop.} $H\lhd G$，则存在群同态 $\pi:G\to G/H$：$x\mapsto xH$，且 $\ker(\pi)=H$.}
\handline{72}{904}{832}{\textbf{Pf.} $\forall x,y\in G$，$(xH)\cdot(yH)=(xy)H$，$\therefore\pi(x)\cdot\pi(y)=\pi(xy)$.}
\handline{72}{944}{711}{$x\in\ker(\pi)\Longleftrightarrow xH=H\Longleftrightarrow x\in H$，$\therefore\ker(\pi)=H$.}
\end{handpage}

% ===== Line-layout page-068.tex =====
\begin{handpage}{68}
\handline{72}{124}{830}{\textbf{Ex.} \circled{1} 在交换群中所有子群都是正规的.}
\handline{72}{154}{782}{\circled{2} $G_1=\{1,c,c^2,c^3\}$，$H_1=\{1,c^2\}$，$G_1/H_1=\{H_1,cH_1\}$.}
\handline{72}{184}{747}{$G_2=\{1,a,b,ab\}$，$H_2=\{1,a\}$，$G_2/H_2=\{H_2,bH_2\}$.}
\handline{72}{216}{743}{$G_1\not\cong G_2$，但 $H_1\cong H_2$ 且 $G_1/H_1\cong G_2/H_2$. \ann{（这个现象与同调代数有关）}}
\handline{72}{249}{789}{\circled{3} $G=S_3$，$H=\{1,r,r^2\}$ 是正规的，$G/H=\{H,sH\}\cong\Z_2$，而 $H\cong\Z_3$.}
\handline{72}{284}{747}{令 $G'=\Z_6\cong\Z_2\times\Z_3$，$G'/\Z_3=\Z_2$.}
\handline{72}{314}{750}{$G'$ 交换而 $G$ 不交换，$\therefore G\not\cong G'$，但 $G/H\cong G'/\Z_3$.}
\handline{72}{349}{755}{$G$ 是可解群\ann{（这类群与 Galois 理论密切相关）}.}
\handline{72}{385}{801}{\circled{4} $X=G/H$ 为齐性空间. 取 $G=\operatorname{SO}(3)$，$H=\operatorname{SO}(2)$.}
\handline{72}{423}{761}{将 $H$ 以 $Q\mapsto\begin{pmatrix}1&0\\0&Q\end{pmatrix}$ 的方式嵌入 $G$，并视为 $G$ 的子群.}
\handline{72}{460}{795}{但 $H$ 不是正规子群，此时 $G/H$ 是二维射影平面\ann{（这是齐性空间给出的由来）}.}
\handline{72}{498}{810}{\circled{5} $Z(G)=\{h\in G\mid\forall g\in G,\ gh=hg\}$ 称为 $G$ 的中心.}
\handline{72}{533}{769}{$Z(G)$ 是正规子群.}
\handline{72}{572}{775}{令 $G_1=G/Z(G)$，$G_2=G_1/Z(G_1)$，$\ldots$，由此得到高阶中心的概念.}
\handline{72}{610}{220}{$\varphi:G\to\operatorname{Aut}(G)$}
\handline{72}{645}{181}{$g\mapsto\mu_g$}
\handline{295}{610}{548}{其中 $\mu_g:G\to G$}
\handline{376}{645}{428}{$x\mapsto gxg^{-1}$.}
\handline{72}{683}{771}{$\operatorname{Inn}(G):=\im(\varphi)$ 称为内自同构群，}
\handline{72}{721}{775}{且 $G/Z(G)\cong\operatorname{Inn}(G)$.}
\handline[AnnotationBlue]{72}{757}{768}{$\varphi$ 是同态显然. $g\in\ker(\varphi)\Longleftrightarrow\mu_g=\id$.}
\handline[AnnotationBlue]{72}{797}{741}{$\Longleftrightarrow\forall x\in G,\ gxg^{-1}=x\Longleftrightarrow\forall x\in G,\ gx=xg\Longleftrightarrow g\in Z(G)$.}
\handline[AnnotationBlue]{72}{831}{772}{$\therefore\ker(\varphi)=Z(G)$.}
\handline[AnnotationBlue]{72}{875}{776}{由同态基本定理 $\Rightarrow G/Z(G)\cong\operatorname{Inn}(G)$.}
\end{handpage}

% ===== Line-layout page-069.tex =====
\begin{handpage}{69}
\handline{72}{126}{766}{\circled{6} $\forall a,b\in G$，$[a,b]:=a^{-1}b^{-1}ab$，称为交换子或换位子.}
\handline{72}{156}{762}{\textbf{Prop.} \circled{1} $[a,b]^{-1}=(a^{-1}b^{-1}ab)^{-1}=b^{-1}a^{-1}ba=[b,a]$.}
\handline{72}{187}{697}{\circled{2} $g[a,b]g^{-1}=ga^{-1}b^{-1}abg^{-1}=(ga^{-1}g^{-1})(gb^{-1}g^{-1})(gag^{-1})(gbg^{-1})=[gag^{-1},gbg^{-1}]$.}
\handline[AnnotationBlue]{72}{219}{668}{交换子的共轭等于共轭的交换子.}
\handline{72}{251}{771}{用 $[G,G]$ 表示所有交换子构成的集合，易验证它是子群.}
\handline{72}{283}{743}{$[G,G]=\{[a_1,b_1]^{d_1}\cdots[a_n,b_n]^{d_n}\mid n\in\omega,\ d_i\in\{\pm1\},\ a_i,b_i\in G\}$.}
\handline{72}{315}{744}{由于 $[a,b]^{-1}=[b,a]$，$\therefore[G,G]=\{[a_1,b_1]\cdots[a_n,b_n]\mid n\in\omega,\ a_i,b_i\in G\}$.}
\handline{72}{348}{750}{又由于 $g[a,b]g^{-1}=[gag^{-1},gbg^{-1}]\Rightarrow g[G,G]g^{-1}\subseteq[G,G]$.}
\handline{72}{381}{754}{$\therefore[G,G]\lhd G$.}
\handline{72}{411}{754}{而 $G/[G,G]$ 是一个交换群，称为群 $G$ 的交换化.}
\handline[AnnotationBlue]{72}{444}{750}{$\forall x\in(ab)[G,G]$，$x=ab[a_1,b_1]\cdots[a_n,b_n]=(baa^{-1}b^{-1})ab[a_1,b_1]\cdots[a_n,b_n]$}
\handline[AnnotationBlue]{72}{478}{610}{$=ba(a^{-1}b^{-1}ab)[a_1,b_1]\cdots[a_n,b_n]=ba[a,b][a_1,b_1]\cdots[a_n,b_n]=(ba)[G,G]$.}
\handline{72}{516}{782}{$\therefore(ab)[G,G]\subseteq(ba)[G,G]$. 同理 $(ba)[G,G]\subseteq(ab)[G,G]$，$\therefore(ba)[G,G]=(ab)[G,G]$.}
\handline{72}{554}{758}{$\therefore\forall a,b\in G$，$(a[G,G])\cdot(b[G,G])=(b[G,G])(a[G,G])$. $\therefore G/[G,G]$ 交换.}
\handline[AnnotationBlue]{72}{593}{757}{$G=S_3$，$H=[G,G]=\{1,r,r^2\}$，$G/H=\{H,sH\}$ 是交换群.}
\handline{72}{633}{803}{\circled{7} 任意多个正规子群的交仍是正规子群.}
\handline{72}{677}{764}{$S\subseteq G$ 为 $G$ 的子集，$\langle S\rangle:=\bigcap_{\substack{S\subseteq H\\H\lhd G}}H$，称为 $S$ 生成的正规子群.}
\handline{72}{739}{809}{\circled{8} 任何同态的 $\ker$ 都是正规子群.}
\handline{72}{776}{754}{$\varphi:G\to G'$ 是同态. 设 $H=\ker(\varphi)$.}
\handline{72}{816}{757}{只需证 $\forall g\in G$，$gHg^{-1}\subseteq H$.}
\handline{72}{854}{753}{$\forall h\in H$，$\varphi(h)=e$.}
\handline{72}{891}{697}{$\therefore\varphi(ghg^{-1})=\varphi(g)\varphi(h)\varphi(g^{-1})=\varphi(g)e\cdot(\varphi(g))^{-1}=e$.}
\handline{72}{931}{700}{$\therefore ghg^{-1}\in\ker(\varphi)=H$，$\therefore gHg^{-1}\subseteq H$.}
\handline{72}{974}{749}{$\therefore H=\ker(\varphi)\lhd G$.}
\end{handpage}

% ===== Line-layout page-070.tex =====
\begin{handpage}{70}
\handline{72}{119}{843}{\circled{9} $H\le G$，$N_G(H):=\{g\in G\mid gHg^{-1}=H\}=\{g\in G\mid gH=Hg\}$}
\handline{72}{150}{792}{称为 $H$ 的正规化子.（若 $H\lhd G$，则 $N_G(H)=G$.）}
\handline[AnnotationBlue]{72}{181}{799}{如果 $H$ 不是正规子群，降低要求：把满足 $gHg^{-1}=H$ 的 $g$ 收集在一起.}
\handline{72}{213}{794}{显然 $N_G(H)\le G$ 且 $H\lhd N_G(H)$.}
\handline[NoteRed]{72}{247}{794}{$N_G(H)$ 是 $G$ 中包含 $H$ 作为正规子群的最大子群.}
\handline{72}{280}{796}{$C_G(H):=\{g\in G\mid\forall h\in H,\ gh=hg\}$，称为 $H$ 的中心化子.}
\handline{72}{314}{856}{\textbf{Prop.} $C_G(H)\lhd N_G(H)$.}
\handline{72}{348}{850}{\textbf{Pf.} 即证 $\forall n\in N_G(H)$，$nC_G(H)n^{-1}\subseteq C_G(H)$.}
\handline{72}{379}{766}{即证 $\forall n\in N_G(H)$，$\forall c\in C_G(H)$，$ncn^{-1}\in C_G(H)$.}
\handline{72}{412}{759}{$\forall h\in H$，由 $n\in N_G(H)$，$\therefore nH=Hn$，$\therefore\exists h'\in H$ s.t. $hn=nh'\Rightarrow n^{-1}h=h'n^{-1}$.}
\handline{72}{444}{768}{由 $c\in C_G(H)$，$\therefore ch'=h'c$.}
\handline{72}{484}{736}{$\therefore h(ncn^{-1})=nh'cn^{-1}=nch'n^{-1}=(ncn^{-1})h$.}
\handline{72}{523}{756}{$\therefore ncn^{-1}\in C_G(H)$.}
\handline{72}{565}{856}{\circled{10} $G=S_n$，$H=A_n\longrightarrow n$ 元偶置换. 则 $|G|=n!$，$|H|=\frac12n!$，$A_n\lhd S_n$.}
\handline{72}{607}{856}{\textbf{Prop.} 若 $|G/H|=2$ 则 $H\lhd G$.}
\handline{72}{645}{856}{\textbf{Pf.} 由于 $|G/H|=2$，$\therefore\exists x\in G$，$x\notin H$，s.t. $G/H=\{H,xH\}$.}
\handline{72}{685}{798}{由于陪集要么相等，要么不交，$\therefore G=H\sqcup xH$.}
\handline{72}{727}{802}{而 $Hx\ne H$，但 $Hx\subseteq G\Rightarrow Hx\subseteq xH$. 同理 $xH\subseteq Hx$，$\therefore xH=Hx$，$\therefore H\lhd G$.}
\end{handpage}

% ===== Line-layout page-071.tex =====
\begin{handpage}{71}
\handline{72}{77}{255}{\textbf{Thm.} \ann{(Lagrange)}}
\handline{72}{111}{350}{设 $|G|<\infty$，$H\le G$，则 $|H|\mid|G|$.}
\handline{72}{145}{310}{\textbf{Pf.} $\forall x\in G$，$x\notin H$，$|H|=|xH|$.}
\handline{72}{174}{320}{而陪集要么相等，要么不交.}
\handline{72}{205}{491}{所以 $G=\bigsqcup_{xH\in G/H}(xH)$，所以 $|G|=|H|\cdot|G/H|$，所以 $|H|\mid|G|$.}
\handline{72}{261}{620}{\textbf{Def.} $|G|<\infty$，$H\le G$，$[G:H]:=|G/H|$ 称为 $H$ 在 $G$ 中的指标.}
\handline[AnnotationBlue]{72}{292}{247}{显然 $|G/H|\mid|G|$}
\handline{72}{322}{821}{\textbf{Def.} $|G|<\infty$，$a\in G$. 使得 $a^n=e$ 的最小正整数 $n$ 称为 $a$ 的阶数或周期，记为 $o(a)$}
\handline[AnnotationBlue]{72}{355}{332}{首先 $o(a)$ 必存在. $1,a,a^2,\ldots$}
\handline[AnnotationBlue]{72}{385}{587}{由于 $|G|<\infty$，所以 $\exists s,t\in\omega$ s.t. $a^s=a^t$，不妨设 $s>t$.}
\handline[AnnotationBlue]{72}{416}{440}{所以 $a^{s-t}=e$. 令 $r=s-t\in\omega$，所以 $a^r=e$.}
\handline[AnnotationBlue]{72}{449}{354}{由 $\omega$ 的良序性 $\Rightarrow o(a)$ 存在.}
\handline{72}{486}{414}{\textbf{Cor.} $|G|<\infty$，$a\in G$，则 $a^{|G|}=e$.}
\handline{72}{518}{502}{\textbf{Pf.} 设 $r=o(a)$，则 $\{e,a,\ldots,a^{r-1}\}\le G$.}
\handline{72}{550}{450}{所以 $r\mid|G|$，所以 $a^{|G|}=(a^r)^{|G|/r}=e^{|G|/r}=e$}
\handline[AnnotationBlue]{72}{585}{573}{数论中的 Euler 定理、Fermat 小定理可由此推出.}
\handline{72}{622}{410}{\textbf{Prop.} $|G|<\infty$，$a\in G$，则 $o(a)\mid|G|$.}
\handline{72}{656}{652}{\textbf{Pf.} 作带余除法 $|G|=bo(a)+r\Rightarrow a^r=e$，$0\le r<o(a)$}
\handline{72}{690}{438}{由 $o(a)$ 的最小性 $\Rightarrow r=0\Rightarrow o(a)\mid|G|$.}
\handline[NoteRed]{72}{727}{559}{事实上，$|G|$ 可改为任意一个 $n\in\omega$，s.t. $a^n=e$，有 $o(a)\mid n$}
\handline{72}{765}{450}{\textbf{Cor.} 若 $|G|=p$ 为素数，则 $G\cong\Z_p$}
\handline{72}{806}{829}{\textbf{Pf.} 设 $a\in G$ 且 $a\ne e$. 于是 $o(a)\mid p\Rightarrow o(a)=p$，而 $\{e,a,\ldots,a^{p-1}\}\subseteq G$，$|G|=p$}
\handline{72}{841}{368}{所以 $G=\{e,a,\ldots,a^{p-1}\}\cong\Z_p$}
\handline[NoteRed]{72}{883}{542}{注意：$\{e,a,\ldots,a^{o(a)-1}\}$ 一定互不相同.}
\handline[NoteRed]{72}{919}{639}{否则 $a^s=a^t$，$0\le t<s\le o(a)-1\Rightarrow0<s-t<o(a)$}
\handline[NoteRed]{72}{950}{445}{$a^{s-t}=e$ 与 $o(a)$ 的最小性矛盾！}
\end{handpage}

% ===== Line-layout page-072.tex =====
\begin{handpage}{72}
\handtitle{72}{85}{350}{$\S3.4$ 同态基本定理}
\handline{72}{116}{290}{\textbf{Thm.} \ann{（第一同构定理）}}
\handline{72}{145}{440}{$\varphi:G\to H$ 为群同态，则 $G/\ker\varphi\cong\im(\varphi)$}
\handline{72}{174}{419}{特别地，若 $\varphi$ 为满同态，则 $G/\ker(\varphi)\cong H$.}
\handline{72}{208}{260}{\textbf{Thm.} \ann{（对应定理）}}
\handline{72}{237}{296}{$\varphi:G\to H$ 为满同态，则}
\handline{72}{266}{283}{\circled{1} 若 $G'\le G$，则 $\varphi(G')\le H$}
\handline{72}{294}{477}{\circled{2} 若 $H'\le H$，则 $\varphi^{-1}(H')\le G$ 且 $\ker(\varphi)\subseteq\varphi^{-1}(H')$}
\handline{72}{326}{542}{\circled{3} $\mathcal G:\{G'\mid G'\le G\text{ 且 }\ker\varphi\le G'\}\to\{H'\mid H'\le H\}$，}
\handline{72}{355}{208}{$G'\mapsto\varphi(G')$}
\handline{72}{387}{206}{则 $\mathcal G$ 是双射}
\handline{72}{420}{528}{\circled{4} 若 $G'\lhd G$，则 $H'=\varphi(G')\lhd H$ 且 $G/G'\cong H/H'$.}
\handline[AnnotationBlue]{72}{449}{321}{这两个定理与环版本类似.}
\handline{72}{482}{490}{\textbf{Def.} 设 $H,K\le G$，$HK:=\{hk\mid h\in H,k\in K\}$.}
\handline{72}{514}{273}{$HK$ 不一定是 $G$ 的子群}
\handline[AnnotationBlue]{72}{545}{736}{如 $G=S_3$，$H=\{1,s\}$，$K=\{1,sr\}$，$\Rightarrow HK=\{1,s,r,sr\}$ 不是 $S_3$ 的子群}
\handline{72}{579}{307}{\textbf{Prop.} 若 $K\lhd G$，则 $HK\le G$}
\handline{72}{609}{267}{且 $K\lhd HK$，$H\cap K\lhd H$.}
\handline{72}{641}{596}{\textbf{Pf.} \circled{1} $\forall a,b\in HK$，$a=h_1k_1$，$b=h_2k_2$，$h_1,h_2\in H$，$k_1,k_2\in K$.}
\handline{72}{684}{753}{$a\cdot b^{-1}=h_1k_1\cdot k_2^{-1}h_2^{-1}=\underbrace{h_1h_2^{-1}}_{h_3}(\underbrace{h_2k_1h_2^{-1}}_{k_3})(\underbrace{h_2k_2^{-1}h_2^{-1}}_{k_4})=h_3\underbrace{k_3k_4}_{k_5}=h_3k_5\in HK$.}
\handline{72}{728}{137}{所以 $HK\le G$.}
\handline{72}{752}{536}{\circled{2} 显然，$K\le HK$. $\forall g\in HK\subseteq G\Rightarrow gK=Kg$，所以 $K\lhd HK$}
\handline{72}{784}{605}{\circled{3} 显然，$H\cap K\le H$. $\forall h\in H$，下证 $h(H\cap K)\subseteq(H\cap K)h$.}
\handline{72}{818}{329}{$\forall g\in H\cap K\Rightarrow g\in H$，$g\in K$.}
\handline{72}{850}{432}{由于 $hK=Kh$，所以 $\exists k\in K$ s.t. $hg=kh$.}
\handline{72}{879}{354}{而 $k=hgh^{-1}\in H$，所以 $k\in H\cap K$}
\handline{72}{915}{665}{所以 $hg=(hgh^{-1})h=kh\in(H\cap K)h$，所以 $h(H\cap K)\subseteq(H\cap K)h$}
\handline{72}{950}{253}{所以 $H\cap K\lhd H$.}
\end{handpage}

% ===== Line-layout page-073.tex =====
\begin{handpage}{73}
\handline{72}{115}{315}{\textbf{Thm.} \ann{（第二同构定理）}}
\handline{72}{146}{356}{若 $K\lhd G$，$H\le G$，则 $H/(H\cap K)\cong HK/K$.}
\handline{72}{176}{496}{\textbf{Pf.} 由前一个命题 $\Rightarrow H/(H\cap K)$、$HK/K$ 是良定义的}
\handline{72}{205}{255}{定义 $\varphi:H\to HK/K$}
\handline{72}{229}{126}{$h\mapsto hK$.}
\handline{72}{263}{526}{则 \circled{1} $\varphi(h_1h_2)=h_1h_2K=(h_1K)(h_2K)=\varphi(h_1)\varphi(h_2)$}
\handline{72}{291}{201}{\circled{2} 满射：显然.}
\handline{72}{319}{490}{\circled{3} $h\in\ker\varphi\Longleftrightarrow\varphi(h)=K=hK\Longleftrightarrow h\in H\cap K$.}
\handline{72}{352}{565}{所以由第一同构定理，$H/\ker\varphi\cong HK/K$，即 $H/(H\cap K)\cong HK/K$.}
\handline{72}{387}{328}{\textbf{Thm.} \ann{（第三同构定理）}}
\handline{72}{417}{555}{若 $K\lhd G$，$K\le H\le G$，则 $K\lhd H$. 从而 $G/K$、$H/K$ 良定义.}
\handline{72}{448}{499}{则 $H\lhd G\Longleftrightarrow H/K\lhd G/K$ 且 $(G/K)/(H/K)\cong G/H$.}
\handline{72}{482}{488}{\textbf{Pf.} “$\Rightarrow$” 若 $H\lhd G$，定义 $\rho:G/K\to G/H$}
\handline{72}{510}{143}{$gK\mapsto gH$}
\handline{72}{546}{490}{由于 $K\subseteq H$，故 $\rho$ 是良定义的，且显然是满同态.}
\handline{72}{577}{509}{$gK\in\ker\rho\Rightarrow gH=H\Longleftrightarrow g\in H$，所以 $gK\in H/K$.}
\handline{72}{610}{252}{所以 $\ker\rho\subseteq H/K$.}
\handline{72}{643}{452}{$\forall gK\in H/K\Rightarrow g\in H\Rightarrow\rho(gK)=gH=H$}
\handline{72}{674}{440}{所以 $H/K\subseteq\ker\rho$，所以 $\ker\rho=H/K$.}
\handline{72}{707}{407}{由于 $\ker\rho$ 一定正规，所以 $H/K\lhd G/K$.}
\handline{72}{741}{427}{由第一同构定理 $\Rightarrow(G/K)/(H/K)\cong G/H$.}
\handline{72}{775}{716}{“$\Leftarrow$” 若 $H/K\lhd G/K$，考虑 $\pi_{G/K}:G\to G/K$，$\bar\pi:G/K\to(G/K)/(H/K)$.}
\handline{72}{810}{722}{则 $\bar\pi,\pi_{G/K}$ 为满同态. 令 $\sigma=\bar\pi\circ\pi_{G/K}:G\to(G/K)/(H/K)$，则 $\sigma$ 也是满同态.}
\handline{72}{844}{566}{$g\in\ker\sigma\Longleftrightarrow\bar\pi(\pi_{G/K}(g))=\bar\pi(gK)=gK(H/K)=H/K$}
\handline{72}{879}{555}{$\Longleftrightarrow gK\in H/K\Longleftrightarrow g\in H$，所以 $\ker\sigma=H$.}
\handline{72}{913}{370}{由于 $\ker\sigma$ 一定正规，所以 $H\lhd G$}
\handline{72}{950}{490}{由第一同构定理 $\Rightarrow G/H\cong(G/K)/(H/K)$.}
\end{handpage}

% ===== Line-layout page-074.tex =====
\begin{handpage}{74}
\handline{72}{89}{651}{正规子群的正规子群不一定是正规子群，即正规性不具有传递性.}
\handline{72}{120}{480}{取 $G=D_8=\{\begin{smallmatrix}e,r,r^2,r^3,\\s,sr,sr^2,sr^3\end{smallmatrix}\}$，设 $H=\{1,s,r^2,sr^2\}$}
\handline{72}{150}{485}{由于 $|G/H|=2$，所以 $H\lhd G$. 取 $K=\{1,s\}$，则 $K\lhd H$.}
\handline{72}{178}{514}{但 $rKr^{-1}=\{1,sr^2\}\ne K$，所以 $K$ 不是 $G$ 的正规子群.}
\handline{72}{209}{394}{\textbf{Def.} \circled{1} 设 $G$ 是有限群，若 $G$ 有子群链.}
\handline{72}{236}{305}{$G=G_0>G_1>G_2>\cdots>G_k=K$}
\handline{72}{265}{451}{每个 $G_i$ 是 $G_{i-1}$ 的正规子群.“$<$”表示真子群.}
\handline{72}{293}{323}{则称 $K$ 为 $k$-次正规子群}
\handline{72}{325}{282}{\circled{2} 若有 $G$ 的子群链}
\handline{72}{353}{389}{$G=G_0\ge G_1\ge\cdots\ge G_r=\{e\}$，$G_i\lhd G_{i-1}$}
\handline{72}{382}{332}{则称其为 $G$ 的次正规子群链.}
\handline{72}{413}{428}{\circled{3} 若 $G$ 只有平凡正规子群，则称 $G$ 为单群.}
\handline{72}{445}{570}{对于一个 $G$ 的次正规子群链，若 $S_i=G_{i-1}/G_i$ 均为单群.}
\handline{72}{475}{521}{则称这个子群链为合成群列. $S_i$ 称为合成因子.}
\handline{72}{506}{525}{若 $\forall i$，$G_i\ne G_{i+1}$，则称此时的合成列是极小的.}
\handline{72}{535}{430}{此时子群链的长度称为合成列的长度.}
\handline{72}{569}{482}{\textbf{Ex.} \circled{1} 有限交换单群一定是 $\Z_p$，$p$ 为素数.}
\handline{72}{599}{372}{\circled{2} $G=D_8$，$H,K$ 同最上面的例子.}
\handline{72}{629}{554}{所以 $G>H>K>\{e\}$，而 $|G/H|=|H/K|=|K/\{e\}|=2$}
\handline{72}{659}{452}{所以 $G/H\cong H/K\cong K/\{e\}\cong\Z_2$ 均为单群.}
\handline{72}{692}{405}{所以 $G>H>K>\{e\}$ 为极小合成群列.}
\handline{72}{725}{630}{\circled{3} $G=\Z_6\cong\Z_2\times\Z_3$. $\quad G>\Z_3>\{e\}$，$S_1=\Z_2$，$S_2=\Z_3$}
\handline{72}{756}{390}{$G>\Z_2>\{e\}$，$S_1=\Z_3$，$S_2=\Z_2$}
\handline{72}{791}{566}{\circled{4} $G$ 只有为有限群时，才可定义合成列. 若 $G$ 为无限群.}
\handline{72}{822}{389}{比如 $G=\Z$，$G_1=p\Z$，$G_2=p^2\Z$，……}
\handline{72}{854}{713}{$G>G_1>G_2,\ldots$，而 $\forall i\in\omega$，$S_i=\Z_p$ 为单群，合成列无止境.}
\handline{72}{888}{421}{而选取不同的素数 $p$，合成因子也不同.}
\handline[AnnotationBlue]{72}{921}{635}{可以类比不是唯一分解的环，如 $\C[x^r\mid r\in\Q,r>0]$.}
\handline[AnnotationBlue]{72}{950}{308}{分解可以无限进行下去.}
\end{handpage}

% ===== Line-layout page-075.tex =====
\begin{handpage}{75}
\handline{72}{143}{288}{\textbf{Thm.} \ann{(Jordan--H\"older)}}
\handline{72}{171}{733}{有限群 $G$ 的极小合成列有相同的长度，且合成因子在不计次序以及同构的意义下唯一.}
\handline{72}{198}{545}{\textbf{Pf.} 设 $G$ 有两个极小合成列. $G=G_0>G_1>\cdots>G_s=\{e\}$ \circled{1}}
\handline{72}{224}{245}{$=H_0>H_1>\cdots>H_t=\{e\}$. \circled{2}}
\handline{72}{255}{337}{要证 $s=t$，且 $G_0/G_1,G_1/G_2,\ldots$}
\handline{72}{283}{402}{与 $H_0/H_1,H_1/H_2,\ldots$ 在调序后是同构的.}
\handline{72}{315}{553}{对 $s$ 进行归纳. $s=1$ 时，$G$ 只有平凡正规子群，即 $G$ 是单群.}
\handline{72}{344}{364}{故只有一条合成列 $G>\{e\}$，结论成立.}
\handline{72}{375}{334}{假设 $s-1$ 成立，下证 $s$ 时成立.}
\handline{72}{407}{691}{Case I：若 $G_1=H_1$，则由归纳假设，对 $G_1$ 和 $H_1$ 使用 $s-1$ 的情形，则结论成立.}
\handline{72}{438}{477}{Case II：若 $G_1\ne H_1$，由于 $G_0/G_1$、$H_0/H_1$ 是单群}
\handline{72}{467}{390}{所以 $G_1,H_1$ 分别是 $G_0,H_0$ 的极大正规子群.}
\handline{72}{498}{675}{令 $G'=G_1\cdot H_1$，则由于 $G_1,H_1$ 正规，所以 $G'\lhd G$，且 $G_1\lhd G'$，$H_1\lhd G'$.}
\handline{72}{529}{350}{所以 $G_1\lhd G'\lhd G$，$H_1\lhd G'\lhd G$.}
\handline{72}{560}{398}{而 $G_1\ne H_1$，结合极大性 $\Rightarrow G'=G$}
\handline{72}{591}{386}{设 $G_1\cap H_1=N_2$，则 $N_2\lhd G_1$，$N_2\lhd H_1$}
\handline{72}{620}{434}{所以 $G_0/G_1=G_1H_1/G_1\cong H_1/(G_1\cap H_1)=H_1/N_2$}
\handline{72}{649}{398}{$H_0/H_1=G_1H_1/H_1\cong G_1/(G_1\cap H_1)=G_1/N_2$}
\handline{72}{683}{477}{取 $N_2$ 的极小合成列. $N_2>N_3>\cdots>N_r>\{e\}$.}
\handline{72}{712}{408}{而 $G=G_0>G_1>N_2>\cdots>N_r>\{e\}$ \circled{1}$'$}
\handline{72}{741}{351}{$=H_0>H_1>N_2>\cdots>N_r>\{e\}$. \circled{2}$'$}
\handline{72}{777}{617}{对于 \circled{1} 和 \circled{1}$'$ 利用 Case I $\Rightarrow r=s$ 且合成因子调序后同构}
\handline{72}{809}{581}{同理，对于 \circled{2} 和 \circled{2}$'$ $\Rightarrow r=t$ 且合成因子调序后同构.}
\handline{72}{840}{445}{最后对于 \circled{1}$'$ 和 \circled{2}$'$ 只需看前两个合成因子.}
\handline{72}{871}{366}{$G_0/G_1,G_1/N_2$ 和 $H_0/H_1,H_1/N_2$}
\handline{72}{908}{773}{而 $G_0/G_1\cong H_1/N_2$，$H_0/H_1\cong G_1/N_2$. 故 \circled{1}$'$、\circled{2}$'$ 在调序后合成因子同构.}
\handline{72}{941}{449}{所以 \circled{1} 和 \circled{2} 的合成因子在调序后同构.}
\end{handpage}

% ===== Line-layout page-076.tex =====
\begin{handpage}{76}
\handline[AnnotationBlue]{72}{109}{778}{有了这个定理作为保障，任意一个有限群 $G$ 会对应一个合成因子 $\{S_1,\ldots,S_r\}$.}
\handline[AnnotationBlue]{72}{144}{409}{其中 $S_1,\ldots,S_r$ 为有限单群.}
\handline[AnnotationBlue]{72}{177}{470}{于是研究有限群的结构可以分两步走：}
\handline[AnnotationBlue]{72}{208}{317}{\circled{1} 找出所有有限单群}
\handline[AnnotationBlue]{72}{239}{425}{\circled{2} 给定合成因子，找出所有的 $G$.}
\handline[AnnotationBlue]{72}{273}{561}{若给定 $S_1,\ldots,S_r$，则 $|G|=|S_1|\cdots|S_r|$，$G$ 的阶可确定}
\handline[AnnotationBlue]{72}{305}{561}{但如何让这些有限单群拼成一个群，仍未解决.}
\handline[AnnotationBlue]{72}{341}{812}{即使 $r=2$，$G>H>\{e\}$，$H,G/H$ 是单群，给定 $H$ 和 $G/H$，去找 $G$，依然很困难.}
\handline[AnnotationBlue]{72}{376}{393}{这类问题称为群扩张问题.}
\handline[AnnotationBlue]{72}{415}{614}{\circled{1} 现在已经解决了. Gorenstein 是这一领域的领军人物.}
\handline[AnnotationBlue]{72}{447}{364}{有限单群只有以下几类：}
\handline[AnnotationBlue]{72}{481}{316}{\circled{1} 素数阶循环群 $\Z_p$.}
\handline{72}{513}{592}{\ann{\circled{2} 交错群 $A_n$（$n\ge5$）.}\redann{$A_5$ 是最小的非交换单群.}}
\handline[AnnotationBlue]{72}{548}{310}{\circled{3} 有限李型单群}
\handline[AnnotationBlue]{72}{581}{300}{\circled{4} $26$ 个散在单群}
\end{handpage}

% ===== Line-layout page-077.tex =====
\begin{handpage}{77}
\handtitle{72}{77}{240}{$\S3.5$ 群作用}
\handline{72}{105}{380}{\textbf{Def.} \circled{1} $G$：群，$X$：集合. $\mu:G\times X\to X$}
\handline{72}{129}{155}{$(g,x)\mapsto g*x$}
\handline{72}{161}{633}{满足 $1^\circ$ $e*x=x$；$2^\circ$ $g_1*(g_2*x)=(g_1g_2)*x$，则称 $\mu$ 为 $G$ 在 $X$ 上的作用}
\handline[AnnotationBlue]{72}{190}{260}{以下将 $g*x$ 简写为 $gx$}
\handline{72}{220}{398}{\circled{2} $g\in G$，$\rho_g:X\to X$，$x\mapsto gx$，称为平移}
\handline{72}{250}{546}{\circled{3} $x\in X$，$Gx:=\{gx\mid g\in G\}$，称为 $x$ 在 $G$ 作用下的轨道.}
\handline[AnnotationBlue]{72}{279}{532}{轨道要与陪集区分，但类似地，轨道要么相等，要么不交.}
\handline{72}{311}{713}{\circled{4} $\operatorname{Stab}(x):=\{g\in G\mid gx=x\}$，称为 $x$ 在 $G$ 中的稳定化子（或迷向子群）}
\handline[AnnotationBlue]{72}{341}{242}{显然 $\operatorname{Stab}(x)\le G$.}
\handline{72}{372}{665}{\circled{5} 若 $\operatorname{Stab}(x)=G$，即 $Gx=\{x\}$，则称 $x$ 为 $G$ 作用下的不动点.}
\handline{72}{405}{400}{\textbf{Lemma.} 存在双射：$G/\operatorname{Stab}(x)\to Gx$}
\handline{72}{437}{439}{\textbf{Pf.} 设 $H=\operatorname{Stab}(x)$. $\varphi:G/H\to Gx$}
\handline{72}{460}{138}{$gH\mapsto gx$}
\handline{72}{497}{553}{\circled{1} 良定：若 $g_1H=g_2H$，即 $g_2^{-1}g_1\in H\Rightarrow g_2^{-1}g_1x=x$}
\handline{72}{527}{502}{$\varphi(g_1H)=g_1x=g_2\cdot g_2^{-1}g_1x=g_2x=\varphi(g_2x)$}
\handline{72}{561}{221}{\circled{2} 满射：显然.}
\handline{72}{592}{666}{\circled{3} 单射：若 $g_1x=g_2x\Rightarrow(g_1^{-1}g_2)x=g_1^{-1}(g_2x)=g_1^{-1}(g_1x)=x$}
\handline{72}{623}{400}{所以 $g_1^{-1}g_2\in H=\operatorname{Stab}(x)$，所以 $g_1H=g_2H$.}
\handline{72}{657}{292}{所以 $\varphi$ 是一个双射.}
\handline{72}{696}{669}{\textbf{Cor.} $|Gx|=|G/\operatorname{Stab}(x)|=[G:\operatorname{Stab}(x)]=\dfrac{|G|}{|\operatorname{Stab}(x)|}$，且 $|Gx|\mid|G|$.}
\handline{72}{729}{383}{\textbf{Def.} 设 $O=\{Gx\mid x\in X\}\subseteq\mathcal P(X)$.}
\handline{72}{759}{380}{$\pi:X\to O$，$x\mapsto Gx$，则 $\pi$ 是满射.}
\handline{72}{789}{599}{由选择公理，存在选择函数 $\theta:O\to X$ 为单射，且 $\pi\circ\theta=\id_O$.}
\handline{72}{826}{524}{设 $C=\im(\theta)$，则 $C$ 与每个轨道 $Gx$ 有且仅有一个交点.}
\handline{72}{855}{254}{$C$ 称为一个截面.}
\handdiagram{81}{873}{412}{132}{
\draw (0,0) ellipse (3 and 1.15);
\foreach \i/\x in {1/-2,2/-.65,3/.65,4/2}{
\draw[AnnotationBlue] (\x-.2,-.9)..controls(\x-.3,0)and(\x+.1,.5)..(\x+.35,.95);
\node[AnnotationBlue] at(\x,1.45){$Gx_{\i}$};
\draw (\x,0) circle(.07);
\draw[->] (\x,-.1)--(\x-.3,-1.3);
\node at(\x-.3,-1.55){$\theta(Gx_{\i})$};}
\draw[NoteRed,thick](-2.7,-.25)..controls(-1,.6)and(1,.5)..(2.5,-.25);
\node[NoteRed]at(3.3,-.2){$C$};
}
\end{handpage}

% ===== Line-layout page-078.tex =====
\begin{handpage}{78}
\handline{72}{113}{640}{\textbf{Prop.}\quad 若 $x_1,x_2\in Gx$，则 $\exists g_0\in G$ s.t. $\operatorname{Stab}(x_2)=g_0\operatorname{Stab}(x_1)g_0^{-1}$.}
\handline{72}{138}{705}{从而 $|\operatorname{Stab}(x_1)|=|\operatorname{Stab}(x_2)|$ 且 $[G:\operatorname{Stab}(x_1)]=[G:\operatorname{Stab}(x_2)]$，$|Gx_1|=|Gx_2|$.}
\handline[AnnotationBlue]{72}{165}{500}{（当然这个数量关系也可直接由 $Gx_1=Gx_2=Gx$ 得出.}
\handline[AnnotationBlue]{72}{189}{725}{这一命题说明同一轨道上元素的稳定化子在相差一个共轭的意义下是唯一的）}
\handline{72}{221}{330}{\textbf{Pf.}\quad 设 $x_1=g_1x,x_2=g_2x$}
\handline{72}{249}{500}{$\therefore x_2=(g_2g_1^{-1}g_1)x=(g_2g_1^{-1})(g_1x)=g_0x_1$}
\handline{72}{279}{330}{$\forall g\in\operatorname{Stab}(x_2)\iff gx_2=x_2$.}
\handline{72}{309}{680}{$\iff g(g_0x_1)=g_0x_1\iff(g_0^{-1}gg_0)(x_1)=x_1\iff g_0^{-1}gg_0\in\operatorname{Stab}(x_1)$.}
\handline{72}{338}{605}{$\iff g\in g_0\operatorname{Stab}(x_1)g_0^{-1}$，$\therefore\operatorname{Stab}(x_2)=g_0\operatorname{Stab}(x_1)g_0^{-1}$.}
\handline{72}{370}{510}{\textbf{Cor.}\quad$\displaystyle|\bar X|=\sum_{x\in C}[G:\operatorname{Stab}(x)]$，$C$ 是一个截面.}
\handline{72}{402}{610}{\textbf{Pf.}\quad$\displaystyle\bar X=\coprod_{x\in C}Gx$，$\displaystyle\therefore|\bar X|=\sum_{x\in C}|Gx|=\sum_{x\in C}[G:\operatorname{Stab}(x)]$.}
\handline{72}{455}{515}{\textbf{Def.}\quad 考虑 $G$ 在自身上的共轭作用.$\mu:G\times G\to G$}
\handline{72}{481}{170}{$(g,x)\longmapsto gxg^{-1}$.}
\handline{72}{517}{570}{\circled{1} 则此时 $x$ 的轨道 $Gx=\{gxg^{-1}\mid g\in G\}$ 称为 $x$ 的共轭类.}
\handline{72}{549}{765}{\circled{2} $x$ 的稳定化子 $\operatorname{Stab}(x)=\{g\in G\mid gxg^{-1}=x\}=C(x)$ 称为 $x$ 的中心化子.}
\handline{72}{583}{680}{\circled{3} $Gx=\{x\}$ 称为平凡共轭类，若 $|Gx|\geq2$，则称为非平凡共轭类.}
\handline{72}{615}{675}{若 $x$ 是不动点，即 $Gx=\{x\}$，即 $\forall g\in G,gxg^{-1}=x$，即 $x\in Z(G)$.}
\handline{72}{647}{635}{$\therefore G$ 在自身上的共轭作用的不动点集即为 $G$ 的中心 $Z(G)$}
\handline{72}{683}{610}{取 $C$ 为一个截面，则 $C$ 与每个轨道 $Gx$ 有且仅有一个交点.}
\handline{72}{715}{375}{$\therefore Z(G)\subseteq C$.令 $\Gamma=C\setminus Z(G)$.}
\handline{72}{747}{500}{则 $\Gamma$ 即为非平凡共轭类的代表元构成的集合}
\handline{72}{783}{640}{于是有 $\displaystyle|G|=\sum_{x\in C}[G:\operatorname{Stab}(x)]=|Z(G)|+\sum_{x\in\Gamma}[G:C(x)]$.}
\handline[AnnotationBlue]{72}{813}{365}{此方程称为 $G$ 的类方程.}
\handline{72}{845}{465}{\textbf{Ex.}\quad 若 $|G|=p^m$，$p$ 素，则 $G$ 称为 $p$ 群.}
\handline{72}{882}{795}{$\displaystyle\therefore p^m=|Z(G)|+\sum_{x\in\Gamma}[G:C(x)]$.由 $C(x)\subsetneq G\Rightarrow[G:C(x)]>1$，而 $[G:C(x)]\mid p^m$}
\handline{72}{923}{655}{$\therefore p\mid[G:C(x)]\Rightarrow p\mid|Z(G)|\Rightarrow Z(G)$ 中的元素不仅 $e$}
\handline[NoteRed]{72}{962}{360}{即 $p$ 群一定有非平凡的中心元.}
\end{handpage}

% ===== Line-layout page-079.tex =====
\begin{handpage}{79}
\handline{72}{149}{320}{\textbf{Lemma.}\quad\ann{(Burnside's Lemma)}}
\handline{72}{175}{520}{设 $G$ 作用在 $\bar X$ 上有 $n$ 个轨道，则 $\displaystyle n=\frac1{|G|}\sum_{g\in G}|\bar X^g|$}
\handline{72}{211}{390}{其中 $\bar X^g:=\{x\in\bar X\mid gx=x\}$，$\forall g\in G$.}
\handline{72}{238}{520}{\textbf{Pf.}\quad 固定 $g$，满足 $gx=x$ 的 $x$ 的个数为 $|\bar X^g|$.}
\handline{72}{266}{650}{$g$ 跑遍 $G$ 时 $\displaystyle\sum_{g\in G}|\bar X^g|$ 即为满足 $gx=x$ 的元素对 $(g,x)$ 的个数.}
\handline{72}{299}{640}{另一方面，先固定 $x$，则满足 $gx=x$ 的 $g$ 的个数为 $|\operatorname{Stab}(x)|$.}
\handline{72}{327}{450}{$\displaystyle\therefore\sum_{g\in G}|\bar X^g|=\sum_{x\in\bar X}|\operatorname{Stab}(x)|=\sum_{x\in\bar X}\frac{|G|}{|Gx|}$.}
\handline{72}{360}{255}{取截面 $C$，则 $|C|=n$.}
\handline{72}{396}{605}{$\displaystyle\therefore\sum_{x\in\bar X}\frac{|G|}{|Gx|}=\sum_{x\in C}\left(\sum_{y\in Gx}\frac{|G|}{|Gx|}\right)=\sum_{x\in C}\left(|Gx|\cdot\frac{|G|}{|Gx|}\right)=n|G|$.}
\handline{72}{446}{250}{$\displaystyle\therefore n=\frac1{|G|}\sum_{g\in G}|\bar X^g|$}
\handline[AnnotationBlue]{72}{481}{690}{Burnside 是第一位写有限群论教材的.这个引理就在那本书中.}
\handline[AnnotationBlue]{72}{514}{735}{但实际上 Cauchy 和 Frobenius 也证明过，有人称之为 Cauchy 引理.}
\handline[AnnotationBlue]{72}{548}{615}{Burnside 提出了 Burnside 猜想：奇数阶群一定可解.}
\handline[AnnotationBlue]{72}{581}{520}{该猜想由 Feit--Thompson 解决.（1963）}
\end{handpage}

% ===== Line-layout page-080.tex =====
\begin{handpage}{80}
\handtitle{72}{116}{280}{§3.6 Sylow 定理}
\handline{72}{144}{480}{Sylow 定理可视为 Lagrange 定理的反问题.}
\handline{72}{173}{480}{即若 $d\mid|G|$，是否存在 $H\leq G$ 满足 $|H|=d$}
\handline{72}{200}{665}{一般来说它不对，但 Sylow 定理表明 $d=p^k$，$p$ 素时命题成立.}
\handline{72}{227}{200}{\textbf{Lemma.}\quad\ann{(Cauchy)}}
\handline{72}{257}{410}{设 $G$ 为有限交换群，$p$ 是素数，$p\mid|G|$.}
\handline{72}{285}{680}{则 $G$ 中有 $p$ 阶元 $a$.于是 $\langle a\rangle=\{e,a,\ldots,a^{p-1}\}$ 是 $G$ 的一个 $p$ 阶子群.}
\handline{72}{317}{460}{\textbf{Pf.}\quad 对 $|G|$ 作归纳.$|G|=1,2:$ trivial}
\handline{72}{348}{530}{设对阶数 $<|G|$ 的交换群结论成立，下面考虑 $G$.}
\handline{72}{376}{385}{$\forall a\in G,a\ne e$.设 $a$ 的阶数为 $m$.}
\handline{72}{406}{530}{\circled{1} 若 $p\mid m$，令 $b=a^r$，则 $b^p=a^{pr}=e$，$b$ 即为所求.}
\handline{72}{439}{470}{\circled{2} 若 $p\nmid m$，设 $H=\langle a\rangle=\{e,a,\ldots,a^{m-1}\}$.}
\handline{72}{470}{420}{由于 $G$ 交换，$\therefore$ 任意子群都正规，$\therefore H\triangleleft G$}
\handline{72}{503}{640}{于是可定义商群 $G/H$，$|G/H|=\dfrac{|G|}{m}$.而 $p\nmid m$ 且 $p$ 素 $\Rightarrow p\mid\dfrac{|G|}{m}$.}
\handline{72}{534}{445}{而 $\dfrac{|G|}{m}<|G|$，由归纳假设 $G/H$ 有 $p$ 阶元}
\handline{72}{565}{465}{$\therefore\exists bH\in G/H$ s.t. $(bH)^p=H\Rightarrow b^p\in H$.}
\handline{72}{602}{755}{设 $b$ 的阶数为 $n$，则 $(bH)^n=b^nH=H=e_{G/H}$，$\therefore p\mid n$，化为第\circled{1}种情形.}
\handline{72}{635}{255}{取 $c=b^{n/p}$ 即为所求.}
\end{handpage}

% ===== Line-layout page-081.tex =====
\begin{handpage}{81}
\handline{72}{128}{210}{\textbf{Thm.}\quad\ann{(Sylow-I)}}
\handline{72}{157}{505}{$G$ 为有限群，$p$ 素，$p^k\mid|G|$，则 $G$ 必有 $p^k$ 阶子群}
\handline{72}{184}{350}{\textbf{Pf.}\quad 对 $|G|$ 归纳，$|G|=1:$ trivial}
\handline{72}{212}{405}{设阶数 $<|G|$ 时结论成立，下证 $|G|$ 时}
\handline{72}{241}{440}{考虑类方程 $\displaystyle|G|=|Z(G)|+\sum_{x\in\Gamma}[G:C(x)]$}
\handline{72}{270}{470}{\circled{1} 若 $p\nmid|Z|$，则 $\exists x\in\Gamma$ s.t. $p\nmid[G:C(x)]$.}
\handline{72}{300}{460}{但 $p^k\mid|G|=[G:C(x)]\cdot|C(x)|\Rightarrow p^k\mid|C(x)|$.}
\handline{72}{329}{405}{而 $C(x)$ 为 $G$ 的真子群，由归纳假设}
\handline{72}{359}{440}{$C(x)$ 有一个 $p^k$ 阶子群，同时也是 $G$ 的子群}
\handline{72}{390}{380}{\circled{2} 若 $p\mid|Z(G)|$，而 $Z(G)$ 为交换群}
\handline{72}{420}{490}{由 Cauchy 引理，$Z(G)$ 有一个 $p$ 阶子群 $H=\langle a\rangle$.}
\handline{72}{451}{630}{由于 $H\subseteq Z(G)$，故 $H$ 中元素与 $G$ 中任意元素可交换，$\therefore H\triangleleft G$.}
\handline{72}{484}{375}{考虑 $G'=G/H$，则 $|G'|<|G|$.}
\handline{72}{512}{300}{而 $|H|=p$，$\therefore p^{k-1}\mid|G'|$.}
\handline{72}{542}{420}{由归纳假设，$G'$ 中存在 $p^{k-1}$ 阶子群 $K'$.}
\handline{72}{573}{505}{由子群对应定理，$G$ 中存在子群 $K$ 与之对应}
\handline{72}{608}{495}{即存在唯一子群 $K$ s.t. $H\triangleleft K$，$K/H=K'$}
\handline{72}{642}{340}{$\therefore|K|=|H|\cdot|K/H|=p\cdot p^{k-1}=p^k$.}
\end{handpage}

% ===== Line-layout page-082.tex =====
\begin{handpage}{82}
\handline{72}{105}{400}{\textbf{Def.}\quad\circled{1} $|H|=p^k$，$p$ 素，则 $H$ 称为 $p$ 群}
\handline{72}{133}{675}{\circled{2} 若 $|G|<\infty$，$p$ 素数，$p$ 在 $|G|$ 中的阶为 $m$，即 $p^m\mid|G|$ 但 $p^{m+1}\nmid|G|$.}
\handline{72}{160}{420}{则 $G$ 中的 $p^m$ 阶子群称为 $p$-Sylow 子群}
\handline{72}{192}{210}{\textbf{Thm.}\quad\ann{(Sylow-II)}}
\handline{72}{218}{500}{$|G|<\infty$，$p$ 素，则 $G$ 的任意两个 $p$-Sylow 子群共轭.}
\handline{72}{248}{570}{即 $\forall P_1,P_2$ 为 $p$-Sylow 子群，$\exists g\in G$ s.t. $P_2=gP_1g^{-1}$.}
\handline{72}{278}{620}{\textbf{Pf.}\quad 设 $P$ 为 $p$-Sylow 子群，$H$ 为 $p$ 子群.$|H|=p^k$，$\forall k\leq m$.}
\handline{72}{307}{390}{考虑群作用 $\mu:H\times(G/P)\longrightarrow G/P$}
\handline{72}{335}{215}{$(k,gP)\longmapsto kgP$}
\handline{72}{369}{390}{$|P|=p^m\Rightarrow|G/P|=\dfrac{|G|}{p^m}$ 且 $p\nmid|G/P|$.}
\handline{72}{398}{250}{考虑 $\mu$ 的轨道 $H(gP)$}
\handline{72}{430}{500}{若 $gP$ 是不动点，即 $H(gP)=\{gP\}\iff|H(gP)|=1$}
\handline{72}{460}{680}{若 $gP$ 不是不动点，由于 $|H(gP)|\mid|H|=p^k\Rightarrow|H(gP)|=p^l$，$1\leq l\leq k$.}
\handline{72}{490}{425}{由于 $p\nmid|G/P|$，$\displaystyle|G/P|=\sum_{gP\in C}|H(gP)|$.}
\handline{72}{523}{555}{$\therefore$ 必 $\exists gP\in G/P$ s.t. $|H(gP)|=1$，即 $H(gP)=\{gP\}$.}
\handline{72}{554}{510}{$\therefore\forall h\in H$，$(h,gP)=hgP=gP\Rightarrow g^{-1}hg\in P$}
\handline{72}{586}{365}{$\therefore g^{-1}Hg\subseteq P\iff H\subseteq gPg^{-1}$.}
\handline{72}{618}{415}{当 $k=m$ 即 $H$ 也为 $p$-Sylow 子群时，}
\handline{72}{655}{595}{$|H|=p^m=|P|=|gPg^{-1}|$，而 $H\subseteq gPg^{-1}\Rightarrow H=gPg^{-1}$.}
\handline{72}{689}{855}{\textbf{Cor.}\quad$|G|<\infty$，$p$ 素，$P$ 为 $p$-Sylow 子群，则 $P$ 正规 $\iff P$ 是唯一的 $p$-Sylow 子群.}
\handline{72}{720}{545}{\textbf{Pf.}\quad“$\Rightarrow$”若 $P\triangleleft G$，设 $Q$ 是另一个 $p$-Sylow 子群}
\handline{72}{757}{615}{则 $\exists g\in G$ s.t. $Q=gPg^{-1}$，而 $gPg^{-1}=P$，$\therefore P=Q$.}
\handline{72}{787}{400}{“$\Leftarrow$”若 $P$ 是唯一的 $p$-Sylow 子群，}
\handline{72}{827}{620}{则 $\forall g\in G$，$gPg^{-1}$ 也为一个子群，且 $|gPg^{-1}|=|P|$}
\handline{72}{859}{345}{$\therefore gPg^{-1}$ 为 $p$-Sylow 子群}
\handline{72}{889}{325}{由唯一性 $\Rightarrow gPg^{-1}=P$}
\handline{72}{924}{150}{$\therefore P\triangleleft G$.}
\end{handpage}

% ===== Line-layout page-083.tex =====
\begin{handpage}{83}
\handline{72}{97}{195}{\textbf{Thm.}\quad\ann{(Sylow-III)}}
\handline{72}{126}{700}{$|G|<\infty$，$p$ 是素数.设 $G$ 的 $p$-Sylow 子群的个数为 $n_p$，$P$ 为 $p$-Sylow 子群}
\handline{72}{153}{330}{则 $n_p\mid m$ 且 $n_p\equiv1\pmod p$.}
\handline{72}{181}{670}{其中 $m=\dfrac{|G|}{p^k}$，$k$ 是 $p$ 在 $|G|$ 中的阶，即 $p\nmid m$.而 $m=\dfrac{|G|}{p^k}=\dfrac{|G|}{|P|}=[G:P]$.}
\handline{72}{209.049}{460}{\textbf{Pf.}\quad\circled{1} 设 $G$ 的所有 $p$-Sylow 子群构成的集合为 $S$}
\handline{72}{233}{310}{考虑群作用 $\mu:G\times S\longrightarrow S$}
\handline{72}{259}{190}{$(g,H)\longmapsto gHg^{-1}$.}
\handline{72}{289}{645}{$\forall P$ 为 $p$-Sylow 子群，即 $P\in S$.考虑 $\operatorname{Stab}(P)=\{g\in G\mid gPg^{-1}=P\}=N_G(P)$.}
\handline{72}{320}{335}{即 $P$ 在 $G$ 中的正规化子.}
\handline{72}{347}{550}{考虑 $P$ 的轨道 $G(P)$，则结合 Sylow-II 可知 $|G(P)|=n_p$.}
\handline{72}{376}{475}{$\therefore n_p=|G(P)|=[G:\operatorname{Stab}(P)]=[G:N_G(P)]$}
\handline{72}{404}{660}{而 $P\leq N_G(P)\leq G$，$\therefore[G:N_G(P)][N_G(P):P]=[G:P]$，$\therefore n_p\mid m$}
\handline{72}{437}{420}{\circled{2} 考虑另一个群作用 $\sigma:P\times(G/P)\longrightarrow G/P$.}
\handline{72}{466}{220}{$(h,gP)\longmapsto hgP$.}
\handline{72}{502}{485}{若 $gP$ 是 $\sigma$ 的不动点，即 $\operatorname{Orb}_P(gP)=\{gP\}$.}
\handline{72}{535}{470}{$\Rightarrow\forall h\in P$，$hgP=gP$，$\therefore g^{-1}hgP=P$.}
\handline{72}{567}{535}{$\therefore g^{-1}hg\in P\Rightarrow h\in gPg^{-1}$，$\therefore P\subseteq gPg^{-1}$.}
\handline{72}{597}{530}{而 $|P|=|gPg^{-1}|$，$\therefore P=gPg^{-1}$，$\therefore g\in N_G(P)$.}
\handline{72}{630}{635}{反之若 $g\in N_G(P)$，则 $P=gPg^{-1}\Rightarrow gP$ 是 $\sigma$ 的不动点.}
\handline{72}{662}{470}{$\therefore\sigma$ 的不动点的个数即为 $[N_G(P):P]$.}
\handline{72}{696}{760}{若 $gP$ 不是不动点，则 $|\operatorname{Orb}_P(gP)|\mid|P|=p^k\Rightarrow|\operatorname{Orb}_P(gP)|=p^l$，$1\leq l\leq k$.}
\handline{72}{731}{265}{而 $\displaystyle|G/P|=\sum_{O\text{ 为轨道}}|O|$}
\handline{72}{766}{590}{两边同时 $\bmod p\Rightarrow m\equiv[N_G(P):P]\pmod p$}
\handline{72}{806}{750}{$m\cdot n_p\equiv[N_G(P):P]\cdot n_p=[N_G(P):P]\cdot[G:N_G(P)]=[G:P]=m\pmod p$}
\handline{72}{849}{445}{而 $(m,p)=1$，$\therefore n_p\equiv1\pmod p$.}
\end{handpage}

% ===== Line-layout page-084.tex =====
\begin{handpage}{84}
\handtitle{72}{106}{270}{Sylow 定理的应用：}
\handline{72}{136}{510}{\circled{1} 设 $p,q$ 为素数，$p<q$，若 $|G|=pq$，则 $G$ 不是单群.}
\handline{72}{165}{605}{考虑 $G$ 的 $q$-Sylow 子群，则 $n_q\mid p$ 且 $n_q\equiv1\pmod q\Rightarrow n_q=1$.}
\handline{72}{193}{505}{即 $G$ 只有一个 $q$-Sylow 子群，故正规，$\therefore G$ 不是单群.}
\handline{72}{221}{555}{\circled{2} $2024$ 阶群是可解群.$2024=2^3\times11\times23$.$|G|=2024$.}
\handline{72}{251}{610}{从较大的因子考虑，$n_{23}\mid2^3\times11$，且 $n_{23}\equiv1\pmod{23}\Rightarrow n_{23}=1$}
\handline{72}{282}{495}{故 $G$ 只有一个 $23$-Sylow 子群 $P_{23}$，故 $P_{23}\triangleleft G$.}
\handline{72}{310}{325}{令 $G_1=G/P_{23}$，$|G_1|=2^3\times11$.}
\handline{72}{342}{725}{$n_{11}\mid8$ 且 $n_{11}\equiv1\pmod{11}\Rightarrow n_{11}=1$，$\therefore G_1$ 只有一个 $11$-Sylow 子群 $P_{11}$ 且正规.}
\handline{72}{372}{335}{令 $G_2=G_1/P_{11}$，$|G_2|=8$.}
\handline[AnnotationBlue]{72}{401}{420}{（若 $|G|=p^n$，$H\leq G$，且 $|H|=p^{n-1}$，则 $H\triangleleft G$）}
\handline{72}{431}{720}{由 Sylow-I $\Rightarrow G_2$ 有 $4$ 阶子群 $P_4$ 且正规，而 $P_4$ 又有 $2$ 阶子群 $P_2$ 且正规.}
\handline{72}{460}{585}{$\therefore G_2\triangleright P_4\triangleright P_2\triangleright\{e\}$，每一步的商群都是 $\Z_2$，故 $G_2$ 可解.}
\handline{72}{492}{590}{结合命题\ann{（若 $N\triangleleft G$，则 $G$ 可解 $\iff N$ 和 $G/N$ 均可解）}}
\handline{72}{523}{430}{以及素数阶循环群均可解，$\Rightarrow G$ 可解.}
\handline{72}{554}{315}{\circled{3} $|G|=56$，则 $G$ 不是单群.}
\handline{72}{588}{515}{$56=7\times2^3$.$n_7\mid8$ 且 $n_7\equiv1\pmod7\Rightarrow n_7=1$ 或 $8$}
\handline{72}{622}{515}{$1^\circ$ 若 $n_7=1$，则 $7$-Sylow 子群正规，$G$ 不是单群}
\handline{72}{658}{800}{$2^\circ$ 若 $n_7=8$，由 Sylow-II $\Rightarrow$ 它们之间相互共轭，且要么相等，要么交集为 $\{e\}$.}
\handline{72}{696}{690}{于是这 $8$ 个 $7$ 阶子群中贡献出 $1+8\times6=49$ 个元素，还剩 $7$ 个元素.}
\handline{72}{730}{455}{而 $G$ 有 $2$-Sylow 子群 $H$，则 $|H|=8$.}
\handline{72}{766}{705}{由于 $H$ 中元素的阶一定为 $8$ 的因子，故一定不在 $8$ 个 $7$ 阶群中.}
\handline{72}{801}{545}{故 $H\setminus\{e\}$ 中的 $7$ 个元素恰为剩下的 $7$ 个.}
\handline{72}{836}{595}{$\therefore2$-Sylow 子群只有一个，故正规，$\therefore G$ 不是单群.}
\end{handpage}

% ===== Line-layout page-085.tex =====
\begin{handpage}{85}
\handline{72}{145}{809}{\circled{4} \textbf{Prop.} 若 $|G|=p^2$，$p$ 素，则 $G$ 交换.}
\handline{72}{176}{792}{\textbf{Pf.} 考虑类方程 $|G|=|Z(G)|+\sum_{y\in T}[G:C(y)]$. $C$ 为一个截面，$T=C-Z(G)$.}
\handline{72}{204}{718}{由 Sylow-II 的中间过程 $\Rightarrow|Z(G)|>1$.}
\handline{72}{235}{707}{而 $p\mid[G:C(y)]$，$\therefore p\mid|Z(G)|$，$\therefore Z(G)=p$ 或 $p^2$.}
\handline{72}{266}{719}{若 $Z(G)=p^2$ 则 $G=Z(G)$ 为交换群.}
\handline{72}{296}{719}{若 $Z(G)=p$，则 $Z(G)\subsetneq G$，$\therefore\exists a\in G$，但 $a\notin Z(G)$.}
\handline{72}{328}{712}{考虑 $a$ 的中心化子 $C(a)$，则 $Z(G)\cup\{a\}\subseteq C(a)$.}
\handline{72}{357}{693}{而 $|C(a)|>|Z(G)|=p$，$\therefore|C(a)|=p^2$.}
\handline{72}{394}{719}{$\therefore C(a)=G\Rightarrow G$ 交换，$\therefore Z(G)=G$，矛盾！}
\handline{72}{428}{723}{$\therefore$ 综上 $G$ 为交换群.}
\handline{72}{462}{759}{若 $|G|=11^2\times13^2$ 则 $G$ 交换.}
\handline{72}{515}{728}{$\begin{cases}n_{11}\mid13^2,\\n_{11}\equiv1\pmod{11},\end{cases}\quad\begin{cases}n_{13}\mid11^2,\\n_{13}\equiv1\pmod{13},\end{cases}\quad\Rightarrow n_{11}=n_{13}=1$.}
\handline{72}{555}{749}{$\therefore G$ 有且仅有一个 $11$-Sylow 子群、$13$-Sylow 子群 $H,K$.}
\handline{72}{587}{743}{且 $H\lhd G$，$K\lhd G$，$|H|=11^2$，$|K|=13^2$.}
\handline{72}{623}{746}{$\therefore H,K$ 是交换群.}
\handline{72}{656}{743}{利用元素的阶进行分析 $\Rightarrow H\cap K=\{e\}$.}
\handline{72}{696}{747}{由于 $H,K$ 正规，$\therefore HK\le G$. 而 $|HK|=\frac{|H|\cdot|K|}{|H\cap K|}=11^2\times13^2=|G|$.}
\handline{72}{734}{746}{$\therefore HK=G$. 下面只需证 $\forall x\in H$，$y\in K$，$xy=yx$ 即可.}
\handline{72}{771}{765}{姚慕生 \S2.3 习题4：\ann{若 $N\lhd G$，$H\lhd G$，$N\cap H=\{e\}$，则 $\forall x\in N$，$y\in H$，$xy=yx$.}}
\handline[AnnotationBlue]{72}{807}{744}{$\forall x\in N$，$y\in H$，考虑交换子 $[x,y]=x^{-1}y^{-1}xy$.}
\handline[AnnotationBlue]{72}{847}{737}{$x^{-1}\cdot(y^{-1}xy)$，由正规 $\Rightarrow y^{-1}xy\in N$，$\therefore x^{-1}(y^{-1}xy)\in N$.}
\handline[AnnotationBlue]{72}{885}{745}{$(x^{-1}y^{-1}x)\cdot y$，由正规 $\Rightarrow x^{-1}y^{-1}x\in H$，$\therefore(x^{-1}y^{-1}x)y\in H$.}
\handline[AnnotationBlue]{72}{924}{750}{$\therefore[x,y]\in N\cap H=\{e\}$，$\therefore xy=yx$.}
\end{handpage}

% ===== Line-layout page-086.tex =====
\begin{handpage}{86}
\handline{72}{92}{802}{\circled{5} $|G|=108=2^2\cdot3^3$ 则 $G$ 不是单群.}
\handline{72}{119}{798}{\textbf{Pf.} $n_3\mid4$，$n_3\equiv1\pmod3\Rightarrow n_3=1$ 或 $4$.}
\handline{72}{148}{726}{$n_3=1$：trivial. 下面考虑 $n_3=4$.}
\handline{72}{178}{732}{$G$ 有 $4$ 个 $3$-Sylow 子群，即 $27$ 阶群.}
\handline{72}{205}{722}{设 $H,K$ 为不同 $3$-Sylow 子群，则 $|H|=|K|=27$.}
\handline{72}{235}{715}{$HK\subseteq G\Rightarrow|HK|\le108$. \ann{（$HK$ 只是一个集合，由于 $H,K$ 均不正规，$HK$ 不一定是子群.）}}
\handline{72}{270}{722}{$|HK|=\frac{|H||K|}{|H\cap K|}=\frac{27^2}{|H\cap K|}\le108\Rightarrow|H\cap K|\ge\frac{27}4$，$\therefore|H\cap K|\ge7$.}
\handline{72}{303}{724}{而 $H\cap K$ 一定为 $H,K$ 的子群，$\therefore|H\cap K|\mid|H|=27$，但 $|H\cap K|<27$.}
\handline{72}{334}{729}{$\therefore|H\cap K|=9$.}
\handline{72}{366}{732}{由于 $|H\cap K|=3^2$，$|H|=|K|=3^3$，$H\cap K\le H,K$.}
\handline{72}{397}{731}{$\therefore H\cap K\lhd H$，$H\cap K\lhd K$.}
\handline{72}{428}{732}{设 $N$ 是 $H\cap K$ 在 $G$ 中的正规化子}
\handline{72}{465}{734}{则 $H\le N$，$K\le N$，$HK\subseteq N$. 而 $|HK|=\frac{|H|\cdot|K|}{|H\cap K|}=\frac{27^2}{9}=81$.}
\handline{72}{502}{717}{$\therefore|N|\ge81$. 而 $N\le G$，$|N|\mid|G|=108$，$\therefore|N|=|G|=108$.}
\handline{72}{534}{716}{$\therefore N=G\Rightarrow H\cap K\lhd G$，$\therefore G$ 不是单群.}
\handline[AnnotationBlue]{72}{567}{741}{本题得到 $4$ 个 $3$-Sylow 子群，故一定不正规；}
\handline[AnnotationBlue]{72}{602}{736}{但我们证明了两两之交是正规的. \redann{（但并不是任意 $p$-Sylow 子群之交都正规.）}}
\handline{72}{639}{829}{\textbf{Thm.} 最小的非交换单群是 $60$ 阶群（$A_5$，$5$ 元偶置换全体）.}
\handline[AnnotationBlue]{72}{674}{755}{此定理的证明只能 Case by Case，可作为练习.}
\handline[AnnotationBlue]{72}{705}{756}{只有 $24$ 和 $36$ 阶群的方法没有提到.}
\handline[AnnotationBlue]{72}{740}{752}{利用姚慕生 \S2.7 习题2：}
\handline[AnnotationBlue]{72}{780}{756}{$|G|<\infty$，$H\le G$. 若 $[G:H]=n$，则 $H$ 必含有 $G$ 的正规子群 $K$ s.t. $[G:K]\mid n!$.}
\handline[AnnotationBlue]{72}{817}{758}{若 $|G|\nmid n!$，则 $G$ 有一个非平凡正规子群.}
\handline{72}{854}{755}{$|G|=24=2^3\times3$.}
\handline{72}{890}{757}{$n_2\mid3$，$n_2\equiv1\pmod2\Rightarrow n_2=1$ 或 $3$.}
\handline{72}{930}{762}{$n_2=3$ 时，有 $8$ 阶子群 $H$. 而 $[G:H]=3$，$24\nmid3!$.}
\handline{72}{973}{765}{$\therefore G$ 有一个非平凡正规子群，$\therefore G$ 不是单群. $36$ 阶群同理.}
\end{handpage}

% ===== Line-layout page-087.tex =====
\begin{handpage}{87}
\handtitle{72}{109}{776}{\S3.7 对称群}
\handline{72}{140}{778}{\textbf{Def.} $X$：集合.}
\handline{72}{168}{734}{\circled{1} $S(X):=\{f\colon X\to X\mid f\text{ 是双射}\}$.}
\handline{72}{195}{706}{$(S(X),\circ,\id)$ 称为 $X$ 的对称群.}
\handline{72}{226}{739}{\circled{2} 若 $G\le S(X)$，则 $G$ 称为 $X$ 上的置换群.}
\handline{72}{256}{740}{\circled{3} $G$：群，$X$：集合，$\mu\colon G\times X\to X$ 为一个群作用}
\handline{72}{288}{710}{则 $(G,X)$ 这个具有群作用的系统称为变换群.}
\handline{72}{321}{797}{\textbf{Thm.} \ann{（Cayley）}}
\handline{72}{353}{776}{任意群 $G$ 都同构于某个置换群.}
\handline{72}{385}{798}{\textbf{Pf.} 考虑左平移 $L_g\colon G\times G\to G$}
\handline{72}{413}{507}{$(g,h)\mapsto gh$.}
\handline{72}{453}{696}{若 $L_g(h_1)=L_g(h_2)$ 则 $gh_1=gh_2\Rightarrow h_1=h_2$，$\therefore L_g$ 是单射.}
\handline{72}{489}{693}{$\forall h\in G$，考虑 $g^{-1}h$，则 $L_g(g^{-1}h)=h$，$\therefore L_g$ 是满射.}
\handline{72}{522}{696}{$\therefore L_g\in S(G)$.}
\handline{72}{555}{230}{令 $\varphi\colon G\to S(G)$}
\handline{72}{585}{128}{$g\mapsto L_g$.}
\handline{331}{552}{432}{$\varphi(g_1g_2)=L_{g_1g_2}=L_{g_1}\circ L_{g_2}=\varphi(g_1)\varphi(g_2)$.}
\handline{72}{622}{684}{若 $\varphi(g_1)=\varphi(g_2)$，即 $L_{g_1}=L_{g_2}$.}
\handline{72}{654}{684}{$\therefore\forall h\in G$，$g_1h=g_2h\Rightarrow g_1=g_2$.}
\handline{72}{688}{693}{$\therefore\varphi$ 是单同态. $\therefore G\cong\varphi(G)$.}
\handline{72}{721}{696}{而 $\varphi(G)\le S(G)$ 为置换群.}
\handline{72}{762}{811}{下面研究有限集的置换群. 取 $X=\{1,2,\ldots,n\}$，记 $S_n=S(X)$.}
\handline{72}{798}{803}{$S_n$ 称为 $n$ 阶对称群，则 $|S_n|=n!$.}
\handline[AnnotationBlue]{72}{839}{805}{以 $n=3$ 为例，$\sigma\in S_3$，$\sigma=\begin{pmatrix}1&2&3\\2&1&3\end{pmatrix}\Longleftrightarrow\begin{cases}\sigma(1)=2,\\\sigma(2)=1,\\\sigma(3)=3.\end{cases}$}
\handline[AnnotationBlue]{72}{910}{801}{但这种表达置换的书写方式不够简洁.}
\handline[AnnotationBlue]{72}{952}{793}{下面我们力求找到置换的某些性质来简化书写.}
\end{handpage}

% ===== Line-layout page-088.tex =====
\begin{handpage}{88}
\handline{72}{114}{823}{$\forall\sigma\in S_n$，考虑 $H=\langle\sigma\rangle=\{e,\sigma,\sigma^2,\ldots,\sigma^{m-1}\}$，$m$ 为 $\sigma$ 的阶.}
\handline{72}{143}{812}{考虑群作用 $H\times X\to X$，$(\sigma,i)\mapsto\sigma(i)$，$1\le i\le n$.}
\handline{72}{174}{809}{将 $X$ 分解为一些轨道的并.}
\handline{72}{202}{806}{$\forall i\in X$，$H_i=\{i,\sigma(i),\sigma^2(i),\ldots\}$.}
\handline{72}{233}{800}{设 $k$ 是使得 $\sigma^k(i)=i$ 的最小正整数.}
\handline{72}{262}{796}{则 $i,\sigma(i),\ldots,\sigma^{k-1}(i)$ 互不相同，否则与 $k$ 的最小性矛盾.}
\handline{72}{293}{802}{$\therefore H_i=\{i,\sigma(i),\ldots,\sigma^{k-1}(i)\}$.}
\handline{72}{323}{799}{而 $\operatorname{Stab}(i)=\{\varphi\in H\mid\varphi(i)=i\}$.}
\handline{72}{353}{792}{首先 $\sigma^k\in\operatorname{Stab}(i)$；$\forall1\le j\le k-1$，$\sigma^j\notin\operatorname{Stab}(i)$. $\therefore\operatorname{Stab}(i)$ 中的元素形如 $\sigma^{kj}$，$j\in\N$.}
\handline{72}{386}{801}{由于 $|\operatorname{Stab}(i)|\cdot|H_i|=|H|$，$\therefore|\operatorname{Stab}(i)|=m/k$，设为 $l$，则 $m=kl$.}
\handline{72}{416}{802}{从而 $\operatorname{Stab}(i)=\{e,\sigma^k,\ldots,\sigma^{k(l-1)}\}$.}
\handline{72}{449}{795}{对于 $i$ 的轨道 $H_i=\{i,\sigma(i),\ldots,\sigma^{k-1}(i)\}$.}
\handline{72}{481}{794}{设 $\sigma^t(i)=i_{t+1}$，$i=i_1$，$0\le t\le k-1$，则 $\sigma^k(i)=i_1$.}
\handline{72}{514}{790}{于是 $H_i$ 可记为 $(i_1,i_2,\ldots,i_k)$.}
\handline{72}{545}{794}{$\therefore(\sigma(i_1),\sigma(i_2),\ldots,\sigma(i_k))=(i_2,i_3,\ldots,i_k,i_1)$.}
\handline{72}{579}{788}{也就是 $\sigma$ 每在 $H_i$ 上作用一次，每个 $i_t$ 向前平移一次.}
\handline{72}{611}{794}{我们把 $(i_1,\ldots,i_k)$ 称为一个循环.}
\handline{72}{645}{798}{将 $i$ 取遍 $1\sim n$，由于轨道之间要么相等，要么不交，故 $\sigma\in S_n$ 可与下式一一对应：}
\handline{72}{681}{791}{$(i_1,\ldots,i_{k_1})(i_{k_1+1},\ldots,i_{k_1+k_2})\cdots(i_{k_1+k_2+\cdots+k_{l-1}+1},\ldots,i_{k_1+\cdots+k_l})$. \ann{$(*)$}}
\handline{72}{713}{786}{而 $(i_1,i_2,\ldots,i_n)$ 是 $1\sim n$ 的一个排列.}
\handline{72}{745}{782}{$k_1,k_2,\ldots,k_l$ 称为每个循环的长度.}
\handline[AnnotationBlue]{72}{782}{796}{以后就将 $(*)$ 式作为一个置换的表达方式，称为循环表达式.}
\handline[AnnotationBlue]{72}{817}{784}{若某个循环的长度为 $1$，$(i)$，说明 $i$ 是 $\sigma$ 的不动点，省略不写.}
\handline[AnnotationBlue]{72}{852}{786}{也就是说在 $\sigma$ 的表达式中只保留 $k_j\ge2$ 的部分.}
\handline{72}{886}{856}{\textbf{Def.} 若 $\sigma$ 的循环表达式中只有一个长度大于 $1$ 的循环，}
\handline{72}{919}{779}{则 $\sigma$ 称为一个循环置换.}
\handline{72}{955}{780}{即 $\sigma=(i_1,\ldots,i_k)$. 显然 $\forall j\notin\{i_1,\ldots,i_k\}$ 则 $\sigma(j)=j$.}
\handline{72}{985}{752}{$\sigma^{-1}=(i_k,i_{k-1},\ldots,i_1)$.}
\end{handpage}

% ===== Line-layout page-089.tex =====
\begin{handpage}{89}
\handline{72}{130}{794}{\textbf{Prop.} 若 $\sigma=(i_1,\ldots,i_k)$，$\tau=(j_1,\ldots,j_l)$，}
\handline{72}{160}{723}{$\{i_1,\ldots,i_k\}\cap\{j_1,\ldots,j_l\}=\varnothing$，则 $\sigma\tau=\tau\sigma$.}
\handline{72}{188}{792}{\textbf{Thm.} 设 $\sigma,\tau$ 同上，则 $\sigma\tau$ 在 $S_n$ 中的阶为 $\operatorname{lcm}(k,l)$.}
\handline{72}{219}{785}{\textbf{Pf.} 用 $o(\sigma)$ 表示 $\sigma$ 的阶，则 $o(\sigma)=k$，$o(\tau)=l$.}
\handline{72}{248}{707}{$\sigma,\tau$ 交换，$\therefore\forall m\in\N$，$(\sigma\tau)^m=\sigma^m\tau^m$.}
\handline{72}{277}{710}{下证 $(\sigma\tau)^m=e\Longleftrightarrow\sigma^m=\tau^m=e$.}
\handline{72}{309}{712}{“$\Leftarrow$” trivial. “$\Rightarrow$”：}
\handline{72}{338}{699}{若 $(\sigma\tau)^m=e$ 则 $\sigma^m\tau^m=e$.}
\handline{72}{369}{707}{$\therefore\sigma^m=\tau^{-m}$. 而 $\sigma^m\in\langle\sigma\rangle$，$\tau^{-m}\in\langle\tau\rangle$.}
\handline{72}{400}{705}{$\sigma$ 与 $\tau$ 中的元素不相交，$\Rightarrow\langle\sigma\rangle\cap\langle\tau\rangle=\{e\}$.}
\handline{72}{432}{710}{$\therefore\sigma^m=\tau^{-m}=e$.}
\handline{72}{465}{715}{$\therefore(\sigma\tau)^m=e\Longleftrightarrow\sigma^m=\tau^m=e$.}
\handline{72}{500}{718}{设 $o(\sigma\tau)=t$，则 $(\sigma\tau)^t=e\Longleftrightarrow\sigma^t=\tau^t=e\Longleftrightarrow k\mid t,\ l\mid t$.}
\handline{72}{536}{718}{由 $t$ 的最小性 $\Rightarrow t=\operatorname{lcm}(k,t)$.}
\handline{72}{573}{804}{于是 $\forall\sigma\in S_n$，要求 $o(\sigma)$：先写出 $\sigma$ 的循环表达式.}
\handline{72}{607}{757}{得到所有的循环长度 $k_1,\ldots,k_l$. 则 $o(\sigma)=\operatorname{lcm}(k_1,\ldots,k_l)$.}
\handline{72}{649}{837}{\textbf{Prop.} $\forall\sigma\in S_n$，设 $\tau=(i_1,\ldots,i_k)$，则 $\sigma^{-1}\tau\sigma=(\sigma(i_1),\ldots,\sigma(i_k))$.}
\handline{72}{685}{835}{\textbf{Pf.} 考虑 $\forall j\in X$，LHS 及 RHS 在 $j$ 上的作用.}
\handline{72}{717}{752}{\circled{1} 若 $j\in\{\sigma(i_1),\ldots,\sigma(i_k)\}$，即 $\exists1\le m\le k$ s.t. $\sigma(i_m)=j$.}
\handline{72}{752}{770}{LHS：$j\sigma^{-1}\tau\sigma=i_m\tau\sigma=i_{m+1}\sigma=\sigma(i_{m+1})$.}
\handline{72}{787}{770}{RHS：由于 $j=\sigma(i_m)$，$\therefore\mathrm{RHS}(j)=\sigma(i_{m+1})$.}
\handline{72}{823}{767}{\circled{2} 若 $j\notin\{\sigma(i_1),\ldots,\sigma(i_k)\}\Rightarrow\sigma^{-1}(j)\notin\{i_1,\ldots,i_k\}$.}
\handline{72}{866}{772}{LHS：$j\sigma^{-1}\tau\sigma=\sigma^{-1}(j)\tau\sigma=\sigma^{-1}(j)\sigma=j$.}
\handline{72}{902}{781}{RHS：$\mathrm{RHS}(j)=j$.}
\handline{72}{949}{779}{$\therefore\sigma^{-1}\tau\sigma=(\sigma(i_1),\ldots,\sigma(i_k))$.}
\handline[AnnotationBlue]{72}{982}{775}{于是一个循环经过 $S_n$ 中元素的共轭作用，循环长度不变.}
\end{handpage}

% ===== Line-layout page-090.tex =====
\begin{handpage}{90}
\handline{72}{103}{791}{下面利用共轭作用来证明一个组合恒等式：$\textstyle\sum_{t=1}^{n}\ \sum_{\substack{0\le m_i\le n\\m_1+2m_2+\cdots+tm_t=n}}\frac1{m_1!m_2!\cdots m_t!\,1^{m_1}2^{m_2}\cdots t^{m_t}}=1$.}
\handline[AnnotationBlue]{72}{139}{785}{首先利用 Taylor 公式. $\dfrac1{1-x}=1+x+x^2+\cdots$.}
\handline{72}{173}{797}{而 $\textstyle\frac1{1-x}=\exp\left(\log\frac1{1-x}\right)=\exp\left(\sum_{k=1}^\infty\frac{x^k}k\right)=\prod_{k=1}^\infty\exp\left(\frac{x^k}k\right)=\prod_{k=1}^\infty\left(\sum_{m_k=0}^\infty\frac1{m_k!}\cdot\frac{x^{km_k}}{k^{m_k}}\right)$}
\handline{72}{222}{769}{$\textstyle=1+\sum_{n=1}^\infty\left(\sum_{t=1}^n\ \sum_{\substack{0\le m_i\le n\\m_1+2m_2+\cdots+tm_t=n}}\frac1{m_1!m_2!\cdots m_t!\,1^{m_1}2^{m_2}\cdots t^{m_t}}\right)x^n$. 从而得证.}
\handline[AnnotationBlue]{72}{259}{788}{两个反函数先 Taylor 展开后再复合，往往会得到一些组合恒等式.}
\handline{72}{294}{798}{首先，$\forall\tau=\tau_1\cdots\tau_s$，其中每个 $\tau_i$ 是长度为 $l_i$ 的循环.}
\handline{72}{323}{741}{则 $\forall\sigma\in S_n$，$\sigma^{-1}\tau\sigma=(\sigma^{-1}\tau_1\sigma)(\sigma^{-1}\tau_2\sigma)\cdots(\sigma^{-1}\tau_s\sigma)$.}
\handline{72}{351}{740}{其中每个 $\sigma^{-1}\tau_i\sigma$ 的长度也为 $l_i$.}
\handline{72}{385}{799}{反之，若先给定 $\tau=\tau_1\cdots\tau_s$，$\widehat\tau=\widehat\tau_1\cdots\widehat\tau_s$，每个 $\tau_i,\widehat\tau_i$ 的长度为 $l_i$，}
\handline{72}{420}{746}{则可构造一个 $\sigma$，使得 $\widehat\tau=\sigma^{-1}\tau\sigma$.}
\handline{72}{455}{190}{构造示意如下：}
\handdiagram{518}{440}{410}{60}{
\node at(-1.2,1){$\tau:$};\node at(-1.2,0){$\widehat\tau:$};
\node(a)at(0,1){$i_1$};\node(b)at(1.8,1){$i_{l_1}$};\node(c)at(3.3,1){$i_{l_1+1}$};\node(d)at(5.5,1){$i_{l_1+l_2}$};
\node(aa)at(0,0){$j_1$};\node(bb)at(1.8,0){$j_{l_1}$};\node(cc)at(3.3,0){$j_{l_1+1}$};\node(dd)at(5.5,0){$j_{l_1+l_2}$};
\node at(.9,1){$\cdots$};\node at(.9,0){$\cdots$};\node at(4.4,1){$\cdots$};\node at(4.4,0){$\cdots$};
\node at(-.5,1){$($};\node at(2.4,1){$)($};\node at(6.1,1){$)\cdots$};\node at(-.5,0){$($};\node at(2.4,0){$)($};\node at(6.1,0){$)\cdots$};
\draw[->,NoteRed](a)--(aa);\draw[->,NoteRed](b)--(bb);\draw[->,NoteRed](c)--(cc);\draw[->,NoteRed](d)--(dd);}
\handline{72}{525}{799}{于是 $S_n$ 中的共轭类由 $l_1,l_2,\ldots,l_s$ 这些循环长度唯一决定.}
\handline{72}{559}{782}{不妨设 $l_1\ge l_2\ge\cdots\ge l_s$. 于是得到一张 Young 图.}
\handline{72}{593}{777}{以 $n=13$ 为例，循环长度分别为 $4,3,2,2,1,1$.}
\handdiagram{708}{607}{220}{209}{
\foreach \len [count=\row from 0] in {4,3,2,2,1,1}{\foreach \col in {1,...,\len}{\draw(\col-1,-\row) rectangle(\col,-\row-1);}\node[text=AnnotationBlue] at(-.6,-\row-.5){$\len$};}}
\handline{72}{630}{612}{一张 Young 图就决定了一个共轭类.}
\handline{72}{663}{600}{只需往左图中填入 $1\sim13$，即可得到共轭类中的所有元素.}
\handline{72}{695}{607}{但会有重复：}
\handline{72}{729}{600}{\circled{1} 相同行中，轮换对称重复. 若该行长度为 $l$，}
\handline{72}{760}{596}{则有 $l$ 个轮换为同一循环.}
\handline{72}{802}{600}{\circled{2} 不同行，但长度相同. 若有 $k$ 行长度相同，}
\handline{72}{841}{598}{则 $k$ 行上下顺序不同会有 $k!$ 种重复.}
\handline{72}{883}{799}{设 $l_1,l_2,\ldots,l_s$ 中有 $m_1$ 个 $1$，$m_2$ 个 $2$，$\ldots$，$m_t$ 个 $t$. 则 $1\le t\le n$.}
\handline{72}{929}{856}{$\textstyle\therefore n!=|S_n|=\sum_{\text{共轭类}}(\text{共轭类中元素个数})=\sum_{t=1}^n\ \sum_{\substack{0\le m_i\le n\\m_1+2m_2+\cdots+tm_t=n}}\frac{n!}{m_1!\cdots m_t!\,1^{m_1}2^{m_2}\cdots t^{m_t}}$.}
\handline{72}{990}{806}{$\textstyle\therefore\sum_{t=1}^n\ \sum_{\substack{0\le m_i\le n\\m_1+2m_2+\cdots+tm_t=n}}\frac1{m_1!\cdots m_t!\,1^{m_1}2^{m_2}\cdots t^{m_t}}=1$.}
\end{handpage}

% ===== Line-layout page-091.tex =====
\begin{handpage}{91}
\handline{72}{105}{795}{\textbf{Prop.} 任意循环可写为对换的乘积. 表示方法不唯一，但所含对换个数的奇偶性不变.}
\handline{72}{134}{794}{\textbf{Pf.} 若 $\sigma=(i_1,\ldots,i_k)$，则 $\sigma=(i_1\ i_2)(i_1\ i_3)(i_1\ i_4)\cdots(i_1\ i_k)$.}
\handline{72}{161}{725}{故任意循环可写为对换的乘积.}
\handline{72}{191}{723}{考虑 $\sigma$ 在 $k$ 阶 Van der Monde 行列式上的作用.}
\handline{72}{240}{707}{$\Delta=\prod_{1\le t<s\le k}(x_{i_s}-x_{i_t})=\begin{vmatrix}1&1&\cdots&1\\x_{i_1}&x_{i_2}&\cdots&x_{i_k}\\\vdots&\vdots&&\vdots\\x_{i_1}^{k-1}&x_{i_2}^{k-1}&\cdots&x_{i_k}^{k-1}\end{vmatrix}$.}
\handline{72}{294}{734}{任一对换可使 $\Delta$ 的符号改变.}
\handline{72}{320}{726}{设 $\sigma$ 有两个对换分解 $\sigma=\tau_1\tau_2\cdots\tau_r=\eta_1\eta_2\cdots\eta_s$.}
\handline{72}{342}{724}{$\forall$ 对换 $\tau$，$\sigma(\Delta)=-\tau\sigma(\Delta)$.}
\handline{72}{376}{715}{$\therefore\sigma(\Delta)=\tau_1\tau_2\cdots\tau_r(\Delta)=(-1)^r\Delta$}
\handline{72}{404}{616}{$=\eta_1\eta_2\cdots\eta_s(\Delta)=(-1)^s\Delta$.}
\handline{72}{436}{686}{而 $\Delta\ne0\Rightarrow(-1)^r=(-1)^s$，$\therefore s\equiv r\pmod2$. 故奇偶性相同.}
\handline{72}{474}{747}{\uline{还可以用逆序数证（见高代）.}}
\handline{72}{513}{817}{\textbf{Def.} 设 $\sigma=\tau_1\cdots\tau_s$，$\tau_i$ 是长度为 $l_i$ 的循环.}
\handline{72}{545}{725}{$\operatorname{sign}(\tau_i):=(-1)^{l_i-1}$，$\operatorname{sing}(\sigma):=(-1)^{l_1+\cdots+l_s-s}$.}
\handline{72}{578}{731}{若 $\operatorname{sign}(\sigma)=1$，则称 $\sigma$ 为偶置换.}
\handline{72}{611}{734}{若 $\operatorname{sign}(\sigma)=-1$，则称 $\sigma$ 为奇置换.}
\handline{72}{648}{738}{$A_n=\{\sigma\mid\operatorname{sign}(\sigma)=1\}$ 称为交错群. \ann{（alternating group）}}
\handline{72}{690}{778}{\circled{1} $|A_n|=\frac12n!=\frac12|S_n|$，即 $[S_n:A_n]=2$，$\therefore A_n\lhd S_n$.}
\handline{72}{726}{805}{\circled{2} 对换是奇置换. 若循环置换长度为偶数，则为奇置换，一定可以分解为奇数个对换.}
\handline{72}{758}{744}{若循环置换长度为奇数，则为偶置换，一定可以分解为偶数个对换.}
\handline{72}{800}{783}{\circled{3} 偶置换成群，但奇置换不成群，因为两个奇置换相乘为偶置换.}
\handline{72}{845}{833}{\textbf{Thm.} $n\ge5$，则 $A_n$ 是单群\ann{（且非交换）}.}
\handline{72}{883}{757}{这一定理在 Galois 理论中也有很重要的作用.}
\handline{72}{919}{758}{下面先看 $n=1,2,3,4$ 以及两个引理.}
\handline[AnnotationBlue]{72}{960}{762}{由于 $n\ge4$ 时，$\sigma=(123)$，$\tau=(124)$.}
\handline[AnnotationBlue]{72}{999}{740}{$\sigma\tau(1)=\tau(2)=4$，$\tau\sigma(1)=\sigma(2)=3$. $\therefore A_n$ 不是交换群.}
\end{handpage}

% ===== Line-layout page-092.tex =====
\begin{handpage}{92}
\begingroup
\newcommand{\pNinetyTwoYoung}{\begin{tikzpicture}[baseline=-.5ex,x=1cm,y=1cm,scale=.13]\draw(0,0)grid(2,2);\end{tikzpicture}}
\handline{72}{94}{178}{$n=1$，$A_1$ 不存在.}
\handline{72}{121}{374}{$n=2$，$A_2=\{e\}$，也可以视为单群}
\handline{72}{149}{577}{$n=3$，$|A_3|=\frac{|S_3|}2=3\quad\therefore A_3\cong\Z_3\quad\therefore$ 是单群但交换.}
\handline{72}{176}{648}{$n=4$\quad 先分析 Young 图.\qquad\circled{1} $4$\quad$(-1)^{4-1}=-1$，为奇置换.}
\handline{72}{204}{570}{考虑 $4$ 的分拆\qquad\circled{2} $3+1$\quad$(-1)^{3-1}\cdot(-1)^{1-1}=1$，为偶置换.}
\handline{72}{232}{660}{从而 $A_4$ 中的元素的 Young 图\quad\circled{3} $2+1+1$\quad$(-1)^{2-1}(-1)^{1-1}(-1)^{1-1}=-1$，为奇置换}
\handline{72}{258}{667}{只有\hspace{49pt}和\hspace{35pt}两种.\qquad\circled{4} $2+2$\quad$(-1)^{2-1}\cdot(-1)^{2-1}=1$，为偶置换.}
\handdiagram{135}{238}{137}{34}{
\draw(0,1)grid(3,2);\draw(0,0)rectangle(1,1);\node at(1,-1){(I)};
\draw(6,0)grid(8,2);\node at(7,-1){(II)};
}
\handline{72}{282}{210}{\circled{5} $1+1+1+1$\quad 为 $e$.}
\handline{72}{315}{485}{而 (I) 共有 $\frac{4!}3=8$ 种，(II) 共有 $\frac{4!}{2!\times2^2}=3$ 种.}
\handline{72}{343}{355}{显然 $K=\{e,\pNinetyTwoYoung\}$ 是 $A_4$ 的子群.}
\handline{72}{371}{643}{而由于共轭保持循环长度.$\therefore\sigma^{-1}\pNinetyTwoYoung\sigma=\pNinetyTwoYoung\subseteq K$.\ $\forall\sigma\in S_4$.}
\handline{72}{403}{340}{$\therefore K\lhd A_4\quad\therefore A_4$ 不是单群.}
\handline[AnnotationBlue]{72}{432}{668}{$K=\{e,(12)(34),(13)(24),(14)(23)\}\cong\Z_2\times\Z_2$，称为 Klein 四元群.}
\handline{72}{465}{313}{Lemma. (I)\quad$A_n$ 可由 $3$ 循环生成.}
\handline{72}{494}{647}{Pf.\quad 由于 $A_n$ 中元素可写为偶数个对换乘积，故只需 $(ij)(kl)$ 由 $3$ 循环生成}
\handline{72}{522}{350}{$1^\circ$ $(ij)=(kl)$\quad 则 $(ij)(kl)=e$}
\handline{72}{553}{507}{$2^\circ$ $|\{i,j\}\cap\{k,l\}|=1$.不妨设 $j=l$.则 $(ij)(kl)=(ikj)$}
\handline{72}{584}{520}{$3^\circ$ $\{i,j\}\cap\{k,l\}=\varnothing$.\ $(ij)(kl)=(ij)(j,k)(kj)(kl)$}
\handline{72}{611}{230}{由 $2^\circ$\quad$=(ikj)(jlk)$}
\handline{72}{650}{496}{Lemma. (II).\quad$n\ge5$ 时，$A_n$ 中的 $3$ 循环相互共轭.}
\handline{72}{682}{478}{Pf.\quad 由于 $S_n$ 中的 $3$ 循环都相互共轭.}
\handline{72}{718}{635}{$\therefore\forall\sigma\in A_n$，$\sigma$ 为 $3$ 循环，$\exists\tau\in S_n\ \text{s.t.}\ \tau^{-1}\sigma\tau=(123)$.}
\handline{72}{748}{230}{若 $\tau\in A_n$，则证毕.}
\handline{72}{783}{437}{若 $\tau\notin A_n$，取 $\widetilde\tau=\tau(45)$.则 $\widetilde\tau\in A_n$.}
\handline{72}{821}{648}{$(\widetilde\tau)^{-1}\sigma\tau=(45)(123)(45)=(123)(45)(45)=(123)$.}
\handline{72}{858}{647}{$\therefore\forall\sigma\in A_n$，$\sigma$ 为 $3$ 循环，$\exists\tau\in A_n\ \text{s.t.}\ \tau^{-1}\sigma\tau=(123)$.}
\handline{72}{889}{378}{即 $A_n$ 中的 $3$ 循环相互共轭.}
\endgroup
\end{handpage}

% ===== Line-layout page-093.tex =====
\begin{handpage}{93}
\handline{72}{111}{265}{Thm.\quad$n\ge5$，$A_n$ 是单群}
\handline{72}{138}{317}{Pf.\quad 设 $K\lhd A_n$ 且 $K\ne\{e\}$}
\handline{72}{163}{169}{下证 $K=A_n$.}
\handline{72}{186}{461}{由引理 (I)，只需证明 $K$ 包含所有 $3$ 循环}
\handline{72}{215}{386}{只需证明 $K$ 至少包含一个 $3$ 循环}
\handline{72}{244}{597}{因为只要 $K$ 包含了一个 $3$ 循环，则这个 $3$ 循环的共轭类}
\handline{72}{269}{368}{包含了所有 $3$ 循环\ann{（由引理 (II)）}}
\handline{72}{297}{533}{又因为 $K$ 正规，从而这个共轭类一定包含于整个 $K$.}
\handline{72}{326}{227}{设 $\sigma\in K$，且 $\sigma\ne e$}
\handline{72}{352}{583}{由于 $\sigma\ne e$ 故 $\sigma$ 在 $\{1,2,\ldots,n\}$ 上作用的不动点数 $d<n$.}
\handline{72}{383}{482}{设 $\sigma$ 的不动点的个数是 $K\setminus\{e\}$ 中最大的}
\handline{72}{410}{271}{下证 $\sigma$ 一定是 $3$ 循环.}
\handline{72}{440}{580}{若否，则 \circled{1} $\sigma$ 的循环分解中只含对换，即 $2$ 循环.}
\handline{72}{469}{552}{由于 $\sigma$ 是偶置换，故至少有 $2$ 个对换.$\therefore\sigma=(12)(34)\rho$}
\handline{72}{499}{455}{\circled{2} $\sigma$ 的循环分解中含有长度 $\ge3$ 的循环.}
\handline{72}{530}{189}{$\sigma=(123\cdots)\rho$.}
\handline{72}{572}{563}{对于 \circled{1}.\quad 设 $\tau=(345)$.考虑 $[\sigma,\tau]=\sigma^{-1}\tau^{-1}\sigma\tau$}
\handline{72}{602}{631}{由于 $K$ 正规.$\therefore\tau^{-1}\sigma\tau\in K\quad\therefore\sigma^{-1}(\tau^{-1}\sigma\tau)=[\sigma,\tau]\in K$.}
\handline{72}{631}{435}{设 $\alpha=[\sigma,\tau]$.下面看 $\alpha$ 的不动点.}
\handline{72}{664}{355}{若 $i$ 是 $\sigma$ 的不动点，且 $i\ge6$}
\handline{72}{698}{665}{则 $\alpha(i)=\sigma^{-1}\tau^{-1}\sigma\tau(i)=i\quad\therefore i$ 也是 $\alpha$ 的不动点.}
\handline{72}{734}{665}{而 $\alpha(1)=\sigma^{-1}\tau^{-1}\sigma\tau(1)=\tau^{-1}\sigma\tau(2)=\sigma\tau(2)=\tau(1)=1$}
\handline{72}{765}{491}{同理 $\alpha(2)=2\quad\therefore1,2$ 也是 $\alpha$ 的不动点.}
\handline{72}{801}{808}{而 $\alpha(3)=5\quad\therefore\alpha\ne e$\redann{（这一步必须要有，因为是 $\sigma$ 的最小性是在 $K\setminus\{e\}$ 中讨论的）}}
\handline{72}{842}{804}{为了比较 $\alpha$ 与 $\sigma$ 的不动点个数，只需比较 $1,2,3,4,5$ 中不动点的个数.}
\handline{72}{883}{822}{而 $\alpha$ 至少 $2$ 个，$\sigma$ 至多 $1$ 个.$\therefore\alpha$ 的不动点更多，与 $\sigma$ 的取法矛盾！}
\handline[NoteRed]{72}{924}{392}{（$1,2,3,4$ 一定不是 $\sigma$ 的不动点）}
\handarrow[NoteRed]{292.0}{902}{289.0}{886}
\end{handpage}

% ===== Line-layout page-094.tex =====
\begin{handpage}{94}
\handline{72}{131}{509}{对于 \circled{2}.由于 $(1234)$ 为奇置换.$\therefore\sigma\ne(1234)$}
\handline{72}{160}{531}{$\therefore\sigma$ 只可能为 $(12345\cdots)\rho$.\quad$\rho$ 可以为 $e$}
\handline{72}{188}{552}{或 $(1234)\rho$.或 $(123)\rho$.\quad$\rho$ 不可以为 $e$.}
\handline{72}{215}{671}{$\therefore\sigma$ 除了 $1,2,3$ 之外，至少还应该动两个点，不妨设为 $4,5$.}
\handline{72}{246}{505}{仍取 $\tau=(345)$\quad$\alpha=[\sigma,\tau]$ 则 $\alpha\in K$}
\handline{72}{277}{324}{若 $i$ 为 $\sigma$ 的不动点，且 $i\ge6$.}
\handline{72}{307}{444}{则 $\alpha(i)=(i)$ 故仍为 $\alpha$ 的不动点.}
\handline{72}{338}{660}{而 $\alpha(2)=\sigma^{-1}\tau^{-1}\sigma\tau(2)=\tau^{-1}\sigma\tau(1)=\sigma\tau(1)=\tau(2)=2$.}
\handline{72}{368}{280}{$\therefore2$ 也为 $\alpha$ 的不动点.}
\handline{72}{399}{651}{而 $\alpha(3)=\sigma^{-1}\tau^{-1}\sigma\tau(3)=\tau^{-1}\sigma\tau(2)=\sigma\tau(2)=\tau(3)=4$}
\handline{72}{432}{126}{$\therefore\alpha\ne e$.}
\handline{72}{459}{418}{而在 $1\sim5$ 中，$2$ 一定为 $\alpha$ 的不动点.}
\handline{72}{491}{379}{但 $1\sim5$ 均不为 $\sigma$ 的不动点.}
\handline{72}{523}{494}{$\therefore\alpha$ 的不动点更多，与 $\sigma$ 的取法矛盾！}
\handline{72}{562}{319}{综上 $\sigma$ 一定是 $3$ 循环.}
\handline{72}{595}{180}{从而 $K=A_n$.}
\handline{72}{633}{177}{$\therefore A_n$ 是单群.}
\end{handpage}

% ===== Line-layout page-095.tex =====
\begin{handpage}{95}
\handtitle{72}{100}{184}{§3.8 交换群}
\handline{72}{128}{258}{Def.\quad$G$：群\quad$S\subseteq G$.}
\handline{72}{158}{598}{$\langle S\rangle:=\bigcap_{\substack{H\le G\\S\subseteq H}}H$ 即包含 $S$ 的最小子群，称为 $S$ 生成的子群.}
\handline{72}{204}{715}{若 $G=\langle S\rangle$.则称 $G$ 由 $S$ 生成，$S$ 称为 $G$ 的生成集，$a\in S$ 称为生成元.}
\handline{72}{240}{390}{若 $|S|<\infty$，则称 $G$ 是有限生成的}
\handline{72}{269}{408}{若 $G=\langle a\rangle$，则称 $G$ 是循环群.}
\handline{72}{300}{528}{Prop.\quad 设 $G=\langle a\rangle$.则若 $|G|=\infty$，则 $G\cong\Z$.}
\handline{72}{329}{483}{\ann{$G$ 中运算为加法}\qquad 若 $|G|=n$，则 $G\cong\Z_n$.}
\handline{72}{359}{611}{Pf.\quad$\phi\colon\Z\to G$\qquad 是自然同态，显然是满同态.}
\handline{72}{383}{254}{$n\mapsto na:=a+a+\cdots+a$}
\handline{72}{449}{707}{$1^\circ$ 若 $\ker\phi=\{0\}$，则 $\phi$ 为单射.$\therefore|G|=\infty$.且 $\phi$ 为同构.$\therefore G\cong\Z$.}
\handline{72}{485}{754}{$2^\circ$ 若 $\ker\phi\ne\{0\}$，由于 $\ker\phi\lhd\Z\quad\therefore\exists n\in\omega\ \text{s.t.}\ \ker\phi=n\Z$.}
\handline{72}{520}{375}{则 $G\cong\Z/\ker\phi=\Z/n\Z=\Z_n$}
\handline[AnnotationBlue]{72}{555}{417}{若 $H\le G$.\quad 当 $|G|=\infty$ 时 $|H|=\infty$}
\handline[AnnotationBlue]{72}{588}{315}{当 $|G|=n$ 时，则 $|H|\mid|G|=n$.}
\handline[AnnotationBlue]{72}{623}{413}{设 $|H|=d$，$n=dk$ 则 $H=\langle ka\rangle$.}
\handline{72}{663}{835}{Thm.\quad 设 $G$ 是有限交换群，则 $G$ 是循环群 $\Longleftrightarrow\max\{o(a)\mid a\in G\}=|G|$.}
\handline{72}{696}{244}{Pf.\quad“$\Rightarrow$” trivial}
\handline{72}{729}{637}{“$\Leftarrow$” 设 $a\in G$.$o(a)$ 是 $G$ 中最大的元素阶.$o(a)=m$.}
\handline{72}{769}{642}{则 $|\langle a\rangle|=m=|G|$ 而 $\langle a\rangle\subseteq G\quad\therefore\langle a\rangle=G$.}
\handline{72}{809}{216}{$\therefore G$ 是循环群.}
\end{handpage}

% ===== Line-layout page-096.tex =====
\begin{handpage}{96}
\handline{72}{96}{285}{下面看循环群的自同构：}
\handline{72}{122}{381}{\circled{1} $\Z$：\quad$\varphi\colon\Z\to\Z$ 设 $\varphi(1)=a$}
\handline{72}{149}{328}{若 $a\ne\pm1$.则 $\varphi(\Z)=a\Z\ne\Z$}
\handline{72}{175}{400}{而 $a=\pm1$ 时 $\varphi(\Z)=\Z$ 且 $\varphi$ 为自同构.}
\handline{72}{204}{255}{$\therefore\operatorname{Aut}(\Z)\cong\{\pm1\}\cong\Z_2$.}
\handline{72}{232}{630}{\circled{2} $\Z_n$：\quad$\varphi\colon\Z_n\to\Z_n$.设 $\varphi(1)=\bar a$.不妨设 $0\le a\le n-1$.}
\handline{72}{261}{480}{若 $\varphi(k)=\bar0$ 则 $k\varphi(1)=k\bar a=\bar0\Rightarrow n\mid ka$.}
\handline{72}{289}{375}{为了使 $\varphi$ 为单射，则需 $\ker\varphi=\{\bar0\}$.}
\handline{72}{319}{322}{即“$\forall k\in\omega$.\ $n\mid ka\Rightarrow n\mid k$”.}
\handline{72}{346}{355}{若 $(a,n)=d\ne1$ 设 $n=dn_1$，$a=da_1$}
\handline{72}{378}{555}{而取 $k=n_1$.则 $n\mid n_1a=dn_1a_1=na_1$，但 $n\nmid k=n_1$.}
\handline{72}{408}{140}{$\therefore(a,n)=1$.}
\handline{72}{438}{427}{另一方面若 $(a,n)=1$，易验证 $\varphi$ 是同构.}
\handline{72}{471}{680}{$\therefore\operatorname{Aut}(\Z_n)\cong(\Z_n)^\times$.即 $\Z_n$ 的单位群.$|\operatorname{Aut}(\Z_n)|=\varphi(n)$}
\handline{72}{505}{649}{Def.\quad$G_1,G_2$ 是群.令 $G=G_1\times G_2$.定义乘法 $G\times G\to G$.}
\handline{72}{535}{339}{$((g_1,h_1),(g_2,h_2))\mapsto(g_1h_1,g_2h_2)$.}
\handline{72}{567}{651}{则 $G$ 是一个群.$e_G=(e_{G_1},e_{G_2})$.称为 $G_1,G_2$ 的外直积.}
\handline[AnnotationBlue]{72}{596}{370}{类似可定义 $n$ 个群的外直积.}
\handline[AnnotationBlue]{72}{631}{609}{利用嵌入映射 $h_1\colon G_1\to G\qquad\quad h_2\colon G_2\to G$}
\handline[AnnotationBlue]{72}{664}{417}{$g\mapsto(g,e_{G_2})\qquad\quad g\mapsto(e_{G_1},g)$}
\handline[AnnotationBlue]{72}{694}{365}{可将 $G_1,G_2$ 视为 $G$ 的子群.}
\handline[AnnotationBlue]{72}{729}{574}{且容易验证 $G_1$ 中的元素与 $G_2$ 中的元素交换}
\handline[AnnotationBlue]{72}{761}{245}{且 $G_1\lhd G$，$G_2\lhd G$.}
\handline[AnnotationBlue]{72}{799}{811}{反过来考虑，若 $G$ 有两个正规子群 $G_1,G_2$，且 $G_1\cap G_2=\{e\}$，$G=G_1G_2$.}
\handline[AnnotationBlue]{72}{830}{624}{且由之前的习题知 $G_1$ 中元素与 $G_2$ 中元素交换.}
\handline{72}{868}{709}{\ann{此时我们希望 $G\cong G_1\times G_2$.}\redann{（事实上这是正确的）}}
\end{handpage}

% ===== Line-layout page-097.tex =====
\begin{handpage}{97}
\handline{72}{106}{641}{Def.\quad\circled{1} $G$：群.$G_1\lhd G$，$G_2\lhd G$.$G=G_1G_2$.\ $G_1\cap G_2=\{e\}$.}
\handline{72}{131}{319}{则称 $G$ 为 $G_1,G_2$ 的内直积.}
\handline{72}{160}{459}{\circled{2} 若 $G_1\lhd G,\ldots,G_n\lhd G$，$G=G_1G_2\cdots G_n$}
\handline{72}{184}{387}{且 $\forall1\le i\le n$.\ $G_i\cap(\prod_{j\ne i}G_j)=\{e\}$.}
\handline{72}{210}{335}{则称 $G$ 为 $G_1,\ldots,G_n$ 的内直积.}
\handline{72}{239}{643}{“$\forall1\le i\le n$.\ $G_i\cap(\prod_{j\ne i}G_j)=\{e\}$”$\Rightarrow$“$\forall i\ne j$，$G_i\cap G_j=\{e\}$”.}
\handline{72}{268}{641}{$\forall i\ne j$，$G_i\cap G_j\subseteq G_i\cap(\prod_{j\ne i}G_j)=\{e\}\Rightarrow G_i\cap G_j=\{e\}$.}
\handline[AnnotationBlue]{72}{294}{611}{事实上，内直积与外直积本质上是等价的，此后不再区分}
\handline{72}{325}{526}{Lemma.\quad$G_1,\ldots,G_n\lhd G$，则 $G$ 是 $G_1,\ldots,G_n$ 的内直积}
\handline{72}{354}{665}{$\Longleftrightarrow$\circled{1} $G=G_1\cdots G_n$\quad\circled{2} $G$ 中元素用 $G_i$ 中元素的乘积表示法唯一.}
\handline{72}{384}{390}{Pf.\quad“$\Rightarrow$” 由于 $G_i\lhd G$，$G_j\lhd G$\quad$i\ne j$}
\handline{72}{414}{630}{$\therefore\forall g_i\in G_i,g_j\in G_j\quad[g_i,g_j]\in G_i\cap G_j=\{e\}\Rightarrow g_ig_j=g_jg_i$}
\handline{72}{444}{502}{若 $g\in G$\quad$g=g_1\cdots g_n=h_1\cdots h_n$.\ $h_i,g_i\in G_i$.}
\handline{72}{477}{426}{$\Rightarrow g'=h_1^{-1}g_1=(h_2g_2^{-1})\cdots(h_ng_n^{-1})$}
\handline{72}{528}{441}{$\therefore g'\in G_1\cap G_2\cdots G_n=\{e\}\Rightarrow h_1=g_1$}
\handline{72}{561}{636}{$\therefore g_2\cdots g_n=h_2\cdots h_n$，重复上述过程可得 $h_i=g_i$（$1\le i\le n$）.}
\handline{72}{598}{454}{“$\Leftarrow$” 若 $g_1=g_2\cdots g_n\in G_1\cap G_2\cdots G_n$.}
\handline{72}{632}{477}{由表示法唯一.$g_1\cdot e\cdots e=e\cdot g_2\cdots g_n$}
\handline{72}{670}{522}{$\Rightarrow g_1=g_2=\cdots=g_n=e\quad\therefore G_1\cap G_2\cdots G_n=\{e\}$}
\handline{72}{701}{386}{同理 $G_i\cap(\prod_{j\ne i}G_j)=\{e\}\quad1\le i\le n$}
\handline{72}{737}{359}{$\therefore G$ 是 $G_1,\ldots,G_n$ 的内直积.}
\handline{72}{778}{836}{Cor.\quad 若 $G$ 是 $G_1,\ldots,G_n$ 内直积，$\widetilde G$ 是 $G_1,\ldots,G_n$ 外直积，则 $G\cong\widetilde G$}
\handline{72}{812}{314}{Pf.\quad 定义 $\varphi\colon G\to\widetilde G$}
\handline{72}{845}{304}{$g_1\cdots g_n\mapsto(g_1,g_2,\ldots,g_n)$}
\handline{72}{887}{765}{由于表示法唯一，故 $\varphi$ 是良定义的.容易验证 $\varphi$ 是同构}
\end{handpage}

% ===== Line-layout page-098.tex =====
\begin{handpage}{98}
\handline{72}{118}{600}{Def.\quad$\Lambda$ 为指标集，$\{G_\alpha\}_{\alpha\in\Lambda}$ 是一族群，则 $\prod_{\alpha\in\Lambda}G_\alpha$ 称为直积.}
\handline{72}{147}{654}{$\{\prod_{\alpha\in\Lambda}g_\alpha\mid$ 只有有限多个 $\alpha$，$g_\alpha\ne e_\alpha\}$ 称为限制直积或直和，记为 $\bigoplus_{\alpha\in\Lambda}G_\alpha$}
\handline[AnnotationBlue]{72}{172}{361}{对于有限个群，直积与直和是一回事.}
\handline[AnnotationBlue]{72}{196}{666}{有限个交换群的直积，一般将其记为直和，因为交换群都是 $\Z$-模.}
\handline[AnnotationBlue]{72}{222}{242}{可以视为 $\Z$-模的直和}
\handline[AnnotationBlue]{72}{250}{415}{有限个非交换群的直积，一般还记为直积.}
\handline[AnnotationBlue]{72}{279}{362}{对于无限个群，直积与直和要作区分.}
\handline[AnnotationBlue]{72}{310}{743}{本节讨论的是有限生成交换群，故接下来的直积都记为直和 $\bigoplus_{i=1}^nG_i$.}
\handline{72}{340}{602}{Thm.\quad$G$ 是循环群，$|G|=mn$.若 $(m,n)=1$.则 $G\cong\Z_m\oplus\Z_n$.}
\handline{72}{371}{391}{Pf.\quad$(m,n)=1$.\ $\exists u,v\in\Z\ \text{s.t.}\ um+vn=1$.}
\handline{72}{398}{313}{设 $G=\langle a\rangle$.则 $a=a^{um+vn}$}
\handline{72}{428}{374}{设 $b=a^m$，$c=a^n$ 则 $b^n=c^m=e$.}
\handline{72}{456}{302}{若 $0<o(b)<n$.设 $o(b)=k$.}
\handline{72}{483}{507}{则 $b^k=a^{mk}=e$.与 $a$ 的阶为 $mn$，$mk<mn$ 矛盾.}
\handline{72}{514}{285}{$\therefore o(b)=n$.同理 $o(c)=m$.}
\handline{72}{548}{441}{令 $G_1=\Z_n$，$G_2=\Z_m$.则 $|G_1|=n$，$|G_2|=m$.}
\handline{72}{581}{369}{下证 $G=G_1G_2$ 且 $G_1\cap G_2=\{e\}$.}
\handline{72}{614}{709}{由于 $G_1\cap G_2\le G_1,G_2\Rightarrow|G_1\cap G_2|\mid n,m$ 而 $(m,n)=1\Rightarrow|G_1\cap G_2|=1$}
\handline{72}{644}{179}{$\therefore G_1\cap G_2=\{e\}$.}
\handline{72}{683}{659}{$\forall g\in G$，$g=a^k=a^{(um+vn)k}=b^{ku}c^{kv}\in G_1G_2\quad\therefore G=G_1G_2$.}
\handline{72}{710}{289}{$\therefore G\cong G_1\times G_2=\Z_m\oplus\Z_n$.}
\handline{72}{749}{614}{Cor.\quad$G$ 是循环群.$|G|=p_1^{\alpha_1}\cdots p_k^{\alpha_k}$ 则 $G\cong\bigoplus_{j=1}^k\Z_{p_j^{\alpha_j}}$.}
\handline{72}{779}{480}{Ex.\quad\ann{$\Z_4$ 是循环群}，但 $\Z_4\not\cong\Z_2\oplus\Z_2$.}
\handline{72}{811}{402}{由于 $\max\{o(a)\mid a\in\Z_2\oplus\Z_2\}=2\ne4$}
\handline{72}{852}{497}{$\therefore\Z_2\oplus\Z_2$ 不是循环群.$\therefore\Z_2\oplus\Z_2\not\cong\Z_4$.}
\end{handpage}

% ===== Line-layout page-099.tex =====
\begin{handpage}{99}
\handline{72}{120}{556}{\textbf{Thm.} $G$：有限生成交换群，则 $G$ 可分解为 $G=C_1\oplus C_2\oplus\cdots\oplus C_s$}
\handline{72}{145}{413}{其中 $C_i$（$1\le i\le s$）要么全部是无限循环群}
\handline{72}{174}{671}{要么存在一个 $1\le t\le s$ s.t. $C_1,\ldots,C_t$ 是有限循环群，$C_{t+1},\ldots,C_s$ 是无限循环群.}
\handline{72}{203}{364}{且设 $|C_i|=m_i$，$1\le i\le t$，则 $m_i\mid m_{i+1}$}
\handline{72}{234}{628}{\textbf{Pf.} 设 $G$ 由 $s$ 个生成元生成，且 $s$ 是生成元集的最小基数，对 $s$ 进行归纳.}
\handline{72}{262}{471}{若 $s=1$，trivial. 下设命题对 $s-1$ 成立，下证 $s$ 时也成立.}
\handline{72}{293}{658}{设 $G=\langle a_1,\ldots,a_s\rangle$，由于 $G$ 是 $\Z$-模，故可定义 $\Z$-线性无关与 $\Z$-线性相关.}
\handline{72}{322}{294}{\circled{1} 若 $a_1,\ldots,a_s$ $\Z$-线性无关.}
\handline{72}{351}{415}{$\forall g\in G$，$\exists x_1,\ldots,x_s\in\Z$ s.t. $g=x_1a_1+\cdots+x_sa_s$.}
\handline{72}{381}{422}{由于 $a_1,\ldots,a_s$ 线性无关，故 $x_1,\ldots,x_s$ 唯一确定.}
\handline{72}{411}{472}{于是 $\varphi:G\to\Z\oplus\cdots\oplus\Z$（$s$ 个 $\Z$）. 是良定义的.}
\handline{72}{437}{176}{$g\mapsto(x_1,\ldots,x_s)$}
\handline{72}{471}{268}{且易验证它是同构.}
\handline{72}{502}{521}{\circled{2} 若存在不全为零的 $x_1,\ldots,x_s\in\Z$，使得 $x_1a_1+\cdots+x_sa_s=0$.}
\handline{72}{533}{468}{不妨设 $x_1,\ldots,x_s$ 中至少有一个为正，否则两边同乘 $(-1)$.}
\handline{72}{565}{635}{考虑如下集合：$\bar X=\{(x_1,\ldots,x_s)\mid\begin{gathered}\text{存在 }G\text{ 的一个生成元集 }\{a_1,\ldots,a_s\},\\x_1a_1+\cdots+x_sa_s=0,\text{ 且至少有个 }x_i>0\end{gathered}\}$}
\handline{72}{601}{338}{由于 $a_1,\ldots,a_s$ 线性相关，故 $\bar X$ 非空.}
\handline{72}{632}{678}{由于 $\bar X$ 中所有正分量坐标构成 $\omega$ 的子集，由 $\omega$ 的良序性，故存在最小值.}
\handline{72}{663}{390}{设最小值为 $m_1$，且不妨设为第一个分量坐标.}
\handline{72}{696}{770}{故 $\exists G$ 的生成元集，s.t. $m_1a_1+\cdots+x_sa_s=0$. 对中每个 $x_i$，$2\le i\le s$，作带余除法，$x_i=q_im_1+r_i$.}
\handline{72}{728}{498}{$\Rightarrow m_1a_1+m_1(q_2a_2+\cdots+q_sa_s)+r_2a_2+\cdots+r_sa_s=0$.}
\handline{72}{758}{485}{$=m_1(a_1+q_2a_2+\cdots+q_sa_s)+r_2a_2+\cdots+r_sa_s=0$.}
\handline{72}{791}{639}{令 $b=a_1+q_2a_2+\cdots+q_sa_s$，则 $b\ne0$. 否则 $a_1=-q_2a_2-\cdots-q_sa_s$.}
\handline{72}{823}{438}{$a_1$ 可由 $a_2,\ldots,a_s$ 生成，与 $s$ 的最小性矛盾.}
\handline{72}{855}{711}{于是 $a_1=b-q_2a_2-\cdots-q_sa_s$，所以 $\{b,a_2,\ldots,a_s\}$ 也构成一个生成元集.}
\handline{72}{887}{690}{$m_1b+r_2a_2+\cdots+r_sa_s=0$. 若 $r_i\ne0$，由带余除法 $\Rightarrow0<r_i\le m_1-1<m_1$}
\handline{72}{920}{680}{由于 $(m_1,r_2,\ldots,r_s)\in\bar X$，这与 $m_1$ 的最小性矛盾！所以 $m_1b=0$.}
\handline{72}{950}{663}{令 $C_1=\langle b\rangle$，由 $m_1$ 的最小性可知，$o(b)=m_1$，所以 $|C_1|=m_1$.}
\end{handpage}

% ===== Line-layout page-100.tex =====
\begin{handpage}{100}
\handline{72}{116}{390}{设 $G_2=\langle a_2,\ldots,a_s\rangle$. 下证 $G=C_1\oplus G_2$.}
\handline{72}{146}{625}{由于 $b,a_2,\ldots,a_s$ 是 $G$ 的生成元，所以 $G=C_1+G_2$. 下证 $C_1\cap G_2=\{0\}$}
\handline{72}{176}{616}{若 $g\in C_1\cap G_2$，则 $g=x_1b=x_2a_2+\cdots+x_sa_s$，$x_1,\ldots,x_s\in\Z$.}
\handline{72}{206}{354}{$\Rightarrow x_1b-x_2a_2-\cdots-x_sa_s=0$}
\handline{72}{236}{382}{由于 $C_1\cong\Z_{m_1}$，不妨设 $0\le x_1\le m_1-1$.}
\handline{72}{266}{616}{则 $(x_1,-x_2,\ldots,-x_s)\in\bar X$，结合 $m_1$ 在 $\bar X$ 中的最小性 $\Rightarrow x_1=0$}
\handline{72}{295}{345}{所以 $g=x_1b=0$，所以 $C_1\cap G_2=\{0\}$}
\handline{72}{325}{214}{所以 $G=C_1\oplus G_2$}
\handline[AnnotationBlue]{72}{357}{703}{最后，为了使用归纳假设，还需证明 $G_2$ 的最小生成元集有 $s-1$ 个元素.}
\handline{72}{387}{589}{若 $G_2$ 由少于 $s-1$ 个生成元生成，则 $G$ 可由少于 $s$ 个生成元生成}
\handline{72}{417}{579}{与 $s$ 的最小性矛盾. 故 $G_2$ 的最小生成元集有 $s-1$ 个元素.}
\handline{72}{447}{365}{于是由归纳假设 $G_2=C_2\oplus\cdots\oplus C_s$}
\handline{72}{477}{573}{要么 $C_i$ 全是 $\Z$，要么 $\exists2\le t\le s$，$C_2,\ldots,C_t$ 为有限循环群}
\handline{72}{510}{573}{$C_{t+1},\ldots,C_s$ 为 $\Z$，且 $m_i=|C_i|$，$2\le i\le t$，满足 $m_2\mid m_3\mid\cdots\mid m_t$.}
\handline{72}{542}{389}{只需在第 \circled{2} 种情形下证明 $m_1\mid m_2$.}
\handline{72}{573}{534}{设 $C_i=\langle b_i\rangle$，$2\le i\le s$，则 $G$ 由 $\langle b_1,b_2,\ldots,b_s\rangle$ 生成.}
\handline{72}{603}{376}{且 $0=m_1b_1+m_2b_2+0\cdot b_3+\cdots+0\cdot b_s$}
\handline{72}{636}{283}{所以 $(m_1,m_2,0,\ldots,0)\in\bar X$.}
\handline{72}{669}{467}{再作带余除法 $m_2=q_2m_1+r_2$，$0\le r_2\le m_1-1$.}
\handline{72}{702}{344}{$\Rightarrow m_1(b_1+q_2b_2)+r_2b_2=0$.}
\handline{72}{736}{791}{则 $b_1+q_2b_2\ne0$，否则 $b_1=-q_2b_2$，$G$ 可由 $\langle b_2,\ldots,b_s\rangle$ 生成，与 $s$ 最小性矛盾！}
\handline{72}{769}{663}{而 $b_1+q_2b_2,b_2$ 可生成 $b_1,b_2$，故 $G$ 可由 $\langle b_1,b_2,\ldots,b_s\rangle$ 生成}
\handline{72}{802}{273}{所以 $(m_1,r_2,0,\ldots,0)\in\bar X$}
\handline{72}{837}{682}{若 $r_2\ne0$，则由带余除法 $0<r_2\le m_1-1<m_1$，与 $m_1$ 最小性矛盾.}
\handline{72}{871}{426}{所以 $r_2=0\Rightarrow m_2=q_2m_1$，所以 $m_1\mid m_2$.}
\end{handpage}

% ===== Line-layout page-101.tex =====
\begin{handpage}{101}
\handline{72}{116}{459}{\textbf{Def.} $G$：交换群. 若存在 $B=\{e_\alpha\mid\alpha\in\Lambda\}\subseteq G$.}
\handline{72}{147}{514}{$\forall g\in G$，$g$ 可被唯一地表示为 $g=\sum_{\alpha\in\Lambda}n_\alpha e_\alpha$，$n_\alpha\in\Z$.}
\handline{72}{177}{618}{其中只有有限个 $n_\alpha$ 不为 $0$. 则称 $G$ 为自由交换群，并称 $B$ 为 $G$ 的一个基.}
\handline{72}{205}{171}{\circled{1} $G=\langle B\rangle$.}
\handline{72}{233}{313}{\circled{2} $B$ 中元素一定 $\Z$-线性无关}
\handline{72}{263}{316}{即对于任意有限个 $e_{\alpha_1},\ldots,e_{\alpha_k}\in B$}
\handline{72}{292}{459}{若 $n_1e_{\alpha_1}+\cdots+n_ke_{\alpha_k}=0$，则 $n_1=n_2=\cdots=n_k=0$.}
\handline{72}{324}{621}{\circled{3} $G\cong\bigoplus_{\alpha\in\Lambda}\Z$，因为表示法唯一，故 $\varphi:G\to\bigoplus_{\alpha\in\Lambda}\Z$ 良定义}
\handline{72}{354}{602}{且易验证 $\varphi$ 是同构.\qquad $\sum_{\alpha\in\Lambda}n_\alpha e_\alpha\mapsto(n_\alpha)_{\alpha\in\Lambda}$}
\handline{72}{386}{716}{\textbf{Lemma. (I)} $G$：交换群，$H\le G$. 若 $G/H\cong\Z^r$，$r\in\omega$，则 $G\cong H\oplus\Z^r$.}
\handline{72}{418}{789}{\textbf{Pf.} 取 $\Z^r$ 的基 $\bar u_1,\ldots,\bar u_r$，它们对应 $u_1,\ldots,u_r\in G$. 其中 $\bar u_i=u_i+H$，$1\le i\le r$.}
\handline{72}{448}{450}{设 $K=\langle u_1,u_2,\ldots,u_r\rangle\le G$. 下证 $G=H\oplus K$.}
\handline{72}{478}{193}{\circled{1} $H\cap K=\{0\}$.}
\handline{72}{509}{639}{若 $x\in H\cap K$，则 $x=m_1u_1+\cdots+m_ru_r$，$m_i\in\Z$，且 $x+H=0+H=H$.}
\handline{72}{540}{657}{所以 $m_1\bar u_1+m_2\bar u_2+\cdots+m_r\bar u_r=0$. 由基的定义 $\Rightarrow m_1=\cdots=m_r=0$，所以 $x=0$.}
\handline{72}{571}{216}{所以 $H\cap K=\{0\}$}
\handline{72}{603}{216}{\circled{2} $H+K=G$.}
\handline{72}{636}{545}{只需证 $G\subseteq H+K$. $\forall x\in G$，考虑 $x+H\in G/H\cong\Z^r$.}
\handline{72}{671}{483}{所以 $\exists m_1,\ldots,m_r\in\Z$ s.t. $x+H=m_1\bar u_1+\cdots+m_r\bar u_r$}
\handline{72}{705}{412}{所以 $\exists y\in H$ s.t. $x+y=m_1u_1+\cdots+m_ru_r$.}
\handline{72}{738}{583}{所以 $x=(-y)+(m_1u_1+\cdots+m_ru_r)\in H+K$，所以 $G\subseteq H\cap K$}
\handline{72}{774}{354}{所以 $G=H\oplus K\cong H\oplus\Z^r$.}
\handline[NoteRed]{72}{809}{471}{反例：\circled{1} 商群不自由，则直和关系不一定存在.}
\handline[NoteRed]{72}{841}{505}{$\Z_4/\Z_2\cong\Z_2$，$\Z_2$ 不自由，而 $\Z_4\not\cong\Z_2\oplus\Z_2$.}
\handline[NoteRed]{72}{876}{402}{\circled{2} $G/H$ 自由不可改成 $H$ 自由.}
\handline[NoteRed]{72}{913}{492}{$G=\Z$，$H=2\Z$，$H$ 自由，但 $G/H\cong\Z_2$ 不自由.}
\handline[NoteRed]{72}{950}{582}{而 $G\not\cong H\oplus\Z_2$，因为 $G$ 中无 $2$ 阶元，$\Z_2$ 中有 $2$ 阶元.}
\end{handpage}

% ===== Line-layout page-102.tex =====
\begin{handpage}{102}
\handline{72}{85}{670}{\textbf{Lemma. (II)} 若交换群 $G$ 有两种直和分解，$G=H\oplus\langle a\rangle=K\oplus\langle b\rangle$.}
\handline{72}{116}{475}{其中 $\langle a\rangle,\langle b\rangle$ 均为无限循环群，则 $H\cong K$.}
\handline{72}{145}{323}{\textbf{Pf.} 考虑 $K/(K\cap H)$、$H/(K\cap H)$.}
\handline{72}{175}{612}{由第二同构定理，$K/(K\cap H)\cong(K+H)/H\le G/H\cong\langle a\rangle\cong\Z$.}
\handline{72}{208}{699}{而由于 $\Z$ 的子群要么为 $\{0\}$，要么也为无限循环群. $H/(K\cap H)$ 同理.}
\handline{72}{240}{380}{\circled{1} 若 $K/(K\cap H)$、$H/(K\cap H)$ 均为 $\{0\}$，}
\handline{72}{269}{279}{$\Rightarrow K=K\cap H=H$.}
\handline{72}{300}{498}{\circled{2} 若 $K/(K\cap H)$、$H/(K\cap H)$ 至少有一个是无限循环群.}
\handline{72}{334}{753}{不妨设 $K/(K\cap H)=\langle\bar u\rangle$，结合引理 (I) 的证明过程 $\Rightarrow K=(K\cap H)\oplus\langle u\rangle$，$u\in K$.}
\handline{72}{367}{589}{于是 $G=(K\cap H)\oplus\langle u\rangle\oplus\langle b\rangle$. 另一方面 $G=H\oplus\langle a\rangle$}
\handline{72}{400}{334}{所以 $G/(K\cap H)\cong\langle u\rangle\oplus\langle b\rangle$}
\handline{72}{434}{701}{考虑 $\varphi:G\to H/(K\cap H)\oplus\langle a\rangle$，则 $\varphi$ 是满同态，且 $\ker\varphi=K\cap H$.}
\handline{72}{467}{603}{$h+na\mapsto(h+K\cap H,na)$，所以 $G/(K\cap H)\cong H/(K\cap H)\oplus\langle a\rangle$.}
\handline{72}{502}{448}{若 $H/(K\cap H)=\{0\}$，则 $\langle u\rangle\oplus\langle b\rangle\cong\langle a\rangle$}
\handline{72}{538}{740}{所以 $\exists m,n\in\omega$ s.t. $u=ma$，$b=na\Rightarrow nu=mb\in\langle u\rangle\cap\langle b\rangle=\{0\}$}
\handline{72}{573}{575}{所以 $nu=mb=0$ 与 $\langle u\rangle,\langle b\rangle$ 为无限阶循环群矛盾！}
\handline{72}{606}{703}{所以 $H/(K\cap H)$ 是无限阶循环群. 同理可设 $H=(K\cap H)\oplus\langle w\rangle$，$w\in H$.}
\handline{72}{642}{642}{定义 $\varphi:H\to K$，可验证，$\varphi$ 是一个同构，所以 $H\cong K$.}
\handline{72}{671}{110}{$w\mapsto u$}
\handline{72}{704}{251}{且 $\varphi|_{K\cap H}=\id_{K\cap H}$.}
\handline[NoteRed]{72}{748}{524}{反例：$G=\Z$，$\langle a\rangle=2\Z$，$\langle b\rangle=3\Z$.}
\handline[NoteRed]{72}{782}{347}{$G/\langle a\rangle\cong\Z_2$，$G/\langle b\rangle\cong\Z_3$，}
\handline[NoteRed]{72}{817}{203}{但 $\Z_2\ne\Z_3$.}
\handline[NoteRed]{72}{855}{620}{这是因为 $(\langle a\rangle,\Z_2)$、$(\langle b\rangle,\Z_3)$ 构不成直和条件.}
\end{handpage}

% ===== Line-layout page-103.tex =====
\begin{handpage}{103}
\handline{72}{93}{326}{\textbf{Thm.} $G$：有限生成交换群}
\handline{72}{119}{401}{若 $G=C_1\oplus C_2\oplus\cdots\oplus C_r\oplus C_{r+1}\oplus\cdots\oplus C_k$}
\handline{72}{144}{338}{$=D_1\oplus D_2\oplus\cdots\oplus D_t\oplus D_{t+1}\oplus\cdots\oplus D_s$}
\handline{72}{174}{543}{其中 $C_1,\ldots,C_r$ 为有限阶循环群，$C_{r+1},\ldots,C_k$ 为无限循环群}
\handline{72}{200}{494}{$D_1,\ldots,D_t$ 为有限阶循环群，$D_{t+1},\ldots,D_s$ 为无限循环群}
\handline{72}{228}{587}{且设 $|C_i|=m_i$，$1\le i\le r$；$|D_j|=n_j$，$1\le j\le s$；$m_1\mid m_2\mid\cdots\mid m_r$，$n_1\mid\cdots\mid n_s$}
\handline{72}{259}{373}{则 $r=t$，$k=s$，$m_i=n_i$，$\forall1\le i\le r$.}
\handline[AnnotationBlue]{72}{284}{717}{$G$ 若有两个无限阶循环群的分解，结合引理 (II)，一个个往外拿，证明是 trivial 的.}
\handline[AnnotationBlue]{72}{311}{432}{$G$ 不可能同时有一个全为无限循环群的分解}
\handline[AnnotationBlue]{72}{339}{696}{以及一个存在有限循环群的分解，因为有限循环群中元素的阶有限.}
\handline[AnnotationBlue]{72}{367}{475}{因此这里只考虑 $G$ 有两个存在有限循环群的分解}
\handline{72}{399}{391}{\textbf{Pf.} 由引理 (II)，不妨设 $k-r=s-t$}
\handline{72}{426}{355}{且 $C_1\oplus\cdots\oplus C_r\cong D_1\oplus\cdots\oplus D_t$. \redann{(*)}}
\handline{72}{454}{359}{只需证 $r=t$ 且 $m_i=n_i$，$\forall1\le i\le r$.}
\handline{72}{483}{565}{由整除，\redann{(*)} 左边阶数最大的为 $m_r$，右边阶数最大的为 $n_t$.}
\handline{72}{512}{246}{再由同构 $\Rightarrow m_r=n_t$.}
\handline{72}{544}{567}{$m_{r-1}C_1\oplus m_{r-1}C_2\oplus\cdots\oplus m_{r-1}C_r\cong m_{r-1}D_1\oplus\cdots\oplus m_{r-1}D_t$.}
\handline{72}{575}{459}{由于 $m_i\mid m_{r-1}$，$\forall1\le i\le r-1$，所以 $m_{r-1}C_i\cong\{0\}$.}
\handline{72}{606}{399}{所以 $m_{r-1}C_r\cong m_{r-1}D_1\oplus\cdots\oplus m_{r-1}D_t$.}
\handline{72}{641}{518}{而 $|m_{r-1}C_r|=\dfrac{|C_r|}{m_{r-1}}=\dfrac{m_r}{m_{r-1}}=\dfrac{n_t}{m_{r-1}}=\dfrac{|D_t|}{m_{r-1}}=|m_{r-1}D_t|$.}
\handline{72}{673}{547}{所以 $m_{r-1}D_i$，$1\le i\le t-1$，不可能包含有更多非零项.}
\handline{72}{704}{508}{再结合直和关系 $\Rightarrow m_{r-1}D_i\cong\{0\}$，$1\le i\le t-1$.}
\handline{72}{738}{466}{特别地，由于 $m_{r-1}D_{t-1}\cong\{0\}$，所以 $n_{t-1}\mid m_{r-1}$}
\handline{72}{770}{284}{同理可证 $m_{r-1}\mid n_{t-1}$.}
\handline{72}{804}{503}{不断重复上述过程 $\Rightarrow r=t$，$m_i=n_i$，$\forall1\le i\le r$.}
\handline[AnnotationBlue]{72}{840}{784}{结合存在性与唯一性，我们得到了有限生成交换群的结构定理.}
\handline{72}{886}{539}{\textbf{Thm.} $G$：有限生成交换群，则 $G=C_1\oplus\cdots\oplus C_r\oplus G_1$.}
\handline{72}{917}{590}{其中 $|C_i|=m_i$，$1\le i\le r$，$C_i$ 为 $m_i$ 阶循环群，且 $m_1\mid\cdots\mid m_r$.}
\handline{72}{950}{597}{$G_1$ 为自由交换群，$G_1\cong\Z^n$，$n\in\omega$，且上述分解唯一.}
\end{handpage}

% ===== Line-layout page-104.tex =====
\begin{handpage}{104}
\handline{72}{88}{650}{\textbf{Def.} $G$：有限生成交换群，$G=C_1\oplus C_2\oplus\cdots\oplus C_r\oplus C_{r+1}\oplus\cdots\oplus C_k$.}
\handline{72}{118}{561}{$|C_i|=m_i\in\omega$，$1\le i\le r$，$C_i$ 为 $m_i$ 阶循环群，且 $m_1\mid m_2\mid\cdots\mid m_r$}
\handline{72}{148}{600}{$C_{r+1},\ldots,C_k$ 为无限阶循环群，上述分解称为不变分解.}
\handline{72}{178}{555}{$k-r$ 称为 $G$ 的秩，$m_1,m_2,\ldots,m_r$ 称为 $G$ 的不变因子组}
\handline{72}{207}{558}{$C_1\oplus\cdots\oplus C_r$ 称为 $G$ 的挠分量 \ann{(Torsion Component)}}
\handline{72}{237}{555}{$C_{r+1}\oplus\cdots\oplus C_k$ 称为 $G$ 的自由分量 \ann{(Free Component)}}
\handline{72}{268}{848}{\textbf{Def.} $G$：交换群，$g\in G$. 若 $\exists n\in\omega$ s.t. $ng=0$，则称 $g$ 为挠元 \ann{(torsion element)}}
\handline{72}{298}{742}{$T(G):=\{g\in G\mid\exists n\ge1\text{ s.t. }ng=0\}$，称为 $G$ 的挠子群 \ann{(torsion subgroup)}}
\handline[AnnotationBlue]{72}{331}{451}{\circled{1} 挠元可类比交换幺环的幂零元，而幂零根.}
\handline[AnnotationBlue]{72}{359}{366}{挠元即为 $G$ 中有限阶元素.}
\handline[AnnotationBlue]{72}{391}{609}{\circled{2} 由有限生成交换群的结构定理，$G$：有限生成交换群.}
\handline[AnnotationBlue]{72}{420}{295}{$T(G)$ 一定是有限交换群.}
\handline[AnnotationBlue]{72}{452}{690}{$\exists G$ 的自由子群 $F$ s.t. $G=T(G)\oplus F$，且 $F$ 在同构意义下唯一.}
\handline[AnnotationBlue]{72}{483}{382}{\circled{3} 秩即为自由分量的分量个数.}
\handline[AnnotationBlue]{72}{515}{715}{\circled{4} 不变因子不是那些循环子群自身，而是 $m_1,m_2,\ldots,m_r$ 这些正整数.}
\handline[AnnotationBlue]{72}{545}{367}{且 $m_1,\ldots,m_r$ 均严格大于 $1$.}
\handline[AnnotationBlue]{72}{577}{592}{不变因子组是一个有序数组 $(m_1,\ldots,m_r)$，满足 $m_1\mid\cdots\mid m_r$.}
\handline[AnnotationBlue]{72}{608}{532}{任给一个有限交换群，可以唯一确定不变因子组.}
\handline[AnnotationBlue]{72}{642}{785}{反之，给定一个不变因子组，可以在同构意义下唯一地给出一个有限交换群.}
\handline[AnnotationBlue]{72}{675}{771}{从秩、不变因子这些术语可以看出：这与 Jordan 标准型及有理标准型理论}
\handline[AnnotationBlue]{72}{706}{300}{有着千丝万缕的联系.}
\handline[AnnotationBlue]{72}{742}{771}{事实上，这两者可以统一在主理想整环上的有限生成模的结构定理之下.}
\handline[AnnotationBlue]{72}{776}{561}{前者考虑 $\Z$-模，后者考虑 $K[x]$-模，$K$ 为一个域.}
\handline[AnnotationBlue]{72}{811}{676}{而 $K[x]$-模恰与 $K$-线性空间上的线性变换全体同构.}
\handline[AnnotationBlue]{72}{846}{658}{而 Jordan 标准型理论中还出现了准素分解、初等因子这些，}
\handline[AnnotationBlue]{72}{881}{531}{同样地在这里也可以类似定义出这些概念.}
\end{handpage}

% ===== Line-layout page-105.tex =====
\begin{handpage}{105}
\handline{72}{127}{461}{\textbf{Thm.} \ann{（准素分解）(Primary Decomposition)}}
\handline{72}{153}{560}{$G$：有限交换群，$|G|=p_1^{\alpha_1}\cdots p_k^{\alpha_k}$，$p_1,\ldots,p_k$ 为两两不同的素数.}
\handline{72}{182}{580}{$P_i$ 是 $p_i$-Sylow 子群，则 $G=P_1\oplus\cdots\oplus P_k$，称为 $G$ 的准素分解.}
\handline{72}{210}{506}{\textbf{Pf.} 由于 $G$ 是交换群，所以所有 $p_i$-Sylow 子群都正规}
\handline{72}{239}{523}{从而所有 $p_i$-Sylow 子群各只有一个，故若分解存在则唯一.}
\handline{72}{268}{353}{对于 $a\in P_1$，则 $o(a)$ 必为 $p_1$ 的幂次.}
\handline{72}{297}{385}{若 $a\in P_2+\cdots+P_k$，则 $a=b_2+\cdots+b_k$.}
\handline{72}{329}{682}{$a^{p_2^{\alpha_2}\cdots p_k^{\alpha_k}}=(b_2+\cdots+b_k)^{p_2^{\alpha_2}\cdots p_k^{\alpha_k}}=0\Rightarrow o(a)\mid p_2^{\alpha_2}\cdots p_k^{\alpha_k}\Rightarrow o(a)=1$，所以 $a=0$.}
\handline{72}{358}{303}{所以 $P_1\cap(P_2+\cdots+P_k)=\{0\}$}
\handline{72}{386}{399}{同理 $P_i\cap(P_1+\cdots+P_{i-1}+P_{i+1}+\cdots+P_k)=\{0\}$.}
\handline{72}{418}{324}{于是可定义 $\widetilde G=P_1\oplus\cdots\oplus P_k$.}
\handline{72}{447}{475}{显然 $\widetilde G\le G$，而 $|\widetilde G|=p_1^{\alpha_1}\cdots p_k^{\alpha_k}=|G|$，所以 $\widetilde G=G$}
\handline{72}{477}{258}{所以 $G=P_1\oplus\cdots\oplus P_k$}
\handline{72}{511}{704}{\textbf{Def.} 上述分解中 $P_i$ 称为 $G$ 的 $p_i$-准素分量 \ann{(primary component)}}
\handline{72}{542}{514}{$p_1^{\alpha_1},\ldots,p_k^{\alpha_k}$ 称为 $G$ 的准素因子 \ann{(primary divisors)}}
\handline{72}{571}{377}{$(p_1^{\alpha_1},\ldots,p_k^{\alpha_k})$ 称为 $G$ 的准素因子组}
\handline[AnnotationBlue]{72}{603}{683}{注意，这一分解与不变分解不同，每个 $p_i$-准素分量不一定是循环群.}
\handline{72}{639}{647}{\textbf{Def.} $n\in\omega$，$1\le\alpha_1\le\alpha_2\le\cdots\le\alpha_k$，$\alpha_1+\alpha_2+\cdots+\alpha_k=n$，称为 $n$ 的分拆.}
\handline{72}{673}{439}{\textbf{Lemma.} $p$ 为素数，$G$ 为 $p^n$ 阶交换群，$n\in\omega$，}
\handline{72}{708}{407}{则 $G$ 的分类与 $n$ 的分拆一一对应.}
\handline{72}{742}{612}{\textbf{Pf.} $G$ 有限交换，所以 $G=C_1\oplus\cdots\oplus C_k$ 为不变分解.}
\handline{72}{774}{310}{其中 $|C_1|\mid|C_2|\mid\cdots\mid|C_k|$.}
\handline{72}{806}{535}{设 $|C_i|=p^{\alpha_i}$，则整除关系等价于 $\alpha_1\le\cdots\le\alpha_k$.}
\handline{72}{840}{330}{$|G|=p^n\Longleftrightarrow\alpha_1+\cdots+\alpha_k=n$.}
\handline{72}{876}{398}{所以 $G$ 的分类与 $n$ 的分拆一一对应.}
\handline[AnnotationBlue]{72}{914}{722}{有了这一引理，可以在准素分解的基础上，将每个准素分量}
\handline[AnnotationBlue]{72}{950}{846}{再进一步进行不变分解，得到更细的分解，使得每个分量均为循环群.}
\end{handpage}

% ===== Line-layout page-106.tex =====
\begin{handpage}{106}
\handline{72}{129}{610}{Cor.\quad$n=p_1^{\alpha_1}\cdots p_r^{\alpha_r}$，则不同的 $n$ 阶交换群的个数为 $p(\alpha_1)\cdots p(\alpha_r)$.}
\handline{72}{154}{553}{其中 $p(k)$，$k\in\N$，表示 $k$ 的分拆的种类数，称为分拆函数.}
\handline{72}{181}{333}{Def.\quad$n=p_1^{\alpha_1}\cdots p_r^{\alpha_r}$.\quad 给定分拆.}
\handline{72}{202}{512}{$\alpha_1$ 的分拆：$1\le\alpha_{1,1}\le\alpha_{1,2}\le\cdots\le\alpha_{1,k_1}\quad\alpha_{1,1}+\cdots+\alpha_{1,k_1}=\alpha_1$}
\handline{72}{228}{525}{$\alpha_2$ 的分拆：$1\le\alpha_{2,1}\le\alpha_{2,2}\le\cdots\le\alpha_{2,k_2}\quad\alpha_{2,1}+\cdots+\alpha_{2,k_2}=\alpha_2$}
\handline{72}{256}{394}{$\vdots\qquad\qquad\qquad\vdots\qquad\qquad\qquad\vdots$}
\handline{72}{286}{538}{$\alpha_r$ 的分拆：$1\le\alpha_{r,1}\le\alpha_{r,2}\le\cdots\le\alpha_{r,k_r}\quad\alpha_{r,1}+\cdots+\alpha_{r,k_r}=\alpha_r$}
\handline{72}{307}{552}{$p_i^{\alpha_{i,j}}\quad1\le i\le r,1\le j\le k_i$\quad 称为初等因子\ann{（elementary divisors）}}
\handline{72}{336}{479}{$(p_i^{\alpha_{i,j}})_{\substack{1\le i\le r\\1\le j\le k_i}}$\quad 这一有序数组称为初等因子组.}
\handline[AnnotationBlue]{72}{367}{460}{这里的有序只需保持同一 $\alpha_i$ 的分拆中的顺序，}
\handline[AnnotationBlue]{72}{383}{736}{以及同一个 $\alpha_i$ 拆出来的要排在一起，而两两不同 $\alpha_i$ 的整体之间的顺序可以任意换.}
\handline[AnnotationBlue]{72}{411}{533}{给定这些初等因子组可唯一确定一个 $n$ 阶交换群 $G$.}
\handline[AnnotationBlue]{72}{441}{571}{$G\cong\Z_{p_1^{\alpha_{1,1}}}\oplus\cdots\oplus\Z_{p_1^{\alpha_{1,k_1}}}\oplus\cdots\oplus\Z_{p_r^{\alpha_{r,1}}}\oplus\cdots\oplus\Z_{p_r^{\alpha_{r,k_r}}}$.}
\handline{72}{472}{460}{Thm.\quad\ann{（初等分解）（Elementary Decomposition）}}
\handline{72}{500}{750}{$G$：有限交换群，则 $G$ 可唯一地分解为 $G\cong\bigoplus_{i=1}^r(\Z_{p_i^{\alpha_{i,1}}}\oplus\cdots\oplus\Z_{p_i^{\alpha_{i,k_i}}})$}
\handline{72}{526}{518}{且 $|G|=p_1^{\alpha_1}\cdots p_r^{\alpha_r}$.\quad$\alpha_{i,1},\ldots,\alpha_{i,k_r}$ 为 $\alpha_i$ 的分拆.}
\handline[AnnotationBlue]{72}{554}{614}{证明很 trivial.先做准素分解，再对准素分量做不变分解.}
\handline[AnnotationBlue]{72}{583}{592}{$p_i^{\alpha_{i,j}}$ 称为 $G$ 的初等因子.$(p_i^{\alpha_{i,j}})$ 称为 $G$ 的初等因子组.}
\handline{72}{616}{467}{Ex.\quad\circled{1} $16$ 阶交换群.$16=2^4$.考虑 $4$ 的分拆.}
\handline{72}{643}{706}{$1+3:\ \Z_2\oplus\Z_8\quad1+1+2:\ \Z_2\oplus\Z_2\oplus\Z_4\quad1+1+1+1\quad\Z_2\oplus\Z_2\oplus\Z_2\oplus\Z_2$}
\handline{72}{675}{419}{$2+2:\ \Z_4\oplus\Z_4\qquad4:\ \Z_{16}$.\quad 共 $5$ 种}
\handline{72}{708}{748}{\circled{2} $12$ 阶交换群.$12=2^2\times3$.共 $2$ 种 $12$ 阶交换群，因为 $p(2)\cdot p(1)=2\times1=2$}
\handline{72}{739}{405}{$1^\circ\quad\Z_2\oplus\Z_2\oplus\Z_3\qquad2^\circ\quad\Z_4\oplus\Z_3$.}
\handline{72}{767}{470}{但都不是不变分解，因为 $2\nmid3$，$3\nmid4$.}
\handline{72}{800}{805}{而 $1^\circ$ 的不变分解为 $\Z_2\oplus\Z_6$.不变因子组为 $2,6$.$2^\circ$ 的不变分解为 $\Z_{12}$.不变因子为 $12$.}
\handline{72}{840}{770}{\circled{3} $\Z_{180}$\quad$180=2^2\cdot3^2\cdot5\quad\therefore\Z_{180}\cong\Z_4\oplus\Z_9\oplus\Z_5$.\redann{（还可由 CRT 得到）}}
\handline{72}{865}{556}{而 $\Z_4,\Z_9$ 不能再分解了.$\therefore$ 初等因子组为 $4,9,5$.}
\handline{72}{894}{644}{而 $\Z_{180}$ 本身就是一个有限循环群，故 $\Z_{180}$ 即为不变分解：}
\handline{72}{930}{182}{不变因子为 $180$.}
\end{handpage}

% ===== Line-layout page-107.tex =====
\begin{handpage}{107}
\handtitle{72}{136}{156}{§3.9 可解群.}
\handline{72}{167}{569}{Def.\quad\circled{1} $G$：群.\quad$G^{(1)}:=[G,G]=\langle[a,b]=a^{-1}b^{-1}ab\mid a,b\in G\rangle$.}
\handline{72}{192}{502}{即 $G$ 的换位子群，称为 $G$ 的导出群.显然 $G^{(1)}\lhd G$}
\handline{72}{220}{459}{令 $G^{(2)}:=[G^{(1)},G^{(1)}]$ 即 $G^{(1)}$ 的换位子群.$\cdots$}
\handline{72}{252}{224}{$G^{(k)}:=[G^{(k-1)},G^{(k-1)}]$}
\handline{72}{280}{500}{$G\triangleright G^{(1)}\triangleright G^{(2)}\triangleright\cdots\triangleright G^{(k)}\triangleright\cdots$ 称为 $G$ 的导出列.}
\handline{72}{305}{480}{\circled{2} 若 $G$ 的导出列终止于 $\{e\}$，则称 $G$ 为可解群.}
\handline{72}{331}{274}{即 $\exists k\in\N\ \text{s.t.}\ G^{(k)}=\{e\}$.}
\handline{72}{364}{555}{\circled{1} 显然，交换群一定是可解群，因为 $G^{(1)}=[G,G]=\{e\}$.}
\handline{72}{392}{485}{\circled{2} 非交换单群一定不是可解群.由于 $G^{(1)}\lhd G$.}
\handline{72}{420}{380}{而 $G$ 是单群 $\Rightarrow G^{(1)}=\{e\}$ 或 $G^{(1)}=G$}
\handline{72}{451}{692}{而 $G$ 非交换 $\Rightarrow G^{(1)}\ne\{e\}\Rightarrow G^{(1)}=G\Rightarrow G^{(k)}=G\ne\{e\}\ \forall k\in\N$.}
\handline{72}{485}{626}{Thm.\quad$G$ 可解 $\Longleftrightarrow$ 存在次正规群列.$G=G_0\triangleright\cdots\triangleright G_r=\{e\}$.}
\handline{72}{515}{388}{其中商因子 $G_i/G_{i+1}$ 均为交换群.}
\handline{72}{546}{784}{Pf.\quad“$\Rightarrow$” $G^{(k)}/G^{(k+1)}=G^{(k)}/[G^{(k)},G^{(k)}]$ 即 $G^{(k)}$ 的交换化，一定是交换群.}
\handline{72}{576}{523}{“$\Leftarrow$” Lemma：$H\lhd G$，若 $G/H$ 交换，则 $[G,G]\subseteq H$.}
\handline{72}{606}{578}{$\forall a,b\in G\quad[a,b]H=a^{-1}b^{-1}abH=(a^{-1}H)(b^{-1}H)(aH)(bH)$}
\handline{72}{640}{493}{$=(a^{-1}H)(aH)(b^{-1}H)(bH)=(a^{-1}H)(Ha)(b^{-1}H)(Hb)$}
\handline{72}{671}{447}{$=(a^{-1}Ha)(b^{-1}Hb)=HH=H.\quad\therefore[a,b]\in H$}
\handline{72}{706}{180}{$\therefore[G,G]\subseteq H$.}
\handline{72}{743}{812}{而 $G_0/G_1$ 交换，$G_1\lhd G_0\Rightarrow[G,G]\subseteq G_1$ 即 $G^{(1)}\subseteq G_1\quad\therefore G^{(2)}\subseteq[G^{(1)},G^{(1)}]\subseteq[G_1,G_1]$}
\handline{72}{777}{574}{$G_1/G_2$ 交换，$G_2\lhd G_1\Rightarrow[G_1,G_1]\subseteq G_2\Rightarrow G^{(2)}\subseteq G_2$.}
\handline{72}{807}{350}{$\cdots\Rightarrow G^{(k)}\subseteq G_k\quad1\le k\le r$.}
\handline{72}{847}{621}{而 $G_r=\{e\}\quad\therefore G^{(r)}\subseteq G_r=\{e\}\quad\therefore G^{(r)}=\{e\}$ 故 $G$ 可解}
\end{handpage}

% ===== Line-layout page-108.tex =====
\begin{handpage}{108}
\handline{72}{92}{350}{Thm.\quad\ann{（可解性质的稳定性）}}
\handline{72}{118}{440}{\circled{1} 可解群的子群和商群也可解.}
\handline{72}{146}{470}{\circled{2} 若 $K\lhd G$，$K,G/K$ 均可解，则 $G$ 可解.}
\handline{72}{175}{620}{Pf.\quad\circled{1} $G$ 可解.若 $H\le G$，则 $H^{(1)}\subseteq G^{(1)}\Rightarrow H^{(k)}\subseteq G^{(k)}$.}
\handline{72}{200}{430}{若 $G^{(k)}=\{e\}\Rightarrow H^{(k)}=\{e\}$，$\therefore H$ 可解.}
\handline{72}{227}{370}{若 $H\lhd G$.商群 $G/H$.}
\handline{72}{254}{530}{考虑典则投影 $\varphi:G\to G/H$.则 $\varphi$ 是满同态.}
\handline{72}{279}{150}{$a\mapsto aH$}
\handline{72}{305}{260}{于是 $G/H=\varphi(G)$.}
\handline{72}{340}{620}{而 $[G/H,G/H]=[\varphi(G),\varphi(G)]=\varphi([G,G])=\varphi(G^{(1)})$.}
\handline{72}{367}{470}{$\Rightarrow$ 可归纳得 $(G/H)^{(k)}=\varphi(G^{(k)})$.}
\handline{72}{397}{635}{若 $G^{(k)}=\{e\}$，则 $(G/H)^{(k)}=\varphi(\{e\})=\{e_{G/H}\}$，$\therefore G/H$ 可解.}
\handline{72}{427}{765}{\circled{2} $K$ 可解 $\Rightarrow\exists m\in\N:K^{(m)}=\{e\}$.$G/K$ 可解 $\Rightarrow\exists n\in\N\ (G/K)^{(n)}=\{e\}$.}
\handline{72}{457}{550}{考虑典则投影 $\varphi:G\to G/K$.则 $\ker(\varphi)=K$.}
\handline{72}{481}{150}{$a\mapsto aK$}
\handline{72}{516}{400}{$\varphi(G^{(n)})=(G/K)^{(n)}=\{e\}$.}
\handline{72}{549}{375}{$\Rightarrow G^{(n)}\subseteq\ker(\varphi)=K$.}
\handline{72}{578}{630}{$\Rightarrow G^{(n+m)}\subseteq K^{(m)}=\{e\}\Rightarrow G^{(n+m)}=\{e\}$，$\therefore G$ 可解.}
\handline[AnnotationBlue]{72}{608}{440}{若 $G$ 是交换单群，则 $G$ 可解.}
\handline[AnnotationBlue]{72}{639}{645}{进一步，若 $G$ 还有限，则 $G$ 一定为素数阶循环群 $\Z_p$.}
\handline[AnnotationBlue]{72}{672}{550}{若否.$|G|=pq$.$p$ 为素数，$q\in\N$ 且 $q>1$.}
\handline[AnnotationBlue]{72}{708}{710}{由 Sylow-I，$G$ 一定有 $p$ 阶子群 $H$，而 $G$ 交换，$\therefore H$ 一定正规.}
\handline[AnnotationBlue]{72}{744}{660}{而 $q>1\Rightarrow H$ 比 $G$ 严格小，与 $G$ 为单群矛盾！}
\handline[AnnotationBlue]{72}{777}{400}{故 $G\cong\Z_p$，$p$ 为素数.}
\handline[AnnotationBlue]{72}{808}{440}{此时 $G$ 的导出列为 $G\triangleright\{e\}$.}
\handline[AnnotationBlue]{72}{846}{735}{$G/\{e\}\cong\Z_p$.由此可见素数阶循环群与可解性有些联系.}
\end{handpage}

% ===== Line-layout page-109.tex =====
\begin{handpage}{109}
\handline{72}{114}{660}{Thm.\quad$G$ 为有限群，则 $G$ 可解 $\iff G$ 的合成因子都是素数阶群.}
\handline{72}{142}{430}{Pf.\quad“$\Leftarrow$”.\quad$G_0\triangleright G_1\triangleright\cdots\triangleright G_k=\{e\}$.}
\handline{72}{167}{420}{$G_i/G_{i+1}\cong\Z_{p_i}$ 是交换群 $\Rightarrow G$ 可解.}
\handline{72}{196}{650}{“$\Rightarrow$” 设 $G$ 有一个合成列.$G=G_0\triangleright\cdots\triangleright G_k=\{e\}$.$G_i/G_{i+1}$ 是单群.}
\handline{72}{226}{440}{由于 $G$ 可解，由可解的稳定性.}
\handline{72}{256}{455}{$\Rightarrow G_i,G_i/G_{i+1}$ 均可解，$\forall0\le i\le k$.}
\handline{72}{284}{490}{由于 $G_i/G_{i+1}$ 是单群，且可解，故一定交换.}
\handline{72}{310}{430}{而有限交换单群一定是素数阶群.}
\handline{72}{338}{510}{$\therefore G_i/G_{i+1}\cong\Z_{p_i}$，$p_i$ 为素数，$0\le i\le k-1$.}
\handline{72}{368}{360}{设 $K\lhd G$.$G/K$.$1=\{e\}$.}
\handline{72}{396}{765}{则有 $1\xrightarrow{e}K\xrightarrow{i}G\xrightarrow{\pi}G/K\xrightarrow{\eta}1$.称为群同态的短正合列.}
\handline[NoteRed]{72}{454}{495}{$e:1\to K$.即从平凡群到 $K$ 的唯一群同态.}
\handline[NoteRed]{72}{481}{140}{$e_1\mapsto e_K$}
\handline[NoteRed]{72}{518}{180}{$i:K\to G$}
\handline[NoteRed]{72}{549}{140}{$k\mapsto k$}
\handline[NoteRed]{72}{582}{220}{$\pi:G\to G/K$}
\handline[NoteRed]{72}{612}{160}{$g\mapsto gK$}
\handline[NoteRed]{72}{645}{500}{$\eta:G/K\to1$.即将 $G/K$ 中所有元素映为 $e$.}
\handline[NoteRed]{72}{678}{170}{$gK\mapsto e$.}
\handdiagram{40}{458}{18}{143}{\draw[AnnotationBlue,line width=.6pt] (0,0) .. controls (.5,-.12) and (.4,-.3) .. (.25,-.6) -- (.25,-1.25) .. controls (.25,-1.45) and (0,-1.4) .. (.25,-1.55) .. controls (0,-1.7) and (.25,-1.7) .. (.25,-1.9) -- (.25,-2.65) .. controls (.25,-2.8) and (.1,-2.95) .. (0,-3.05);}
\handline[AnnotationBlue]{332}{484}{270}{满足：$\ker(i)=\operatorname{Im}(e)=1$}
\handline[AnnotationBlue]{460}{514}{210}{$\ker(\pi)=\operatorname{Im}(i)=K$}
\handline[AnnotationBlue]{413}{543}{240}{$\ker(\eta)=\operatorname{Im}(\pi)=G/K$.}
\handline[AnnotationBlue]{446}{576}{355}{后一个同态的核是前一个同态的像}
\handline{72}{721}{785}{对于可解群.$G$.$G=G_0\triangleright G_1\triangleright\cdots\triangleright G_k=\{e\}$.$G_i/G_{i+1}\cong A_i$.$A_i$ 交换.}
\handline{72}{755}{605}{每一步都有短正合列 $1\to G_{i+1}\to G_i\to A_i\to1$.}
\handline{72}{791}{610}{群扩张：把 $G_{i+1}$ 用交换群 $A_i$ 扩张成了 $G_i$.}
\handline{72}{826}{805}{可解群就可以看成从平凡群出发，每一步都用交换群扩张而成的群.}
\end{handpage}

% ===== Line-layout page-110.tex =====
\begin{handpage}{110}
\handline[AnnotationBlue]{72}{108}{610}{最后介绍一个暂时无法证明但很重要的大定理.}
\handline[AnnotationBlue]{72}{140}{370}{Burnside 可解性定理.}
\handline[AnnotationBlue]{72}{167}{640}{若 $|G|=p^aq^b$.$a,b\in\N$.$p,q$ 为素数.则 $G$ 可解.}
\handline[AnnotationBlue]{72}{199}{570}{它的证明需要用群表示论，借助特征标.}
\handline[AnnotationBlue]{72}{230}{440}{具体可看 Dummit--Foote \S 19.2}
\handline[AnnotationBlue]{72}{261}{520}{Burnside 猜想：奇数阶群为可解群}
\handline[AnnotationBlue]{72}{292}{750}{于 $1963$ 年被 Feit--Thompson 解决.现在叫 Feit--Thompson 定理}
\handline[AnnotationBlue]{72}{325}{300}{论文长达 $255$ 页.}
\end{handpage}

\end{document}
